% Generated by build_chapter.py from the GY1969 files. Do not edit by hand; edit the files and rebuild.
\chapter{The academic years 1969--70 and 1970--71}\label{ch:1969}
This chapter follows all 55 catalogue records, 174--228, through the supplied
leaves for 1969--70 and 1970--71. Source order is decisive: Hewitt's November
1969 lecture precedes the February 1970 monodromy entry, and the finite-field
Ramsey question precedes the next sheaf lecture despite its higher record
number. The opening Buchstaber fragment on P000282 belongs to the preceding
year; the Kirillov leaf on P000360 belongs to the following one. Both are
documented as boundary evidence rather than assigned to adjacent speakers
\citep[catalogue, records 173--229]{mccme2026gelfand}.

The principal continuous thread is Kazhdan's course in schemes, sheaves,
affine cohomology and projective constructions, joined to Griffiths's
periods by monodromy, residues and vanishing cycles. A second thread passes
from McDuff's factors and Sinai's quantum chains to Gelfand's current-group
questions and Perelomov's coherent states. Finite groups, four subspaces,
matrix deformations, real curves, characteristic numbers, torsion and
inverse scattering repeatedly connect these strands. Announcements and
reading programmes remain catalogue entries even when they contain no
identifiable mathematical assertion.

The occurrence index contains 290 items and the problem index
17 definite questions or conjectures. Counts measure editorial
segmentation, not originality. The historical classes are:
\begin{center}
\begin{tabular}{lrrrrrr}
\toprule
Class&A&B&C&D&E&F\\
\midrule
Items&19&8&4&0&231&28\\
\bottomrule
\end{tabular}
\end{center}
Class E dominates: the seminar mixed exposition, verification and
research questions. Chronology matters.
McDuff's continuum of factors had appeared before her October lecture;
Griffiths's period results had circulated before their 1970 journal issue;
Perelomov's lattice completeness paper preceded his April lecture.
Conversely, later general solutions must not be substituted for narrower
seminar constructions. The general modularity theorem, the Gudkov
congruence and analytic--combinatorial torsion are followed through
their later literature in their respective entries.

Two bounded developments are supplied with complete proofs and independent
checks. For current groups, an infinite tensor product of closely paired
evaluations gives a strongly continuous irreducible representation of
$C^1([0,1],SU(2))$, still irreducible on smooth currents, with support
$[0,1]$. A countersequence shows why its topology cannot be weakened to
uniform continuity. Operator-norm-continuous smooth representations have
finite support by the checked classification
\citep{janssensneeb2015norm}; the notebook's unspecified category remains
an explicit qualification. For the four-subspace question, the orbit of
$S(3,2)^{\oplus2}$ is open dense in
$\operatorname{Gr}(4,6)^4$, giving the requested deformation of
$S(5,2)\oplus S(1,2)$. Its mathematical conclusion is settled here;
the exact historical attribution of that elementary proof is not.
Neither development claims novelty.

Two further annotations audit settled material. A faithful Bell-block
state disproves the claim that commuting finite marginals necessarily
have compatible one-site product eigenbases. Exact rational computations
check all four-site marginals and commutators at two parameter values.
Wavelet scaling verifies the positive answer to the product-uncertainty
question, without validating the separate proposed Gram--Schmidt method.
The arbitrary-sequence coherent-state completeness problem remains
qualified: checked partial criteria do not settle all irregular sets,
and an inaccessible September 2026 preprint prevents a certified global
openness or resolution claim as of 9 November 2026.

The 99 corrections distinguish transcription, translation,
ambiguity, provenance and faithful mathematical error. Earlier-year links
are retained where supported; related themes are not merged into one
problem merely because their titles resemble one another. In particular,
the earlier distribution/foliation slogan is a contextual link, while
the Buchstaber carryover keeps its existing identity. Throughout, the
notes, modern reconstruction, publication match and present mathematical
status are separately visible.

\GStalk{174}{1969-09-15}{B. R. Vainberg}{P000283}{Fundamental solutions and singularities; the Bernstein problem}
The dated September entry reports constant-coefficient PDE representations, naming H\"ormander, Borovikov and Grushin, and then the dimension threshold in the Bernstein problem. The unheaded material on P000282 is the continuation of Buchstaber\textquotesingle s earlier record 173, already indexed as 1968-173.1--3 and GP-1968-42; it is not a new Vainberg theorem. The present record ends at the separate heading \ru{Каждан и Бернштейн}.

\begin{GSitem}{1969-174.1}{F}{Superpositions on a characteristic cone}
\GSfield{Kind; source}{construction; \GSpage{P000283}}
\GSfield{The notes}{The notes call \(\int_{P(\xi)=0}f((x,\xi))A(\xi)\,d\xi\) the general solution of the homogeneous equation, with arbitrary \(f,A\), and attribute a singularity description to Borovikov.}
\GSfield{Statement}{If \(P\) is a homogeneous polynomial of degree \(m\), each smooth ridge function \(x\mapsto f(x\cdot\xi)\) with \(P(\xi)=0\) satisfies \(P(\partial_x)f(x\cdot\xi)=0\). Any superposition for which differentiation under the integral is valid is likewise a solution. Exhaustion of all local or global solutions requires a specified real or complex characteristic cycle, solution space and amplitude class.}
\GSfield{Reconstruction}{The checked elementary ridge-function assertion is separated from the unverified completeness claim. For an elliptic real symbol, its real characteristic cone may reduce to the origin, so an ordinary integral over that cone cannot represent all harmonic functions.}
\GSfield{Fate}{Borovikov\textquotesingle s 1959 paper is the named antecedent \citep{borovikov1959fundamental}. The fragment does not identify a precise completeness or singular-support theorem from that paper. The unrestricted phrase \textquotedblleft general solution\textquotedblright{} is therefore not certified.}
\GSfield{Present status}{As of 9 November 2026: the construction has a classical antecedent; the exact completeness and singularity claim intended in the notebook remains unidentified.}
\GSfield{Assessment}{literal statement: false as written; publication match: partial; verification: publisher metadata checked; development: none; confidence: medium}
\end{GSitem}
\begin{GSitem}{1969-174.2}{F}{Homogeneous powers and normalized distributions}
\GSfield{Kind; source}{construction; \GSpage{P000283}}
\GSfield{The notes}{The displayed example has \((x,\xi)^\lambda/\Gamma(\lambda+1)\) under the characteristic-surface integral and refers to \(R_0^+\).}
\GSfield{Statement}{The one-variable distributions \(t_+^\lambda/\Gamma(\lambda+1)\), initially defined for \(\operatorname{Re}\lambda>-1\), have an entire distribution-valued continuation and satisfy \(\partial_t(t_+^\lambda/\Gamma(\lambda+1))=t_+^{\lambda-1}/\Gamma(\lambda)\). Applying such distributions to a linear functional is a standard ingredient of homogeneous fundamental solutions. The notebook does not specify the power branch, sign, cycle or meaning of \(R_0^+\).}
\GSfield{Reconstruction}{The positive-part convention is a modern example, not a recovered subscript or branch from the scan.}
\GSfield{Fate}{The method belongs to the classical construction discussed by \citet{borovikov1959fundamental}; an exact formula for the unnamed operator cannot be matched from this fragment.}
\GSfield{Present status}{As of 9 November 2026: identifiable method, unidentified specific fundamental solution.}
\GSfield{Assessment}{literal statement: not determinable; publication match: thematic only; verification: publisher metadata checked; development: none; confidence: low}
\end{GSitem}
\begin{GSitem}{1969-174.3}{F}{Boundary values of reciprocal symbols}
\GSfield{Kind; source}{construction; \GSpage{P000283}}
\GSfield{The notes}{A Grushin formula is written as a sum of Fourier integrals with denominator \(P(\xi)\pm i0\) and a partition of unity on the sphere.}
\GSfield{Statement}{When the distributional limits \((P\pm i\varepsilon)^{-1}\to T_\pm\) exist, \(P T_\pm=1\); with the Fourier normalization fixed, \(\mathcal F^{-1}T_\pm\) is a fundamental solution of \(P(D)\). A localization partition must sum to one in the required symbol region, with the origin and any singular characteristic points treated separately.}
\GSfield{Fate}{Borovikov\textquotesingle s paper is a verified nearby antecedent, but targeted searches of Grushin\textquotesingle s fundamental-solution work did not identify this exact partition formula. No unrestricted existence theorem for the \(i0\) limit is inferred.}
\GSfield{Present status}{As of 9 November 2026: the formal multiplier mechanism is clear; its intended hypotheses and precise Grushin source are not.}
\GSfield{Assessment}{literal statement: correct with missing hypotheses; publication match: thematic only; verification: unresolved lead; development: none; confidence: medium}
\end{GSitem}
\begin{GSitem}{1969-174.4}{E}{Dimension threshold for entire minimal graphs}
\GSfield{Kind; source}{theorem; \GSpage{P000283}}
\GSfield{The notes}{The notes announce that an entire minimal graph \(z=f(x_1,\ldots,x_{n-1})\) is a plane for \(n\leq8\), and that the assertion fails for \(n\geq9\), naming De Giorgi, Simons and Bombieri.}
\GSfield{Statement}{For a smooth entire solution \(u:\mathbb R^m\to\mathbb R\) of the minimal-surface equation, every graph is affine when \(m\leq7\). Non-affine entire minimal graphs exist for every \(m\geq8\). Thus the ambient dimension threshold is \(m+1=9\).}
\GSfield{Reconstruction}{The variable dimension and the ambient dimension are explicitly distinguished; their one-unit difference explains the threshold in the notebook.}
\GSfield{Fate}{The counterexamples of \citet{bombieridegiorgigiusti1969minimal} settle the remaining high-dimensional direction. This is an announcement of an already obtained result, continuing the earlier seminar occurrence 1966-114.6, rather than a new unsolved question.}
\GSfield{Present status}{As of 9 November 2026: established in the stated scope; no open problem is asserted by this item.}
\GSfield{Assessment}{literal statement: correct; publication match: exact; verification: read the result; development: none; confidence: high; problem \hyperref[prob:GP-1966-13]{GP-1966-13}}
\end{GSitem}

\GStalk{175}{1969-09-15}{D. A. Kazhdan}{P000283-P000284}{Local spectral measures and Artin approximation}
The source heading credits Kazhdan and Bernstein jointly for the spectral-measure discussion. A second Kazhdan note reports Artin\textquotesingle s analytic approximation theorem through the Bourbaki seminar. The latter belongs to this record, before the Fichera heading on P000284.

\begin{GSitem}{1969-175.1}{E}{Pointwise spectral measure and projector kernel}
\GSfield{Kind; source}{definition; \GSpage{P000283-P000284}; D. A. Kazhdan and I. N. Bernstein}
\GSfield{The notes}{For a self-adjoint operator \(A\), a spectral interval \(\Lambda\), and a point \(x_0\), the notes set \(M_{x_0}(\Lambda)=\max|f(x_0)|^2\) over unit vectors in the spectral subspace and introduce its projector kernel.}
\GSfield{Statement}{Let \(A\) be self-adjoint on \(L^2(X,\mu)\) with spectral projections \(E(\Lambda)\). If evaluation at \(x\) is a well-defined bounded functional on \(E(\Lambda)L^2\), then its squared norm is \(\sup_{\|u\|=1,u\in E(\Lambda)L^2}|u(x)|^2\). When \(E(\Lambda)\) has a sufficiently regular reproducing kernel this equals \(K_\Lambda(x,x)\). Compatible point evaluations give the usual positive local spectral measure.}
\GSfield{Fate}{This is the classical spectral-function construction; see \citet[definition of the spectral function]{hormander1968spectral}. Density of continuous vectors alone does not ensure bounded evaluation, existence of a maximum, or a regular projector kernel.}
\GSfield{Present status}{As of 9 November 2026: established in the stated scope; no open problem is asserted by this item.}
\GSfield{Assessment}{literal statement: correct with missing hypotheses; publication match: partial; verification: read the result; development: none; confidence: high}
\end{GSitem}
\begin{GSitem}{1969-175.2}{F}{Spectral convergence under domain exhaustion}
\GSfield{Kind; source}{claim; \GSpage{P000283-P000284}; D. A. Kazhdan and I. N. Bernstein}
\GSfield{The notes}{The heading says convergence of the spectrum as the domains increase, accompanied by nested-domain sketches and the preceding local measure.}
\GSfield{Statement}{To formulate a convergence theorem one must specify domains \(X_j\uparrow X\), the realizations and boundary conditions of the operators, a common Hilbert-space comparison, and whether the conclusion concerns resolvents, spectral projections or local spectral measures. None of these data is recorded sufficiently to select one theorem.}
\GSfield{Fate}{Targeted searches for Kazhdan--Bernstein spectral-measure or domain-exhaustion results, including Russian author and topic variants, did not identify a precise publication. H\"ormander\textquotesingle s spectral-function work is thematic background, not proof of the unnamed exhaustion result.}
\GSfield{Present status}{As of 9 November 2026: no uniquely determined claim is available for a solved/open verdict.}
\GSfield{Assessment}{literal statement: not determinable; publication match: none found; verification: unresolved lead; development: none; confidence: low}
\end{GSitem}
\begin{GSitem}{1969-175.3}{E}{Analytic approximation of formal solutions}
\GSfield{Kind; source}{theorem; \GSpage{P000284}}
\GSfield{The notes}{The Kazhdan paragraph says that an analytic system has a formal parametric solution if and only if it has a convergent one, and that convergent solutions may be chosen arbitrarily close to the formal solution.}
\GSfield{Statement}{For \(F\in\mathbb C\{x_1,\ldots,x_n,y_1,\ldots,y_m\}^r\) and \(\widehat y\in\mathbb C[[x]]^m\) with \(\widehat y(0)=0\) and \(F(x,\widehat y)=0\), every integer \(c\geq1\) admits \(y\in\mathbb C\{x\}^m\) satisfying \(F(x,y)=0\) and \(y-\widehat y\in(x)^c\mathbb C[[x]]^m\).}
\GSfield{Reconstruction}{\textquotedblleft Arbitrarily close\textquotedblright{} means agreement to any prescribed finite jet in the maximal-ideal topology, not convergence of the original formal series.}
\GSfield{Fate}{Published by \citet{artin1968analytic} in December 1968. The Bourbaki exposition is \citet[Thm.~1]{raynaud1969artin}, delivered by Michel Raynaud in June 1969 and printed in 1971. The notebook\textquotesingle s \textquotedblleft Artin in Bourbaki\textquotedblright{} identifies the subject, not the lecturer.}
\GSfield{Present status}{As of 9 November 2026: established in the stated scope; no open problem is asserted by this item.}
\GSfield{Assessment}{literal statement: correct with missing hypotheses; publication match: exact; verification: read the result; development: none; confidence: high; problem \hyperref[prob:GP-1969-03]{GP-1969-03}}
\end{GSitem}

\GStalk{176}{1969-09-22}{Gaetano Fichera}{P000284-P000286}{The Signorini problem and unilateral constraints}
Fichera describes elastic contact with a rigid frictionless support, its variational formulation and a scalar model. The existence problem is weak and variational; classical regularity is a separate issue. The source uses a stress sign compatible with Fichera\textquotesingle s inward-normal convention. The record ends before the dated Kazhdan entry of 29 September on P000286.

\begin{GSitem}{1969-176.1}{E}{Unilateral frictionless elastic contact}
\GSfield{Kind; source}{definition; \GSpage{P000284-P000285}}
\GSfield{The notes}{On \(\Sigma\), either \(u\cdot\nu=0\) with nonnegative normal reaction and zero tangential traction, or \(u\cdot\nu>0\) with zero traction. Loads are given on \(\Sigma^*=\partial A\setminus\Sigma\).}
\GSfield{Statement}{Let a bounded elastic body \(A\subset\mathbb R^d\) have a sufficiently regular boundary decomposed into a contact portion \(\Sigma\) and a loaded portion \(\Sigma^*\). With an explicitly chosen normal and stress convention, write the normal gap as \(g\geq0\) and compressive contact reaction as \(\lambda\geq0\). Frictionless contact is \(g\lambda=0\) and zero tangential reaction; prescribed traction acts on \(\Sigma^*\).}
\GSfield{Reconstruction}{The modern gap/reaction variables avoid silently switching the notebook\textquotesingle s inward normal and stress convention.}
\GSfield{Fate}{This is the Signorini problem treated in Fichera\textquotesingle s 1964 work and summarized before this lecture in \citet[pp.~64--66]{fichera1967unilateral}.}
\GSfield{Present status}{As of 9 November 2026: established in the stated scope; no open problem is asserted by this item.}
\GSfield{Assessment}{literal statement: correct with missing hypotheses; publication match: exact; verification: read the result; development: none; confidence: high}
\end{GSitem}
\begin{GSitem}{1969-176.2}{E}{Strain, constitutive law and equilibrium}
\GSfield{Kind; source}{construction; \GSpage{P000285}}
\GSfield{The notes}{The notebook gives \(\varepsilon_{ik}=\tfrac12(\partial_k u_i+\partial_i u_k)\), \(\sigma_{ik}+\partial W/\partial\varepsilon_{ik}=0\), and the equilibrium and traction equations.}
\GSfield{Statement}{For small-strain linear elasticity, the strain is the symmetric gradient \(\varepsilon(u)\). A positive-definite quadratic energy \(W(x,\varepsilon)\) defines the conventional stress \(\tau=\partial_\varepsilon W\); with outward normal \(n\), equilibrium is \(-\operatorname{div}\tau=f\) and natural boundary traction is \(\tau n=t\). Taking \(\sigma=-\tau\) and \(\nu=-n\) recovers the sign pattern in the notebook.}
\GSfield{Fate}{The constitutive and boundary conventions are explicit in \citet[pp.~64--65]{fichera1967unilateral}. They are pre-existing elastic theory, not a new sign-changing material law.}
\GSfield{Present status}{As of 9 November 2026: established in the stated scope; no open problem is asserted by this item.}
\GSfield{Assessment}{literal statement: correct with missing hypotheses; publication match: exact; verification: read the result; development: none; confidence: high}
\end{GSitem}
\begin{GSitem}{1969-176.3}{E}{Energy minimization and complementarity}
\GSfield{Kind; source}{theorem; \GSpage{P000285}}
\GSfield{The notes}{The proposed energy is the elastic energy minus work of body and surface forces; a classical minimizer in the admissible class is said to solve the contact problem, and conversely.}
\GSfield{Statement}{Let \(V\) be a Sobolev displacement space, \(K\subset V\) its nonempty closed convex set satisfying the trace inequality, and \(J(u)=\tfrac12a(u,u)-\ell(u)\), with \(a\) continuous symmetric positive semidefinite. Then \(u\in K\) minimizes \(J\) exactly when \(a(u,v-u)\geq\ell(v-u)\) for all \(v\in K\). For sufficiently smooth \(u\), integration by parts gives the equilibrium, prescribed-traction and complementarity conditions.}
\GSfield{Reconstruction}{The displacement must be defined throughout \(\overline A\), not merely in \(C^1(\Sigma)\), and the prescribed surface work is integrated over \(\Sigma^*\).}
\GSfield{Fate}{The variational contact formulation predates the lecture; \citet[pp.~65--67]{fichera1967unilateral} gives Fichera\textquotesingle s construction. Classical differentiability is needed only to recover pointwise boundary equations.}
\GSfield{Present status}{As of 9 November 2026: established in the stated scope; no open problem is asserted by this item.}
\GSfield{Assessment}{literal statement: correct with missing hypotheses; publication match: exact; verification: read the result; development: annotated; confidence: high}
\end{GSitem}
\begin{GSitem}{1969-176.4}{E}{Weak existence and rigid-motion compatibility}
\GSfield{Kind; source}{theorem; \GSpage{P000285-P000286}}
\GSfield{The notes}{The classical admissible class may contain no minimizer; the discussion passes to weak minimization.}
\GSfield{Statement}{If \(a\) is coercive on \(V\), \(J\) has a unique minimizer in every nonempty closed convex \(K\). In pure-traction elasticity \(a\) vanishes on rigid displacements, so one needs a quotient or normalization and compatibility of the work functional with admissible rigid directions; strict negativity in unbounded admissible directions is a typical coercivity condition. Positivity of the strain energy alone is insufficient.}
\GSfield{Reconstruction}{The coercive version is stated precisely as a modern sufficient theorem; the notebook omits enough geometric and load data to state Fichera\textquotesingle s most general compatibility theorem verbatim.}
\GSfield{Fate}{\citet[pp.~66--67]{fichera1967unilateral} explicitly treats the necessary work conditions and uniqueness modulo admissible rigid motions. This is an expository report of Fichera\textquotesingle s earlier solution.}
\GSfield{Present status}{As of 9 November 2026: established in the stated scope; no open problem is asserted by this item.}
\GSfield{Assessment}{literal statement: correct with missing hypotheses; publication match: exact; verification: read the result; development: none; confidence: high}
\end{GSitem}
\begin{GSitem}{1969-176.5}{E}{Scalar boundary-obstacle problem}
\GSfield{Kind; source}{construction; \GSpage{P000285-P000286}}
\GSfield{The notes}{A scalar quadratic functional with uniformly positive coefficient matrix is introduced; on \(\Sigma\), \(u\geq0\) and either the value or the conormal derivative vanishes.}
\GSfield{Statement}{Let \(a(x)\) be a bounded measurable symmetric uniformly elliptic matrix on a bounded Lipschitz domain \(A\), let \(K=\{v\in H^1(A):\operatorname{Tr}v\geq0\text{ on }\Sigma\}\), and let \(\ell(v)=\int_A fv+\int_{\Sigma^*}\varphi v\) be continuous. Minimize \(\tfrac12\int_A\nabla v\cdot a\nabla v-\ell(v)\) on \(K\), subject to coercivity or constant-direction compatibility. In a smooth realization the outward conormal derivative on \(\Sigma\) is nonnegative and complementary to \(u\).}
\GSfield{Fate}{The scalar analogue is already given in \citet[pp.~66--68]{fichera1967unilateral}; later thin-obstacle regularity cannot be used to assert unrestricted classical solvability for the mixed boundary setup here.}
\GSfield{Present status}{As of 9 November 2026: established in the stated scope; no open problem is asserted by this item.}
\GSfield{Assessment}{literal statement: correct with missing hypotheses; publication match: exact; verification: read the result; development: none; confidence: high}
\end{GSitem}
\begin{GSitem}{1969-176.6}{E}{Necessary sign of total load}
\GSfield{Kind; source}{claim; \GSpage{P000286}}
\GSfield{The notes}{The notes write \(\int_A f+\int_{\Sigma^*}\varphi\leq0\) as a necessary condition.}
\GSfield{Statement}{In the scalar energy convention of the preceding item, every positive constant is an admissible recession direction. If \(J\) is bounded below on \(K\), then necessarily \(\ell(1)=\int_A f+\int_{\Sigma^*}\varphi\leq0\); otherwise \(J(u+c)\to-\infty\) as \(c\to+\infty\). This condition alone does not supply all existence or uniqueness hypotheses.}
\GSfield{Fate}{The load condition is part of the classical scalar model in \citet[pp.~66--67]{fichera1967unilateral}. Its elementary recession-direction interpretation also resolves the normal-sign convention.}
\GSfield{Present status}{As of 9 November 2026: established in the stated scope; no open problem is asserted by this item.}
\GSfield{Assessment}{literal statement: correct with missing hypotheses; publication match: exact; verification: read the result; development: annotated; confidence: high}
\end{GSitem}
\begin{GSitem}{1969-176.7}{F}{Failure of classical regularity at a contact interface}
\GSfield{Kind; source}{claim; \GSpage{P000286}}
\GSfield{The notes}{For signed data \(f\leq0\), \(\varphi<0\), the notebook invokes the maximum principle and Green\textquotesingle s formula and says that no \(C^1(\overline A)\) solution exists; a jump of the normal derivative is indicated.}
\GSfield{Statement}{The source does not specify a complete domain, decomposition of its boundary, behavior at \(\overline\Sigma\cap\overline{\Sigma^*}\), or an explicit load. Thus the asserted universal nonexistence cannot be reconstructed as a precise theorem. Mixed contact and load interfaces can obstruct classical regularity even when the variational solution exists.}
\GSfield{Fate}{Fichera\textquotesingle s pre-lecture note discusses the distinction between weak existence and regularity \citep[pp.~67--68]{fichera1967unilateral}. It is a partial contextual match; it does not justify quantifying the notebook\textquotesingle s incomplete example over all domains and loads.}
\GSfield{Present status}{As of 9 November 2026: weak existence and regularity theory are established, but this particular nonexistence example is not determined by the recorded data.}
\GSfield{Assessment}{literal statement: not determinable; publication match: partial; verification: read the result; development: none; confidence: medium}
\end{GSitem}

\GStalk{177}{1969-09-29}{D. A. Kazhdan}{P000286-P000287}{Artin approximation and proof reductions}
The analytic approximation theorem is restated in terms of convergent and formal power-series rings. The cusp and reduction to irreducible components provide examples and proof preparation. A Soviet-priority remark names Pokornyi, Rybin and Melamed; its scope must be checked independently of the general theorem.

\begin{GSitem}{1969-177.1}{E}{Formal points and finite-jet approximation}
\GSfield{Kind; source}{theorem; \GSpage{P000286-P000287}}
\GSfield{The notes}{The rings of convergent and formal series are denoted \(A\subset B\), and formal solutions of analytic equations are approximated by convergent solutions.}
\GSfield{Statement}{For the analytic system and formal solution in 1969-175.3, every prescribed finite jet is realized by an exact convergent solution. Equivalently, the analytic solution set is dense in the formal solution set for the maximal-ideal topology; this does not say that every formal solution converges.}
\GSfield{Fate}{The theorem is \citet[Thm.~1]{raynaud1969artin}, reporting \citet{artin1968analytic}; both the 1968 original and June 1969 Bourbaki lecture precede this seminar.}
\GSfield{Present status}{As of 9 November 2026: established in the stated scope; no open problem is asserted by this item.}
\GSfield{Assessment}{literal statement: correct with missing hypotheses; publication match: exact; verification: read the result; development: none; confidence: high; problem \hyperref[prob:GP-1969-03]{GP-1969-03}}
\end{GSitem}
\begin{GSitem}{1969-177.2}{E}{Parametrizing the cusp}
\GSfield{Kind; source}{construction; \GSpage{P000286}}
\GSfield{The notes}{The example is \(x_1^2-x_2^3=0\), with \(x_1=y^3\), \(x_2=y^2\).}
\GSfield{Statement}{The analytic map \(y\mapsto(y^3,y^2)\) satisfies the cusp equation identically and gives its normalization near the origin. The coordinate-ring map \(\mathbb C\{x_1,x_2\}/(x_1^2-x_2^3)\to\mathbb C\{y\}\) sends \(x_1\) to \(y^3\) and \(x_2\) to \(y^2\).}
\GSfield{Reconstruction}{The normalization interpretation is explanatory; the notebook records only the parametrization.}
\GSfield{Fate}{This is a classical analytic-curve example illustrating the algebra used in the approximation proof, not a new parametrization claim; compare the local-ring setting of \citet[Section II]{raynaud1969artin}.}
\GSfield{Present status}{As of 9 November 2026: established in the stated scope; no open problem is asserted by this item.}
\GSfield{Assessment}{literal statement: correct; publication match: thematic only; verification: read the result; development: none; confidence: high}
\end{GSitem}
\begin{GSitem}{1969-177.3}{E}{Reduction to irreducible components}
\GSfield{Kind; source}{method; \GSpage{P000286-P000287}}
\GSfield{The notes}{The argument proposes induction in dimension and writing the defining ideal as an intersection of prime ideals.}
\GSfield{Statement}{Let \(R\) be a Noetherian analytic local ring and \(\phi:R\to\mathbb C[[y_1,\ldots,y_s]]\) a local homomorphism. Its kernel is prime and contains a minimal prime of \(R\), so the formal point factors through an irreducible component. Passing first to \(R_{\mathrm{red}}\) removes nilpotents; a radical ideal has a finite intersection decomposition into minimal prime ideals.}
\GSfield{Reconstruction}{The domain property of the target and the passage to the reduced analytic germ are supplied explicitly.}
\GSfield{Fate}{This is a standard reduction in the proof summarized by \citet[Section II]{raynaud1969artin}. An arbitrary ideal is not itself the intersection of its minimal primes unless it is radical.}
\GSfield{Present status}{As of 9 November 2026: established in the stated scope; no open problem is asserted by this item.}
\GSfield{Assessment}{literal statement: correct with missing hypotheses; publication match: exact; verification: read the result; development: none; confidence: high}
\end{GSitem}
\begin{GSitem}{1969-177.4}{F}{Earlier Soviet work on formal implicit functions}
\GSfield{Kind; source}{remark; \GSpage{P000287}}
\GSfield{The notes}{The notes cite \ru{Покорный и Рыбкин}, Uspekhi 15(4), 1960, p.~169, and a work of Melamed in the Siberian Mathematical Journal, suggesting an earlier result.}
\GSfield{Statement}{The first citation is V.~V. Pokornyi and P.~P. Rybin, on stabilization of a scalar formal implicit-function procedure. The second is identifiable as V.~B. Melamed\textquotesingle s 1965 paper. For its analytic equation in a complex Banach space, with compact linear part and one complex parameter, Theorem 1 says that all formal Puiseux solution branches converge near zero exactly when their set is finite; Theorem 2 gives a finite nonempty set of convergent branches if the zero solution at parameter zero is isolated.}
\GSfield{Reconstruction}{The precise earlier statements are supplied from the papers. Their identification does not recover an unstated equivalence of hypotheses or conclusions with Artin approximation.}
\GSfield{Fate}{\citet[pp.~169--172]{pokornyirybin1960formal} was received in December 1958; \citet[Theorems 1--2, p.~175]{melamed1965formal} was received 13 January 1964 and published in 1965. Both are verified earlier contributions. Their stated scalar or one-parameter convergence conclusions do not identify the notebook\textquotesingle s suggested priority comparison with the arbitrary finite-jet approximation theorem of \citet{artin1968analytic}.}
\GSfield{Present status}{As of 9 November 2026: both earlier references are identified and their relevant statements checked; the broad priority comparison remains too imprecise to certify.}
\GSfield{Assessment}{literal statement: not determinable; publication match: partial; verification: read the result; development: none; confidence: medium}
\end{GSitem}

\GStalk{178}{1969-09-29}{L. D. Faddeev}{P000287-P000289}{Quantization, prequantum operators and the rotation sphere}
The lecture connects unitary representations to Hamiltonian mechanics, beginning with symplectic forms and the Van Hove operator construction. A sphere orbit and stereographic coordinates illustrate the reduction to holomorphic functions. Several normalizations and global integrality assumptions are absent; the proposed weight must not be mistaken for a proof that a measure exists only at integer parameters. The announced next report on Schmid is a plan, not an additional theorem in this record.

\begin{GSitem}{1969-178.1}{E}{Symplectic phase space and Liouville measure}
\GSfield{Kind; source}{definition; \GSpage{P000287}}
\GSfield{The notes}{The notes take a \(2n\)-dimensional phase space with nondegenerate \(\Omega=d\omega\) and invariant measure \(\Omega^n\).}
\GSfield{Statement}{A symplectic manifold is a smooth real \(2n\)-manifold with a closed nondegenerate real two-form \(\Omega\); its Liouville volume is \(\Omega^n/n!\). Exactness \(\Omega=d\omega\) is an additional hypothesis or a local choice. In particular a compact symplectic sphere has no globally defined primitive.}
\GSfield{Reconstruction}{The exact/local distinction is necessary because the lecture subsequently takes a sphere as its example.}
\GSfield{Fate}{This is classical symplectic structure underlying the orbit method \citep{kirillov1962unitary}; the line-bundle treatment in \citet[Sections 2 and 4]{kostant1970quantization} makes the global replacement for a primitive explicit.}
\GSfield{Present status}{As of 9 November 2026: established in the stated scope; no open problem is asserted by this item.}
\GSfield{Assessment}{literal statement: correct with missing hypotheses; publication match: exact; verification: read the result; development: none; confidence: high}
\end{GSitem}
\begin{GSitem}{1969-178.2}{E}{Hamiltonian vector fields and the Poisson bracket}
\GSfield{Kind; source}{theorem; \GSpage{P000287-P000288}}
\GSfield{The notes}{The correspondence \(f\mapsto X_f\) preserves brackets, and these vector fields are described as antisymmetric operators in \(L^2(M)\).}
\GSfield{Statement}{Fix \(\iota_{X_f}\Omega=-df\) and \(\{f,g\}=X_f g\). Then \([X_f,X_g]=X_{\{f,g\}}\), and each Hamiltonian flow preserves Liouville volume. Consequently \(X_f\) is formally skew-symmetric on compactly supported smooth functions. A globally defined unitary flow additionally requires a complete Hamiltonian vector field; operator self-adjointness requires a domain statement.}
\GSfield{Fate}{\citet[Proposition 4.1.1 and Section 4.3]{kostant1970quantization} records the Poisson-algebra homomorphism and explains why global integration is a separate requirement. The construction itself precedes the talk.}
\GSfield{Present status}{As of 9 November 2026: established in the stated scope; no open problem is asserted by this item.}
\GSfield{Assessment}{literal statement: correct with missing hypotheses; publication match: exact; verification: read the result; development: none; confidence: high}
\end{GSitem}
\begin{GSitem}{1969-178.3}{E}{Van Hove prequantization operator}
\GSfield{Kind; source}{construction; \GSpage{P000288}}
\GSfield{The notes}{The notebook writes \(Z_f=X_f+\alpha(\omega(X_f)-f)\), attributes the construction to Van Hove (1951), and says that the bracket relation remains valid.}
\GSfield{Statement}{On an exact symplectic chart with the preceding convention, \(\rho_\alpha(f)=X_f+\alpha(\omega(X_f)-f)\), acting on smooth complex functions, satisfies \([\rho_\alpha(f),\rho_\alpha(g)]=\rho_\alpha(\{f,g\})\). For real \(f\), formal skew-symmetry requires \(\alpha\in i\mathbb R\). Globally, the corresponding operators act on sections of a Hermitian line bundle with compatible connection, when its curvature class is integral.}
\GSfield{Reconstruction}{Signs are fixed by the displayed Hamiltonian convention; scalar multiplication and differentiation are understood on a common smooth domain.}
\GSfield{Fate}{Van Hove\textquotesingle s 1951 paper is verified bibliographically \citep{vanhove1951relations}. The full geometric operator statement is read in \citet[Theorem 4.3.1]{kostant1970quantization}. This is a prequantization construction, not an unrestricted quantization of all observables in an irreducible Hilbert-space model.}
\GSfield{Present status}{As of 9 November 2026: established in the stated scope; no open problem is asserted by this item.}
\GSfield{Assessment}{literal statement: correct with missing hypotheses; publication match: exact; verification: read the result; development: none; confidence: high}
\end{GSitem}
\begin{GSitem}{1969-178.4}{E}{Dependence on the prequantization parameter}
\GSfield{Kind; source}{claim; \GSpage{P000288}}
\GSfield{The notes}{The operators for different values of \(\alpha\) are said to give inequivalent representations, and the language also refers to the Lie algebra of vector fields.}
\GSfield{Statement}{As representations of the Poisson algebra containing constants, \(\rho_\alpha(1)=-\alpha I\), so distinct values of \(\alpha\) cannot be intertwined by a nonzero linear isomorphism. For \(\alpha\ne0\), this representation does not factor through Hamiltonian vector fields: a constant Hamiltonian has zero vector field but acts nontrivially.}
\GSfield{Fate}{The precise distinction is the central extension of Hamiltonian vector fields by constants, described in \citet[Section 4.1]{kostant1970quantization}. The elementary central-character argument explains the notebook\textquotesingle s inequivalence statement.}
\GSfield{Present status}{As of 9 November 2026: established in the stated scope; no open problem is asserted by this item.}
\GSfield{Assessment}{literal statement: correct with missing hypotheses; publication match: exact; verification: read the result; development: annotated; confidence: high}
\end{GSitem}
\begin{GSitem}{1969-178.5}{E}{Polarization and the free-particle example}
\GSfield{Kind; source}{method; \GSpage{P000288}}
\GSfield{The notes}{The notes say that quantum states should depend on half the phase variables, and contrast the classical kinetic energy \(p^2/(2m)\) with the second derivative in the position variable.}
\GSfield{Statement}{For \(T^*\mathbb R\), the Schr\"odinger realization uses \(L^2(\mathbb R,dq)\), \(\widehat p=-i\hbar\,d/dq\), and \(\widehat H=-(\hbar^2/2m)d^2/dq^2\), with its usual self-adjoint domain \(H^2(\mathbb R)\). A polarization is the geometric choice selecting an appropriate subspace of sections; an arbitrary choice of half the coordinates is not a coordinate-invariant definition.}
\GSfield{Reconstruction}{The mass is assumed positive and the real line is boundaryless; other configuration spaces require different self-adjoint boundary conditions.}
\GSfield{Fate}{\citet[Remarks 4.3.1--4.3.2]{kostant1970quantization} separates prequantization and polarized quantization. The free-particle model is standard prior quantum mechanics, not a new theorem of this seminar.}
\GSfield{Present status}{As of 9 November 2026: established in the stated scope; no open problem is asserted by this item.}
\GSfield{Assessment}{literal statement: correct with missing hypotheses; publication match: partial; verification: read the result; development: none; confidence: high}
\end{GSitem}
\begin{GSitem}{1969-178.6}{E}{Stereographic sphere and holomorphic generators}
\GSfield{Kind; source}{construction; \GSpage{P000288}}
\GSfield{The notes}{The rotation orbit is the sphere \(x_1^2+x_2^2+x_3^2=a^2\). In stereographic coordinates the notes display \(\Omega=a\,dz\wedge d\bar z/(1+|z|^2)^2\), \(\omega=(a/i)(z\,d\bar z-\bar z\,dz)/(1+|z|^2)\), and a decomposition of a rotation generator into commuting holomorphic and antiholomorphic pieces.}
\GSfield{Statement}{The sphere is \(\mathbb{CP}^1\). On a complex chart its real invariant area form is a real multiple of \(i\,dz\wedge d\bar z/(1+|z|^2)^2\); the coefficient is fixed by the total symplectic area. Holomorphic differential operators and their conjugates commute when their coefficients depend only on \(z\) and \(\bar z\), respectively. The two displayed forms in the notebook do not satisfy \(d\omega=\Omega\): direct differentiation gives \(d\omega=(2a/i)\,dz\wedge d\bar z/(1+|z|^2)^2\).}
\GSfield{Reconstruction}{The derivative is an editorial check of the recorded formula, not a recovered replacement coefficient for the original speaker.}
\GSfield{Fate}{The geometric representation mechanism is the rank-one Borel--Weil model \citep[Theorem 4]{serre1954borelweil}. The local coefficients written in the notebook require a normalization correction before they define that model.}
\GSfield{Present status}{As of 9 November 2026: established in the stated scope; no open problem is asserted by this item.}
\GSfield{Assessment}{literal statement: false as written; publication match: partial; verification: read the result; development: annotated; confidence: high}
\end{GSitem}
\begin{GSitem}{1969-178.7}{E}{Integrality of the sphere representation}
\GSfield{Kind; source}{claim; \GSpage{P000289}}
\GSfield{The notes}{The page says that a measure making the operators self-adjoint exists only for integral \(a\), writes a weight \((1+|z|^2)^{-a}\), and obtains a finite space of analytic polynomials.}
\GSfield{Statement}{The global condition is that \([\Omega/(2\pi\hbar)]\) be an integral cohomology class. For its nonnegative degree \(k\) on \(\mathbb{CP}^1\), the holomorphic sections of \(\mathcal O(k)\) have dimension \(k+1\) and realize the irreducible \(SU(2)\) representation of highest weight \(k\); it descends to \(SO(3)\) precisely when \(k\) is even. A positive weighted measure \((1+|z|^2)^{-a}d^2z\) itself exists for every real \(a\), and is finite when \(a>1\).}
\GSfield{Reconstruction}{The source parameter \(a\) is not identified with \(k\) until the missing factors of \(i,2,\hbar\) and the area-density convention are fixed. In the stated flat area measure, \(z^j\) is square-integrable exactly when \(a>j+1\); this check prevents an unsupported dimension formula.}
\GSfield{Fate}{Integral line bundles and their irreducible representations are already in \citet[Theorem 4]{serre1954borelweil}; the prequantum integrality criterion is \citet[Proposition 2.1.1]{kostant1970quantization}. Thus integrality belongs to a global equivariant quantization, not to bare existence of a measure.}
\GSfield{Present status}{As of 9 November 2026: established in the stated scope; no open problem is asserted by this item.}
\GSfield{Assessment}{literal statement: false as written; publication match: partial; verification: read the result; development: annotated; confidence: high}
\end{GSitem}

\GStalk{179}{1969-10-06}{Dusa McDuff}{P000289-P000290}{Factors and uncountably many type II1 examples}
McDuff introduces von Neumann factors through projection equivalence, their type classification, traces and the hyperfinite example, then states her continuum family of nonisomorphic finite factors. The two relevant Annals papers were published in September 1969, before this October lecture; authorship alone is not a reason to place the result in class A.

\begin{GSitem}{1969-179.1}{E}{Von Neumann algebras and factors}
\GSfield{Kind; source}{definition; \GSpage{P000289}}
\GSfield{The notes}{A weakly closed self-adjoint operator algebra \(R\) is considered; it is a factor when its center consists of scalars.}
\GSfield{Statement}{A von Neumann algebra on a complex Hilbert space \(H\) is a unital weak-operator closed \(*\)-subalgebra \(M\subset B(H)\). It is a factor if \(Z(M)=M\cap M^{\prime}=\mathbb CI\). Its abstract algebra and its particular representation on \(H\) should be distinguished.}
\GSfield{Fate}{This is the classical Murray--von Neumann framework; a checked modern statement is \citet[Chapters 1 and 4]{jones2015vonneumann}.}
\GSfield{Present status}{As of 9 November 2026: established in the stated scope; no open problem is asserted by this item.}
\GSfield{Assessment}{literal statement: correct with missing hypotheses; publication match: exact; verification: read the result; development: none; confidence: high}
\end{GSitem}
\begin{GSitem}{1969-179.2}{E}{Matrix factors and type I tensor decompositions}
\GSfield{Kind; source}{theorem; \GSpage{P000289}}
\GSfield{The notes}{For matrix algebras the notes write a tensor decomposition and describe minimal projections.}
\GSfield{Statement}{If a factor \(M\subset B(H)\) contains a nonzero minimal projection, there are Hilbert spaces \(K,L\) and a unitary \(H\simeq K\otimes L\) such that \(M=B(K)\otimes I_L\) and \(M^{\prime}=I_K\otimes B(L)\). In finite dimensions this is the ordinary matrix-factor structure. The tensor description does not apply to arbitrary type II or III factors.}
\GSfield{Reconstruction}{The minimal-projection hypothesis restricts the tensor factorization to type I.}
\GSfield{Fate}{The type I characterization is classical; see \citet[Theorem 4.2.1]{jones2015vonneumann}.}
\GSfield{Present status}{As of 9 November 2026: established in the stated scope; no open problem is asserted by this item.}
\GSfield{Assessment}{literal statement: correct with missing hypotheses; publication match: exact; verification: read the result; development: none; confidence: high}
\end{GSitem}
\begin{GSitem}{1969-179.3}{E}{Central decomposition into factors}
\GSfield{Kind; source}{theorem; \GSpage{P000289}}
\GSfield{The notes}{The general algebra is described as a direct integral of factors.}
\GSfield{Statement}{For a von Neumann algebra acting on a separable Hilbert space, central decomposition yields a measurable field of factor algebras over a standard measure space, with the original algebra their direct integral, up to the usual almost-everywhere equivalence. The center becomes the diagonal \(L^\infty\) algebra.}
\GSfield{Fate}{This is von Neumann\textquotesingle s reduction theory, published in 1949; \citet[Chapter V, Section 1]{connes1994noncommutative} gives the classical theorem and its role. It is a direct integral, not a discrete sum in general.}
\GSfield{Present status}{As of 9 November 2026: established in the stated scope; no open problem is asserted by this item.}
\GSfield{Assessment}{literal statement: correct with missing hypotheses; publication match: exact; verification: read the result; development: none; confidence: high}
\end{GSitem}
\begin{GSitem}{1969-179.4}{E}{Projections generate a von Neumann algebra}
\GSfield{Kind; source}{theorem; \GSpage{P000289-P000290}}
\GSfield{The notes}{The factor is said to be generated by its projection operators.}
\GSfield{Statement}{For every self-adjoint \(x\in M\), all its spectral projections lie in \(M\), and step-function spectral approximants converge to \(x\). Since every element is a linear combination of two self-adjoint elements, the von Neumann algebra is generated by its projections.}
\GSfield{Fate}{This is the spectral-theorem foundation of projection classification, recalled in \citet[Chapters 1 and 4]{jones2015vonneumann}.}
\GSfield{Present status}{As of 9 November 2026: established in the stated scope; no open problem is asserted by this item.}
\GSfield{Assessment}{literal statement: correct; publication match: exact; verification: read the result; development: none; confidence: high}
\end{GSitem}
\begin{GSitem}{1969-179.5}{E}{Murray--von Neumann comparison and dimension}
\GSfield{Kind; source}{definition; \GSpage{P000290}}
\GSfield{The notes}{Projection equivalence, comparison and addition of orthogonal projections are used to define a dimension scale.}
\GSfield{Statement}{For projections \(p,q\in M\), write \(p\sim q\) when some partial isometry \(v\in M\) satisfies \(v^*v=p\), \(vv^*=q\), and write \(p\precsim q\) when \(p\) is equivalent to a subprojection of \(q\). In a factor any two projections are comparable. A normal dimension function is invariant under equivalence and additive on orthogonal sums, with normalization dependent on the type.}
\GSfield{Fate}{The precise comparison statement is \citet[Definition 6.1.1 and Theorem 6.1.8]{jones2015vonneumann}; it is classical material predating the lecture.}
\GSfield{Present status}{As of 9 November 2026: established in the stated scope; no open problem is asserted by this item.}
\GSfield{Assessment}{literal statement: correct with missing hypotheses; publication match: exact; verification: read the result; development: none; confidence: high}
\end{GSitem}
\begin{GSitem}{1969-179.6}{E}{Factor types and projection-dimension ranges}
\GSfield{Kind; source}{theorem; \GSpage{P000290}}
\GSfield{The notes}{The alternatives \(I\), \(II_1\), \(II_\infty\), \(III\) are distinguished by existence of minimal projections and dimension values.}
\GSfield{Statement}{For factors on a separable Hilbert space, the standard projection-dimension scales are \(\{0,1,\ldots,n\}\) for \(I_n\), \(\{0,1,2,\ldots,\infty\}\) for \(I_\infty\), \([0,1]\) for normalized \(II_1\), \([0,\infty]\) for \(II_\infty\), and \(\{0,\infty\}\) for \(III\). The type III scale records equivalence of nonzero projections in the countably decomposable factor setting; it is not a finite trace.}
\GSfield{Fate}{The classical ranges and the distinction between finiteness of the identity and mere order type are explicitly set out in \citet[Section 6.1, pp.~35--37]{jones2015vonneumann}.}
\GSfield{Present status}{As of 9 November 2026: established in the stated scope; no open problem is asserted by this item.}
\GSfield{Assessment}{literal statement: correct with missing hypotheses; publication match: exact; verification: read the result; development: none; confidence: high}
\end{GSitem}
\begin{GSitem}{1969-179.7}{E}{Trace, determinant and the limits of matrix analogy}
\GSfield{Kind; source}{claim; \GSpage{P000290}}
\GSfield{The notes}{The finite factor is said to possess trace, determinant and a minimax principle, with the suggestion that matrix arguments continue.}
\GSfield{Statement}{A \(II_1\) factor has a unique faithful normal tracial state \(\tau\), normalized by \(\tau(1)=1\). For invertible \(x\), the Fuglede--Kadison determinant is \(\Delta(x)=\exp\tau(\log|x|)\), a positive real number. Spectral distribution and variational analogues exist, but a general self-adjoint element need not have a discrete list of eigenvalues, and the notebook supplies no particular minimax formula.}
\GSfield{Fate}{The determinant is \citet[Section 2]{fugledekadison1952determinant}; the trace is reviewed in \citet[Section 6.1]{jones2015vonneumann}. These verify the identified structures, not the unlimited phrase that all finite-matrix arguments apply.}
\GSfield{Present status}{As of 9 November 2026: established in the stated scope; no open problem is asserted by this item.}
\GSfield{Assessment}{literal statement: correct with missing hypotheses; publication match: partial; verification: read the result; development: none; confidence: high}
\end{GSitem}
\begin{GSitem}{1969-179.8}{E}{Uniqueness of the separable hyperfinite II1 factor}
\GSfield{Kind; source}{theorem; \GSpage{P000290}}
\GSfield{The notes}{The approximately finite factor is described as unique.}
\GSfield{Statement}{Among \(II_1\) factors with separable predual, the hyperfinite factor is unique up to isomorphism. Here hyperfinite means generated in the weak operator topology by an increasing sequence of finite-dimensional \(*\)-subalgebras. The claim is not uniqueness of all finite factors.}
\GSfield{Fate}{The uniqueness is the classical Murray--von Neumann theorem, recalled with proof sketch in \citet[Theorem 15.1.1]{jones2015vonneumann}.}
\GSfield{Present status}{As of 9 November 2026: established in the stated scope; no open problem is asserted by this item.}
\GSfield{Assessment}{literal statement: correct with missing hypotheses; publication match: exact; verification: read the result; development: none; confidence: high}
\end{GSitem}
\begin{GSitem}{1969-179.9}{F}{Named earlier nonisomorphic factor constructions}
\GSfield{Kind; source}{remark; \GSpage{P000290}}
\GSfield{The notes}{The notebook lists Schwartz, Ching, Sakai, Dixmier, Lance, Zeller-Meier, Stepin and Golodets in a compressed history of examples.}
\GSfield{Statement}{The list by itself contains no specified algebra, invariant or cardinality assertion for each author. It cannot be expanded into separate priority theorems without additional source evidence.}
\GSfield{Fate}{McDuff\textquotesingle s author bibliography identifies Ching, Dixmier--Lance, Sakai and Zeller-Meier as antecedents to her 1969 work \citep{mcduff1969uncountably}. That is contextual confirmation; it does not identify every construction in this notebook list.}
\GSfield{Present status}{As of 9 November 2026: the historical list is preserved; no unrecorded specific theorem is assigned to every name.}
\GSfield{Assessment}{literal statement: not determinable; publication match: partial; verification: indirect; development: none; confidence: medium}
\end{GSitem}
\begin{GSitem}{1969-179.10}{E}{A continuum of pairwise nonisomorphic II1 factors}
\GSfield{Kind; source}{theorem; \GSpage{P000290}}
\GSfield{The notes}{McDuff announces \(\mathfrak c\) nonisomorphic factors of type \(II_1\).}
\GSfield{Statement}{There are \(2^{\aleph_0}\) pairwise nonisomorphic type \(II_1\) factors acting on separable Hilbert space.}
\GSfield{Fate}{This is exactly the principal result of \citet{mcduff1969uncountably}, Annals 90(2), pp.~372--377. Its author abstract was read and the September 1969 issue date verified; the lecture is 6 October 1969. Hence the chronological class is E, although the speaker is the author. The full proof was not directly read in this task.}
\GSfield{Present status}{As of 9 November 2026: established in the stated scope; no open problem is asserted by this item.}
\GSfield{Assessment}{literal statement: correct; publication match: exact; verification: indirect; development: none; confidence: high}
\end{GSitem}

\GStalk{180}{1969-10-13}{I. M. Gelfand}{P000291-P000292}{Current groups, support, and Fock constructions}
The main questions concern irreducible unitary representations of \(SU(2)\)-valued functions on an interval and the possibility of interval support. Evaluation representations, tensor products, exponential Hilbert space and abelian characters serve as comparisons. The notebook specifies neither regularity of functions nor a group topology; this omission materially changes the questions. Araki\textquotesingle s checked paper uses bounded measurable functions and a one-parameter continuity convention, not an implicit answer in the smooth category.

\begin{GSitem}{1969-180.1}{C}{Classifying irreducible representations of an interval current group}
\GSfield{Kind; source}{question; \GSpage{P000291}}
\GSfield{The notes}{For matrices \(u(t)\in SU(2)\), \(0\leq t\leq1\), the notes ask for irreducible unitary representations. No differentiability class or topology is written.}
\GSfield{Statement}{After specifying a topological group \(G\) of \(SU(2)\)-valued maps on \([0,1]\), classify its strongly continuous irreducible unitary representations up to unitary equivalence. The choices \(C^0\), \(C^1\), \(C^\infty\), bounded measurable maps, and operator-norm versus strong continuity of the representation define different problems.}
\GSfield{Reconstruction}{The topological categories are supplied for comparison, not attributed to the 1969 lecturer.}
\GSfield{Fate}{\citet[Theorem 4.11; group theorem in the arXiv version, Theorem 4.15]{janssensneeb2015norm} classify the norm-continuous irreducible representations in the smooth compact-semisimple setting by finite evaluation tensors. Their preprint circulated 11 February 2013, the paper appeared online in December 2014 and in print in 2015. This is a proper subcategory of all strongly continuous representations. \citet[Chapter 1]{janssensneeb2024positive} impose positive-energy assumptions absent from the notebook; \citet[Definitions 2.1--2.3]{araki1970factorizable} use a different measurable current-group category.}
\GSfield{Present status}{As of 9 November 2026: strong results exist in precisely restricted categories; the notebook\textquotesingle s unrestricted classification is not a single well-posed problem, and the checked sources do not supply a complete strong-continuity classification.}
\GSfield{Assessment}{literal statement: correct with missing hypotheses; publication match: partial; verification: read the result; development: none; confidence: high; problem \hyperref[prob:GP-1969-01]{GP-1969-01}}
\GSfield{Worth developing}{Worth developing only after fixing regularity and continuity, because evaluation classification in the norm-continuous category does not control the strongly continuous category.}
\end{GSitem}
\begin{GSitem}{1969-180.2}{C}{An irreducible representation with interval support}
\GSfield{Kind; source}{question; \GSpage{P000291}}
\GSfield{The notes}{The support is the smallest closed \(F\) such that \(u=1\) on a neighborhood of \(F\) implies \(T_u=I\). The explicit question is whether an irreducible representation can have a whole interval as support.}
\GSfield{Statement}{For a chosen function group and continuity convention, ask for an irreducible unitary representation which is nontrivial on the subgroup of functions supported in every nonempty open subinterval. This is a precise full-support formulation of the displayed question. Norm-continuous smooth-current representations have finite support, whereas strongly continuous representations require separate analysis.}
\GSfield{Reconstruction}{The direct local-subgroup criterion avoids presupposing existence of a smallest closed support for an entirely unspecified class of maps.}
\GSfield{Fate}{The finite-support obstruction follows from \citet[Theorem 4.11]{janssensneeb2015norm}. The \(SU(2)\) energy irreducibility theorem of \citet{ismagilov1976unitary} assumes base dimension at least five and cannot answer an interval question; \citet[p.~74]{ismagilov1981segment} proves reducibility of the interval energy representation. The chapter develops an explicit paired-evaluation infinite tensor construction for \(C^1\) maps, also irreducible on smooth maps, and checks full interval support. It does not assert an exact archival topology or novelty.}
\GSfield{Present status}{As of 9 November 2026: the bounded \(C^1\)/smooth strong-continuity existence variant is proved in this chapter; no verdict is asserted for every possible interpretation of the unspecified 1969 group.}
\GSfield{Assessment}{literal statement: correct with missing hypotheses; publication match: partial; verification: read the result; development: developed in this chapter; confidence: high; problem \hyperref[prob:GP-1969-02]{GP-1969-02}}
\GSfield{Worth developing}{Worth developing in a fixed topology, because the contrast between finite norm-continuous support and full strongly continuous support is both explicit and mathematically consequential.}
\end{GSitem}
\begin{GSitem}{1969-180.3}{E}{Single-point evaluation representations}
\GSfield{Kind; source}{construction; \GSpage{P000291}}
\GSfield{The notes}{The simplest example evaluates \(u\) at \(t_0\) and applies an irreducible representation of \(SU(2)\).}
\GSfield{Statement}{For \(t_0\in[0,1]\) and the irreducible \(SU(2)\) representation \(V_\ell\) of highest weight \(\ell\geq0\), \(\pi(u)=V_\ell(u(t_0))\) is irreducible because the evaluation homomorphism is onto. It is continuous for each usual \(C^r\) topology. For \(\ell>0\) its support is \(\{t_0\}\); the trivial representation has empty support.}
\GSfield{Fate}{This is the elementary evaluation construction, included as the basic family in \citet{janssensneeb2015norm}. It was available before the lecture; the later theorem classifies a specified category, rather than introducing this example.}
\GSfield{Present status}{As of 9 November 2026: established in the stated scope; no open problem is asserted by this item.}
\GSfield{Assessment}{literal statement: correct with missing hypotheses; publication match: exact; verification: read the result; development: none; confidence: high}
\end{GSitem}
\begin{GSitem}{1969-180.4}{E}{Finite tensor products at distinct points}
\GSfield{Kind; source}{construction; \GSpage{P000291}}
\GSfield{The notes}{The notes tensor evaluation representations at finitely many points and describe them as irreducible.}
\GSfield{Statement}{At pairwise distinct \(t_1,\ldots,t_m\), the joint evaluation from \(C^r([0,1],SU(2))\) onto \(SU(2)^m\) is surjective for \(r=0,1,\ldots,\infty\). Thus the tensor product of irreducible factor representations is irreducible, with support the finite subset carrying nontrivial factors. Repeated evaluation points must first be combined and decomposed as \(SU(2)\) tensor products.}
\GSfield{Reconstruction}{Distinctness, removal of trivial factors and the regularity convention are made explicit.}
\GSfield{Fate}{The finite evaluation family is precisely the one in \citet{janssensneeb2015norm}. The interpolation and irreducibility argument is classical and does not depend on the later completeness theorem.}
\GSfield{Present status}{As of 9 November 2026: established in the stated scope; no open problem is asserted by this item.}
\GSfield{Assessment}{literal statement: correct with missing hypotheses; publication match: exact; verification: read the result; development: none; confidence: high}
\end{GSitem}
\begin{GSitem}{1969-180.5}{F}{The countably many evaluation points remark}
\GSfield{Kind; source}{claim; \GSpage{P000291}}
\GSfield{The notes}{The source says that the construction cannot be done for a countable set, followed by a parenthesis that it can be done for a convergent sequence.}
\GSfield{Statement}{An infinite tensor product requires choices of reference vectors and convergence conditions for the product unitaries; continuity and irreducibility are additional issues. Cardinality alone is not an obstruction, and the recorded parenthesis provides too little information to recover the intended construction or exclusion theorem.}
\GSfield{Fate}{\citet[Section 4]{araki1970factorizable} analyzes discrete and continuous tensor contributions in a specified measurable setting. The chapter\textquotesingle s paired-evaluation construction also demonstrates why a bare prohibition for countable sets is untenable without a category and tensor-sector convention.}
\GSfield{Present status}{As of 9 November 2026: the exact intended claim is unidentified; one cannot infer a general nonexistence theorem from this sentence.}
\GSfield{Assessment}{literal statement: not determinable; publication match: partial; verification: read the result; development: none; confidence: medium}
\end{GSitem}
\begin{GSitem}{1969-180.6}{E}{Exponential Hilbert space and coherent vectors}
\GSfield{Kind; source}{construction; \GSpage{P000291-P000292}}
\GSfield{The notes}{The notes define \(e^K\) as the sum of symmetric tensor powers and write exponential vectors with normalization \(\sqrt{p!}\).}
\GSfield{Statement}{For a complex Hilbert space \(K\), let \(\Gamma_s(K)=\bigoplus_{p=0}^{\infty}\operatorname{Sym}^pK\), where the zeroth summand is \(\mathbb C\). Define \(e(\phi)=\bigoplus_{p\geq0}\phi^{\otimes p}/\sqrt{p!}\). Then \(\langle e(\phi),e(\psi)\rangle=\exp\langle\phi,\psi\rangle\), and these vectors span a dense subspace.}
\GSfield{Reconstruction}{The tensor factor repeated in degree p is phi; the handwritten final factor and the distinction between the degree index and vector must not alter the coherent-vector definition.}
\GSfield{Fate}{This is classical symmetric Fock space, used explicitly in \citet[Section 4]{araki1970factorizable}; his paper was submitted in June 1969 and revised in November, and builds on earlier exponential-space work cited there.}
\GSfield{Present status}{As of 9 November 2026: established in the stated scope; no open problem is asserted by this item.}
\GSfield{Assessment}{literal statement: correct with missing hypotheses; publication match: exact; verification: read the result; development: none; confidence: high}
\end{GSitem}
\begin{GSitem}{1969-180.7}{E}{Fock factorization over orthogonal sums}
\GSfield{Kind; source}{theorem; \GSpage{P000292}}
\GSfield{The notes}{The identity \(e^{K_1\oplus K_2}=e^{K_1}\otimes e^{K_2}\) is used to motivate a continuous product.}
\GSfield{Statement}{There is a canonical unitary \(\Gamma_s(K_1\oplus K_2)\simeq\Gamma_s(K_1)\otimes\Gamma_s(K_2)\) carrying \(e(\phi_1\oplus\phi_2)\) to \(e(\phi_1)\otimes e(\phi_2)\). It follows by comparing exponential inner products and using their dense spans; direct-integral versions require a specified measurable field and measure.}
\GSfield{Fate}{The exponential-space factorization is the framework of \citet[Sections 3--4]{araki1970factorizable}. The displayed finite-sum identity is established independently of any irreducibility claim for a current group.}
\GSfield{Present status}{As of 9 November 2026: established in the stated scope; no open problem is asserted by this item.}
\GSfield{Assessment}{literal statement: correct with missing hypotheses; publication match: exact; verification: read the result; development: none; confidence: high}
\end{GSitem}
\begin{GSitem}{1969-180.8}{E}{A reducible direct-integral Fock representation}
\GSfield{Kind; source}{construction; \GSpage{P000292}}
\GSfield{The notes}{The notes pass to \(K=\int^\oplus K(t)\,dt\) and explicitly observe that the representation constructed in Fock space is reducible.}
\GSfield{Statement}{If \(\rho(u)\) is the decomposable unitary action on a nonzero direct-integral one-particle space, its second quantization \(\Gamma_s(\rho(u))=\bigoplus_{p\geq0}\operatorname{Sym}^p\rho(u)\) preserves each particle-number summand. In particular the vacuum line is invariant, so this representation is reducible whenever a nonzero positive-degree summand is present.}
\GSfield{Reconstruction}{The exact action is not written fully, so the number-preserving interpretation is stated as the version verified by the recorded direct-integral example.}
\GSfield{Fate}{Araki\textquotesingle s more general formulas include cocycle displacement terms \citep[Section 4, Lemma 4.2]{araki1970factorizable}; they must not be identified with the simpler number-preserving construction. The notebook\textquotesingle s stated reducibility is correct for the latter.}
\GSfield{Present status}{As of 9 November 2026: established in the stated scope; no open problem is asserted by this item.}
\GSfield{Assessment}{literal statement: correct with missing hypotheses; publication match: partial; verification: read the result; development: none; confidence: high}
\end{GSitem}
\begin{GSitem}{1969-180.9}{E}{Abelian characters from linear functionals}
\GSfield{Kind; source}{construction; \GSpage{P000292}}
\GSfield{The notes}{For \(\alpha(t)=\exp x(t)\), a character is written as an exponential of a linear functional or generalized function; delta mass gives point evaluation.}
\GSfield{Statement}{On the additive real space \(C^\infty([0,1],\mathbb R)\), every continuous real linear functional \(T\) gives the unitary character \(f\mapsto\exp(iT(f))\). It descends along \(f\mapsto e^{if}\) to the smooth \(U(1)\) current group on the interval exactly when \(T(1)\in\mathbb Z\), because two lifts differ by a constant multiple of \(2\pi\). Point evaluation corresponds to a delta functional. For the uniform \(C^0\) topology, continuous functionals are finite signed measures, rather than arbitrary distributions.}
\GSfield{Reconstruction}{The factor i, real-valuedness and integral-period condition are supplied to make the formula a unitary group character; the notebook uses informal complex exponential notation.}
\GSfield{Fate}{The abelian factorizable setting is discussed in \citet[Sections 1--2 and 9]{araki1970factorizable}. The interval lifting condition and the topology-dependent dual clarify the notebook\textquotesingle s schematic exponential formula; no classification of nonabelian irreducibles follows from it.}
\GSfield{Present status}{As of 9 November 2026: established in the stated scope; no open problem is asserted by this item.}
\GSfield{Assessment}{literal statement: correct with missing hypotheses; publication match: partial; verification: read the result; development: annotated; confidence: high}
\end{GSitem}

\GStalk{183}{1969-11-17}{Edwin Hewitt}{P000293}{Sharp Fourier inequalities on compact groups}
Hewitt discusses Fourier transforms with matrix-valued coefficients, Plancherel and Hausdorff--Young, and the functions giving equality. The preliminary list of names concerning singular values contains no actual inequality. This leaf precedes the 1970 heading and the February Kazhdan lecture in the manuscript, notwithstanding the catalogue record numbers.

\begin{GSitem}{1969-183.1}{F}{Unspecified singular-value inequalities}
\GSfield{Kind; source}{remark; \GSpage{P000293}}
\GSfield{The notes}{The heading \ru{О s-числах} is followed by Gantmacher, Wielandt, Berezin--Gelfand, Weyl and Ky Fan, with a reference to Trudy MMO and semisimple groups.}
\GSfield{Statement}{Singular values of a compact operator are the eigenvalues of its absolute value, counted with multiplicity. The notebook does not state which inequality, representation-theoretic construction or extension to semisimple groups is intended.}
\GSfield{Fate}{The checked Fourier-inequality literature confirms the use of Schatten norms, but does not identify an individual theorem from this list. No result is assigned solely on the basis of the named authors.}
\GSfield{Present status}{As of 9 November 2026: the precise singular-value claim remains unidentified; the names and publication hint are retained as a research lead.}
\GSfield{Assessment}{literal statement: not determinable; publication match: none found; verification: unresolved lead; development: none; confidence: low}
\end{GSitem}
\begin{GSitem}{1969-183.2}{E}{Weighted Schatten norms for the compact-group Fourier transform}
\GSfield{Kind; source}{definition; \GSpage{P000293}}
\GSfield{The notes}{The displayed \(p\)-norm is a sum over irreducibles weighted by \(d_\sigma\), but lacks a readable final exponent and a closing argument parenthesis.}
\GSfield{Statement}{For a compact group \(G\) with normalized Haar measure, let \(\widehat G\) denote the unitary irreducibles and \(\widehat f(\sigma)=\int_G f(g)\sigma(g)^*\,dg\). For \(1\leq p<\infty\), put \(\|A\|_{\ell^p(\widehat G)}=(\sum_\sigma d_\sigma\operatorname{Tr}|A_\sigma|^p)^{1/p}\); for \(p=\infty\), use \(\sup_\sigma\|A_\sigma\|_{\mathrm{op}}\). The sum is interpreted over the finite-support partial sums if the dual is uncountable.}
\GSfield{Reconstruction}{Normalized Haar measure, Schatten rather than entrywise norms, and the exponent \(1/p\) are made explicit.}
\GSfield{Fate}{These are the standard compact-group Fourier norms used in the compact-group case described by \citet[pp.~293--294]{russo1974fourier} and in \citet{hewittross1970abstract}. The missing root is a mathematical emendation, not an established reading of the handwriting.}
\GSfield{Present status}{As of 9 November 2026: established in the stated scope; no open problem is asserted by this item.}
\GSfield{Assessment}{literal statement: false as written; publication match: exact; verification: indirect; development: none; confidence: high}
\end{GSitem}
\begin{GSitem}{1969-183.3}{E}{Peter--Weyl--Plancherel identity}
\GSfield{Kind; source}{theorem; \GSpage{P000293}}
\GSfield{The notes}{The notebook states \(\|\widehat f\|_2^2=\int_G|f|^2\,d\lambda\).}
\GSfield{Statement}{For a compact group with Haar probability measure and \(f\in L^2(G)\), \(\sum_{\sigma\in\widehat G}d_\sigma\operatorname{Tr}(\widehat f(\sigma)^*\widehat f(\sigma))=\|f\|_2^2\). Equivalently, the Fourier transform extends to a unitary onto the Hilbert direct sum of the matrix spaces with these weights.}
\GSfield{Fate}{This is the classical compact-group Plancherel theorem. The exact probability-Haar formula was checked in \citet[Theorem 5.1, p.~6]{sakellaridisndcompact}; the older book treatment is \citet[Theorem 28.43]{hewittross1970abstract}. Neither later exposition is asserted to be the origin of the theorem.}
\GSfield{Present status}{As of 9 November 2026: established in the stated scope; no open problem is asserted by this item.}
\GSfield{Assessment}{literal statement: correct with missing hypotheses; publication match: exact; verification: read the result; development: none; confidence: high}
\end{GSitem}
\begin{GSitem}{1969-183.4}{E}{Compact-group Riemann--Lebesgue lemma}
\GSfield{Kind; source}{theorem; \GSpage{P000293}}
\GSfield{The notes}{For every \(f\in L^1(G)\) and \(\varepsilon>0\), only finitely many irreducibles have Fourier coefficient operator norm larger than \(\varepsilon\).}
\GSfield{Statement}{For a compact group and \(f\in L^1(G)\), the field \(\widehat f(\sigma)\) vanishes at infinity in operator norm on the discrete dual: \(\{\sigma:\|\widehat f(\sigma)\|_{\mathrm{op}}>\varepsilon\}\) is finite for every \(\varepsilon>0\). Approximation by finite linear combinations of matrix coefficients proves the assertion, using \(\|\widehat f(\sigma)\|_{\mathrm{op}}\leq\|f\|_1\).}
\GSfield{Reconstruction}{The short approximation proof is supplied to make the precise operator-norm conclusion explicit.}
\GSfield{Fate}{The finite-matrix-coefficient density statement was checked in \citet[Proposition 4.5, p.~5]{sakellaridisndcompact}. Together with the displayed \(L^1\)-to-operator-norm estimate it gives the recorded conclusion: first approximate by a finite Fourier polynomial, then bound every remaining coefficient by the approximation error. This is classical Peter--Weyl theory.}
\GSfield{Present status}{As of 9 November 2026: established in the stated scope; no open problem is asserted by this item.}
\GSfield{Assessment}{literal statement: correct; publication match: exact; verification: read the result; development: annotated; confidence: high}
\end{GSitem}
\begin{GSitem}{1969-183.5}{E}{Kunze--Hausdorff--Young inequality}
\GSfield{Kind; source}{theorem; \GSpage{P000293}}
\GSfield{The notes}{The notes attribute \(\|\widehat f\|_{p^{\prime}}\leq\|f\|_p\), \(1<p<2\), to Kunze and label it Hausdorff--Young.}
\GSfield{Statement}{Let \(G\) be compact with Haar probability measure and \(1\leq p\leq2\), \(p^{\prime}=p/(p-1)\). Then \(\|\widehat f\|_{\ell^{p^{\prime}}(\widehat G)}\leq\|f\|_{L^p(G)}\), with the \(p=1\) target interpreted in operator norm. The coefficient norm is the weighted Schatten norm, not a scalar norm on matrix entries.}
\GSfield{Fate}{Kunze published the unimodular-group extension in 1958, as explicitly documented in \citet[Introduction and reference 8]{russo1974fourier}. Its compact specialization is the theorem in the notes; this also continues the earlier seminar discussion 1967-137.11.}
\GSfield{Present status}{As of 9 November 2026: established in the stated scope; no open problem is asserted by this item.}
\GSfield{Assessment}{literal statement: correct with missing hypotheses; publication match: exact; verification: read the result; development: none; confidence: high}
\end{GSitem}
\begin{GSitem}{1969-183.6}{E}{Equality on the circle}
\GSfield{Kind; source}{theorem; \GSpage{P000293}}
\GSfield{The notes}{The Hardy--Littlewood theorem is stated as equality only for \(f(x)=\alpha e^{inx}\).}
\GSfield{Statement}{For \(1<p<2\) and nonzero \(f\in L^p(\mathbb T)\), equality in Hausdorff--Young holds exactly when \(f(x)=c e^{inx}\) almost everywhere, for some nonzero \(c\in\mathbb C\) and \(n\in\mathbb Z\). The zero function is an additional trivial equality case; at the endpoint \(p=2\), all functions give equality.}
\GSfield{Fate}{The circle statement is the classical Hardy--Littlewood equality case, discussed as the antecedent of the compact-group result in \citet[pp.~293--294]{russo1974fourier}. It follows as well from the compact-group characterization, because the circle has no proper open subgroup.}
\GSfield{Present status}{As of 9 November 2026: established in the stated scope; no open problem is asserted by this item.}
\GSfield{Assessment}{literal statement: correct with missing hypotheses; publication match: exact; verification: read the result; development: none; confidence: high}
\end{GSitem}
\begin{GSitem}{1969-183.7}{A}{Compact-group Hausdorff--Young extremizers}
\GSfield{Kind; source}{theorem; \GSpage{P000293}}
\GSfield{The notes}{The general-group examples are continuous multiplicative functions on an open subgroup \(G_0\), equal to one at the identity and extended by zero outside \(G_0\).}
\GSfield{Statement}{For a compact group \(G\), \(1<p<2\), and nonzero \(f\in L^p(G)\), equality holds exactly for functions \(f(x)=c\chi(x_0^{-1}x)\mathbf1_{x_0H}(x)\) almost everywhere, where \(H\) is an open subgroup, \(\chi:H\to U(1)\) is a continuous character, \(x_0\in G\), and \(c\ne0\). Openness makes \(H\) compact; translation and scalar multiplication preserve the equality.}
\GSfield{Reconstruction}{The complete classification, including almost-everywhere equality, translates and scalar multiples, is supplied from the published result.}
\GSfield{Fate}{The compact-group result appears in \citet[Section 43]{hewittross1970abstract}, published in 1970 after this November 1969 lecture. \citet[Theorem 1 and pp.~293--294]{russo1974fourier} states the broader unimodular theorem and expressly credits Hewitt--Ross for the compact case. The exact theorem was read in Russo; the original book was checked through publisher metadata and this explicit attribution. The notebook provides subgroup representatives and omits translates and scalar multiples from the final line.}
\GSfield{Present status}{As of 9 November 2026: the compact-group equality classification is established. The verified chronology is the 17 November 1969 seminar report followed by the 1970 book; no earlier circulation date or absolute priority is inferred from the checked evidence.}
\GSfield{Assessment}{literal statement: correct with missing hypotheses; publication match: exact; verification: indirect; development: none; confidence: high}
\end{GSitem}
\begin{GSitem}{1969-183.8}{E}{Interpolation behind Fourier norm inequalities}
\GSfield{Kind; source}{method; \GSpage{P000293}}
\GSfield{The notes}{The closing sentence says that analogues of the Riesz--Thorin theorems exist.}
\GSfield{Statement}{The Fourier map has endpoint bounds \(L^1(G)\to\ell^\infty(\widehat G)\) and \(L^2(G)\to\ell^2(\widehat G)\). Interpolation in the weighted Schatten spaces gives the compact Hausdorff--Young estimate at intermediate exponents. The sentence does not identify any additional interpolation theorem or a different operator.}
\GSfield{Reconstruction}{The endpoint spaces and resulting interpolation application are supplied as an interpretation of the brief remark.}
\GSfield{Fate}{This is the standard interpolation mechanism in the Hausdorff--Young theory discussed by \citet[Introduction]{russo1974fourier} and \citet[Section 43]{hewittross1970abstract}. Only this explicit application is reconstructed; no unspecified stronger analogue is claimed.}
\GSfield{Present status}{As of 9 November 2026: established in the stated scope; no open problem is asserted by this item.}
\GSfield{Assessment}{literal statement: correct with missing hypotheses; publication match: partial; verification: indirect; development: none; confidence: high}
\end{GSitem}

\GStalk{181}{1970}{}{P000294}{1970}
The catalogue year heading introduces the February 1970 entry. It carries no separate mathematical assertion and does not begin a new academic-year chapter. The preceding Hewitt leaf remains part of the 1969--1970 academic year.


\GStalk{182}{1970-02-23}{D. A. Kazhdan}{P000294-P000295}{Griffiths: monodromy and regularity of periods}
Kazhdan explains the cohomology local system of an algebraic family over a punctured curve, its Gauss--Manin connection, the monodromy theorem and regular-singular expansions of periods. The parenthetical assertion that monodromy itself is unipotent needs a finite-base-change qualification. Griffiths\textquotesingle s published survey records an August 1968 invited lecture and submission in September 1969, although its print issue is March 1970.

\begin{GSitem}{1969-182.1}{E}{Quasi-unipotence and the monodromy exponent}
\GSfield{Kind; source}{theorem; \GSpage{P000294}}
\GSfield{The notes}{For monodromy \(T:H^i(V_s)\to H^i(V_s)\), the notes state \((T^N-1)^{i+1}=0\) for some \(N\), then call \(T\) unipotent.}
\GSfield{Statement}{For a smooth projective family of complex algebraic varieties over a punctured smooth algebraic curve, local monodromy on \(H^i\) has roots of unity as eigenvalues, and for some positive integer \(N\), \((T^N-I)^{i+1}=0\). Thus \(T\) is quasi-unipotent; after the base change \(z=w^N\), monodromy is unipotent with nilpotency index at most \(i+1\).}
\GSfield{Reconstruction}{Smoothness and projectivity over the punctured curve are sufficient modern hypotheses. The correction concerns the separate parenthesis, not the correctly written power of monodromy.}
\GSfield{Fate}{\citet[Theorem 3.1, pp.~235--236]{griffiths1970periods} states this theorem and discusses the proofs of Landman, Clemens and Grothendieck. The survey was submitted 18 September 1969 and published March 1970, after an invited lecture on 29 August 1968; the result predates the February 1970 seminar report.}
\GSfield{Present status}{As of 9 November 2026: established in the stated scope; no open problem is asserted by this item.}
\GSfield{Assessment}{literal statement: false as written; publication match: exact; verification: read the result; development: none; confidence: high}
\end{GSitem}
\begin{GSitem}{1969-182.2}{E}{The cohomology local system and Gauss--Manin bundle}
\GSfield{Kind; source}{construction; \GSpage{P000294}}
\GSfield{The notes}{The bundle \(W\to S\) has algebraic-curve base; \(\pi_1(S)\) acts on fibre cohomology. The associated analytic vector bundle is trivial with parallel transport on the universal cover.}
\GSfield{Statement}{For a smooth proper morphism \(\pi:W\to S\) of complex algebraic varieties, \(\mathcal L=R^i\pi_*\mathbb C\) is a finite-rank local system. Its analytic bundle \(E=\mathcal L\otimes_{\mathbb C}\mathcal O_S\) has the canonical flat Gauss--Manin connection, with horizontal sections \(\mathcal L\). Pullback to a simply connected covering trivializes the local system and its flat bundle.}
\GSfield{Fate}{This is the classical construction recalled in \citet[Section 2, pp.~231--233]{griffiths1970periods}. The notebook\textquotesingle s phrase algebraic bundle is replaced by an explicit smooth proper family, not an arbitrary algebraic map.}
\GSfield{Present status}{As of 9 November 2026: established in the stated scope; no open problem is asserted by this item.}
\GSfield{Assessment}{literal statement: correct with missing hypotheses; publication match: exact; verification: read the result; development: none; confidence: high}
\end{GSitem}
\begin{GSitem}{1969-182.3}{E}{Picard--Fuchs equations in a local frame}
\GSfield{Kind; source}{method; \GSpage{P000294}}
\GSfield{The notes}{The first abbreviated line reads \(\partial/\partial z=A(z)\); the next complete line says every flat section satisfies \(f^{\prime}(z)=A(z)f(z)\), and conversely.}
\GSfield{Statement}{In a local frame of a flat holomorphic vector bundle on a complex curve, a horizontal section with coefficient column \(u(z)\) satisfies a first-order system \(u^{\prime}(z)=A(z)u(z)\), with the sign determined by the connection-matrix convention. Conversely every solution determines a horizontal section. Changes of frame transform the coefficient matrix by the usual gauge formula.}
\GSfield{Reconstruction}{The vector of coefficients and the connection sign convention are specified; the abbreviated first line is not treated as an equality of an operator with a matrix.}
\GSfield{Fate}{The local connection equation is standard and is used explicitly in \citet[Section 4]{griffiths1970periods}. The complete equation in the notebook resolves the earlier operator abbreviation.}
\GSfield{Present status}{As of 9 November 2026: established in the stated scope; no open problem is asserted by this item.}
\GSfield{Assessment}{literal statement: correct with missing hypotheses; publication match: exact; verification: read the result; development: none; confidence: high}
\end{GSitem}
\begin{GSitem}{1969-182.4}{E}{Periods pair moving cycles with cohomology}
\GSfield{Kind; source}{construction; \GSpage{P000294-P000295}}
\GSfield{The notes}{The notes transport a cycle \(\phi_z\subset V_z\) and write \((f_i(z),\phi_z)\), described as the value of cohomology on homology and a vector function.}
\GSfield{Statement}{A locally constant family of homology classes \(\gamma_z\in H_i(V_z,\mathbb C)\) is horizontal for the dual Gauss--Manin connection. For a local cohomology frame \(\omega_j(z)\), the period functions \(\langle\omega_j(z),\gamma_z\rangle\) satisfy the induced dual differential system. They are not individually constant unless the corresponding cohomology classes also are horizontal.}
\GSfield{Reconstruction}{The flat cycle lies in the homology local system; period coordinates are obtained by pairing with the chosen cohomology frame.}
\GSfield{Fate}{The homology--cohomology pairing and period system are the construction in \citet[Sections 2 and 4]{griffiths1970periods}. This identifies the recorded cycle example while keeping the bundle and its dual distinct.}
\GSfield{Present status}{As of 9 November 2026: established in the stated scope; no open problem is asserted by this item.}
\GSfield{Assessment}{literal statement: correct with missing hypotheses; publication match: exact; verification: read the result; development: none; confidence: high}
\end{GSitem}
\begin{GSitem}{1969-182.5}{E}{Regular singularity of the Gauss--Manin connection}
\GSfield{Kind; source}{theorem; \GSpage{P000295}}
\GSfield{The notes}{The theorem attributed to Griffiths says that a basis can be chosen with \(zA(z)\) holomorphic at the puncture.}
\GSfield{Statement}{For the Gauss--Manin connection of a smooth projective algebraic family over a punctured curve, there is a logarithmic extension across each puncture: in a local holomorphic frame of that extension, the connection matrix has at most a simple pole. Equivalently, starting with an algebraic meromorphic frame, the system is regular singular. An arbitrary holomorphic gauge on the punctured disk can introduce worse apparent growth.}
\GSfield{Reconstruction}{The basis is a frame of an extension, or an admissible meromorphic frame; the assertion is not invariant under unrestricted holomorphic gauges on a punctured disk.}
\GSfield{Fate}{\citet[Theorem 4.3, p.~238]{griffiths1970periods} states regularity and explains its relation to the monodromy theorem. The article was submitted in September 1969; publication in March 1970 does not make this a post-seminar discovery.}
\GSfield{Present status}{As of 9 November 2026: established in the stated scope; no open problem is asserted by this item.}
\GSfield{Assessment}{literal statement: correct with missing hypotheses; publication match: exact; verification: read the result; development: none; confidence: high}
\end{GSitem}
\begin{GSitem}{1969-182.6}{E}{Puiseux--logarithmic local expansions}
\GSfield{Kind; source}{theorem; \GSpage{P000295}}
\GSfield{The notes}{The displayed flat-section series has fractional powers \(z^{n/N}\) and finitely many powers of \(\log z\); its lower-index handwriting is uncertain and is transcribed as \(n\leq0\).}
\GSfield{Statement}{After choosing a branch on a sector and a meromorphic frame, regular singularity with quasi-unipotent monodromy gives a convergent expansion \(\sum_{n\geq n_0}z^{n/N}\sum_{j=0}^{i}c_{nj}(\log z)^j\) for horizontal degree-\(i\) sections, for some integer \(n_0\) and positive integer \(N\). The coefficients are vectors. There are only finitely many negative powers after ramification, and logarithmic degree is bounded by the monodromy nilpotency index minus one.}
\GSfield{Reconstruction}{The finite principal part, sector, coefficient vectors and direction of summation are explicit additions needed for a correct local expansion.}
\GSfield{Fate}{The factorization of a fundamental solution into a meromorphic factor and a monodromy exponential is described by \citet[pp.~238--239]{griffiths1970periods}; it gives exactly this expansion under the stated hypotheses. The index is mathematically reconstructed rather than claimed as a certain decipherment.}
\GSfield{Present status}{As of 9 November 2026: established in the stated scope; no open problem is asserted by this item.}
\GSfield{Assessment}{literal statement: correct with missing hypotheses; publication match: exact; verification: read the result; development: none; confidence: high}
\end{GSitem}
\begin{GSitem}{1969-182.7}{E}{Moderate growth of periods near a puncture}
\GSfield{Kind; source}{theorem; \GSpage{P000295}}
\GSfield{The notes}{The lemma bounds any flat section by \(C|z|^{-a}|\log z|^b\).}
\GSfield{Statement}{In a meromorphic frame of the regular-singular Gauss--Manin bundle, on each sufficiently small sector with a fixed bounded range of argument, every horizontal section satisfies \(\|u(z)\|\leq C|z|^{-a}(1+|\log z|)^b\), for suitable constants depending on the section and sector. A freely chosen holomorphic frame with an essential singularity would invalidate this statement.}
\GSfield{Fate}{The moderate-growth characterization is explained in \citet[pp.~238--239]{griffiths1970periods}, and follows directly from the preceding Puiseux--log expansion. This is a theorem in the specified geometric and gauge category, not a universal estimate for flat analytic connections.}
\GSfield{Present status}{As of 9 November 2026: established in the stated scope; no open problem is asserted by this item.}
\GSfield{Assessment}{literal statement: correct with missing hypotheses; publication match: exact; verification: read the result; development: none; confidence: high}
\end{GSitem}

\GStalk{186}{1970-09-07}{}{P000296}{Programme for the 1970--1971 seminar}
The opening programme assigns Kazhdan the relation between formal and algebraic geometry, Arnold work of Griffiths, Mityagin and Henkin their own work, and Moishezon work of Artin. It also lists finite-group \(K\)-theory with Brown, Tate, Thompson, Swan and Katz, Hironaka, and for the next semester zeta functions, \(p\)-adic harmonic analysis and finite groups. These are proposed topics and names, not stated theorems or open conjectures. The programme physically precedes the Kazhdan lecture on the same leaf.


\GStalk{184}{1970-09-07}{D. A. Kazhdan}{P000296-P000299}{Foundations of schemes and quasi-coherent sheaves}
Kazhdan begins an expository course on coherent-sheaf cohomology with spectra, localization, sheaves, fibre products and quasi-coherent modules. A Kirillov intervention explains the gluing construction of products. The source uses the older distinction between preschemes and separated schemes; the modern terminology is made explicit. The entry ends at the dated 28 September heading on P000299, rather than including that later lecture.

\begin{GSitem}{1970-184.1}{E}{Affine spectra as ringed spaces}
\GSfield{Kind; source}{definition; \GSpage{P000296}}
\GSfield{The notes}{A scheme is introduced as a topological space with a structure sheaf of rings; for an affine scheme, its points are prime ideals of a commutative ring with identity.}
\GSfield{Statement}{For a commutative unital ring \(A\), \(\operatorname{Spec}A\) consists of its proper prime ideals. Equipped with the Zariski topology and structure sheaf, it is an affine scheme; its stalk at \(\mathfrak p\) is the local ring \(A_{\mathfrak p}\). The zero ring has empty spectrum. A scheme is a locally ringed space covered by such affine models.}
\GSfield{Fate}{This is the established Grothendieck framework of \citet{grothendieck1960ega1}; its current formulation is checked in \citet[Section 26.5]{stacks2026foundations}. It continues the seminar foundations of 1965-85.1 and 1968-163.1--2.}
\GSfield{Present status}{As of 9 November 2026: established in the stated scope; no open problem is asserted by this item.}
\GSfield{Assessment}{literal statement: correct with missing hypotheses; publication match: exact; verification: read the result; development: none; confidence: high}
\end{GSitem}
\begin{GSitem}{1970-184.2}{E}{Zariski topology and principal-open basis}
\GSfield{Kind; source}{problem; \GSpage{P000296-P000297}}
\GSfield{The notes}{For an ideal \(I\), \(U_I\) comprises primes not containing \(I\). The exercises ask that these sets form a topology and that the sets \(U_f\) form a basis.}
\GSfield{Statement}{With \(D(I)=\{\mathfrak p:I\not\subseteq\mathfrak p\}\), one has \(D(I)=\bigcup_{f\in I}D(f)\), \(D(f)\cap D(g)=D(fg)\), \(D(0)=\varnothing\), and \(D(1)=\operatorname{Spec}A\). These identities prove that the \(D(I)\) are precisely the opens of a topology and that the principal opens form a basis.}
\GSfield{Fate}{These are the standard elementary spectrum identities, checked in \citet[Lemma 10.17.2]{stacks2026foundations}, within the foundational theory already published in \citet{grothendieck1960ega1}. The displayed identities resolve both exercises.}
\GSfield{Present status}{As of 9 November 2026: established in the stated scope; no open problem is asserted by this item.}
\GSfield{Assessment}{literal statement: correct; publication match: exact; verification: read the result; development: annotated; confidence: high}
\end{GSitem}
\begin{GSitem}{1970-184.3}{E}{The spectrum is a T0 space}
\GSfield{Kind; source}{problem; \GSpage{P000297}}
\GSfield{The notes}{The exercise asks for the weakest separation axiom: of two points, one can be separated from the other by a neighborhood.}
\GSfield{Statement}{For distinct prime ideals \(\mathfrak p,\mathfrak q\), choose an element in their symmetric difference. Its principal open contains exactly one of the two points. Thus \(\operatorname{Spec}A\) is \(T_0\); it is generally neither \(T_1\) nor Hausdorff.}
\GSfield{Fate}{This is an immediate consequence of the spectrum topology in \citet[Section 10.17]{stacks2026foundations}. It is a classical exercise, and no stronger separation property is inferred.}
\GSfield{Present status}{As of 9 November 2026: established in the stated scope; no open problem is asserted by this item.}
\GSfield{Assessment}{literal statement: correct; publication match: exact; verification: read the result; development: annotated; confidence: high}
\end{GSitem}
\begin{GSitem}{1970-184.4}{E}{Localization and the equality of fractions}
\GSfield{Kind; source}{construction; \GSpage{P000297}}
\GSfield{The notes}{The localization uses a multiplicative set \(S\) excluding zero and identifies \(a/s\) and \(a^{\prime}/s^{\prime}\) when some \(s^{\prime\prime}\in S\) kills \(s^{\prime}a-sa^{\prime}\). The exercise asks for equivalence of this relation.}
\GSfield{Statement}{For a multiplicative set \(S\subset A\) containing \(1\), the relation \((a,s)\sim(a^{\prime},s^{\prime})\) iff \(t(s^{\prime}a-sa^{\prime})=0\) for some \(t\in S\) is reflexive, symmetric and transitive; fraction addition and multiplication make the quotient \(S^{-1}A\). If \(0\in S\), this is the zero ring. For modules the same formula replaces \(a\) by a module element.}
\GSfield{Reconstruction}{The conventional condition that 1 belongs to S, and the harmless zero-ring case needed by the next construction, are supplied.}
\GSfield{Fate}{This is the standard localization construction used throughout \citet{grothendieck1960ega1} and \citet[Section 10.9]{stacks2026foundations}. Allowing the zero ring is necessary to include \(A_f\) when \(f\) is nilpotent and its principal open is empty.}
\GSfield{Present status}{As of 9 November 2026: established in the stated scope; no open problem is asserted by this item.}
\GSfield{Assessment}{literal statement: correct with missing hypotheses; publication match: exact; verification: read the result; development: none; confidence: high}
\end{GSitem}
\begin{GSitem}{1970-184.5}{E}{The structure sheaf from localizations}
\GSfield{Kind; source}{construction; \GSpage{P000297}}
\GSfield{The notes}{The structure sheaf is specified by \(\Gamma(U_f,\mathcal O)=A_{\{f^n\}}\).}
\GSfield{Statement}{There is a sheaf of rings \(\mathcal O_X\) on \(X=\operatorname{Spec}A\) with \(\Gamma(D(f),\mathcal O_X)=A_f\), compatible with localization restriction maps. Its stalk at \(\mathfrak p\) is \(A_{\mathfrak p}\), so it is a sheaf of local rings at points, and \(\Gamma(X,\mathcal O_X)=A\).}
\GSfield{Fate}{The result is foundational scheme theory; the explicit construction and section and stalk identifications are checked in \citet[Lemma 26.5.4]{stacks2026foundations}.}
\GSfield{Present status}{As of 9 November 2026: established in the stated scope; no open problem is asserted by this item.}
\GSfield{Assessment}{literal statement: correct with missing hypotheses; publication match: exact; verification: read the result; development: none; confidence: high}
\end{GSitem}
\begin{GSitem}{1970-184.6}{E}{The sheaf equalizer and its two restriction arrows}
\GSfield{Kind; source}{definition; \GSpage{P000297}}
\GSfield{The notes}{The diagram uses circled-plus signs, but its caption says product of sets. The question asks why there are two arrows to the sections on intersections.}
\GSfield{Statement}{A presheaf \(F\) of sets is a sheaf if, for every open cover \(U=\bigcup U_\alpha\), the map \(F(U)\to\prod_\alpha F(U_\alpha)\) identifies \(F(U)\) with the equalizer of two maps to \(\prod_{\alpha,\beta}F(U_\alpha\cap U_\beta)\). Their \((\alpha,\beta)\)-components restrict \(s_\alpha\) and \(s_\beta\), respectively; the target open is the same, but the input component differs.}
\GSfield{Reconstruction}{The product meaning follows the prose caption; the overwritten source index in the second arrow is not treated as a secure alpha.}
\GSfield{Fate}{This is the standard gluing and uniqueness axiom, checked in \citet[Section 6.7]{stacks2026foundations}. Products, not direct sums, permit arbitrary compatible local families. The written caption supports reading the symbol as an informal product notation.}
\GSfield{Present status}{As of 9 November 2026: established in the stated scope; no open problem is asserted by this item.}
\GSfield{Assessment}{literal statement: correct with missing hypotheses; publication match: exact; verification: read the result; development: annotated; confidence: high}
\end{GSitem}
\begin{GSitem}{1970-184.7}{E}{Stalks and canonical sheafification}
\GSfield{Kind; source}{problem; \GSpage{P000297}}
\GSfield{The notes}{After mentioning stalks, the exercise asks for a unique sheaf with the same stalks as any presheaf.}
\GSfield{Statement}{For a presheaf \(F\) of sets or abelian groups, there is a sheaf \(aF\) and a canonical morphism \(F\to aF\) inducing an isomorphism on every stalk. It is characterized by \(\operatorname{Hom}(aF,G)\simeq\operatorname{Hom}(F,G)\) for each sheaf \(G\), and is unique up to unique isomorphism compatible with this morphism. Merely prescribing abstract stalks without the specialization and local-section data does not determine a sheaf.}
\GSfield{Reconstruction}{Uniqueness is relative to the canonical map or universal property, not an assertion about a bare collection of stalks.}
\GSfield{Fate}{Canonical sheafification is classical; \citet[Lemmas 6.17.1--6.17.3]{stacks2026foundations} gives the construction and universal property. The seminar had already recorded it in 1968-152.4; it recurs below in 1970-187.7.}
\GSfield{Present status}{As of 9 November 2026: established in the stated scope; no open problem is asserted by this item.}
\GSfield{Assessment}{literal statement: correct with missing hypotheses; publication match: exact; verification: read the result; development: none; confidence: high}
\end{GSitem}
\begin{GSitem}{1970-184.8}{E}{Constructing a sheaf from a basis}
\GSfield{Kind; source}{method; \GSpage{P000297}}
\GSfield{The notes}{The remark says that to define a sheaf it is enough to know a presheaf on a basis.}
\GSfield{Statement}{A sheaf is uniquely determined by its values and restrictions on a topological basis. Conversely, basis data extend to a sheaf with those basis values when they satisfy the sheaf gluing and uniqueness conditions for covers of basis opens by basis opens. An arbitrary presheaf on a basis can be sheafified, but its original values need not be preserved.}
\GSfield{Fate}{The precise construction is \citet[Section 6.30]{stacks2026foundations}; it is the established method used in the structure-sheaf construction. The extra basis sheaf condition is essential to the literal converse.}
\GSfield{Present status}{As of 9 November 2026: established in the stated scope; no open problem is asserted by this item.}
\GSfield{Assessment}{literal statement: correct with missing hypotheses; publication match: exact; verification: read the result; development: none; confidence: high}
\end{GSitem}
\begin{GSitem}{1970-184.9}{E}{Principal-open inclusions, radicals and restriction maps}
\GSfield{Kind; source}{theorem; \GSpage{P000297-P000298}}
\GSfield{The notes}{The notes write \(U_f\subset U_g\) iff a power of \(f\) is divisible by \(g\), refer to the intersection of prime ideals, and then give the induced fraction map.}
\GSfield{Statement}{For \(f,g\in A\), \(D(f)\subseteq D(g)\) iff \(f\in\sqrt{(g)}\), equivalently \(f^m=gh\) for some \(m\geq1\), \(h\in A\). Then the restriction \(A_g\to A_f\) sends \(a/g^k\) to \(ah^k/f^{mk}\). The algebraic fact is \(\sqrt I=\bigcap_{\mathfrak p\supseteq I}\mathfrak p\); the intersection of all primes is the nilradical \(\sqrt{(0)}\).}
\GSfield{Reconstruction}{The radical is explicitly the radical of an ideal; the unqualified intersection of all prime ideals means the nilradical.}
\GSfield{Fate}{These are standard localization and radical identities, checked in \citet[Lemmas 10.17.2 and 26.5.1]{stacks2026foundations}. The mismatched quantifier and exponent in the source are normalized to a single integer.}
\GSfield{Present status}{As of 9 November 2026: established in the stated scope; no open problem is asserted by this item.}
\GSfield{Assessment}{literal statement: correct with missing hypotheses; publication match: exact; verification: read the result; development: none; confidence: high}
\end{GSitem}
\begin{GSitem}{1970-184.10}{E}{Local affine models and the older word prescheme}
\GSfield{Kind; source}{definition; \GSpage{P000298}}
\GSfield{The notes}{A prescheme is defined as a ringed space locally isomorphic to an affine scheme.}
\GSfield{Statement}{In current terminology a scheme is a locally ringed space admitting an open cover by affine schemes. The older word prescheme used here imposes no separatedness; the notebook reserves scheme for the later diagonal-closed condition.}
\GSfield{Fate}{This is historical terminology from the early EGA framework \citep{grothendieck1960ega1}. The modern definition is checked in \citet[Section 26.5]{stacks2026foundations}. The change is terminological, not a defect in the underlying construction.}
\GSfield{Present status}{As of 9 November 2026: established in the stated scope; no open problem is asserted by this item.}
\GSfield{Assessment}{literal statement: correct; publication match: exact; verification: read the result; development: none; confidence: high}
\end{GSitem}
\begin{GSitem}{1970-184.11}{E}{Principal opens are affine schemes}
\GSfield{Kind; source}{problem; \GSpage{P000298}}
\GSfield{The notes}{The exercise asks that \((U_f,\mathcal O)\) be affine for every \(f\in A\).}
\GSfield{Statement}{The canonical map \(\operatorname{Spec}A_f\to\operatorname{Spec}A\) is an isomorphism of locally ringed spaces onto \(D(f)\) equipped with the restricted structure sheaf. In particular this also covers the empty scheme when \(A_f=0\).}
\GSfield{Fate}{This is part of the basic affine-spectrum theorem, checked in \citet[Lemmas 10.17.6 and 26.5.4]{stacks2026foundations}. It supplies the local compatibility required in the definition of a scheme.}
\GSfield{Present status}{As of 9 November 2026: established in the stated scope; no open problem is asserted by this item.}
\GSfield{Assessment}{literal statement: correct; publication match: exact; verification: read the result; development: none; confidence: high}
\end{GSitem}
\begin{GSitem}{1970-184.12}{E}{Affine fibre products and tensor products of rings}
\GSfield{Kind; source}{construction; \GSpage{P000298}}
\GSfield{The notes}{The product is written \(\operatorname{Spec}A\times\operatorname{Spec}B=\operatorname{Spec}(A\otimes B)\), followed by an exercise on general preschemes.}
\GSfield{Statement}{For commutative \(R\)-algebras \(A,B\), the fibre product over \(\operatorname{Spec}R\) is \(\operatorname{Spec}A\times_{\operatorname{Spec}R}\operatorname{Spec}B\simeq\operatorname{Spec}(A\otimes_RB)\). It satisfies the categorical fibre-product universal property. An unqualified product of schemes is over \(\operatorname{Spec}\mathbb Z\).}
\GSfield{Fate}{The affine formula and universal property are checked in \citet[Lemma 26.17.2]{stacks2026foundations}; they are foundational scheme theory already available in \citet{grothendieck1960ega1}. The base ring cannot be omitted when working over a fixed field or base scheme.}
\GSfield{Present status}{As of 9 November 2026: established in the stated scope; no open problem is asserted by this item.}
\GSfield{Assessment}{literal statement: correct with missing hypotheses; publication match: exact; verification: read the result; development: none; confidence: high}
\end{GSitem}
\begin{GSitem}{1970-184.13}{E}{The empty product of the two residue-field spectra}
\GSfield{Kind; source}{claim; \GSpage{P000298}}
\GSfield{The notes}{The boxed example says \(\operatorname{Spec}\mathbb Z_2\times\operatorname{Spec}\mathbb Z_3\) is empty.}
\GSfield{Statement}{Interpreting \(\mathbb Z_2=\mathbb Z/2\mathbb Z\) and \(\mathbb Z_3=\mathbb Z/3\mathbb Z\), the product over \(\operatorname{Spec}\mathbb Z\) is empty: the tensor product of the rings is zero, because both 2 and 3 annihilate it and \(3-2=1\). This is not the Cartesian product of the underlying one-point sets.}
\GSfield{Reconstruction}{The interpretation of the two subscripted rings and the base are made explicit.}
\GSfield{Fate}{The statement follows directly from the affine fibre-product formula \citep[Lemma 26.17.2]{stacks2026foundations}. The source notation is read as residue fields, not as rings of 2-adic and 3-adic integers.}
\GSfield{Present status}{As of 9 November 2026: established in the stated scope; no open problem is asserted by this item.}
\GSfield{Assessment}{literal statement: correct with missing hypotheses; publication match: exact; verification: read the result; development: annotated; confidence: high}
\end{GSitem}
\begin{GSitem}{1970-184.14}{E}{Kirillov: gluing products from affine charts}
\GSfield{Kind; source}{method; \GSpage{P000298}; A. A. Kirillov}
\GSfield{The notes}{The framed Kirillov intervention covers \(X,Y\) by affine opens, forms their products, and refines their overlaps; its lemma says products of open inclusions are open inclusions.}
\GSfield{Statement}{Given schemes \(X,Y\) over \(S\), choose affine covers over affine opens of \(S\). The affine products \(U_\alpha\times_S V_\beta\) glue along their canonical overlap products. Base change preserves open immersions, so these pieces form an open cover of \(X\times_SY\), and the affine universal properties glue uniquely to the global one. Both cover indices range independently.}
\GSfield{Fate}{The precise gluing statement and open-immersion property are \citet[Lemmas 26.17.3--26.17.4 and 26.17.6]{stacks2026foundations}. This explains the improvisation without interpreting the illegible union subscript as a new mathematical restriction.}
\GSfield{Present status}{As of 9 November 2026: established in the stated scope; no open problem is asserted by this item.}
\GSfield{Assessment}{literal statement: correct with missing hypotheses; publication match: exact; verification: read the result; development: none; confidence: high}
\end{GSitem}
\begin{GSitem}{1970-184.15}{E}{Separatedness and the affine diagonal exercise}
\GSfield{Kind; source}{definition; \GSpage{P000298}}
\GSfield{The notes}{A prescheme is called a scheme if its diagonal in the product is closed; the exercise asks that an affine scheme satisfy this condition.}
\GSfield{Statement}{A morphism \(X\to S\) is separated when \(\Delta_{X/S}:X\to X\times_SX\) is a closed immersion. Every affine morphism is separated: for \(X=\operatorname{Spec}A\) over \(\operatorname{Spec}R\), the diagonal is induced by the surjective multiplication map \(A\otimes_RA\to A\). For a scheme over \(S\), its diagonal is an immersion, so closed image is an equivalent criterion.}
\GSfield{Fate}{The modern definition and affine proof are \citet[Lemmas 26.21.1--26.21.2 and Definition 26.21.3]{stacks2026foundations}. The lecture uses an older naming convention in which separated schemes alone are called schemes.}
\GSfield{Present status}{As of 9 November 2026: established in the stated scope; no open problem is asserted by this item.}
\GSfield{Assessment}{literal statement: correct with missing hypotheses; publication match: exact; verification: read the result; development: annotated; confidence: high}
\end{GSitem}
\begin{GSitem}{1970-184.16}{E}{Quasi-coherent modules by local free presentations}
\GSfield{Kind; source}{definition; \GSpage{P000298-P000299}}
\GSfield{The notes}{A sheaf of \(\mathcal O_X\)-modules is quasi-coherent if locally it is a cokernel \(\mathcal O^{(S)}\to\mathcal O^{(T)}\), with arbitrary sets of summands.}
\GSfield{Statement}{On a scheme \(X\), an \(\mathcal O_X\)-module \(\mathcal F\) is quasi-coherent iff locally it admits an exact presentation \(\mathcal O_U^{(S)}\to\mathcal O_U^{(T)}\to\mathcal F|_U\to0\), where the free sheaves are direct sums indexed by arbitrary sets. On an affine open \(U=\operatorname{Spec}A\), this is equivalent to \(\mathcal F|_U\simeq\widetilde M\) for an \(A\)-module.}
\GSfield{Fate}{The equivalence is checked in \citet[Lemma 26.7.4 and Section 17.10]{stacks2026foundations}. No finiteness of the index sets is part of quasi-coherence itself.}
\GSfield{Present status}{As of 9 November 2026: established in the stated scope; no open problem is asserted by this item.}
\GSfield{Assessment}{literal statement: correct; publication match: exact; verification: read the result; development: none; confidence: high}
\end{GSitem}
\begin{GSitem}{1970-184.17}{E}{An example of a non-quasi-coherent module sheaf}
\GSfield{Kind; source}{problem; \GSpage{P000299}}
\GSfield{The notes}{The exercise asks for a sheaf that is not quasi-coherent.}
\GSfield{Statement}{Let \(X=\operatorname{Spec}k[t]\), \(U=D(t)\), and \(j:U\hookrightarrow X\). Extension by zero \(j_!\mathcal O_U\), considered as an \(\mathcal O_X\)-module, is nonzero but has no nonzero global section: a nonzero regular function on the irreducible open \(U\) has support all of \(U\), whose closure contains the missing origin. The affine equivalence would force any quasi-coherent sheaf with zero global module to vanish. Thus \(j_!\mathcal O_U\) is a counterexample.}
\GSfield{Reconstruction}{The notebook contains only the question; the explicit affine-line example and verification are supplied by the companion.}
\GSfield{Fate}{The exercise is classical; the supplied example is a direct application of extension by zero \citep[Lemma 6.31.6]{stacks2026foundations} and the affine equivalence \citep[Lemma 26.7.5]{stacks2026foundations}. It is an explanatory solution constructed here, not a new research result.}
\GSfield{Present status}{As of 9 November 2026: established in the stated scope; no open problem is asserted by this item.}
\GSfield{Assessment}{literal statement: correct; publication match: exact; verification: read the result; development: annotated; confidence: high}
\end{GSitem}
\begin{GSitem}{1970-184.18}{E}{The affine module--sheaf equivalence}
\GSfield{Kind; source}{theorem; \GSpage{P000299}}
\GSfield{The notes}{The notes call \(\mathcal F\mapsto\Gamma(\operatorname{Spec}A,\mathcal F)\) an isomorphism from quasi-coherent sheaves to modules, and say global sections determine the sheaf.}
\GSfield{Statement}{For a commutative ring \(A\), the functors \(M\mapsto\widetilde M\) and \(\mathcal F\mapsto\Gamma(\operatorname{Spec}A,\mathcal F)\) are mutually quasi-inverse equivalences between \(A\)-modules and quasi-coherent \(\mathcal O_{\operatorname{Spec}A}\)-modules. The identifications \(M\simeq\Gamma(\widetilde M)\) and \(\widetilde{\Gamma(\mathcal F)}\simeq\mathcal F\) are natural; literal equality of the categories is unnecessary.}
\GSfield{Reconstruction}{The word isomorphism is interpreted as categorical equivalence with the canonical natural isomorphisms.}
\GSfield{Fate}{This is standard affine scheme theory, checked in \citet[Lemma 26.7.5]{stacks2026foundations}, with antecedents in \citet{serre1955coherent} and its general-ring formulation in \citet{grothendieck1960ega1}. It is developed again in record 190.}
\GSfield{Present status}{As of 9 November 2026: established in the stated scope; no open problem is asserted by this item.}
\GSfield{Assessment}{literal statement: correct with missing hypotheses; publication match: exact; verification: read the result; development: none; confidence: high}
\end{GSitem}
\begin{GSitem}{1970-184.19}{E}{Vanishing of positive-degree quasi-coherent cohomology on affines}
\GSfield{Kind; source}{theorem; \GSpage{P000299}}
\GSfield{The notes}{The separate parenthesis says that for affine spaces the cohomology is zero, without degrees or coefficients.}
\GSfield{Statement}{If \(X=\operatorname{Spec}A\) and \(\mathcal F\) is quasi-coherent, then \(H^q(X,\mathcal F)=0\) for every \(q>0\). Degree zero is the generally nonzero module of global sections. The assertion concerns affine schemes and quasi-coherent coefficients, not arbitrary sheaves on an affine topological space.}
\GSfield{Reconstruction}{Positive degree and quasi-coherent coefficients are indispensable qualifications.}
\GSfield{Fate}{The classical affine acyclicity theorem is checked in \citet[Lemma 30.2.2]{stacks2026foundations}, following \citet{serre1955coherent} and \citet{grothendieck1957homological}. Exactness of global sections on quasi-coherent modules is related to the affine equivalence, but the notebook\textquotesingle s unqualified equivalence phrase is not itself a proof of sheaf-cohomological vanishing.}
\GSfield{Present status}{As of 9 November 2026: established in the stated scope; no open problem is asserted by this item.}
\GSfield{Assessment}{literal statement: false as written; publication match: exact; verification: read the result; development: none; confidence: high}
\end{GSitem}

\GStalk{185}{1970-09-07}{I. M. Gelfand}{P000296}{Reading advice: Serre and the linear-algebra minimum}
An intervention inside Kazhdan\textquotesingle s introductory bibliography recommends learning Serre\textquotesingle s books on finite groups and on Lie groups and Lie algebras, as part of the required linear-algebra background. The nearby reference to first and second parts is preserved as a reading hint; the note states no mathematical proposition requiring a separate occurrence.


\GStalk{187}{1970-09-28}{D. A. Kazhdan}{P000299-P000301}{Sheaf images, sheafification and cohomology}
The lecture constructs direct and inverse images of sheaves, using the adjunction to characterize inverse image and a two-point example to explain sheafification. It then recalls the abelian category of sheaves, enough injectives, derived sections and the use of acyclic covers. The sign saying that the next meeting is skipped does not supply a separate date for the following cohomology reminder; its placement within this record is retained with that qualification.

\begin{GSitem}{1970-187.1}{E}{Direct image of a sheaf}
\GSfield{Kind; source}{construction; \GSpage{P000299}}
\GSfield{The notes}{For \(f:X\to Y\), the notes define \(\Gamma(V,f_*\mathcal F)=\Gamma(f^{-1}V,\mathcal F)\) and ask to check the sheaf axiom. For a one-point target they recover global sections.}
\GSfield{Statement}{If \(f:X\to Y\) is continuous and \(\mathcal F\) a sheaf of sets or abelian groups, then \((f_*\mathcal F)(V)=\mathcal F(f^{-1}V)\), with inherited restrictions, is a sheaf. Inverse images of an open cover cover its inverse image, so the gluing axiom is inherited. For \(Y\) a point its value is \(\Gamma(X,\mathcal F)\).}
\GSfield{Fate}{This is the classical direct-image functor, checked in \citet[Lemma 6.21.1]{stacks2026foundations}; the verification is the elementary exercise posed in the notebook.}
\GSfield{Present status}{As of 9 November 2026: established in the stated scope; no open problem is asserted by this item.}
\GSfield{Assessment}{literal statement: correct with missing hypotheses; publication match: exact; verification: read the result; development: annotated; confidence: high}
\end{GSitem}
\begin{GSitem}{1970-187.2}{E}{Sections on the empty open set}
\GSfield{Kind; source}{claim; \GSpage{P000299}}
\GSfield{The notes}{The parenthesis says that when \(f^{-1}(V)\) is empty, its set of sections is empty.}
\GSfield{Statement}{For a sheaf of sets, the section set on the empty open is a singleton, by the sheaf axiom for the empty cover. For a sheaf of abelian groups it is the zero group, whose underlying set has one element. Accordingly, \((f_*\mathcal F)(V)\) is terminal, not an empty set, when \(f^{-1}(V)=\varnothing\).}
\GSfield{Fate}{The correction follows immediately from the sheaf definition \citep[Section 6.7]{stacks2026foundations}. The scan repeats the word empty, so this is a faithful mathematical error rather than a translation error.}
\GSfield{Present status}{As of 9 November 2026: established in the stated scope; no open problem is asserted by this item.}
\GSfield{Assessment}{literal statement: false as written; publication match: exact; verification: read the result; development: annotated; confidence: high}
\end{GSitem}
\begin{GSitem}{1970-187.3}{E}{Inverse image characterized by adjunction}
\GSfield{Kind; source}{definition; \GSpage{P000299-P000300}}
\GSfield{The notes}{The source writes \(f^*\mathcal G\) and a natural bijection \(\operatorname{Hom}(f^*\mathcal G,\mathcal F)=\operatorname{Hom}(\mathcal G,f_*\mathcal F)\).}
\GSfield{Statement}{For continuous \(f:X\to Y\), the inverse-image functor on sheaves of sets or abelian groups, denoted \(f^{-1}\), is left adjoint to \(f_*\): \(\operatorname{Hom}_X(f^{-1}\mathcal G,\mathcal F)\simeq\operatorname{Hom}_Y(\mathcal G,f_*\mathcal F)\), naturally in both arguments. This determines inverse image up to canonical isomorphism. For modules on ringed spaces, the usual \(f^*\) additionally tensors with \(\mathcal O_X\).}
\GSfield{Reconstruction}{Modern notation f inverse is used for the functor denoted f star in this topological discussion.}
\GSfield{Fate}{The adjunction is checked in \citet[Section 6.21]{stacks2026foundations}. The distinction between sheaf inverse image and module pullback prevents transferring the bare topological formula to a different category.}
\GSfield{Present status}{As of 9 November 2026: established in the stated scope; no open problem is asserted by this item.}
\GSfield{Assessment}{literal statement: correct with missing hypotheses; publication match: exact; verification: read the result; development: none; confidence: high}
\end{GSitem}
\begin{GSitem}{1970-187.4}{E}{The inverse-image presheaf and its sheafification}
\GSfield{Kind; source}{construction; \GSpage{P000300}}
\GSfield{The notes}{The construction uses the direct limit of \(\mathcal G(V)\) over neighborhoods \(V\supset f(U)\), and asks to show that its associated sheaf is inverse image.}
\GSfield{Statement}{For open \(U\subseteq X\), define \(P(U)=\varinjlim_{V\supseteq f(U)}\mathcal G(V)\), where \(V\) runs through open subsets of \(Y\), ordered by reverse inclusion. Restrictions make \(P\) a presheaf, and \(f^{-1}\mathcal G=aP\). Its stalk at \(x\) is canonically \(\mathcal G_{f(x)}\); sheafification is essential for the adjunction on sheaves.}
\GSfield{Fate}{The construction and stalk statement are \citet[Lemmas 6.21.3--6.21.4 and the following definition]{stacks2026foundations}. The source formula is correct once the open-neighborhood indexing and associated sheaf are retained.}
\GSfield{Present status}{As of 9 November 2026: established in the stated scope; no open problem is asserted by this item.}
\GSfield{Assessment}{literal statement: correct with missing hypotheses; publication match: exact; verification: read the result; development: none; confidence: high}
\end{GSitem}
\begin{GSitem}{1970-187.5}{E}{Why the two-point pullback presheaf is not a sheaf}
\GSfield{Kind; source}{claim; \GSpage{P000300}}
\GSfield{The notes}{Two isolated points map to one point. The construction gives the same value \(\mathcal G\) on each singleton and on the whole two-point space, which the source labels not a sheaf.}
\GSfield{Statement}{Let \(X=\{x_1,x_2\}\) be discrete, \(Y=\{y\}\), and let the sheaf on \(Y\) have value a set \(A\) with at least two elements. The inverse-image presheaf has value \(A\) on every nonempty open, but its value on \(X\) cannot glue arbitrary independent choices at the two points. Its sheafification has \(A\times A\) on \(X\), the locally constant sheaf with fibre \(A\). The example is trivial for a singleton \(A\).}
\GSfield{Fate}{This is the standard distinction between a constant presheaf and constant sheaf, checked in \citet[Example 6.17.4]{stacks2026foundations}. It explains the notebook\textquotesingle s emphatic need for sheafification.}
\GSfield{Present status}{As of 9 November 2026: established in the stated scope; no open problem is asserted by this item.}
\GSfield{Assessment}{literal statement: correct with missing hypotheses; publication match: exact; verification: read the result; development: annotated; confidence: high}
\end{GSitem}
\begin{GSitem}{1970-187.6}{E}{Morphisms of sheaves as natural transformations}
\GSfield{Kind; source}{definition; \GSpage{P000300}}
\GSfield{The notes}{The footnote defines a morphism by maps on the section sets for every open, commuting with restrictions. One arrow is overwritten and resembles a two-headed arrow.}
\GSfield{Statement}{A morphism \(\phi:\mathcal F\to\mathcal G\) over \(X\) is a family \(\phi_U:\mathcal F(U)\to\mathcal G(U)\) such that restriction and \(\phi\) commute for every inclusion of opens. In the group or module category, these are group or module homomorphisms. They need not be bijections; isomorphisms are those with an inverse natural transformation.}
\GSfield{Fate}{This is the standard presheaf and sheaf morphism definition \citep[Definition 6.3.1 and Definition 6.7.1]{stacks2026foundations}. The clear prose and separate right-pointing arrow determine the intended general morphism; the overwritten arrow does not impose invertibility.}
\GSfield{Present status}{As of 9 November 2026: established in the stated scope; no open problem is asserted by this item.}
\GSfield{Assessment}{literal statement: correct with missing hypotheses; publication match: exact; verification: read the result; development: none; confidence: high}
\end{GSitem}
\begin{GSitem}{1970-187.7}{E}{The universal property of the associated sheaf}
\GSfield{Kind; source}{theorem; \GSpage{P000301}}
\GSfield{The notes}{The corrected Hom display is followed by a canonical map from the presheaf to its associated object, but the word presheaf is mistakenly repeated for both objects.}
\GSfield{Statement}{For a presheaf \(P\), sheafification is a sheaf \(aP\) with a morphism \(P\to aP\), such that every presheaf morphism \(P\to G\) into a sheaf factors uniquely through \(aP\). Equivalently, \(\operatorname{Hom}_{\mathrm{Sh}}(aP,G)\simeq\operatorname{Hom}_{\mathrm{PSh}}(P,G)\) naturally. This makes sheafification left adjoint to inclusion of sheaves into presheaves.}
\GSfield{Reconstruction}{The second object is a sheaf, as required by the heading and Hom universal property; the cancelled preliminary formula is not used.}
\GSfield{Fate}{The exact statement is \citet[Lemma 6.17.3]{stacks2026foundations}; the explicit canonical map in this occurrence supplies the qualification omitted in the earlier same-stalks exercise, 1970-184.7.}
\GSfield{Present status}{As of 9 November 2026: established in the stated scope; no open problem is asserted by this item.}
\GSfield{Assessment}{literal statement: false as written; publication match: exact; verification: read the result; development: none; confidence: high}
\end{GSitem}
\begin{GSitem}{1970-187.8}{E}{Sheaves of abelian groups form an abelian category}
\GSfield{Kind; source}{theorem; \GSpage{P000301}}
\GSfield{The notes}{Lemma 1 says that the category of sheaves of abelian groups on \(X\) is abelian.}
\GSfield{Statement}{For every topological space \(X\), the category \(\operatorname{Sh}(X,\mathbf{Ab})\) is abelian. Kernels are computed on each open, whereas cokernels are obtained by sheafifying the presheaf cokernel; a sequence is exact iff its sequences of stalks are exact. Sections on an arbitrary open need not preserve surjectivity.}
\GSfield{Fate}{This is a basic theorem of sheaf homological algebra, established in the framework of \citet{grothendieck1957homological} and checked in \citet[Lemma 17.3.1, applied to the constant sheaf of integers]{stacks2026foundations}. It is an expository lemma, not a research claim of 1970.}
\GSfield{Present status}{As of 9 November 2026: established in the stated scope; no open problem is asserted by this item.}
\GSfield{Assessment}{literal statement: correct; publication match: exact; verification: read the result; development: none; confidence: high}
\end{GSitem}
\begin{GSitem}{1970-187.9}{E}{Enough injective sheaves}
\GSfield{Kind; source}{theorem; \GSpage{P000301}}
\GSfield{The notes}{Lemma 2 says that every sheaf embeds in an injective sheaf.}
\GSfield{Statement}{For every topological space \(X\), the abelian category of sheaves of abelian groups has enough injectives. Thus every sheaf admits an injective resolution. One construction embeds each stalk into an injective abelian group and maps the sheaf into the product of the corresponding direct-image sheaves at the points.}
\GSfield{Fate}{This is the classical existence theorem used by \citet{grothendieck1957homological}; the pointwise construction is checked in \citet[Section 19.4]{stacks2026foundations}. The claim is about abelian-group sheaves, not an unqualified assertion for sheaves of sets.}
\GSfield{Present status}{As of 9 November 2026: established in the stated scope; no open problem is asserted by this item.}
\GSfield{Assessment}{literal statement: correct with missing hypotheses; publication match: exact; verification: read the result; development: none; confidence: high}
\end{GSitem}
\begin{GSitem}{1970-187.10}{E}{Cohomology as derived global sections}
\GSfield{Kind; source}{definition; \GSpage{P000301}}
\GSfield{The notes}{The notes call \(\Gamma\) left semiexact and define \(H^i\) as its derived functors, contrasting them with Cech calculations.}
\GSfield{Statement}{Global sections \(\Gamma(X,-):\operatorname{Sh}(X,\mathbf{Ab})\to\mathbf{Ab}\) are left exact. For an injective resolution \(\mathcal F\to I^\bullet\), define \(H^q(X,\mathcal F)=H^q(\Gamma(X,I^\bullet))=R^q\Gamma(X,\mathcal F)\). The resulting groups are independent of the resolution up to canonical isomorphism and have the usual long exact sequence.}
\GSfield{Fate}{This is the derived-functor definition of sheaf cohomology from \citet{grothendieck1957homological}; its precise formulation is checked in \citet[Section 20.2]{stacks2026foundations}. The notebook correctly separates the abstract definition from methods of computing it.}
\GSfield{Present status}{As of 9 November 2026: established in the stated scope; no open problem is asserted by this item.}
\GSfield{Assessment}{literal statement: correct; publication match: exact; verification: read the result; development: none; confidence: high}
\end{GSitem}
\begin{GSitem}{1970-187.11}{E}{Cech cohomology on acyclic covers}
\GSfield{Kind; source}{theorem; \GSpage{P000301}}
\GSfield{The notes}{The final lines mention a topology basis with vanishing Cech groups and say that actual cohomology agrees with Cech computations on acyclic covers.}
\GSfield{Statement}{Let \(\mathcal U=\{U_\alpha\}\) cover \(X\), and let \(\mathcal F\) be an abelian sheaf. If every nonempty finite intersection \(U_{\alpha_0}\cap\cdots\cap U_{\alpha_p}\) is \(\mathcal F\)-acyclic in positive derived-functor degrees, then \(\check H^q(\mathcal U,\mathcal F)\simeq H^q(X,\mathcal F)\). A basis criterion using Cech vanishing must in addition specify covers cofinal under refinement and stability of their intersections; one unspecified vanishing display is insufficient.}
\GSfield{Reconstruction}{The stated acyclic-cover theorem is a precise sufficient interpretation; the notebook does not record every assumption of its separate basis criterion.}
\GSfield{Fate}{The precise cover comparison and basis criterion are \citet[Section 20.11, especially Lemmas 20.11.6 and 20.11.9]{stacks2026foundations}. For quasi-coherent sheaves, affine intersections give the familiar application \citep[Lemma 30.2.6]{stacks2026foundations}, building on \citet{grothendieck1957homological}.}
\GSfield{Present status}{As of 9 November 2026: established in the stated scope; no open problem is asserted by this item.}
\GSfield{Assessment}{literal statement: correct with missing hypotheses; publication match: exact; verification: read the result; development: none; confidence: high}
\end{GSitem}

\GStalk{188}{1970-10-05}{I. M. Gelfand}{P000302}{An incomplete question about second-order elliptic systems}
The short note asks about solutions of a second-order elliptic system on a manifold, explicitly without lower-order terms. It then mentions the Laplace operator, counting solutions modulo two, topology and the Dirichlet problem. No operator, equation, boundary data or intended notion of nontrivial solution is specified, so the passage does not identify a definite modern open problem. The date and speaker come from the catalogue; the leaf itself has no date before this fragment.

\begin{GSitem}{1970-188.1}{F}{Solvability of an unspecified elliptic system}
\GSfield{Kind; source}{question; \GSpage{P000302}}
\GSfield{The notes}{The source says \ru{Эллипт. сист. 2-го порядка} on a manifold \ru{без младших членов}, then asks whether it has solutions.}
\GSfield{Statement}{A determinate solvability question requires the operator or nonlinear system, the domain and boundary conditions, the right-hand side, the solution class and the exclusion or inclusion of trivial solutions. Absence of lower-order terms does not supply these missing data. For a homogeneous linear system, zero is always a solution.}
\GSfield{Fate}{Targeted English and Russian searches combining Gelfand, elliptic systems without lower-order terms and parity of solutions yielded no exact match. \citet[Section 2, pp.~119--123]{gelfand1960elliptic} studies homotopy invariants and indices of specified elliptic boundary problems, and discusses operators without lower-order terms. It is thematic background; it does not determine the missing operator or make the fragment an index-theorem question.}
\GSfield{Present status}{As of 9 November 2026: the mathematical question is underdetermined; no solved or open status is assigned.}
\GSfield{Assessment}{literal statement: not determinable; publication match: thematic only; verification: read the result; development: none; confidence: low}
\end{GSitem}
\begin{GSitem}{1970-188.2}{F}{The unexplained Laplace comparison}
\GSfield{Kind; source}{claim; \GSpage{P000302}}
\GSfield{The notes}{The next sentence says that for the Laplace operator it is known that there are no solutions.}
\GSfield{Statement}{Without the omitted qualifier, this is not a valid general statement about harmonic functions: constants solve the homogeneous scalar Laplace equation, and zero solves homogeneous Dirichlet data. The source may intend a nontriviality or boundary requirement, but it supplies none that can be reconstructed uniquely.}
\GSfield{Fate}{The Laplace sentence cannot be identified with a particular nonexistence theorem from the available context. Its failure as a bare statement does not license attributing a specific false theorem to the speaker.}
\GSfield{Present status}{As of 9 November 2026: an essential condition is missing; the intended nonexistence assertion remains unidentified.}
\GSfield{Assessment}{literal statement: not determinable; publication match: none found; verification: unresolved lead; development: none; confidence: low}
\end{GSitem}
\begin{GSitem}{1970-188.3}{F}{Parity of solutions and an unspecified topological principle}
\GSfield{Kind; source}{question; \GSpage{P000302}}
\GSfield{The notes}{The remaining lines mention the number of solutions modulo two, that in principle the maximum is topology, and a zero-th solution of the Dirichlet problem.}
\GSfield{Statement}{To define a parity count one needs a finite or suitably compact zero-dimensional solution set, or an appropriate proper Fredholm or finite-dimensional degree construction, with regularity or multiplicity conventions. None is stated. The words topology and Dirichlet do not determine such an invariant or establish its stability.}
\GSfield{Fate}{\citet[Section 2, especially pp.~120--123]{gelfand1960elliptic} gives a related topological research programme, including the Fredholm index as kernel dimension minus adjoint-kernel dimension. This is not a count of individual solutions modulo two. No exact parity theorem was located in the checked sources; the notebook direction remains unidentified.}
\GSfield{Present status}{As of 9 November 2026: neither a parity theorem nor a well-posed parity problem can be extracted without additional archival information.}
\GSfield{Assessment}{literal statement: not determinable; publication match: thematic only; verification: read the result; development: none; confidence: low}
\end{GSitem}

\GStalk{189}{1970-10-05}{D. A. Kazhdan}{P000302}{Regular elements and the adjoint quotient}
The fragment calls itself a report on one particular work but names no author. It discusses regular and subregular elements of a complex semisimple group and the quotient by conjugation. The centralizer numbers and the use of rank require correction. Steinberg\textquotesingle s 1965 theorem matches the identifiable quotient criterion under its hypotheses; the notebook does not state a theorem on ADE surface singularities or transverse slices.

\begin{GSitem}{1970-189.1}{E}{Regularity is minimal centralizer dimension}
\GSfield{Kind; source}{definition; \GSpage{P000302}}
\GSfield{The notes}{A regular element is said to have centralizer of rank \(n-1\); the group is described as complex semisimple, of type \(A_n\) or \(D_n\).}
\GSfield{Statement}{For a connected complex semisimple algebraic group of rank \(r\), an element is regular when \(\dim Z_G(g)=r\), the minimal possible centralizer dimension. It is dimension, not rank of the centralizer, that detects regularity. In \(SL_n\), \(r=n-1\), whereas a group of type \(A_n\) has rank \(n\); the notebook leaves its integer convention unresolved.}
\GSfield{Reconstruction}{The group rank r replaces the unspecified matrix-size or Dynkin-index convention.}
\GSfield{Fate}{This is the definition in \citet[p.~49]{steinberg1965regular}, published in 1965. The mismatch of centralizer rank and dimension is a faithful source error; replacing the ambiguous n by rank r makes the assertion precise.}
\GSfield{Present status}{As of 9 November 2026: established in the stated scope; no open problem is asserted by this item.}
\GSfield{Assessment}{literal statement: false as written; publication match: exact; verification: read the result; development: none; confidence: high}
\end{GSitem}
\begin{GSitem}{1970-189.2}{E}{Subregular centralizers have dimension rank plus two}
\GSfield{Kind; source}{definition; \GSpage{P000302}}
\GSfield{The notes}{The source assigns a subregular element centralizer dimension \(n-2\) and adds two coinciding eigenvalues.}
\GSfield{Statement}{In the usual complex semisimple setting, subregular elements have centralizer dimension \(r+2\), the next dimension after regularity. For a semisimple element of \(SL_n\) with exactly one eigenvalue of multiplicity two and all other eigenvalues distinct, the centralizer dimension is \(4+(n-2)-1=n+1=(n-1)+2\). This eigenvalue description is a type-A example, not a uniform definition for type D.}
\GSfield{Reconstruction}{The standard subregular dimension and the checked type-A example are given separately; the archival n minus two is not retained as a valid centralizer formula.}
\GSfield{Fate}{The minimum-dimension characterization and structure of irregular elements are part of the classical setting of \citet[Sections 1 and 5]{steinberg1965regular}. The explicit block-centralizer calculation verifies the type-A correction directly; no later singularity theorem is inserted into the fragment.}
\GSfield{Present status}{As of 9 November 2026: established in the stated scope; no open problem is asserted by this item.}
\GSfield{Assessment}{literal statement: false as written; publication match: partial; verification: read the result; development: annotated; confidence: high}
\end{GSitem}
\begin{GSitem}{1970-189.3}{E}{Regularity and smoothness of the adjoint quotient}
\GSfield{Kind; source}{theorem; \GSpage{P000302}}
\GSfield{The notes}{The map \(G\to H/W\), with \(H\) a Cartan subgroup and \(W\) its Weyl group, is followed by regular element if and only if the map is regular there.}
\GSfield{Statement}{For a connected simply connected complex semisimple group \(G\) of rank \(r\), the adjoint quotient \(\chi:G\to H/W\simeq\mathbb A^r\) is smooth at \(g\) iff \(g\) is regular. Equivalently the differentials of its fundamental character coordinates are independent there. Here regular for the map must mean nonsingular or submersive, not merely algebraic regularity, which holds everywhere.}
\GSfield{Reconstruction}{Simply connectedness and the smooth/submersive meaning of regular map are explicit sufficient hypotheses and terminology.}
\GSfield{Fate}{\citet[Theorem 1.5, Theorem 8.1 and Remark 8.7]{steinberg1965regular} supplies the exact criterion and its simply connected qualification. This verifies the identifiable statement, but does not prove that Steinberg is the unnamed particular work reported in 1970; no stronger bibliographical identification is made.}
\GSfield{Present status}{As of 9 November 2026: established in the stated scope; no open problem is asserted by this item.}
\GSfield{Assessment}{literal statement: correct with missing hypotheses; publication match: exact; verification: read the result; development: none; confidence: high}
\end{GSitem}

\GStalk{191}{1970-10-12}{I. M. Gelfand}{P000302}{Rota's finite-field Ramsey conjecture}
The dated lower block formulates a Ramsey assertion for subspaces over a fixed finite field, notes the known line case, and translates that case into projective language. The number of colors is suppressed in the notebook; it must be fixed before the ambient dimension is chosen. This is the vector-space Ramsey conjecture, not Rota's matroid-basis conjecture.

\begin{GSitem}{1970-191.1}{B}{Monochromatic subspaces over a finite field}
\GSfield{Kind; source}{conjecture; \GSpage{P000302}}
\GSfield{The notes}{The notes fix \(F_q\), integers \(\ell<k\), and assert existence of \(N\) such that a partition of the \(\ell\)-subspaces of \(F_q^N\) contains a \(k\)-subspace whose \(\ell\)-subspaces lie in one class.}
\GSfield{Statement}{For every prime power \(q\), every integers \(r\geq1\) and \(1\leq\ell<k\), there is \(N=N(q,r,\ell,k)\) such that every map \(c:\operatorname{Gr}_{\ell}(F_q^N)\to\{1,\ldots,r\}\) is constant on \(\operatorname{Gr}_{\ell}(E)\) for some \(k\)-dimensional \(E\leq F_q^N\).}
\GSfield{Reconstruction}{The explicit finite color parameter \(r\), and its position before \(N\), are supplied. Without that order of quantifiers, coloring every subspace differently disproves the literal reading.}
\GSfield{Fate}{Established by \citet[abstract and pp.~119--120]{grahamleebrothschild1972ramsey}; the announcement appeared in January 1972. The current search did not establish an earlier public circulation date.}
\GSfield{Present status}{As of 9 November 2026: Solved; the finite-field, fixed-finite-color formulation is essential.}
\GSfield{Assessment}{literal statement: correct with missing hypotheses; publication match: exact; verification: publisher metadata checked; development: none; confidence: high; problem \hyperref[prob:GP-1970-20]{GP-1970-20}}
\end{GSitem}
\begin{GSitem}{1970-191.2}{E}{The line case and its projective interpretation}
\GSfield{Kind; source}{remark; \GSpage{P000302}}
\GSfield{The notes}{\ru{Для \(\ell=1\) доказана.} A projective line is sought in one of two classes of projective points.}
\GSfield{Statement}{For fixed \(q,r,k\), sufficiently large projective space \(\mathbf P^{N-1}(F_q)\), colored with \(r\) colors, contains a monochromatic projective \((k-1)\)-space. Vector-space dimension is one greater than projective dimension.}
\GSfield{Fate}{This is the \(\ell=1\) specialization of \citet[abstract]{grahamleebrothschild1972ramsey}. The notebook itself attests that this specialization was already known to the seminar; this run did not identify the exact earlier announcement used by Gelfand.}
\GSfield{Present status}{As of 9 November 2026: Established special case; the archival assertion about prior proof is preserved without assigning an unverified pre-1970 reference.}
\GSfield{Assessment}{literal statement: correct with missing hypotheses; publication match: partial; verification: indirect; development: none; confidence: medium; problem \hyperref[prob:GP-1970-20]{GP-1970-20}}
\end{GSitem}

\GStalk{190}{1970-10-12}{D. A. Kazhdan}{P000303-P000304}{Quasi-coherent sheaves and their global modules on affines}
Kazhdan returns to the module sheaf, the equivalence of local definitions of quasi-coherence and recovery from global sections. The proof uses finite principal-open covers and the canonical localization map. The finite exponents in the presentation diagram need arbitrary index sets in the general quasi-coherent case. The record continues only through the localization-map paragraph at the top of P000304, ending before the Thompson heading; the Rota material at the bottom of P000302 belongs to a separate catalogue record.

\begin{GSitem}{1970-190.1}{E}{The sheaf associated to a module}
\GSfield{Kind; source}{construction; \GSpage{P000303}}
\GSfield{The notes}{The notes recall module localization by fractions and set \(\Gamma(U_f,\mathcal F_M)=M_f\); listeners are asked to verify that this is a sheaf.}
\GSfield{Statement}{For a commutative ring \(A\) and any \(A\)-module \(M\), the localization rule \(D(f)\mapsto M_f\), with canonical restrictions, satisfies the basis sheaf condition and determines a unique sheaf \(\widetilde M\). Its stalk at \(\mathfrak p\) is \(M_{\mathfrak p}\), and its global sections are \(M\). Fractions satisfy \(m/s=m^{\prime}/s^{\prime}\) iff some \(t\in S\) kills \(s^{\prime}m-sm^{\prime}\).}
\GSfield{Fate}{The construction is the established affine module-sheaf construction, checked in \citet[Lemma 26.5.4 and Lemma 26.7.1]{stacks2026foundations}. It extends the localization material of 1970-184.4--5 to arbitrary modules.}
\GSfield{Present status}{As of 9 November 2026: established in the stated scope; no open problem is asserted by this item.}
\GSfield{Assessment}{literal statement: correct with missing hypotheses; publication match: exact; verification: read the result; development: none; confidence: high}
\end{GSitem}
\begin{GSitem}{1970-190.2}{E}{Equivalent definitions of quasi-coherence}
\GSfield{Kind; source}{problem; \GSpage{P000303}}
\GSfield{The notes}{The initial definition uses local cokernels of free sheaves; the second uses local isomorphisms to \(\mathcal F_M\). An exercise asks for equivalence.}
\GSfield{Statement}{For a sheaf of modules on a scheme, the following local conditions are equivalent: it is associated to a module on an affine open neighborhood; it admits a presentation there by direct sums of copies of the structure sheaf. Arbitrary modules admit free presentations, and localization is exact, so a module presentation sheafifies to a local free presentation; conversely an affine free-sheaf map comes from the corresponding module map.}
\GSfield{Fate}{This is \citet[Definition 17.10.1 and Lemmas 17.10.5, 26.7.1 and 26.7.4]{stacks2026foundations}, already implicit in \citet{grothendieck1960ega1}. The equivalence is an expository exercise, not a new problem of 1970.}
\GSfield{Present status}{As of 9 November 2026: established in the stated scope; no open problem is asserted by this item.}
\GSfield{Assessment}{literal statement: correct with missing hypotheses; publication match: exact; verification: read the result; development: none; confidence: high}
\end{GSitem}
\begin{GSitem}{1970-190.3}{E}{Free presentations need not have finite rank}
\GSfield{Kind; source}{method; \GSpage{P000303}}
\GSfield{The notes}{The proposed solution displays \(A^m\to A^n\to M\) and \(\mathcal O^m\to\mathcal O^n\to\mathcal F_M\), then reverses the construction by taking sections.}
\GSfield{Statement}{For arbitrary \(M\), use an exact sequence \(A^{(I)}\to A^{(J)}\to M\to0\) with possibly infinite sets \(I,J\), and the corresponding sheaf sequence. Finite \(m,n\) are justified when \(M\) is finitely presented, a stronger assumption than quasi-coherence. In the converse, form the cokernel of the module map and identify its sheaf locally; do not assume right exactness of arbitrary section functors as a premise.}
\GSfield{Reconstruction}{The letters m and n are replaced by index sets for the unrestricted theorem. The proof order avoids using the desired affine exactness circularly.}
\GSfield{Fate}{Exactness of the associated-sheaf functor and the free-presentation characterization are checked in \citet[Lemma 26.5.4 and Section 26.7]{stacks2026foundations}. The arbitrary-index convention had been stated explicitly in the 7 September lecture but is suppressed in this diagram.}
\GSfield{Present status}{As of 9 November 2026: established in the stated scope; no open problem is asserted by this item.}
\GSfield{Assessment}{literal statement: correct with missing hypotheses; publication match: exact; verification: read the result; development: annotated; confidence: high}
\end{GSitem}
\begin{GSitem}{1970-190.4}{E}{Every quasi-coherent sheaf on an affine is a module sheaf}
\GSfield{Kind; source}{theorem; \GSpage{P000303}}
\GSfield{The notes}{The main theorem says every quasi-coherent \(\mathcal F\) on \(\operatorname{Spec}A\) has the form \(\mathcal F_M\); its proof sets \(M=\Gamma(X,\mathcal F)\).}
\GSfield{Statement}{For \(X=\operatorname{Spec}A\) and quasi-coherent \(\mathcal F\), the canonical morphism \(\widetilde{\Gamma(X,\mathcal F)}\to\mathcal F\) is an isomorphism. Thus the local module presentations globalize uniquely up to the canonical module identification; no finite-generation, Noetherian or separatedness hypothesis beyond affineness is required.}
\GSfield{Fate}{This is the affine equivalence of \citet[Lemma 26.7.5]{stacks2026foundations}, first introduced in record 184 and classical by \citet{grothendieck1960ega1}. The choice of the global-section module is part of the theorem, rather than an arbitrary existence assertion.}
\GSfield{Present status}{As of 9 November 2026: established in the stated scope; no open problem is asserted by this item.}
\GSfield{Assessment}{literal statement: correct; publication match: exact; verification: read the result; development: none; confidence: high}
\end{GSitem}
\begin{GSitem}{1970-190.5}{E}{Quasi-compactness of affine spectra}
\GSfield{Kind; source}{theorem; \GSpage{P000303}}
\GSfield{The notes}{The proof recalls that every affine scheme is quasi-compact, with a parenthesis about absence of separation, and chooses a finite subcover.}
\GSfield{Statement}{Every open cover of \(\operatorname{Spec}A\) has a finite subcover. Refine to principal opens \(D(f_i)\); covering means the \(f_i\) generate the unit ideal, so a finite sum \(\sum a_if_i=1\) selects a finite subcover. This topological quasi-compactness does not assert that the underlying space is Hausdorff.}
\GSfield{Fate}{This is checked in \citet[Lemma 10.17.8]{stacks2026foundations}. The parenthesis distinguishes quasi-compactness from a Hausdorff compactness convention; it does not deny the scheme-theoretic separatedness of affine schemes.}
\GSfield{Present status}{As of 9 November 2026: established in the stated scope; no open problem is asserted by this item.}
\GSfield{Assessment}{literal statement: correct; publication match: exact; verification: read the result; development: annotated; confidence: high}
\end{GSitem}
\begin{GSitem}{1970-190.6}{E}{The canonical localization map and clearing denominators}
\GSfield{Kind; source}{method; \GSpage{P000303-P000304}}
\GSfield{The notes}{After choosing a finite cover with \(\mathcal F|_{D(f_i)}=\widetilde M_i\), the notes construct \(M_f\to\Gamma(D(f),\mathcal F)\) because \(1/f\) is a section of the structure sheaf there.}
\GSfield{Statement}{For any module sheaf \(\mathcal F\), restriction \(M=\Gamma(X,\mathcal F)\to\Gamma(D(f),\mathcal F)\) factors canonically through \(M_f\), since multiplication by \(f\) is invertible on \(D(f)\). For quasi-coherent \(\mathcal F\) on affine \(X\), a finite principal-open cover with local module descriptions permits a common power of \(f\) to clear all denominators and compatibility defects; this proves the map is an isomorphism and hence the main theorem.}
\GSfield{Reconstruction}{The missing remainder of the proof is summarized using the standard denominator-clearing argument; the source itself stops after defining the natural map.}
\GSfield{Fate}{The localization identities and proof of global recovery are checked in \citet[Lemmas 26.7.1--26.7.5]{stacks2026foundations}. The construction of the map is valid before quasi-coherence; its being an isomorphism uses quasi-coherence and the finite cover.}
\GSfield{Present status}{As of 9 November 2026: established in the stated scope; no open problem is asserted by this item.}
\GSfield{Assessment}{literal statement: correct with missing hypotheses; publication match: exact; verification: read the result; development: none; confidence: high}
\end{GSitem}

\GStalk{192}{1970-10-12}{I. M. Gelfand}{P000304-P000305}{Thompson quadratic pairs and Conway's lattice group}
Gelfand reports Thompson's quadratic-pair classification. Kirillov then relays Thompson's account of Conway's sporadic group through a 24-dimensional lattice. The passage needs two distinctions: groups of Lie type are broader than classical groups, and the full automorphism group of the Leech lattice is a central double cover of the simple group. The reading list below this discussion is not a mathematical theorem.

\begin{GSitem}{1970-192.1}{E}{Quadratic pairs}
\GSfield{Kind; source}{definition; \GSpage{P000304}}
\GSfield{The notes}{\((G,M)\): \(M\) is a faithful irreducible \(F_p\)-representation and \(G\) is generated by elements satisfying \(M(g-1)^2=0\).}
\GSfield{Statement}{A quadratic pair consists of a nontrivial finite group \(G\), an odd prime \(p\), and a faithful irreducible finite-dimensional \(F_pG\)-module \(M\) such that \(G=\langle Q\rangle\), where \(Q=\{g\ne1:(g-I)^2M=0\}\). Such generators have Jordan blocks of sizes at most two, with at least one size-two block.}
\GSfield{Reconstruction}{Finiteness of the group, finite-dimensionality, exclusion of the identity as a generator, and the odd-prime convention are supplied from the published definition.}
\GSfield{Fate}{Thompson's Nice ICM lecture circulated in 1970 and appeared in 1971; see \citet[p.~375, Definition]{thompson1971quadratic}. The later treatment is \citet[Hypothesis~1.0]{chermak2004quadratic}.}
\GSfield{Present status}{As of 9 November 2026: Established with the hypotheses in the statement.}
\GSfield{Assessment}{literal statement: correct with missing hypotheses; publication match: exact; verification: indirect; development: none; confidence: high}
\end{GSitem}
\begin{GSitem}{1970-192.2}{E}{The central-product and Lie-type classification for odd primes at least five}
\GSfield{Kind; source}{theorem; \GSpage{P000304}}
\GSfield{The notes}{For \(p>5\), the notes say that \(G\) is only a classical group, listing \(S,O\), and adding \ru{кроме \(E_8\)}.}
\GSfield{Statement}{For a quadratic pair with \(p\geq5\), the structural conclusion involves a central product of perfect groups with simple central quotients and a corresponding tensor decomposition of the module. In the quasisimple case the simple quotient is of Lie type in characteristic \(p\), from the families listed in Thompson's theorem, including exceptional families \(G_2,F_4,E_6,E_7\), but excluding \(E_8\). It is not a conclusion that \(G\) itself is a single classical simple group.}
\GSfield{Reconstruction}{The central-product reduction and exceptional Lie types are restored from the checked Thompson theorem; the literal classical-group slogan is not retained as a theorem.}
\GSfield{Fate}{\citet[p.~375, Central product theorem and Main theorem]{thompson1971quadratic} gives precisely the missing structural qualifications. It was reported at Nice before the October seminar, although the proceedings were printed in 1971.}
\GSfield{Present status}{As of 9 November 2026: Established structural classification; the notebook's word classical and its omission of central products are not literal equivalents of the theorem.}
\GSfield{Assessment}{literal statement: false as written; publication match: partial; verification: indirect; development: none; confidence: high}
\end{GSitem}
\begin{GSitem}{1970-192.3}{B}{The small-characteristic continuation of quadratic-pair theory}
\GSfield{Kind; source}{question; \GSpage{P000304}}
\GSfield{The notes}{\ru{Остались хар-ки 2 и 3.}}
\GSfield{Statement}{The definite odd-characteristic continuation is to classify faithful irreducible quadratic pairs over \(F_3\). The characteristic-two phrase does not specify what additional restrictions distinguish the intended problem: in characteristic two every involution satisfies \((g-I)^2=0\).}
\GSfield{Reconstruction}{The characteristic-three problem is separated from the underspecified characteristic-two slogan; no claim of solving every characteristic-two classification problem is made.}
\GSfield{Fate}{For characteristic three, \citet[Theorems~A--B; introduction]{chermak2004quadratic} treats the component case using the finite-simple-group classification and points to the earlier component-free treatment. The opened preprint is dated February 2002; publication was 2004. Its introduction explains that the original program was displaced by a different route to the finite-simple-group classification.}
\GSfield{Present status}{As of 9 November 2026: The odd-characteristic quadratic-pair program has a classification-based treatment. The unqualified characteristic-two slogan is not assigned an exact solved theorem.}
\GSfield{Assessment}{literal statement: correct with missing hypotheses; publication match: partial; verification: read the result; development: none; confidence: medium; problem \hyperref[prob:GP-1970-21]{GP-1970-21}}
\end{GSitem}
\begin{GSitem}{1970-192.4}{E}{The rootless even unimodular lattice in dimension twenty-four}
\GSfield{Kind; source}{construction; \GSpage{P000304-P000305}; A. A. Kirillov}
\GSfield{The notes}{Kirillov describes \(L\subset\mathbf R^{24}\) with covolume one, integral pairings, even squared norms, and minimum squared norm four, and says it is unique.}
\GSfield{Statement}{Up to Euclidean isometry there is a unique positive-definite even unimodular lattice of rank 24 with no norm-two vectors: the Leech lattice. Its minimum over nonzero vectors is four. The zero vector must be excluded from the minimum.}
\GSfield{Reconstruction}{Evenness, unimodularity, positive definiteness, rank 24, and absence of roots delimit the standard uniqueness statement.}
\GSfield{Fate}{\citet[Theorem~1.1 and proof]{brouwerleech} reproduces Conway's uniqueness proof. The original 1969 characterization is identified there; this is prior work, not an unpublished seminar conjecture.}
\GSfield{Present status}{As of 9 November 2026: Established; uniqueness is up to isometry and the minimum is over nonzero lattice vectors.}
\GSfield{Assessment}{literal statement: correct with missing hypotheses; publication match: exact; verification: read the result; development: none; confidence: high}
\end{GSitem}
\begin{GSitem}{1970-192.5}{E}{Conway's simple group requires a central quotient}
\GSfield{Kind; source}{claim; \GSpage{P000304}; A. A. Kirillov}
\GSfield{The notes}{The notes describe another finite simple group, Conway's, as the automorphism group of this lattice.}
\GSfield{Statement}{For the Leech lattice \(\Lambda\), \(\operatorname{Aut}(\Lambda)\) contains the central inversion \(-I\). The simple Conway group is \(\mathrm{Co}_1=\operatorname{Aut}(\Lambda)/\{\pm I\}\), not the full automorphism group.}
\GSfield{Reconstruction}{The central quotient by \(\{\pm I\}\) separates the full automorphism group from the simple Conway group.}
\GSfield{Fate}{The lattice automorphism group and its order occur in \citet[Theorem~1.1]{brouwerleech}; \citet[introduction, pp.~2--3]{chermak2004quadratic} explicitly identifies it as \(2\cdot\mathrm{Co}_1\). These confirm the missing central quotient independently of the notebook.}
\GSfield{Present status}{As of 9 November 2026: Established after inserting the quotient by central inversion.}
\GSfield{Assessment}{literal statement: false as written; publication match: partial; verification: read the result; development: none; confidence: high}
\end{GSitem}
\begin{GSitem}{1970-192.6}{E}{Divisibility of the signature for even unimodular lattices}
\GSfield{Kind; source}{remark; \GSpage{P000305}}
\GSfield{The notes}{\ru{Сигнатура дел. на 8 у полож. опр. форм.}}
\GSfield{Statement}{An even unimodular positive-definite integral lattice has rank, hence signature, divisible by eight. Positivity alone is insufficient: a one-dimensional positive form has signature one.}
\GSfield{Reconstruction}{The surrounding even-unimodular lattice context is imposed; the unrestricted positive-definite-form reading is false.}
\GSfield{Fate}{The lattice hypotheses in the surrounding passage supply the intended scope; \citet[Theorem~1.2(i)]{brouwerleech} states rank divisibility for even unimodular lattices.}
\GSfield{Present status}{As of 9 November 2026: Classical theorem under evenness and unimodularity; false if read as an assertion about all positive-definite forms.}
\GSfield{Assessment}{literal statement: correct with missing hypotheses; publication match: exact; verification: read the result; development: none; confidence: high}
\end{GSitem}

\GStalk{193}{1970-10-19}{I. M. Gelfand}{P000305}{A proposed non-quasi-coherent sheaf and a localization reminder}
This short intervention starts at the dated Gelfand heading and ends before Bernstein's example. Its proposed point-supported sheaf lacks enough module-structure data to justify the advertised counterexample.

\begin{GSitem}{1970-193.1}{F}{The proposed point-supported counterexample}
\GSfield{Kind; source}{construction; \GSpage{P000305}}
\GSfield{The notes}{A sheaf is said to have fibre \(A\) at one point \(x\in\operatorname{Spec}A\), zero elsewhere, and sections \(A\) exactly on opens containing \(x\); it is called non-quasi-coherent.}
\GSfield{Statement}{The formula defines a skyscraper sheaf of abelian groups. To make it an \(\mathcal O_X\)-module one must specify an \(\mathcal O_{X,x}=A_{\mathfrak p}\)-module structure on its value. Moreover, its stalk support is the closure of \(\{x\}\), so the description of its fibres requires \(x\) to be closed. A usual skyscraper \(i_*\kappa(x)\) at a closed point is quasi-coherent, and therefore does not give the claimed counterexample.}
\GSfield{Fate}{No exact well-posed non-quasi-coherent example is identifiable from these data. The affine module/sheaf correspondence in \citet[affine case]{serre1955coherent} explains why the standard closed-point skyscraper cannot support the claim. The missing action cannot be supplied from the word fibre alone.}
\GSfield{Present status}{As of 9 November 2026: The intended example remains unidentified from the source; this is an identification limit, not a presently open theorem.}
\GSfield{Assessment}{literal statement: not determinable; publication match: none found; verification: unresolved lead; development: none; confidence: high}
\end{GSitem}
\begin{GSitem}{1970-193.2}{E}{Localization kills elements annihilated by the denominator}
\GSfield{Kind; source}{remark; \GSpage{P000305}}
\GSfield{The notes}{The boxed margin calculation is \(fm=0\Rightarrow m/f=mf/f^2=0\).}
\GSfield{Statement}{For any commutative ring \(A\), \(A\)-module \(M\), and \(f\in A\), an element \(m/f^r\in M_f\) is zero if and only if \(f^Nm=0\) for some \(N\geq0\). The displayed calculation is the case \(fm=0\).}
\GSfield{Fate}{This elementary localization fact underlies the denominator-clearing proof of affine acyclicity, given in \citet[Lemma~30.2.1]{stacks2026affine}. It is not an injectivity assertion for \(M\to M_f\).}
\GSfield{Present status}{As of 9 November 2026: Established with the hypotheses in the statement.}
\GSfield{Assessment}{literal statement: correct; publication match: exact; verification: read the result; development: none; confidence: high}
\end{GSitem}

\GStalk{194}{1970-10-19}{I. N. Bernstein}{P000305-P000306}{Regular functions on affine varieties and the module/sheaf correspondence}
Bernstein's example is followed through its ideal-sheaf sequence, the supported-section gluing argument, and the presentation argument for quasi-coherent sheaves. The distinction between one global denominator and a locally regular function is made explicitly in the manuscript.

\begin{GSitem}{1970-194.1}{E}{A globally defined rational expression on an affine variety is polynomial modulo its ideal}
\GSfield{Kind; source}{theorem; \GSpage{P000305-P000306}}
\GSfield{The notes}{If \(P/Q\) is defined everywhere on the subvariety corresponding to \(I\subset\mathbf C[x_1,\ldots,x_n]\), then \(P/Q=R\pmod I\) for a polynomial \(R\).}
\GSfield{Statement}{Let \(A=\mathbf C[x_1,\ldots,x_n]\), and suppose \(Q\) has no zero on \(V(I)\). Then \((I,Q)=A\) by the Nullstellensatz; hence \(Q\) is invertible in \(A/I\) and \(P/Q\) is the image of a polynomial. More generally, the ring of global regular functions on the affine scheme \(\operatorname{Spec}(A/I)\) is \(A/I\).}
\GSfield{Reconstruction}{For the single-fraction proof, defined means that the chosen denominator \(Q\) is nowhere zero, not merely that a removable singularity can be canceled locally.}
\GSfield{Fate}{This is a standard affine consequence of the module/sheaf correspondence in \citet[affine case]{serre1955coherent}; the explicit Nullstellensatz proof is resumed on P000309.}
\GSfield{Present status}{As of 9 November 2026: Established with the hypotheses in the statement.}
\GSfield{Assessment}{literal statement: correct with missing hypotheses; publication match: exact; verification: read the result; development: none; confidence: high}
\end{GSitem}
\begin{GSitem}{1970-194.2}{E}{Ideal sheaves and the closed affine subscheme}
\GSfield{Kind; source}{construction; \GSpage{P000305-P000306}}
\GSfield{The notes}{An ideal \(J\subset A\) gives \(\mathcal F_J\subset\mathcal O_X\); the quotient sheaf is associated with \(A/J\), and its support is the subvariety.}
\GSfield{Statement}{Sheafification on \(X=\operatorname{Spec}A\) is exact. Thus \(0\to\widetilde J\to\mathcal O_X\to\widetilde{A/J}\to0\) is exact and \(\widetilde{A/J}=i_*\mathcal O_{\operatorname{Spec}(A/J)}\) for the closed immersion \(i\). Prime ideals, not only complex closed points, are included in \(X\).}
\GSfield{Fate}{Classical affine scheme theory, as in \citet[affine case]{serre1955coherent}; arbitrary-ring affine acyclicity is verified in \citet[Lemma~30.2.2]{stacks2026affine}.}
\GSfield{Present status}{As of 9 November 2026: Established with the hypotheses in the statement.}
\GSfield{Assessment}{literal statement: correct; publication match: exact; verification: read the result; development: none; confidence: high}
\end{GSitem}
\begin{GSitem}{1970-194.3}{E}{Gluing a supported section before taking global sections}
\GSfield{Kind; source}{method; \GSpage{P000306}}
\GSfield{The notes}{The notes glue \(P/Q\) on the locus where \(Q\ne0\) to zero on the complement of the subvariety.}
\GSfield{Statement}{In \(\widetilde{A/I}\), which vanishes away from \(V(I)\), the two sections agree on the overlap. Their glued section over all of \(\operatorname{Spec}A\) lies in \(\Gamma(X,\widetilde{A/I})=A/I\). This gluing must be performed in the quotient sheaf, not in \(\mathcal O_X\).}
\GSfield{Reconstruction}{The sheaf in which the two local sections are compared is made explicit.}
\GSfield{Fate}{The argument is a concrete use of the affine correspondence, not a new extension theorem; compare \citet[affine case]{serre1955coherent}.}
\GSfield{Present status}{As of 9 November 2026: Established with the hypotheses in the statement.}
\GSfield{Assessment}{literal statement: correct with missing hypotheses; publication match: exact; verification: read the result; development: none; confidence: high}
\end{GSitem}
\begin{GSitem}{1970-194.4}{E}{Recovering a quasi-coherent sheaf from its global module}
\GSfield{Kind; source}{theorem; \GSpage{P000306}}
\GSfield{The notes}{A presentation by sums of \(\mathcal O_X\) is compared with the corresponding module presentation, concluding \(\mathcal F=\mathcal F_M\).}
\GSfield{Statement}{For every affine scheme \(X=\operatorname{Spec}A\), the functors \(M\mapsto\widetilde M\) and \(\mathcal F\mapsto\Gamma(X,\mathcal F)\) are inverse equivalences between \(A\)-modules and quasi-coherent \(\mathcal O_X\)-modules. Presentations can use arbitrary direct sums; finite presentations require the corresponding finite-presentation hypothesis.}
\GSfield{Fate}{The coherent affine case predates the seminar in \citet{serre1955coherent}; the standard arbitrary-ring formulation is compatible with \citet[Lemma~30.2.2 and its preceding section]{stacks2026affine}.}
\GSfield{Present status}{As of 9 November 2026: Established with the hypotheses in the statement.}
\GSfield{Assessment}{literal statement: correct with missing hypotheses; publication match: exact; verification: read the result; development: none; confidence: high}
\end{GSitem}

\GStalk{195}{1970-10-19}{Ya. G. Sinai}{P000307-P000308}{States of a one-dimensional quantum spin chain}
The lecture constructs the quasi-local matrix algebra, states and their density matrices, diagonal classical states, finite-temperature Gibbs limits, and the GNS representation. The existence and factor questions are revisited in records 197 and 200.

\begin{GSitem}{1970-195.1}{E}{The quasi-local spin algebra}
\GSfield{Kind; source}{construction; \GSpage{P000307}}
\GSfield{The notes}{Two-dimensional spaces \(H_i\), tensor products \(H_I\), matrix rings \(\mathcal A(I)\), and inclusions by tensoring with the identity are introduced.}
\GSfield{Statement}{For finite \(I\subset\mathbf Z\), set \(H_I=\bigotimes_{i\in I}\mathbf C^2\) and \(\mathcal A(I)=B(H_I)\). The compatible embeddings \(a\mapsto a\otimes1\) give the unital UHF \(C^*\)-algebra \(\mathcal A=\overline{\bigcup_I\mathcal A(I)}^{\|\cdot\|}\). In particular \(\dim H_I=2^{|I|}\), not \(2|I|\).}
\GSfield{Fate}{This uniformly hyperfinite setup is used in \citet[section~2]{powers1967representations}. It was standard before the lecture.}
\GSfield{Present status}{As of 9 November 2026: Established with the hypotheses in the statement.}
\GSfield{Assessment}{literal statement: correct; publication match: exact; verification: read the result; development: none; confidence: high}
\end{GSitem}
\begin{GSitem}{1970-195.2}{E}{States and local density matrices}
\GSfield{Kind; source}{definition; \GSpage{P000307}}
\GSfield{The notes}{A positive functional \(\varphi\) with \(\varphi(1)=1\) restricts to \(\operatorname{Tr}(a\rho_I)\) on each local matrix algebra.}
\GSfield{Statement}{A state is a positive complex-linear functional of norm one. Every restriction to \(\mathcal A(I)\) has a unique density matrix \(\rho_I\geq0\), \(\operatorname{Tr}\rho_I=1\); consistency means \(\operatorname{Tr}_{J\setminus I}\rho_J=\rho_I\) for \(I\subset J\). Conversely a compatible family defines a state on \(\mathcal A\).}
\GSfield{Fate}{The finite-dimensional trace representation and the UHF state construction are the framework of \citet[section~2]{powers1967representations}.}
\GSfield{Present status}{As of 9 November 2026: Established with the hypotheses in the statement.}
\GSfield{Assessment}{literal statement: correct; publication match: exact; verification: read the result; development: none; confidence: high}
\end{GSitem}
\begin{GSitem}{1970-195.3}{E}{Product-diagonal states and cylinder measures}
\GSfield{Kind; source}{construction; \GSpage{P000307}}
\GSfield{The notes}{In a fixed sitewise basis the eigenvalues of the density matrices are treated as joint probabilities.}
\GSfield{Statement}{Choose an orthonormal basis at each site. Compatible density matrices diagonal in all resulting tensor bases are equivalent to probability measures on \(\{1,2\}^{\mathbf Z}\). For a finite word \(\alpha\), its probability is \(\varphi(\bigotimes_{i\in I}|e^i_{\alpha_i}\rangle\langle e^i_{\alpha_i}|)\). It is not the value of \(\varphi\) on a tensor product of vectors.}
\GSfield{Reconstruction}{Rank-one projections replace the notebook's vector expression so that the state is applied to an element of its algebra.}
\GSfield{Fate}{This is the commutative cylinder-measure construction inside the diagonal subalgebra of the UHF algebra. Product states are examined in \citet[section~4]{powers1967representations}. Record 197 attempts a broader commutation criterion, which fails without additional hypotheses.}
\GSfield{Present status}{As of 9 November 2026: Established with the hypotheses in the statement.}
\GSfield{Assessment}{literal statement: correct with missing hypotheses; publication match: partial; verification: read the result; development: none; confidence: high}
\end{GSitem}
\begin{GSitem}{1970-195.4}{E}{Finite-volume Gibbs functionals}
\GSfield{Kind; source}{construction; \GSpage{P000308}}
\GSfield{The notes}{A translated nearest-neighbor operator is summed to form a Hamiltonian, and normalized traces of \(ae^{-\beta H_n}\) are taken.}
\GSfield{Statement}{For a self-adjoint nearest-neighbor interaction \(h\in\mathcal A(\{0,1\})\), real finite \(\beta\), and a finite interval \(\Lambda\), put \(H_\Lambda=\sum_{\{j,j+1\}\subset\Lambda}\tau_j(h)\). The state is \(\varphi_\Lambda(a)=\operatorname{Tr}(ae^{-\beta H_\Lambda})/\operatorname{Tr}(e^{-\beta H_\Lambda})\) on \(\mathcal A(\Lambda)\).}
\GSfield{Reconstruction}{Self-adjointness of the local interaction, real finite inverse temperature, and the free finite-volume boundary convention are supplied to make the normalized trace a state.}
\GSfield{Fate}{The finite-range spin-chain Gibbs construction is treated in \citet{araki1969gibbs}. Positivity requires a self-adjoint Hamiltonian; no arbitrary matrix \(h\) suffices.}
\GSfield{Present status}{As of 9 November 2026: Established with the hypotheses in the statement.}
\GSfield{Assessment}{literal statement: correct with missing hypotheses; publication match: exact; verification: publisher metadata checked; development: none; confidence: high}
\end{GSitem}
\begin{GSitem}{1970-195.5}{E}{The one-dimensional Gibbs limit}
\GSfield{Kind; source}{theorem; \GSpage{P000308}}
\GSfield{The notes}{The notes assert convergence \(\varphi_n(a)\to\varphi(a)\) for every quasi-local observable.}
\GSfield{Statement}{For a translation-invariant finite-range self-adjoint interaction on a one-dimensional quantum spin lattice, the finite-temperature thermodynamic limit with the usual free finite-volume boundary convention exists. The resulting state is primary; the finite-dimensional local state spaces and finite interaction range are part of the assertion.}
\GSfield{Reconstruction}{The hypotheses are those of the identified Araki theorem; the assertion is not extended beyond finite-range one-dimensional quantum spin systems.}
\GSfield{Fate}{\citet[abstract]{araki1969gibbs}, published in 1969, states existence of the infinite-volume primary Gibbs state, before Sinai's lecture. This is an exact match for the finite-range one-dimensional situation, not an assertion about arbitrary dimension or long-range interactions.}
\GSfield{Present status}{As of 9 November 2026: Established for the stated one-dimensional finite-range finite-temperature system.}
\GSfield{Assessment}{literal statement: correct with missing hypotheses; publication match: exact; verification: publisher metadata checked; development: none; confidence: high}
\end{GSitem}
\begin{GSitem}{1970-195.6}{E}{The analogy with Markov chains}
\GSfield{Kind; source}{remark; \GSpage{P000308}}
\GSfield{The notes}{The Gibbs construction is called a noncommutative analogue of Markov chains.}
\GSfield{Statement}{The nearest-neighbor energy resembles a one-step Markov model with Gibbs weights. For noncommuting quantum interactions this analogy alone does not assert a conditional-independence identity, a transition operator, or the modern quantum Markov property.}
\GSfield{Fate}{\citet{araki1969gibbs} supplies a rigorous quantum Gibbs theory. No exact quantum Markov theorem is stated or proved in the notebook passage; the publication match is thematic.}
\GSfield{Present status}{As of 9 November 2026: A useful structural analogy, with no additional Markov theorem inferred.}
\GSfield{Assessment}{literal statement: correct; publication match: thematic only; verification: publisher metadata checked; development: none; confidence: high}
\end{GSitem}
\begin{GSitem}{1970-195.7}{E}{The GNS Hilbert space and generated von Neumann algebra}
\GSfield{Kind; source}{construction; \GSpage{P000308}}
\GSfield{The notes}{\((a,b)=\varphi(b^*a)\); factor by null vectors, complete, and let the algebra act by left multiplication.}
\GSfield{Statement}{Let \(N_\varphi=\{a:\varphi(a^*a)=0\}\). Complete \(\mathcal A/N_\varphi\) in the inner product \(\langle[a],[b]\rangle=\varphi(b^*a)\). Left multiplication gives a bounded representation \(\pi_\varphi\), since \(\|\pi_\varphi(x)\|\leq\|x\|\), and \(M_\varphi=\pi_\varphi(\mathcal A)^{\prime\prime}\). The null space is a left ideal; it need not be a two-sided ideal.}
\GSfield{Fate}{The construction and factor representations are used throughout \citet{powers1967representations}; the notebook cites the classical Gelfand--Naimark tradition. This is prior foundational theory.}
\GSfield{Present status}{As of 9 November 2026: Established with the hypotheses in the statement.}
\GSfield{Assessment}{literal statement: correct; publication match: exact; verification: read the result; development: none; confidence: high}
\end{GSitem}
\begin{GSitem}{1970-195.8}{E}{When is the GNS algebra a factor, and what type occurs?}
\GSfield{Kind; source}{question; \GSpage{P000308}}
\GSfield{The notes}{The notes ask when the generated ring is a factor and emphasize factors of type III.}
\GSfield{Statement}{The factor question is whether \(Z(M_\varphi)=\mathbf C1\). It is separate from the type of a factor. States of the same UHF algebra can yield type I, type II, or type III factors; the phrase principally type III is a motivation, not a theorem about every state.}
\GSfield{Fate}{\citet[Theorem~2.5]{powers1967representations} gives the factor criterion later explained in record 200; section~4 exhibits the different types for product states. In particular the tracial product state yields type \(\mathrm{II}_1\).}
\GSfield{Present status}{As of 9 November 2026: The factor criterion and important examples were established before the seminar; no universal type-III assertion is valid.}
\GSfield{Assessment}{literal statement: correct with missing hypotheses; publication match: exact; verification: read the result; development: none; confidence: high; problem \hyperref[prob:GP-1970-22]{GP-1970-22}}
\end{GSitem}

\GStalk{196}{1970-10-26}{D. A. Kazhdan}{P000309-P000311}{Affine acyclicity by localization and Čech cochains}
The proof works with a finite distinguished-open cover. It includes both an explicit denominator-clearing attempt and Bernstein's annihilator argument; the manuscript itself acknowledges that localization equalities need more care.

\begin{GSitem}{1970-196.1}{E}{Vanishing of higher cohomology on an affine scheme}
\GSfield{Kind; source}{theorem; \GSpage{P000309}}
\GSfield{The notes}{For an affine scheme \(X\) and a quasi-coherent sheaf \(\mathcal F\), \(H^q(X,\mathcal F)=0\), \(q>0\).}
\GSfield{Statement}{For every commutative ring \(A\) and every \(A\)-module \(M\), \[H^q(\operatorname{Spec}A,\widetilde M)=0\qquad(q>0).\] Neither Noetherianity nor finite generation of \(M\) is required.}
\GSfield{Fate}{The coherent algebraic-geometric theorem is in \citet{serre1955coherent}. The exact arbitrary-ring formulation and proof are \citet[Lemma~30.2.2]{stacks2026affine}.}
\GSfield{Present status}{As of 9 November 2026: Established with the hypotheses in the statement.}
\GSfield{Assessment}{literal statement: correct; publication match: exact; verification: read the result; development: none; confidence: high}
\end{GSitem}
\begin{GSitem}{1970-196.2}{E}{The Nullstellensatz proof of Bernstein's rational-function example}
\GSfield{Kind; source}{method; \GSpage{P000309}}
\GSfield{The notes}{For generators \(F_i\) of the ideal and a nowhere-zero \(Q\), find \(L_0Q+\sum_iL_iF_i=P\).}
\GSfield{Statement}{If \((Q,F_1,\ldots,F_r)\) has no common zero over an algebraically closed field, it is the unit ideal. Multiplying a representation of one by \(P\) produces the displayed identity and \(P/Q\equiv L_0\pmod I\).}
\GSfield{Fate}{This proves the single-denominator instance of 1970-194.1 by elementary affine algebra; it does not need the cohomology theorem subsequently being proved. Compare \citet[affine case]{serre1955coherent}.}
\GSfield{Present status}{As of 9 November 2026: Established with the hypotheses in the statement.}
\GSfield{Assessment}{literal statement: correct; publication match: exact; verification: read the result; development: none; confidence: high}
\end{GSitem}
\begin{GSitem}{1970-196.3}{E}{The localization Čech complex}
\GSfield{Kind; source}{construction; \GSpage{P000309-P000310}}
\GSfield{The notes}{For \(X_i=D(f_i)\) covering \(\operatorname{Spec}A\), cochains have components in \(M_{f_{i_0}\cdots f_{i_q}}\), with alternating restriction differential.}
\GSfield{Statement}{Choose \(f_1,\ldots,f_r\) generating the unit ideal. The augmented complex starts \(M\to\prod_i M_{f_i}\to\prod_{i<j}M_{f_if_j}\to\cdots\). With increasing indices, \((dc)_{i_0\cdots i_{q+1}}=\sum_{s=0}^{q+1}(-1)^s c_{i_0\cdots\widehat{i_s}\cdots i_{q+1}}\) after restriction.}
\GSfield{Fate}{\citet[Lemma~30.2.1]{stacks2026affine} gives exactness for this finite principal cover and identifies the finite localization argument.}
\GSfield{Present status}{As of 9 November 2026: Established with the hypotheses in the statement.}
\GSfield{Assessment}{literal statement: correct with missing hypotheses; publication match: exact; verification: read the result; development: none; confidence: high}
\end{GSitem}
\begin{GSitem}{1970-196.4}{E}{Alternating cochains and the failed single-chart contraction}
\GSfield{Kind; source}{method; \GSpage{P000310}}
\GSfield{The notes}{The notes pass to skew-symmetric cochains; for a one-cocycle they first try one distinguished index, then note that \(m_{1i}\) may have an unwanted denominator.}
\GSfield{Statement}{The normalized alternating Čech complex computes the same cohomology as the full ordered-index complex. A contraction based on a fixed chart is valid after localizing where that chart covers everything, but \(m_{1i}\in M_{f_1f_i}\) need not already lie in \(M_{f_i}\). The latter obstruction invalidates a global single-chart primitive.}
\GSfield{Fate}{The localization contraction in \citet[proof of Lemma~30.2.1]{stacks2026affine} makes precisely this distinction. The notebook records the failed step as part of a proof search, not a completed proof.}
\GSfield{Present status}{As of 9 November 2026: Established with the hypotheses in the statement.}
\GSfield{Assessment}{literal statement: correct; publication match: partial; verification: read the result; development: none; confidence: high}
\end{GSitem}
\begin{GSitem}{1970-196.5}{E}{Uniform denominator clearing and a partition of unity in the ring}
\GSfield{Kind; source}{method; \GSpage{P000311}}
\GSfield{The notes}{Write \(m_{ij}=\bar m_{ij}/(f_i^Nf_j^N)\), choose \(\sum_iL_if_i^N=1\), and sum over \(i\).}
\GSfield{Statement}{After increasing a common power enough to clear all denominators and all finitely many cocycle equalities, a ring identity \(\sum_i L_if_i^N=1\) gives a primitive. For the convention \(m_{ij}=n_i-n_j\), one can set \(n_i=\sum_k L_k f_k^N m_{ik}\) in \(M_{f_i}\), after the clearing step. Equality in a localization may require an additional power before it holds in an earlier module.}
\GSfield{Fate}{This is the standard algebraic proof behind \citet[Lemma~30.2.1]{stacks2026affine}. The corrected formula removes an extraneous denominator from one displayed intermediate equality on P000311; the later lines of the manuscript have the intended sum.}
\GSfield{Present status}{As of 9 November 2026: Established with the hypotheses in the statement.}
\GSfield{Assessment}{literal statement: correct with missing hypotheses; publication match: exact; verification: read the result; development: annotated; confidence: high}
\end{GSitem}
\begin{GSitem}{1970-196.6}{E}{Bernstein's annihilator proof of affine acyclicity}
\GSfield{Kind; source}{method; \GSpage{P000311}; I. N. Bernstein}
\GSfield{The notes}{Each class is killed by a power of every \(f_i\), and those powers still generate the unit ideal.}
\GSfield{Statement}{Localizing the Čech complex at \(f_i\) makes its cover contain the whole affine scheme and hence contracts it. For any given higher-cohomology class \(c\), there are powers with \(f_i^{N_i}c=0\). Finiteness of the cover gives a common \(N\); since \((f_1^N,\ldots,f_r^N)=A\), one obtains \(c=0\).}
\GSfield{Fate}{\citet[proof of Lemma~30.2.1]{stacks2026affine} presents the same localization-to-vanishing mechanism. The argument is valid for arbitrary modules.}
\GSfield{Present status}{As of 9 November 2026: Established with the hypotheses in the statement.}
\GSfield{Assessment}{literal statement: correct; publication match: exact; verification: read the result; development: none; confidence: high}
\end{GSitem}

\GStalk{197}{1970-10-26}{Ya. G. Sinai}{P000312-P000313}{Commuting densities, classical measures, and a failed diagonalization claim}
The quantum-state construction is repeated, now with partial traces. The assertion that mutually commuting embedded marginals admit a fixed sitewise product eigenbasis fails in the presence of degeneracies; an explicit faithful counterexample is supplied in the chapter's audit annotation.

\begin{GSitem}{1970-197.1}{E}{Tensor-product local systems and compatible density operators}
\GSfield{Kind; source}{construction; \GSpage{P000312}}
\GSfield{The notes}{The first display has \(H_I=\bigoplus_{i\in I}H_i\), but the basis, algebra embeddings and partial trace formulas use tensor products.}
\GSfield{Statement}{The intended local Hilbert space is \(H_I=\bigotimes_{i\in I}H_i\), with \(\mathcal A(J)=\mathcal A(I)\otimes\mathcal A(J\setminus I)\) and \(\operatorname{Tr}_{J\setminus I}\rho_J=\rho_I\). This restates 1970-195.1--2.}
\GSfield{Fate}{The UHF framework in \citet[section~2]{powers1967representations} fixes the notation. The direct-sum sign is a faithful manuscript slip, not a valid alternative construction.}
\GSfield{Present status}{As of 9 November 2026: Established with the hypotheses in the statement.}
\GSfield{Assessment}{literal statement: false as written; publication match: exact; verification: read the result; development: annotated; confidence: high}
\end{GSitem}
\begin{GSitem}{1970-197.2}{E}{Commuting finite marginals need not have a common product eigenbasis}
\GSfield{Kind; source}{claim; \GSpage{P000312-P000313}}
\GSfield{The notes}{A state is called of measure type when all embedded \(\rho_I\) commute. The next page asserts a common tensor basis of eigenvectors of the one-site densities, explicitly noting that the argument is obvious for simple spectra.}
\GSfield{Statement}{Product diagonality implies commutation. The converse is valid if every one-site density has simple spectrum: commuting with each local one-site density forces diagonalization in their product basis. Without that hypothesis, pairwise commutation of all embedded finite marginals is insufficient, even when every finite marginal is strictly positive.}
\GSfield{Reconstruction}{The notebook's commutation definition and the stronger product-diagonal property are kept separate. A repair by simple one-site spectra is explicitly conditional.}
\GSfield{Fate}{The implication fails by the faithful Bell-pair construction checked in Section~\ref{dev:bell}. The product-diagonal construction remains valid; \citet[section~4]{powers1967representations} treats product-state representations but does not supply the false general implication.}
\GSfield{Present status}{As of 9 November 2026: False without the missing nondegeneracy or another hypothesis ensuring a product eigenbasis; it is not an open problem.}
\GSfield{Assessment}{literal statement: false as written; publication match: partial; verification: read the result; development: annotated; confidence: high}
\end{GSitem}
\begin{GSitem}{1970-197.3}{E}{Consistent eigenvalue tables and the infinite product probability space}
\GSfield{Kind; source}{construction; \GSpage{P000313}}
\GSfield{The notes}{The sums \(\sum_{\alpha_n}\mu(\alpha_1,\ldots,\alpha_n)=\mu(\alpha_1,\ldots,\alpha_{n-1})\) define a cylinder measure on \(\prod_i\mathbf Z_2\); Haar measure is an example.}
\GSfield{Statement}{When a fixed product basis really diagonalizes the densities, partial traces give consistent finite probability tables. The resulting Borel probability measure on \(\{0,1\}^{\mathbf Z}\) is unique. The product of equal weights is the Haar probability measure on the compact product group \(\prod_{i\in\mathbf Z}\mathbf Z/2\mathbf Z\).}
\GSfield{Fate}{This is the valid classical part of 1970-195.3. For the Haar example the corresponding quantum state is tracial, so \citet[section~4]{powers1967representations} identifies a type \(\mathrm{II}_1\), rather than type III, factor.}
\GSfield{Present status}{As of 9 November 2026: Established with the hypotheses in the statement.}
\GSfield{Assessment}{literal statement: correct with missing hypotheses; publication match: partial; verification: read the result; development: none; confidence: high}
\end{GSitem}
\begin{GSitem}{1970-197.4}{E}{GNS representation and the role of type III factors}
\GSfield{Kind; source}{construction; \GSpage{P000313}}
\GSfield{The notes}{The positive form, quotient, completion, and left multiplication are repeated, followed by the statement that type III factors are important.}
\GSfield{Statement}{The GNS construction in 1970-195.7 applies to all these states. The study of type III factors motivates selecting nontracial infinite product and Gibbs states; it does not classify every state constructed on P000312--313 as type III.}
\GSfield{Fate}{\citet[section~4]{powers1967representations} constructs type III product-state factors and distinguishes the pure and tracial endpoints. This supplies the prior mathematical background without converting a motivational remark into a universal theorem.}
\GSfield{Present status}{As of 9 November 2026: Established with the hypotheses in the statement.}
\GSfield{Assessment}{literal statement: correct; publication match: partial; verification: read the result; development: none; confidence: high; problem \hyperref[prob:GP-1970-22]{GP-1970-22}}
\end{GSitem}

\GStalk{198}{1970-11-02}{I. M. Gelfand}{P000314-P000317}{Four subspaces: classification, normal forms, and deformation}
Gelfand reports joint work with V. A. Ponomarev. The graph examples and most displayed normal forms match the 1971 Doklady announcement very closely. The defect sign, missing odd-dimensional defect-zero systems, and broad stability assertions require separate audit; the classification itself is established.

\begin{GSitem}{1970-198.1}{A}{The category of ordered four-subspace systems}
\GSfield{Kind; source}{definition; \GSpage{P000314}}
\GSfield{The notes}{\(S=(P;E_1,E_2,E_3,E_4)\), with ordered subspaces and morphisms carrying each \(E_i\) into its counterpart; direct sums and permutations are introduced.}
\GSfield{Statement}{Objects consist of a finite-dimensional vector space \(P\) over \(k\) and four ordered subspaces. A morphism \(u:S\to S^{\prime}\) satisfies \(u(E_i)\subset E_i^{\prime}\); an isomorphism is an invertible \(u\) carrying each subspace onto the corresponding subspace. Direct sums are componentwise. Permuting labels is a separate action, not part of ordered isomorphism.}
\GSfield{Fate}{Exactly \citet[section~1, p.~762]{gelfandponomarev1971quadruples}. The announcement was received 21 October 1970 and printed in 1971.}
\GSfield{Present status}{As of 9 November 2026: Established with the hypotheses in the statement.}
\GSfield{Assessment}{literal statement: correct; publication match: exact; verification: read the result; development: none; confidence: high}
\end{GSitem}
\begin{GSitem}{1970-198.2}{A}{The four-armed quiver and its injectivity condition}
\GSfield{Kind; source}{construction; \GSpage{P000314}}
\GSfield{The notes}{The category is said to be equivalent to diagrams with four arrows directed into one central vertex.}
\GSfield{Statement}{Four-subspace systems identify with representations of the inward-oriented affine \(D_4\) quiver in which all four arrow maps are injective. Arbitrary representations also include kernels at outer vertices and hence are not literally the same category. The diagrammatic language unifies several matrix classification problems.}
\GSfield{Reconstruction}{Injectivity of the four arrows is supplied; it is forced by calling their domains subspaces.}
\GSfield{Fate}{The subsequent quiver framework is explicit in \citet[introduction and final example]{bernsteingelfandponomarev1973coxeter}; compare \citet[section~1]{gelfandponomarev1971quadruples}.}
\GSfield{Present status}{As of 9 November 2026: Established with the hypotheses in the statement.}
\GSfield{Assessment}{literal statement: correct with missing hypotheses; publication match: exact; verification: read the result; development: none; confidence: high}
\end{GSitem}
\begin{GSitem}{1970-198.3}{A}{Encoding similarity of an operator by four subspaces}
\GSfield{Kind; source}{construction; \GSpage{P000314}}
\GSfield{The notes}{In \(E\oplus E\), take the two coordinate summands, the graph of \(A:E\to E\), and the diagonal.}
\GSfield{Statement}{The quadruple \(S_A=(E\oplus E;E\oplus0,0\oplus E,\operatorname{graph}A,\Delta_E)\) is isomorphic to \(S_B\) exactly when \(A\) and \(B\) are similar. The first two subspaces force a block diagonal map and the diagonal forces the two blocks to agree. Direct sum and indecomposability correspond to those of operators.}
\GSfield{Fate}{The identical construction is \citet[section~1, example~1]{gelfandponomarev1971quadruples}; the equivalence is also stated in \citet[p.~2]{enomotowatatani2004subspaces}.}
\GSfield{Present status}{As of 9 November 2026: Established with the hypotheses in the statement.}
\GSfield{Assessment}{literal statement: correct; publication match: exact; verification: read the result; development: none; confidence: high}
\end{GSitem}
\begin{GSitem}{1970-198.4}{A}{Opposite maps, matrix pencils, and linear relations}
\GSfield{Kind; source}{construction; \GSpage{P000314}}
\GSfield{The notes}{The notes mention \(A:E\to F\), \(B:F\to E\), Kronecker's problem, and relations regarded as arbitrary graphs.}
\GSfield{Statement}{In \(E\oplus F\), take the two coordinate subspaces together with \(\{(x,Ax)\}\) and \(\{(By,y)\}\). Two parallel maps \(E\to F\) instead give two ordinary graphs and encode a matrix pencil. Replacing graphs by arbitrary subspaces of \(E\oplus F\) encodes linear relations; a relation need not be the graph of a single-valued map.}
\GSfield{Reconstruction}{The three related encodings are written separately rather than treating the word graph as imposing single-valuedness.}
\GSfield{Fate}{These are exactly the five examples in \citet[section~1, pp.~762--763; section~4]{gelfandponomarev1971quadruples}. The article explains that the four-subspace problem includes, but is broader than, pairs of linear relations.}
\GSfield{Present status}{As of 9 November 2026: Established with the hypotheses in the statement.}
\GSfield{Assessment}{literal statement: correct; publication match: exact; verification: read the result; development: none; confidence: high}
\end{GSitem}
\begin{GSitem}{1970-198.5}{E}{Four lines and the cross-ratio parameter}
\GSfield{Kind; source}{construction; \GSpage{P000314-P000316}}
\GSfield{The notes}{Four lines in a plane have one invariant \(\lambda\), with its transformations under relabeling.}
\GSfield{Statement}{For four distinct ordered lines in a two-dimensional vector space over an algebraically closed field, normalize three projective points to \(0,1,\infty\). The fourth gives \(\lambda\notin\{0,1\}\). Relabeling gives the six values \(\lambda,1-\lambda,1/\lambda,1/(1-\lambda),\lambda/(\lambda-1),(\lambda-1)/\lambda\), with coincidences at special values.}
\GSfield{Fate}{This is the \(k=1\) parameter family in \citet[section~2]{gelfandponomarev1971quadruples}. The P000316 list, and even the compressed list in the announcement, omit one of the six cross-ratio transforms; the earlier P000314 generators recover it.}
\GSfield{Present status}{As of 9 November 2026: Established with the hypotheses in the statement.}
\GSfield{Assessment}{literal statement: correct with missing hypotheses; publication match: partial; verification: read the result; development: annotated; confidence: high}
\end{GSitem}
\begin{GSitem}{1970-198.6}{A}{Indecomposable classification and the five possible defects}
\GSfield{Kind; source}{theorem; \GSpage{P000315}}
\GSfield{The notes}{Every system decomposes into indecomposables of defect \(0,\pm1,\pm2\); the boxed definition is \(2\dim P-\sum_i\dim E_i\).}
\GSfield{Statement}{Over an algebraically closed field of characteristic zero, every finite four-subspace system is a direct sum of the canonical indecomposables in the Gelfand--Ponomarev list. With the published convention \(\delta(S)=\sum_i\dim E_i-2\dim P\), indecomposable defects belong to \(\{-2,-1,0,1,2\}\). The notebook's boxed definition is \(-\delta\), while its subsequent labels use \(\delta\).}
\GSfield{Fate}{\citet[Claims~1--2 and Theorem~1, pp.~764--765]{gelfandponomarev1971quadruples} gives the classification; \citet[Theorems~9.1--9.2]{enomotowatatani2004subspaces} reproduces the complete list. The original received date is before this lecture, but printing is 1971.}
\GSfield{Present status}{As of 9 November 2026: Classified; the defect sign must be normalized before using the listed labels.}
\GSfield{Assessment}{literal statement: correct with missing hypotheses; publication match: exact; verification: read the result; development: annotated; confidence: high}
\end{GSitem}
\begin{GSitem}{1970-198.7}{A}{The two even-dimensional defect-one normal forms}
\GSfield{Kind; source}{construction; \GSpage{P000315}}
\GSfield{The notes}{For \(P=\langle e_1,\ldots,e_k,f_1,\ldots,f_k\rangle\), the first system has \(E_3=\langle e_2+f_1,\ldots,e_k+f_{k-1}\rangle\); the second adjoins \(e_1,f_k\).}
\GSfield{Statement}{In both forms \(E_1=\langle e_i\rangle\), \(E_2=\langle f_i\rangle\), and \(E_4=\langle e_i+f_i\rangle\). The first has \(\dim E_3=k-1\) and published defect \(-1\), the second \(\dim E_3=k+1\) and defect \(+1\). Label permutations give the other distinguished-index forms.}
\GSfield{Fate}{The displayed bases match \citet[section~2, p.~763]{gelfandponomarev1971quadruples}. Only the notebook's preceding definition reverses the defect sign.}
\GSfield{Present status}{As of 9 November 2026: Established with the hypotheses in the statement.}
\GSfield{Assessment}{literal statement: correct with missing hypotheses; publication match: exact; verification: read the result; development: annotated; confidence: high}
\end{GSitem}
\begin{GSitem}{1970-198.8}{A}{The exceptional and one-parameter even-dimensional defect-zero forms}
\GSfield{Kind; source}{construction; \GSpage{P000315-P000316}}
\GSfield{The notes}{One form takes \(E_3=\langle e_1,e_2+f_1,\ldots,e_k+f_{k-1}\rangle\); the other has a Jordan-chain parameter \(\lambda\ne0,1\).}
\GSfield{Statement}{Keep \(E_1,E_2,E_4\) as in 1970-198.7. The exceptional form has the stated \(E_3\); its permutations give \(S_{i,j}(2k,0)\). The family \(S(2k,0;\lambda)\) has \(E_3=\langle e_1+\lambda f_1,e_2+f_1+\lambda f_2,\ldots,e_k+f_{k-1}+\lambda f_k\rangle\), \(\lambda\in k\setminus\{0,1\}\). The excluded parameters belong to exceptional label families, not to additional generic families.}
\GSfield{Fate}{\citet[section~2, p.~763 and Claim~1]{gelfandponomarev1971quadruples} gives these exact forms and their permutation rules.}
\GSfield{Present status}{As of 9 November 2026: Established with the hypotheses in the statement.}
\GSfield{Assessment}{literal statement: correct; publication match: exact; verification: read the result; development: none; confidence: high}
\end{GSitem}
\begin{GSitem}{1970-198.9}{A}{Odd-dimensional defect-two forms and duality}
\GSfield{Kind; source}{construction; \GSpage{P000316}}
\GSfield{The notes}{The notes write \(S(2k-1,2)\) explicitly and obtain \(S(2k-1,-2)\) by taking orthogonals.}
\GSfield{Statement}{For \(k\geq2\), in dimension \(2k-1\), the displayed four subspaces each have dimension \(k\), so their published defect is \(+2\). Their annihilators in \(P^*\) have dimension \(k-1\), hence defect \(-2\). Dual annihilators are the field-independent meaning of orthogonals; a chosen inner product may identify them with orthogonal complements. The one-dimensional endpoints are handled separately.}
\GSfield{Fate}{The bases coincide with \citet[section~2, p.~764]{gelfandponomarev1971quadruples}; see also \citet[section~9, forms~8--9]{enomotowatatani2004subspaces}.}
\GSfield{Present status}{As of 9 November 2026: Established with the hypotheses in the statement.}
\GSfield{Assessment}{literal statement: correct with missing hypotheses; publication match: exact; verification: read the result; development: none; confidence: high}
\end{GSitem}
\begin{GSitem}{1970-198.10}{A}{The notebook incorrectly excludes odd-dimensional defect zero}
\GSfield{Kind; source}{claim; \GSpage{P000316}}
\GSfield{The notes}{\ru{def 0 в неч.-мерном не бывает.}}
\GSfield{Statement}{Odd-dimensional indecomposables of defect zero do exist. For example \((k;k,0,k,0)\) in ambient dimension one is indecomposable and has defect zero. More generally the published list contains \(S_{i,j}(2k-1,0)\). The manuscript omits these forms, together with its full odd-dimensional defect-one list.}
\GSfield{Fate}{\citet[Claim~2, p.~764]{gelfandponomarev1971quadruples} explicitly includes the missing odd families. The counterexample follows directly from the definition, with no unresolved classification issue.}
\GSfield{Present status}{As of 9 November 2026: The literal exclusion is false; the published classification supplies the omitted families.}
\GSfield{Assessment}{literal statement: false as written; publication match: partial; verification: read the result; development: annotated; confidence: high}
\end{GSitem}
\begin{GSitem}{1970-198.11}{E}{Parameters and local rigidity are not interchangeable}
\GSfield{Kind; source}{claim; \GSpage{P000316-P000317}}
\GSfield{The notes}{General-position systems are said to have no parameter; parameter-free indecomposables are said always to be stable under deformation.}
\GSfield{Statement}{For fixed dimensions, local rigidity means that the \(\mathrm{GL}(P)\)-orbit is open in the product of Grassmannians. Four distinct lines in a plane form a general-position family with varying cross ratio. A finite exceptional list does not by itself imply open orbits: exceptional defect-zero systems can deform into the parameter family.}
\GSfield{Fate}{The explicit families and exceptional specializations in \citet[section~2]{gelfandponomarev1971quadruples} already disprove the unrestricted parameter/general-position slogan. The regular-versus-irregular distinction in \citet[final example]{bernsteingelfandponomarev1973coxeter} is the appropriate later framework; it does not identify every parameter-free case as rigid.}
\GSfield{Present status}{As of 9 November 2026: False in the stated generality; rigidity requires an orbit or deformation calculation.}
\GSfield{Assessment}{literal statement: false as written; publication match: partial; verification: read the result; development: annotated; confidence: high}
\end{GSitem}
\begin{GSitem}{1970-198.12}{C}{A direct sum of rigid systems need not be rigid}
\GSfield{Kind; source}{question; \GSpage{P000317}}
\GSfield{The notes}{\(S(5,2)\) and \(S(1,2)\) are called stable, but their sum deforms to \(S(3,2)\oplus S(3,2)\); a simple proof is requested.}
\GSfield{Statement}{Using the published defect convention, all three terms are the defect-two canonical systems; both sums have ambient dimension six and four subspaces of dimension four. The intended assertion is that arbitrarily small deformations of \(S(5,2)\oplus S(1,2)\) have the isomorphism type \(S(3,2)^{\oplus2}\). Thus the latter is the generic nearby orbit; the former lies in its orbit closure.}
\GSfield{Reconstruction}{Work over \(\mathbf C\) with ordered subspaces and analytic deformations in the product of Grassmannians. The supplied proof is independently checked reconstruction, not a claim of historical novelty.}
\GSfield{Fate}{The canonical systems are exactly those of \citet[section~2]{gelfandponomarev1971quadruples}. The opened classification texts did not locate this precise degeneration. Section~\ref{dev:four-subspaces} proves it: \(\operatorname{End}S(3,2)=\mathbf C\), so the orbit of \(S(3,2)^{\oplus2}\) has dimension \(36-4=32=\dim\operatorname{Gr}(4,6)^4\) and is open dense. A generic line in a Grassmann chart through the special sum supplies an analytic family. Class C here records the partial historical literature match followed by a completed bounded development, not present mathematical openness.}
\GSfield{Present status}{As of 9 November 2026: the precise deformation over the complex numbers is proved in this chapter. No mathematical obstruction remains for this bounded question; only the exact historical publication of its elementary proof was not verified.}
\GSfield{Assessment}{literal statement: correct with missing hypotheses; publication match: partial; verification: read the result; development: developed in this chapter; confidence: high; problem \hyperref[prob:GP-1970-23]{GP-1970-23}}
\GSfield{Worth developing}{Worth a bounded reconstruction: an orbit-dimension argument gives the requested elementary deformation proof, while the exact historical attribution remains unverified.}
\end{GSitem}
\begin{GSitem}{1970-198.13}{E}{Uniqueness of decomposition into indecomposables}
\GSfield{Kind; source}{theorem; \GSpage{P000317}}
\GSfield{The notes}{The last parenthesis says that decomposition is unique if the indecomposables are distinct and that the number is always the same.}
\GSfield{Statement}{The finite-dimensional category has the Krull--Schmidt property: the multiset of indecomposable isomorphism classes, with multiplicities, is unique up to permutation. Distinctness of summands is unnecessary. Existence follows by decreasing ambient dimension; uniqueness uses finite-dimensional endomorphism algebras and the usual local-endomorphism argument.}
\GSfield{Fate}{The direct-sum classification in \citet[Theorem~1]{gelfandponomarev1971quadruples}, viewed in the representation category of \citet{bernsteingelfandponomarev1973coxeter}, is a classical Krull--Schmidt setting.}
\GSfield{Present status}{As of 9 November 2026: Established with the hypotheses in the statement.}
\GSfield{Assessment}{literal statement: correct with missing hypotheses; publication match: partial; verification: read the result; development: none; confidence: high}
\end{GSitem}

\GStalk{199}{1970-11-02}{D. A. Kazhdan}{P000317-P000319}{A repaired affine proof, a Picard-group exercise, and logarithmic forms}
The Kazhdan block begins after the four-subspace discussion on P000317. The brief Deligne theorem at the top of P000319 precedes Sinai's explicit heading and is recorded here as an unattributed mathematical report, not assigned a new speaker by guesswork.

\begin{GSitem}{1970-199.1}{E}{Decomposing a Čech cocycle into pieces with a local primitive}
\GSfield{Kind; source}{method; \GSpage{P000317-P000318}}
\GSfield{The notes}{Write \(m_{ij}=\sum_km_{ij}^{(k)}\), with \(m_{ik}^{(k)}\in M_{f_i}\) and the cocycle identity already valid in \(M_{f_if_j}\).}
\GSfield{Statement}{For a finite distinguished-open cover, choose \(N\) and lifts \(u_i^{(k)}\in M_{f_i}\) of \(f_k^N m_{ik}\). Increase the common exponent, multiplying the lifts by the same extra powers of \(f_k\), so that all finitely many localization-equality errors vanish and \(u_i^{(k)}-u_j^{(k)}=f_k^N m_{ij}\) in \(M_{f_if_j}\). Choose \(g_k\in A\) with \(\sum_k g_kf_k^N=1\). Then \(n_i=\sum_k g_ku_i^{(k)}\) satisfies \(m_{ij}=n_i-n_j\), proving degree-one affine acyclicity. The lifts are chosen; localization maps need not be injective.}
\GSfield{Fate}{This corrects the proof attempt of record 196 and agrees with the localization argument of \citet[Lemma~30.2.1]{stacks2026affine}. The first one-index trial on P000317 has a sign slip; the final P000318 formula has the consistent sign.}
\GSfield{Present status}{As of 9 November 2026: Established with the hypotheses in the statement.}
\GSfield{Assessment}{literal statement: correct with missing hypotheses; publication match: exact; verification: read the result; development: annotated; confidence: high}
\end{GSitem}
\begin{GSitem}{1970-199.2}{E}{The sheaf of units is not a quasi-coherent module sheaf}
\GSfield{Kind; source}{definition; \GSpage{P000318}}
\GSfield{The notes}{\(\mathcal O^*(U)\) consists of functions avoiding all primes of \(U\), and is called a subsheaf of the structure sheaf.}
\GSfield{Statement}{\(\mathcal O_X^\times\) is the sheaf of invertible regular functions under multiplication. It is a subsheaf of sets and an abelian sheaf under that law, but not an \(\mathcal O_X\)-submodule: zero is not a unit and units are not closed under addition. Affine acyclicity for quasi-coherent modules therefore does not apply to it.}
\GSfield{Fate}{The Picard-group interpretation used in \citet[section~59.69]{stacks2026picardcurves} provides the relevant cohomology. The different algebraic structures are essential to the next exercise.}
\GSfield{Present status}{As of 9 November 2026: Established with the hypotheses in the statement.}
\GSfield{Assessment}{literal statement: correct with missing hypotheses; publication match: partial; verification: read the result; development: none; confidence: high}
\end{GSitem}
\begin{GSitem}{1970-199.3}{E}{Nonzero first cohomology of units on a punctured elliptic curve}
\GSfield{Kind; source}{problem; \GSpage{P000318}}
\GSfield{The notes}{For \(A=\mathbf C[x,y]/(y^2-x^3+1)\), prove \(H^1(\operatorname{Spec}A,\mathcal O^*)\ne0\).}
\GSfield{Statement}{Let \(E\) be the smooth projective completion of \(y^2=x^3-1\), with point at infinity \(O\), and \(U=E\setminus\{O\}\). Then \(H^1_{\mathrm{Zar}}(U,\mathcal O_U^\times)=\operatorname{Pic}(U)\cong\operatorname{Pic}(E)/\mathbf Z[O]\cong\operatorname{Pic}^0(E)\cong E(\mathbf C)\), which is nonzero.}
\GSfield{Fate}{\citet[proof of Lemma~59.69.3]{stacks2026picardcurves} gives the quotient by removed-point classes. Together with the degree decomposition for a pointed elliptic curve, this solves the notebook exercise. No claim that an arbitrary topological torus has this sheaf cohomology is intended.}
\GSfield{Present status}{As of 9 November 2026: A solved standard exercise; affine quasi-coherent vanishing does not cover the multiplicative unit sheaf.}
\GSfield{Assessment}{literal statement: correct; publication match: exact; verification: read the result; development: none; confidence: high; problem \hyperref[prob:GP-1970-24]{GP-1970-24}}
\end{GSitem}
\begin{GSitem}{1970-199.4}{E}{Cohomology of punctured affine space}
\GSfield{Kind; source}{theorem; \GSpage{P000318}}
\GSfield{The notes}{For \(S=\mathbf A^n\setminus\{0\}\), only degrees zero and \(n-1\) can be nonzero, and the latter is infinite-dimensional.}
\GSfield{Statement}{Over a field \(k\), for \(n\geq2\), \(H^0(S,\mathcal O)=k[x_1,\ldots,x_n]\), \(H^q(S,\mathcal O)=0\) for \(q\notin\{0,n-1\}\), and \(H^{n-1}\) is freely spanned over \(k\) by monomials \(x_1^{-a_1}\cdots x_n^{-a_n}\), all \(a_i\geq1\). For \(n=1\), \(H^0=k[x,x^{-1}]\) instead.}
\GSfield{Fate}{The detailed calculation is carried out in record 201. It is the localization-complex calculation underlying \citet[Lemma~30.8.1]{stacks2026projective}; the affine-cover acyclicity needed for it is \citet[Lemma~30.2.2]{stacks2026affine}.}
\GSfield{Present status}{As of 9 November 2026: Established with the hypotheses in the statement.}
\GSfield{Assessment}{literal statement: correct with missing hypotheses; publication match: exact; verification: read the result; development: none; confidence: high; problem \hyperref[prob:GP-1970-25]{GP-1970-25}}
\end{GSitem}
\begin{GSitem}{1970-199.5}{E}{Deligne's logarithmic de Rham description of an open variety}
\GSfield{Kind; source}{theorem; \GSpage{P000319}; Unattributed report of P. Deligne's theorem}
\GSfield{The notes}{A complete smooth \(X\), a transversal codimension-one boundary \(Y\), and local forms \(df/f\) are described; resolution reduces to this situation and the cohomology ring is said to be independent of the resolution.}
\GSfield{Statement}{For a smooth complex algebraic variety \(U\), choose a smooth proper compactification \(j:U\hookrightarrow X\) whose reduced boundary \(D=X\setminus U\) is a simple normal-crossings divisor. Then \(\Omega_X^\bullet(\log D)\) computes \(H^*(U^{\mathrm{an}},\mathbf C)\) by hypercohomology, compatibly with products. The resulting cohomology and the mixed Hodge structure do not depend on the chosen such compactification. Smoothness, characteristic zero, and normal crossings are necessary parts of this formulation.}
\GSfield{Reconstruction}{Hypercohomology, the smooth open variety, and the precise normal-crossings condition replace the compressed word ring. The mixed-Hodge statement is additional published context, not words attributed to the notebook.}
\GSfield{Fate}{\citet[Proposition~3.1.8 and Theorem~3.2.5]{deligne1971hodge} supplies the comparison and independence. The paper was received 15 September 1970 and printed in 1971. The notebook already reports it as Deligne's theorem; the submission date is not treated as a public-release date.}
\GSfield{Present status}{As of 9 November 2026: Established; the notebook records an already circulating theorem, with the first printing checked here in 1971.}
\GSfield{Assessment}{literal statement: correct with missing hypotheses; publication match: exact; verification: read the result; development: none; confidence: high}
\end{GSitem}

\GStalk{200}{1970-11-02}{Ya. G. Sinai}{P000319}{Powers' factor criterion and asymptotic independence}
The whole Sinai block is on P000319 below the separator. Its principal formula has an existential quantifier where a uniform tail-algebra estimate is needed; the original Powers theorem resolves that compression exactly.

\begin{GSitem}{1970-200.1}{E}{The uniform relative-commutant criterion for a factor state}
\GSfield{Kind; source}{theorem; \GSpage{P000319}}
\GSfield{The notes}{The displayed theorem says \(\forall x\,\forall\varepsilon>0\,\exists y\) satisfying the correlation estimate, with \(y\) in a local commutant.}
\GSfield{Statement}{For a state \(\varphi\) of the spin UHF algebra, \(M_\varphi\) is a factor if and only if for every \(x\in\mathcal A\) and \(\varepsilon>0\) there exists \(n\) such that \[|\varphi(xy)-\varphi(x)\varphi(y)|\leq\varepsilon\|y\|\quad\text{for all }y\in\mathcal A([-n,n])^{\prime}\cap\mathcal A.\] The quantifier is over every element of the tail relative commutant, not one chosen \(y\).}
\GSfield{Reconstruction}{The manuscript quantifier is emended to the uniform relative-commutant condition in Powers; the reconstruction is mathematically necessary because the literal condition is vacuous.}
\GSfield{Fate}{\citet[Theorem~2.5 and proof; Russian translation pp.~100--101]{powers1967representations} gives this criterion, after rescaling \(x\) to incorporate \(\varepsilon\). The original theorem is 1967; the Russian translation available to the seminar is 1969.}
\GSfield{Present status}{As of 9 November 2026: Established with the corrected uniform quantifiers. The literal existential condition is vacuous because \(y=0\) always works.}
\GSfield{Assessment}{literal statement: false as written; publication match: exact; verification: read the result; development: annotated; confidence: high; problem \hyperref[prob:GP-1970-22]{GP-1970-22}}
\end{GSitem}
\begin{GSitem}{1970-200.2}{E}{The classical tail-independence form of the criterion}
\GSfield{Kind; source}{construction; \GSpage{P000319}}
\GSfield{The notes}{For measure-type states the estimate becomes \(\sup_B|\mu(A\cap B)-\mu(A)\mu(B)|\to0\).}
\GSfield{Statement}{For the product-diagonal states of 1970-195.3, the equivalent classical condition is: for every fixed cylinder event \(A\), \[\sup_{B\in\sigma(\omega_j:|j|>n)}|\mu(A\cap B)-\mu(A)\mu(B)|\longrightarrow0.\] Uniformity in the distant event is required. For a two-sided lattice, outside \([-n,n]\) replaces a one-sided first-\(n\)-coordinates description.}
\GSfield{Reconstruction}{The two-sided tail sigma-field and the fixed-cylinder quantifier are made explicit; the notebook's measure terminology is not used to bypass its diagonalization error.}
\GSfield{Fate}{This follows by restricting the correlation criterion of \citet[Theorem~2.5]{powers1967representations} to diagonal projections and extending by simple-function approximation. It applies to the genuinely product-diagonal construction, not automatically to the weaker commuting-density definition of record 197.}
\GSfield{Present status}{As of 9 November 2026: Established with the hypotheses in the statement.}
\GSfield{Assessment}{literal statement: correct with missing hypotheses; publication match: partial; verification: read the result; development: none; confidence: high; problem \hyperref[prob:GP-1970-22]{GP-1970-22}}
\end{GSitem}
\begin{GSitem}{1970-200.3}{E}{Product states satisfy the factor criterion}
\GSfield{Kind; source}{theorem; \GSpage{P000319}}
\GSfield{The notes}{Product measures with site probabilities bounded below by a fixed \(c>0\) are given as examples.}
\GSfield{Statement}{Every infinite product state of the tensor-product matrix algebra is a factor state. The lower bound on site probabilities in the notebook is sufficient but unnecessary; pure site states are allowed as well. For an observable supported in \([-n,n]\), correlation with the outside algebra factorizes exactly, and norm approximation handles quasi-local observables.}
\GSfield{Fate}{\citet[Lemma~4.1, Russian translation p.~113]{powers1967representations} proves the product-state assertion. Section~4 distinguishes type I, tracial type \(\mathrm{II}_1\), and nontracial type III examples.}
\GSfield{Present status}{As of 9 November 2026: Established with the hypotheses in the statement.}
\GSfield{Assessment}{literal statement: correct; publication match: exact; verification: read the result; development: none; confidence: high; problem \hyperref[prob:GP-1970-22]{GP-1970-22}}
\end{GSitem}
\begin{GSitem}{1970-200.4}{E}{The finite-temperature one-dimensional Gibbs state is primary}
\GSfield{Kind; source}{theorem; \GSpage{P000319}}
\GSfield{The notes}{The last line says that the state obtained from the Hamiltonian is a factor.}
\GSfield{Statement}{For the self-adjoint finite-range one-dimensional interaction and finite-temperature Gibbs limit of 1970-195.4--5, the GNS von Neumann algebra is a factor. This conclusion does not assert that all states of the algebra, or all infinite-volume Gibbs states in higher dimensions, are factors.}
\GSfield{Fate}{The primary-state conclusion is explicitly in the abstract of \citet{araki1969gibbs}. The notebook is reporting a known result within the stated spin-chain setup.}
\GSfield{Present status}{As of 9 November 2026: Established with the hypotheses in the statement.}
\GSfield{Assessment}{literal statement: correct with missing hypotheses; publication match: exact; verification: publisher metadata checked; development: none; confidence: high; problem \hyperref[prob:GP-1970-22]{GP-1970-22}}
\end{GSitem}

\GStalk{201}{1970-11-16}{D. A. Kazhdan}{P000320-P000323}{Punctured affine space, Koszul exactness, and residue duality}
The manuscript's 16/XI date supplies the ISO date missing from the catalog's underscore-spelled entry. The computation is valid, but the final module-Hom identification is not: the residue pairing produces a restricted linear dual, not \(\operatorname{Hom}_R(R,k)\).

\begin{GSitem}{1970-201.1}{E}{The principal-open cover of punctured affine space}
\GSfield{Kind; source}{construction; \GSpage{P000320}}
\GSfield{The notes}{\(X=\operatorname{Spec}A[x_1,\ldots,x_n]\setminus\{0\}\) is covered by the \(X_i\) where \(x_i\ne0\).}
\GSfield{Statement}{Put \(R=A[x_1,\ldots,x_n]\), \(I=(x_1,\ldots,x_n)\), and \(U=\operatorname{Spec}R\setminus V(I)\). Over a general base ring, \(V(I)\cong\operatorname{Spec}A\) is the zero section, not a single point. The cover is \(D(x_i)\), with sections \(R_{x_i}\); every intersection inverts all corresponding variables.}
\GSfield{Reconstruction}{The zero-section notation makes explicit the set of containing prime ideals already specified in the source on P320.}
\GSfield{Fate}{The localization cover and its cohomology computation are the same monomial complex used in \citet[Lemma~30.8.1]{stacks2026projective}, justified by \citet[Lemma~30.2.2]{stacks2026affine}.}
\GSfield{Present status}{As of 9 November 2026: Established with the hypotheses in the statement.}
\GSfield{Assessment}{literal statement: correct; publication match: exact; verification: read the result; development: none; confidence: high}
\end{GSitem}
\begin{GSitem}{1970-201.2}{E}{Global functions extend across the zero section in dimension at least two}
\GSfield{Kind; source}{theorem; \GSpage{P000320}}
\GSfield{The notes}{The notes identify \(H^0\) with the polynomial ring and explicitly exclude \(n=1\) in the preliminary calculation.}
\GSfield{Statement}{For every commutative ring \(A\) and \(n\geq2\), \(\Gamma(U,\mathcal O_U)=A[x_1,\ldots,x_n]\). For \(n=1\), \(U=D(x_1)\) and the answer is \(A[x_1,x_1^{-1}]\). The polynomial-extension statement follows by intersecting the relevant Laurent subrings, or by monomial Čech exactness.}
\GSfield{Fate}{The affine-cover proof is classical; the graded version is computed in \citet[proof of Lemma~30.8.1]{stacks2026projective}. This continues the earlier punctured-affine examples listed as 1965-58.7 and affine-cover methods 1965-70.6.}
\GSfield{Present status}{As of 9 November 2026: Established with the hypotheses in the statement.}
\GSfield{Assessment}{literal statement: correct; publication match: exact; verification: read the result; development: none; confidence: high; problem \hyperref[prob:GP-1970-25]{GP-1970-25}}
\end{GSitem}
\begin{GSitem}{1970-201.3}{E}{The three-variable one-cocycle calculation is a Koszul syzygy}
\GSfield{Kind; source}{method; \GSpage{P000321}}
\GSfield{The notes}{The \(n=3\) calculation clears denominators and separates polynomial multiples of powers of the coordinates.}
\GSfield{Statement}{The sequence \(x_1^N,\ldots,x_n^N\) is regular in \(A[x_1,\ldots,x_n]\) for every \(N\geq1\), even when \(A\) has zero divisors. Its Koszul complex is exact away from degree zero. Thus a relation among the powered coordinates is generated by the elementary skew relations; after localization this gives the required Čech primitive. In three variables it proves \(H^1(U,\mathcal O)=0\).}
\GSfield{Reconstruction}{The regular-sequence statement is supplied as a precise modern formulation of the polynomial-divisibility argument.}
\GSfield{Fate}{The notebook announces the Koszul complex as the next language for this calculation. The monomial proof in \citet[proof of Lemma~30.8.1]{stacks2026projective} verifies the same vanishing without depending on the inconsistent intermediate polynomial subscripts.}
\GSfield{Present status}{As of 9 November 2026: Established with the hypotheses in the statement.}
\GSfield{Assessment}{literal statement: correct with missing hypotheses; publication match: partial; verification: read the result; development: none; confidence: high}
\end{GSitem}
\begin{GSitem}{1970-201.4}{E}{Vanishing in all intermediate degrees}
\GSfield{Kind; source}{theorem; \GSpage{P000321-P000322}}
\GSfield{The notes}{The argument for \(H^1\) is extended to \(H^q\), except the last possible degree.}
\GSfield{Statement}{For \(n\geq2\), \(H^q(U,\mathcal O_U)=0\) when \(1\leq q\leq n-2\); it is also zero for \(q\geq n\) because the cover has \(n\) members. Splitting the Čech complex by Laurent monomials gives a simplex complex unless every coordinate has negative exponent, which contributes only in degree \(n-1\).}
\GSfield{Fate}{This is the ungraded version of \citet[proof of Lemma~30.8.1]{stacks2026projective}. The complete list of nonzero degrees agrees with the preliminary announcement in record 199.}
\GSfield{Present status}{As of 9 November 2026: Established with the hypotheses in the statement.}
\GSfield{Assessment}{literal statement: correct; publication match: exact; verification: read the result; development: none; confidence: high; problem \hyperref[prob:GP-1970-25]{GP-1970-25}}
\end{GSitem}
\begin{GSitem}{1970-201.5}{E}{The top cohomology as negative Laurent monomials}
\GSfield{Kind; source}{theorem; \GSpage{P000322-P000323}}
\GSfield{The notes}{Top cohomology consists of negative powers of every variable; multiplication by polynomials is followed by throwing away terms with nonnegative powers.}
\GSfield{Statement}{For \(n\geq2\), \[H^{n-1}(U,\mathcal O_U)\cong R_{x_1\cdots x_n}\big/\sum_iR_{x_1\cdots\widehat{x_i}\cdots x_n}\cong\bigoplus_{a_1,\ldots,a_n\geq1}A\,x_1^{-a_1}\cdots x_n^{-a_n}.\] The quotient, not ordinary Laurent multiplication alone, determines its \(R\)-module action.}
\GSfield{Fate}{\citet[Lemma~30.8.1 and its monomial proof]{stacks2026projective} gives the homogeneous pieces of this module. Over a nonzero field it is infinite-dimensional. The case \(n=1\) belongs to degree-zero global sections and must be kept separate.}
\GSfield{Present status}{As of 9 November 2026: Established with the hypotheses in the statement.}
\GSfield{Assessment}{literal statement: correct; publication match: exact; verification: read the result; development: none; confidence: high; problem \hyperref[prob:GP-1970-25]{GP-1970-25}}
\end{GSitem}
\begin{GSitem}{1970-201.6}{E}{Residue pairing with the restricted linear dual}
\GSfield{Kind; source}{construction; \GSpage{P000323}}
\GSfield{The notes}{\(H^{n-1}=\operatorname{Hom}_{R}(R,k)\) is written twice, with \(k\) the evaluation-at-zero module; the proposed pairing takes the coefficient of \((x_1\cdots x_n)^{-1}\).}
\GSfield{Statement}{For \(n\geq2\), over a field \(k\), the pairing \(\langle f,\eta\rangle=\operatorname{coeff}_{x_1^{-1}\cdots x_n^{-1}}(f\eta)\) is nondegenerate and identifies top cohomology with the finite-support, or graded restricted, \(k\)-linear dual \(R^\circ=\bigoplus_{d\geq0}\operatorname{Hom}_k(R_d,k)\), where \(\operatorname{Hom}_k(R_d,k)\) is placed in degree \(-d-n\). Its \(R\)-action is \((r\ell)(f)=\ell(rf)\). It is neither \(\operatorname{Hom}_R(R,k)\cong k\) nor the full algebraic dual \(\prod_d\operatorname{Hom}_k(R_d,k)\).}
\GSfield{Reconstruction}{The incorrect ring-linear Hom is replaced by the restricted linear dual detected directly by the displayed coefficient-extraction pairing; \(n\geq2\) is explicit.}
\GSfield{Fate}{The monomial bases in \citet[Lemma~30.8.1]{stacks2026projective} verify the pairing term by term. The source's reference to Serre duality is a thematic link: this affine local-duality calculation needs the restricted dual and the differential/degree convention stated here.}
\GSfield{Present status}{As of 9 November 2026: The corrected residue-dual identification is established. The displayed \(R\)-linear Hom has the wrong dimension and action.}
\GSfield{Assessment}{literal statement: false as written; publication match: partial; verification: read the result; development: annotated; confidence: high; problem \hyperref[prob:GP-1970-25]{GP-1970-25}}
\end{GSitem}
\begin{GSitem}{1970-201.7}{E}{General base-ring version and the limits of the duality slogan}
\GSfield{Kind; source}{remark; \GSpage{P000323}}
\GSfield{The notes}{The theorem is restated with coefficient ring \(A\), writing \(H^{n-1}=\operatorname{Hom}_{H^0}(H^0,A)\).}
\GSfield{Statement}{The Laurent-quotient and direct-sum description in 1970-201.5 works for every commutative \(A\). Each homogeneous polynomial piece \(R_d\) is finite free, so coefficient extraction gives the corresponding restricted \(A\)-linear dual. The general-base statement does not turn this into an \(R\)-linear Hom to the evaluation-at-zero module.}
\GSfield{Reconstruction}{The same correction is made over the coefficient ring, using finite free homogeneous polynomial pieces instead of a field-only argument.}
\GSfield{Fate}{The arbitrary-ring computation is explicitly allowed in \citet[Lemma~30.8.1]{stacks2026projective}. No field-only dimension argument is used to justify extension of the Laurent formula to general \(A\).}
\GSfield{Present status}{As of 9 November 2026: Established with the hypotheses in the statement.}
\GSfield{Assessment}{literal statement: false as written; publication match: partial; verification: read the result; development: annotated; confidence: high}
\end{GSitem}
\begin{GSitem}{1970-201.8}{F}{Commuting operators and the announced Koszul complex}
\GSfield{Kind; source}{remark; \GSpage{P000322-P000323}}
\GSfield{The notes}{The notes briefly connect the calculation with several commuting operators and promise the Koszul complex next time.}
\GSfield{Statement}{Commuting operators give a module over a polynomial ring, and their Koszul complex is a natural multivariable replacement for a single operator complex. The manuscript gives no precise spectral set, resolvent theorem, or functional calculus to identify beyond that construction.}
\GSfield{Fate}{The polynomial-module/Koszul connection is precise, but no exact additional operator theorem was verified from the abbreviated remark. The item is retained as a methodological indication rather than expanded into an unrecorded multivariable spectral theorem.}
\GSfield{Present status}{As of 9 November 2026: The exact intended operator-theoretic statement remains unspecified; no open problem is inferred.}
\GSfield{Assessment}{literal statement: not determinable; publication match: none found; verification: unresolved lead; development: none; confidence: medium}
\end{GSitem}

\GStalk{202}{1970-11-23}{I. M. Gelfand}{P000324}{Plans for the next semester}
The plan names Arnold and Gabrielov on Griffiths and integrals of rational differential forms over cycles; S. Gelfand on Brauer theory; Sinai on factors; Knorr on a new proof of a theorem of Grauert; and new work of Atiyah on skew Fredholm operators, alongside Gohberg, Krupnik, and Gorin. These are proposed subjects, with no precise new theorem asserted. The catalog dates this record 23 November, whereas the clearly visible manuscript header reads 22/XI. The isolated Arnold heading on P000325, about vector fields in general position, supplies no mathematical statement and is preserved here as an uncatalogued heading between Gorin and Kazhdan.


\GStalk{203}{1970-11-23}{E. A. Gorin}{P000324-P000325}{Symbol algebras, the Arens–Royden index, and Toeplitz operators}
The Gorin block begins at the boxed name in the lower part of P000324 and ends before the isolated Arnold heading on P000325. It passes from a quotient by characters to a group-valued index, then gives the completed-curve Fredholm criterion for piecewise continuous Toeplitz symbols. These are reports on established Banach-algebra and Gohberg--Krupnik results; the manuscript does not identify a separate new Gorin theorem.

\begin{GSitem}{1970-203.1}{E}{The character quotient as a commutative symbol algebra}
\GSfield{Kind; source}{construction; \GSpage{P000324}}
\GSfield{The notes}{For a noncommutative Banach algebra \(\alpha\), set \(J=\bigcap_\varphi\ker\varphi\), with \(\varphi:\alpha\to\mathbf C\), and call \(A=\alpha/J\) the symbol algebra.}
\GSfield{Statement}{Let \(\alpha\) be a unital complex Banach algebra and take all nonzero complex algebra characters. Their kernels are closed, contain every commutator, and have intersection \(J\). If there are characters, \(A=\alpha/J\) is a unital commutative semisimple Banach algebra; if there are none, \(J=\alpha\). For noncommutative \(\alpha\), this character kernel ideal need not be its Jacobson radical.}
\GSfield{Reconstruction}{Unitality and the complex scalar field are specified; the character-kernel quotient is distinguished from a general Jacobson radical.}
\GSfield{Fate}{This is the classical character-quotient construction underlying the application of \citet{arens1963invertible}. The related operator symbol algebra is described concretely in \citet[Theorem~6.1]{gohbergkrupnik1970algebra}; no claim of an exact publication of this introductory formulation by Gorin is made.}
\GSfield{Present status}{As of 9 November 2026: Established with the hypotheses in the statement.}
\GSfield{Assessment}{literal statement: correct with missing hypotheses; publication match: partial; verification: read the result; development: none; confidence: high}
\end{GSitem}
\begin{GSitem}{1970-203.2}{E}{The algebra of all bounded Hilbert-space operators has no scalar characters}
\GSfield{Kind; source}{remark; \GSpage{P000324}}
\GSfield{The notes}{The example says that for the algebra of all bounded operators, \(J=\alpha\): there are no symbols.}
\GSfield{Statement}{For a complex Hilbert space \(H\) with \(\dim H\geq2\), \(B(H)\) has no nonzero scalar character, so its character quotient is zero. The dimension-one exception is \(B(\mathbf C)=\mathbf C\). An arbitrary Banach space is not silently included in this Hilbert-space assertion.}
\GSfield{Reconstruction}{The Hilbert-space interpretation and dimension restriction are supplied; the source does not specify the underlying space.}
\GSfield{Fate}{The claim is an elementary obstruction to this particular symbol construction. In finite dimension matrix units exclude characters; in infinite Hilbert dimension two orthogonal isometries do so. The Fredholm quotient by compact operators is a different construction, and is the operator-algebra context of \citet{gohbergkrupnik1970algebra}.}
\GSfield{Present status}{As of 9 November 2026: Established with the hypotheses in the statement.}
\GSfield{Assessment}{literal statement: correct with missing hypotheses; publication match: thematic only; verification: read the result; development: none; confidence: high}
\end{GSitem}
\begin{GSitem}{1970-203.3}{E}{The group-valued Arens–Royden index}
\GSfield{Kind; source}{theorem; \GSpage{P000324}}
\GSfield{The notes}{\(\mathcal E(A)\) is the invertible group, \(\mathcal E_1(A)\) its identity component, and \(\operatorname{ind}T\in\mathcal E(A)/\mathcal E_1(A)=H^1(\mathfrak M_A,\mathbf Z)\).}
\GSfield{Statement}{For a unital commutative complex Banach algebra \(A\), the Gelfand transform induces \(A^\times/\exp(A)\cong\check H^1(\mathfrak M_A;\mathbf Z)\), and \(\exp(A)\) is the identity component of \(A^\times\). If \(\psi:\alpha\to A\) and \(\psi(T)\) is invertible, its component class is a group-valued index. It need not be a numerical Fredholm index unless the operator extension and its connecting map are specified.}
\GSfield{Reconstruction}{Čech cohomology, the exponential identity component, and the distinction from a numerical Fredholm index are made explicit.}
\GSfield{Fate}{The result is the Arens--Royden theorem, predating the seminar: \citet{arens1963invertible}, and the explicit modern formulation in \citet[introduction, p.~2]{luptonphillipsschochetsmith2007banach}. The notebook itself correctly stresses that the index need not be a number.}
\GSfield{Present status}{As of 9 November 2026: Established; Čech cohomology is the natural theory for a general compact maximal ideal space.}
\GSfield{Assessment}{literal statement: correct with missing hypotheses; publication match: exact; verification: read the result; development: none; confidence: high}
\end{GSitem}
\begin{GSitem}{1970-203.4}{E}{Hardy-space Toeplitz operators with piecewise continuous symbols}
\GSfield{Kind; source}{construction; \GSpage{P000325}}
\GSfield{The notes}{\(T_af=(af)_+\) acts on \(H_2(\Delta)\), the holomorphic functions with square-summable Taylor coefficients; \(\alpha\) is generated by these operators.}
\GSfield{Statement}{Identify \(H^2(\mathbb D)\) with the nonnegative-Fourier subspace of \(L^2(S^1)\), and let \(P_+\) be the orthogonal projection. For \(a\in L^\infty(S^1)\), \(T_a=P_+M_a|_{H^2}\) is bounded with \(\|T_a\|\leq\|a\|_\infty\). The example restricts \(a\) to scalar piecewise continuous symbols and considers the norm-closed algebra generated by them; the set of individual \(T_a\) is not itself generally closed under multiplication.}
\GSfield{Fate}{This is the Toeplitz matrix setup of \citet[p.~91]{gohberg1967toeplitz}. The related quotient-symbol construction is developed in \citet{gohbergkrupnik1970algebra}.}
\GSfield{Present status}{As of 9 November 2026: Established with the hypotheses in the statement.}
\GSfield{Assessment}{literal statement: correct with missing hypotheses; publication match: exact; verification: read the result; development: none; confidence: high}
\end{GSitem}
\begin{GSitem}{1970-203.5}{E}{The completed-symbol curve and the Fredholm index}
\GSfield{Kind; source}{theorem; \GSpage{P000325}}
\GSfield{The notes}{At every discontinuity join the two limiting symbol values by a segment; \(T_a\) is Fredholm exactly when the completed curve avoids zero, and its index is the number of turns.}
\GSfield{Statement}{For a scalar piecewise continuous symbol on the positively oriented circle with finitely many jumps, form the oriented closed curve \(a_2\) by inserting a straight segment from each left limit to its right limit. On \(H^2\), \(T_a\) is Fredholm iff \(0\notin a_2\). With \(\operatorname{ind}T_a=\dim\ker T_a-\dim\operatorname{coker}T_a\), \(\operatorname{ind}T_a=-\operatorname{wind}(a_2,0)\). The straight-segment completion is specific to this Hilbert-space exponent.}
\GSfield{Reconstruction}{The circle orientation, Fredholm-index convention, Hilbert exponent, and finite-jump version read in the original theorem are specified.}
\GSfield{Fate}{Exact match with \citet[Theorem~1, p.~91]{gohberg1967toeplitz}, received 5 November 1966 and printed 1967. Its separate kernel and cokernel formulas fix the sign. The later singular-integral algebra treatment is \citet[Theorem~7.1]{gohbergkrupnik1970algebra}, received 5 June 1969 and printed 1970.}
\GSfield{Present status}{As of 9 November 2026: Established before the lecture. The notebook leaves the winding orientation and Fredholm-index sign convention unstated.}
\GSfield{Assessment}{literal statement: correct with missing hypotheses; publication match: exact; verification: read the result; development: annotated; confidence: high}
\end{GSitem}

\GStalk{204}{1970-11-23}{D. A. Kazhdan}{P000325-P000328}{Projective space, Serre twists, and coherent cohomology}
Kazhdan constructs projective space and its twisting sheaves, states global generation, cohomological finiteness and vanishing, and proves polynomiality of Euler characteristics. The later proofs require a Noetherian base for coherence and finite generation; the Euler characteristic is then specialized to a field. The hyperplane argument needs a section that is a non-zero-divisor on the sheaf.

\begin{GSitem}{1970-204.1}{E}{Projective space as the scalar quotient of the punctured affine space}
\GSfield{Kind; source}{construction; \GSpage{P000325}}
\GSfield{The notes}{\(A^*\) acts on \(A^{n+1}\setminus\{0\}\), and \(A^*=\operatorname{Spec}A[x,y]/(xy-1)\).}
\GSfield{Statement}{Over a ring \(A\), the multiplicative group scheme is \(\mathbf G_{m,A}=\operatorname{Spec}A[t,t^{-1}]\). It acts by simultaneous scalar multiplication on \(U=\mathbf A_A^{n+1}\setminus V(x_0,\ldots,x_n)\); the scheme quotient is \(\mathbf P_A^n\), with \(U\to\mathbf P_A^n\) the associated \(\mathbf G_m\)-torsor. A quotient of sets of \(A\)-valued points is not generally a substitute for this scheme quotient.}
\GSfield{Reconstruction}{The scalar quotient is interpreted as a scheme quotient and torsor, rather than a quotient of the set of points over a general ring.}
\GSfield{Fate}{This is the standard projective-space construction, whose functorial line-bundle interpretation is \citet[Lemma~27.12.3]{stacks2026projconstruction}. It is expository material in the lecture.}
\GSfield{Present status}{As of 9 November 2026: Established with the hypotheses in the statement.}
\GSfield{Assessment}{literal statement: correct with missing hypotheses; publication match: partial; verification: read the result; development: none; confidence: high}
\end{GSitem}
\begin{GSitem}{1970-204.2}{E}{Homogeneous primes and the degree-zero localization sheaf}
\GSfield{Kind; source}{definition; \GSpage{P000325}}
\GSfield{The notes}{Projective points are called homogeneous prime ideals not contained in \((x_0,\ldots,x_n)\); sections on \(P_f\) are \((A[x_0,\ldots,x_n]_f)_0\).}
\GSfield{Statement}{Let \(S=A[x_0,\ldots,x_n]\) be standard graded. Then \(\operatorname{Proj}S\) consists of homogeneous prime ideals \(\mathfrak p\) such that \(S_+\not\subseteq\mathfrak p\). Its basic open \(D_+(f)\), for homogeneous positive-degree \(f\), has coordinate ring \((S_f)_0\). The irrelevant ideal must not be contained in the prime; the source reverses this containment.}
\GSfield{Fate}{The construction and structure sheaf are given in \citet[section~27.8]{stacks2026projconstruction}. The formula for the sections is correct after fixing the prime-ideal condition.}
\GSfield{Present status}{As of 9 November 2026: Established with the hypotheses in the statement.}
\GSfield{Assessment}{literal statement: false as written; publication match: exact; verification: read the result; development: annotated; confidence: high}
\end{GSitem}
\begin{GSitem}{1970-204.3}{E}{Affine charts on projective space}
\GSfield{Kind; source}{theorem; \GSpage{P000326}}
\GSfield{The notes}{\(P_{x_i}\cong A^n\), denoted \(U_i\).}
\GSfield{Statement}{The chart \(D_+(x_i)\) is \(\operatorname{Spec}A[x_0/x_i,\ldots,\widehat{x_i/x_i},\ldots,x_n/x_i]\cong\mathbf A_A^n\). These \(n+1\) charts cover \(\mathbf P_A^n\); on overlaps the coordinate ratios give the gluing maps.}
\GSfield{Fate}{This follows directly from degree-zero localization in \citet[section~27.8]{stacks2026projconstruction}; no hypothesis on \(A\) beyond commutativity is needed.}
\GSfield{Present status}{As of 9 November 2026: Established with the hypotheses in the statement.}
\GSfield{Assessment}{literal statement: correct; publication match: exact; verification: read the result; development: none; confidence: high}
\end{GSitem}
\begin{GSitem}{1970-204.4}{E}{Twisting sheaves and the Hopf-bundle convention}
\GSfield{Kind; source}{definition; \GSpage{P000326}}
\GSfield{The notes}{\(\Gamma(P_f,\mathcal O(k))=(S_f)_k\), \(\mathcal O(1)\) is called the Hopf sheaf, and \(\mathcal F(k)=\mathcal F\otimes\mathcal O(k)\).}
\GSfield{Statement}{The invertible sheaf \(\mathcal O(k)=\widetilde{S(k)}\) has the specified degree-\(k\) sections on basic opens. On complex projective space, the tautological line bundle is \(\mathcal O(-1)\); \(\mathcal O(1)\) is its dual. Since Hopf bundle conventions differ, the degree convention rather than that name fixes the meaning. Tensoring by \(\mathcal O(k)\) is exact.}
\GSfield{Fate}{The standard graded construction is compatible with \citet[Lemma~27.12.3]{stacks2026projconstruction} and the cohomology convention of \citet[Lemma~30.8.1]{stacks2026projective}.}
\GSfield{Present status}{As of 9 November 2026: Established with the hypotheses in the statement.}
\GSfield{Assessment}{literal statement: correct with missing hypotheses; publication match: partial; verification: read the result; development: none; confidence: high}
\end{GSitem}
\begin{GSitem}{1970-204.5}{E}{Coherence and finite presentations}
\GSfield{Kind; source}{definition; \GSpage{P000326}}
\GSfield{The notes}{A coherent sheaf is described locally by \(\mathcal O^J\to\mathcal O^I\to\mathcal F\to0\), with \(I,J\) finite.}
\GSfield{Statement}{On a locally Noetherian scheme, coherence is equivalent to local finite presentation, and kernels of maps between coherent sheaves are coherent. For the following Serre finiteness argument assume \(A\) Noetherian. Over arbitrary noncoherent rings, merely writing a finite presentation is not the general definition of a coherent module.}
\GSfield{Reconstruction}{The missing Noetherian base hypothesis is restored for the finite-presentation and coherent-kernel argument.}
\GSfield{Fate}{The Noetherian hypothesis is explicit in \citet[Lemma~30.14.1]{stacks2026projective}. The classical formulation in \citet{serre1955coherent} works in a Noetherian algebraic setting.}
\GSfield{Present status}{As of 9 November 2026: Established with the hypotheses in the statement.}
\GSfield{Assessment}{literal statement: correct with missing hypotheses; publication match: exact; verification: read the result; development: annotated; confidence: high}
\end{GSitem}
\begin{GSitem}{1970-204.6}{E}{Finiteness and cohomological dimension of coherent sheaves}
\GSfield{Kind; source}{theorem; \GSpage{P000326-P000327}}
\GSfield{The notes}{All \(H^i(P,\mathcal F)\) are finite-type \(A\)-modules, and they vanish for \(i>n\).}
\GSfield{Statement}{If \(A\) is Noetherian and \(\mathcal F\) coherent on \(\mathbf P_A^n\), every \(H^i(\mathbf P_A^n,\mathcal F)\) is a finite \(A\)-module. The vanishing for \(i>n\) follows from the standard affine cover and also holds for arbitrary quasi-coherent sheaves over an arbitrary base ring.}
\GSfield{Fate}{These are the projective finiteness and dimension statements in \citet{serre1955coherent}; their precise ring hypotheses and descending-induction proof are \citet[Lemma~30.14.1]{stacks2026projective}.}
\GSfield{Present status}{As of 9 November 2026: Established with the hypotheses in the statement.}
\GSfield{Assessment}{literal statement: correct with missing hypotheses; publication match: exact; verification: read the result; development: none; confidence: high}
\end{GSitem}
\begin{GSitem}{1970-204.7}{E}{Serre vanishing for large positive twists}
\GSfield{Kind; source}{theorem; \GSpage{P000326-P000328}}
\GSfield{The notes}{For sufficiently large \(k\), \(H^i(P,\mathcal F(k))=0\) for \(i\ne0\).}
\GSfield{Statement}{For \(A\) Noetherian and \(\mathcal F\) coherent on \(\mathbf P_A^n\), there is \(k_0\) such that all \(H^i(\mathbf P_A^n,\mathcal F(k))\) vanish for \(i>0\) and \(k\geq k_0\). The bound can be chosen simultaneously for the finitely many relevant cohomological degrees.}
\GSfield{Fate}{Classical Serre vanishing, \citet{serre1955coherent}; \citet[Lemma~30.14.1]{stacks2026projective} verifies this exact formulation. It was established long before the lecture.}
\GSfield{Present status}{As of 9 November 2026: Established with the hypotheses in the statement.}
\GSfield{Assessment}{literal statement: correct with missing hypotheses; publication match: exact; verification: read the result; development: none; confidence: high}
\end{GSitem}
\begin{GSitem}{1970-204.8}{E}{Global generation after twisting}
\GSfield{Kind; source}{theorem; \GSpage{P000326-P000327}}
\GSfield{The notes}{The main lemma says that \(\mathcal F(k)\) is generated by global sections for every sufficiently large \(k\), hence \(\mathcal O^N\twoheadrightarrow\mathcal F(k)\).}
\GSfield{Statement}{In the same Noetherian coherent setting, \(\mathcal F(k)\) is generated by finitely many global sections for all \(k\gg0\). Equivalently, for some \(m,N\) there is a surjection \(\mathcal O(-m)^N\twoheadrightarrow\mathcal F\); larger twists preserve global generation because \(\mathcal O(1)\) is globally generated.}
\GSfield{Fate}{The finite sum of twists surjecting onto \(\mathcal F\) is \citet[Lemma~30.14.1(1)]{stacks2026projective}. This is Serre's standard projective construction, not a new seminar assertion.}
\GSfield{Present status}{As of 9 November 2026: Established with the hypotheses in the statement.}
\GSfield{Assessment}{literal statement: correct with missing hypotheses; publication match: exact; verification: read the result; development: none; confidence: high}
\end{GSitem}
\begin{GSitem}{1970-204.9}{E}{Clearing a section’s poles along a coordinate hyperplane}
\GSfield{Kind; source}{method; \GSpage{P000326-P000327}}
\GSfield{The notes}{A section \(\alpha\) over \(P_{x_0}\) is said to extend after changing the twist; multiplication by \(x_0\) clears its pole.}
\GSfield{Statement}{Given a local section of a quasi-coherent finite-type sheaf over \(D_+(x_0)\), some power of the coordinate section \(x_0\) turns it into the restriction of a global section of an appropriate positive twist. Applied to finitely many local generators on the finitely many standard charts, a common power yields the global-generation theorem.}
\GSfield{Fate}{The extension-and-clearing step is the standard projective argument used in \citet{serre1955coherent} and in the proof underlying \citet[Lemma~30.14.1(1)]{stacks2026projective}. The word poles is a useful picture, but for a general coherent sheaf the rigorous operation is multiplication and localization.}
\GSfield{Present status}{As of 9 November 2026: Established with the hypotheses in the statement.}
\GSfield{Assessment}{literal statement: correct with missing hypotheses; publication match: partial; verification: read the result; development: none; confidence: high}
\end{GSitem}
\begin{GSitem}{1970-204.10}{E}{Cohomology of the twists from the graded punctured-affine complex}
\GSfield{Kind; source}{theorem; \GSpage{P000327}}
\GSfield{The notes}{\(\bigoplus_{k\in\mathbf Z}H^i(\mathcal O(k))\) is identified with \(H^i(\mathbf A^{n+1}\setminus\{0\},\mathcal O)\).}
\GSfield{Statement}{For \(n\geq1\) and every ring \(A\), \(H^0(\mathbf P_A^n,\mathcal O(k))=S_k\) for \(k\geq0\), while \(H^n\) is freely generated over \(A\) by monomials \(x_0^{-a_0}\cdots x_n^{-a_n}\) with every \(a_i\geq1\) and \(-\sum a_i=k\). All other groups vanish. These are precisely the degree-\(k\) pieces of the punctured-affine Čech complex from record 201. For \(n=0\), \(\mathcal O(k)\) is trivial and has \(H^0=A\) for all \(k\).}
\GSfield{Fate}{The arbitrary-ring monomial computation is \citet[Lemma~30.8.1]{stacks2026projective}; the notebook correctly clarifies that its sum sign means direct sum.}
\GSfield{Present status}{As of 9 November 2026: Established with the hypotheses in the statement.}
\GSfield{Assessment}{literal statement: correct with missing hypotheses; publication match: exact; verification: read the result; development: none; confidence: high}
\end{GSitem}
\begin{GSitem}{1970-204.11}{E}{Descending induction through a coherent kernel}
\GSfield{Kind; source}{method; \GSpage{P000327-P000328}}
\GSfield{The notes}{Use \(0\to\mathcal G\to\mathcal O(-m)^N\to\mathcal F\to0\), and induct downwards in cohomological degree for finiteness and vanishing together.}
\GSfield{Statement}{The kernel \(\mathcal G\) is coherent over a Noetherian base. The long exact sequence bounds \(H^i(\mathcal F)\) between a quotient of the known twist cohomology and a submodule of \(H^{i+1}(\mathcal G)\). Noetherianity ensures that the latter submodule is finite. After a sufficiently large common twist, both outer terms vanish for \(i>0\).}
\GSfield{Fate}{The proof is the one in \citet[proof of Lemma~30.14.1]{stacks2026projective}, and is already in the classical Serre theory. The twist variable must remain the same in all three terms on P000328.}
\GSfield{Present status}{As of 9 November 2026: Established with the hypotheses in the statement.}
\GSfield{Assessment}{literal statement: correct with missing hypotheses; publication match: exact; verification: read the result; development: none; confidence: high}
\end{GSitem}
\begin{GSitem}{1970-204.12}{E}{The Euler characteristic is the Hilbert polynomial}
\GSfield{Kind; source}{theorem; \GSpage{P000326; P000328}}
\GSfield{The notes}{Over a field, \(\chi(\mathcal F(k))=\sum_i(-1)^i\dim H^i(\mathcal F(k))\) is polynomial in \(k\); for large \(k\), it equals \(\dim H^0\).}
\GSfield{Statement}{For a coherent sheaf \(\mathcal F\) on \(\mathbf P_K^n\) over a field, the integer-valued function \(k\mapsto\chi(\mathcal F(k))\), for all integers \(k\), is given by a rational polynomial of degree at most \(\dim\operatorname{Supp}\mathcal F\). For \(k\gg0\), higher cohomology vanishes and this polynomial equals \(\dim_K H^0(\mathcal F(k))\).}
\GSfield{Fate}{This is the classical Euler/Hilbert polynomial consequence of \citet{serre1955coherent}. The source explicitly restricts to a field before defining the dimension-valued Euler characteristic.}
\GSfield{Present status}{As of 9 November 2026: Established with the hypotheses in the statement.}
\GSfield{Assessment}{literal statement: correct with missing hypotheses; publication match: exact; verification: read the result; development: none; confidence: high}
\end{GSitem}
\begin{GSitem}{1970-204.13}{E}{Hyperplane induction requires an injective multiplication map}
\GSfield{Kind; source}{method; \GSpage{P000328}}
\GSfield{The notes}{The notes use \(0\to\mathcal F\xrightarrow{x_i}\mathcal F(1)\to\overline{\mathcal F}\to0\), then finite differences of Euler characteristics.}
\GSfield{Statement}{Over an infinite field, choose a linear form \(\ell\) avoiding the associated points of \(\mathcal F\). Multiplication by \(\ell\) is then injective and its cokernel is supported on the hyperplane. Additivity gives \(\chi(\mathcal F(k+1))-\chi(\mathcal F(k))=\chi(\overline{\mathcal F}(k))\), allowing induction on support dimension. A fixed coordinate \(x_i\) may annihilate a nonzero sheaf and need not be injective. Polynomiality over finite fields follows after a field extension, or by a graded resolution.}
\GSfield{Reconstruction}{The coordinate is replaced by a general linear form regular on the sheaf; the finite-field case is separated from the infinite-field hyperplane choice.}
\GSfield{Fate}{\citet[Lemma~33.35.3 and Remark~33.35.5]{stacks2026projconstruction} gives the general-hyperplane condition explicitly. This supplies the missing hypothesis in the abbreviated proof, while the intended Hilbert-polynomial conclusion remains classical.}
\GSfield{Present status}{As of 9 November 2026: Established with the hypotheses in the statement.}
\GSfield{Assessment}{literal statement: correct with missing hypotheses; publication match: exact; verification: read the result; development: annotated; confidence: high}
\end{GSitem}

\GStalk{205}{}{}{P000329}{1971 calendar-year heading}
The catalog inserts a calendar-year heading before the first February 1971 entries. It is an organizational marker, with no independent mathematical statement. The next surviving dated block on P000329 reads 8/II--71; it remains in the 1970--71 academic-year chapter.


\GStalk{206}{1971-02-08}{I. M. Gelfand}{P000329-P000330}{Bott’s integrability obstruction and analytic continuation programs}
The dated Gelfand passage reports Bott's obstruction to deforming a distribution to a foliation, explains the easier submersion case, and lists proposed Sinai, Gabrielov, and Bernstein subjects. The quartic-oscillator note at the top of P000330 precedes Fedosov's clear heading; its speaker name is cancelled and uncertain, so it is retained here as an unattributed adjacent remark rather than assigned to Kazhdan from the catalog alone.

\begin{GSitem}{1970-206.1}{E}{Can every tangent distribution be deformed to an integrable one?}
\GSfield{Kind; source}{question; \GSpage{P000329}}
\GSfield{The notes}{A tangent subbundle is given, and the notes ask whether an integrable foliation homotopic to it exists; Bott is said to give a negative example.}
\GSfield{Statement}{For a smooth constant-rank subbundle \(E\subset TM\), ask whether it can be joined through subbundles of the same rank to the tangent bundle of a nonsingular smooth foliation. In general the answer is negative: the normal bundle of a foliation satisfies characteristic-class restrictions that a general subbundle quotient need not satisfy. Homotopy of subbundles preserves these normal characteristic classes.}
\GSfield{Fate}{The reference to Global Analysis, volume 16, identifies \citet{bott1970integrability}, printed in 1970 in proceedings of the July 1968 symposium. Bott's later primary exposition explicitly derives the obstruction from characteristic classes in \citet[sections~2 and 4]{bott1976characteristic}. The notebook is reporting this existing obstruction, not posing a still-open universal existence conjecture. Compare the earlier class-F foliation slogan, 1968-149.3, without assuming that its unstated hypotheses coincide.}
\GSfield{Present status}{As of 9 November 2026: The unrestricted question has a negative answer; additional geometric and homotopy hypotheses lead to more refined existence problems not specified here.}
\GSfield{Assessment}{literal statement: correct with missing hypotheses; publication match: exact; verification: read the result; development: none; confidence: high; problem \hyperref[prob:GP-1970-26]{GP-1970-26}}
\end{GSitem}
\begin{GSitem}{1970-206.2}{E}{Bott vanishing concerns the normal Pontryagin ring}
\GSfield{Kind; source}{theorem; \GSpage{P000329}}
\GSfield{The notes}{Pontryagin classes and their products vanish for \(k>2\dim\) of the indicated bundle, whose identity and degree convention are compressed.}
\GSfield{Statement}{Let \(E\subset TM\) be the tangent bundle of a smooth foliation of real codimension \(q\), and let \(\nu=TM/E\). Every homogeneous element of the Pontryagin subring of \(H^*(M;\mathbf R)\) generated by \(p_i(\nu)\) vanishes in cohomological degree \(>2q\). Thus a monomial \(\prod_i p_i(\nu)^{a_i}\) vanishes when \(4\sum_i ia_i>2q\). The assertion concerns the normal bundle, not the tangent bundle of the leaves, and does not assert integral torsion vanishing.}
\GSfield{Reconstruction}{The normal bundle and total real cohomological degree are identified from the subsequent submersion motivation and the checked Bott exposition.}
\GSfield{Fate}{\citet[pp.~125--126 and 135--136]{bott1976characteristic} identifies the normal characteristic homomorphism and the truncation giving this precise bound. The historical source is \citet{bott1970integrability}; its original chapter was checked bibliographically, not read in full.}
\GSfield{Present status}{As of 9 November 2026: Established with the hypotheses in the statement.}
\GSfield{Assessment}{literal statement: correct with missing hypotheses; publication match: exact; verification: read the result; development: annotated; confidence: high; problem \hyperref[prob:GP-1970-26]{GP-1970-26}}
\end{GSitem}
\begin{GSitem}{1970-206.3}{F}{The compressed CP² example cannot be reconstructed as stated}
\GSfield{Kind; source}{claim; \GSpage{P000329}}
\GSfield{The notes}{The next line says that an example already occurs on \(\mathbf CP^2\), using Pontryagin classes.}
\GSfield{Statement}{No rank, normal bundle, or ambient total-space construction is supplied. A literal proper nonsingular real distribution on \(\mathbf CP^2\) is problematic: ranks one and three would force a nonzero vector field despite \(\chi(\mathbf CP^2)=3\); a rank-two splitting would require integers \(a,b\) with \(ab=3\) and \(a^2+b^2=3\), from Euler and first Pontryagin classes, which is impossible. Thus this line cannot identify the asserted homotopy counterexample on that manifold without additional construction or an emendation.}
\GSfield{Reconstruction}{The elementary characteristic-class consistency check explains why the literal manifold name is insufficient; no missing example is invented.}
\GSfield{Fate}{The general Bott obstruction is securely identified by \citet{bott1970integrability} and \citet[section~4]{bott1976characteristic}; those checks do not identify the manuscript's exact abbreviated projective-space example. It is kept separate from the valid theorem rather than silently replaced by an example on a different space.}
\GSfield{Present status}{As of 9 November 2026: The exact example intended by the notebook is unidentified. This is an archival identification limit, not a declaration that the general obstruction problem remains open.}
\GSfield{Assessment}{literal statement: not determinable; publication match: none found; verification: unresolved lead; development: annotated; confidence: medium}
\end{GSitem}
\begin{GSitem}{1970-206.4}{E}{The stronger dimension bound for a foliation by submersion fibres}
\GSfield{Kind; source}{method; \GSpage{P000329}}
\GSfield{The notes}{For \(\pi:Y\to X\), the notes write \(0\to E\to TY\to\pi^*TX\to0\) and use \(H^k(X)=0\) above \(\dim X\).}
\GSfield{Statement}{If \(\pi:Y\to X\) is a smooth submersion with \(E=\ker d\pi\), then \(TY/E\cong\pi^*TX\). Every normal Pontryagin polynomial is pulled back from \(X\), and hence vanishes in cohomological degree \(>\dim_{\mathbf R}X\). A general foliation need not possess a manifold leaf space or such a global submersion, so this argument does not prove the stronger bound for every foliation.}
\GSfield{Fate}{This is the naturality-and-dimension motivation for the weaker universal \(2q\) bound. The normal-bundle interpretation agrees with \citet[section~4]{bott1976characteristic}.}
\GSfield{Present status}{As of 9 November 2026: Established with the hypotheses in the statement.}
\GSfield{Assessment}{literal statement: correct with missing hypotheses; publication match: partial; verification: read the result; development: none; confidence: high}
\end{GSitem}
\begin{GSitem}{1970-206.5}{F}{Planned analytic fundamental solutions and finite-dimensional branching}
\GSfield{Kind; source}{remark; \GSpage{P000329}}
\GSfield{The notes}{Beside Gabrielov and Bernstein, three goals are listed for constant-coefficient equations: analyticity outside finite codimension, continuation throughout complex \(\mathbf C^n\), and a finite-dimensional space of branches.}
\GSfield{Statement}{The passage is a program, not a fully specified theorem. It leaves unstated the class of differential operators, the distributional or holomorphic meaning of a fundamental solution, the exceptional set to be removed, and the paths along which continuation is allowed. The finite-dimensionality phrase may concern the span of branch germs; it does not literally mean that only finitely many branches occur.}
\GSfield{Fate}{The source and adjacent later Bernstein passages were compared with the general analytic-continuation/holonomic research line. No exact primary statement establishing all three assertions with the missing hypotheses was verified. Consequently no universal continuation theorem or new open problem is extracted from this list.}
\GSfield{Present status}{As of 9 November 2026: The intended precise theorem and its exact publication remain unidentified; the three phrases are retained as programmatic content.}
\GSfield{Assessment}{literal statement: not determinable; publication match: none found; verification: unresolved lead; development: none; confidence: medium}
\end{GSitem}
\begin{GSitem}{1970-206.6}{E}{The quartic anharmonic-oscillator spectral problem}
\GSfield{Kind; source}{construction; \GSpage{P000330}; Unattributed adjacent remark}
\GSfield{The notes}{The differential equation is \(-y^{\prime\prime}+x^2y+\beta x^4y=\lambda y\), introduced under a cancelled speaker name.}
\GSfield{Statement}{The natural spectral interpretation is the one-dimensional Schrödinger operator \(H_\beta=-d^2/dx^2+x^2+\beta x^4\) on \(L^2(\mathbf R)\), initially for real \(\beta>0\), with the square-integrable/decaying boundary conditions. Its discrete eigenvalues may then be studied by analytic continuation in complex \(\beta\); the continuation requires specifying sheets and boundary sectors.}
\GSfield{Reconstruction}{The Hilbert-space and decay conditions are supplied from the identified spectral model; the cancelled name does not establish a seminar speaker.}
\GSfield{Fate}{The model agrees after rescaling with \citet[abstract]{benderwu1969anharmonic}, published 25 August 1969, and with the title and scope of \citet{simon1970coupling}, published 1970. The boundary conditions are explicit in Bender--Wu, but omitted in the notebook.}
\GSfield{Present status}{As of 9 November 2026: Established with the hypotheses in the statement.}
\GSfield{Assessment}{literal statement: correct with missing hypotheses; publication match: exact; verification: publisher metadata checked; development: none; confidence: high}
\end{GSitem}
\begin{GSitem}{1970-206.7}{E}{Accumulating branch points and the three-sheeted description}
\GSfield{Kind; source}{claim; \GSpage{P000330}; Unattributed adjacent remark}
\GSfield{The notes}{A provisionally countable number of second-order branch points is said to tend, as \(\beta\to0\), to a third-order branch point; the work is called semi-physical.}
\GSfield{Statement}{The identifiable content is the Bender--Wu analytic continuation of quartic-oscillator levels, with branch singularities accumulating at zero. Zero is not an isolated ordinary algebraic branch point. A three-sheeted global description must be distinguished from a local Puiseux singularity of order three, and the notebook's unrestricted assertion that every relevant branch point is quadratic is not verified by the retrieved abstract.}
\GSfield{Fate}{\citet[abstract]{benderwu1969anharmonic} explicitly gives an infinite accumulation of branch points and uses WKB methods. \citet{simon1970coupling} is the rigorous 1970 follow-up; its publisher search abstract describes global three-sheeted monodromy, but the full paper was not read in this run. Thus the topic is securely matched, while the exact compressed branch-order statement has only a partial match.}
\GSfield{Present status}{As of 9 November 2026: The nonisolated singularity and continuation problem were already studied before the seminar. The literal branch-order slogan is not promoted to an exactly verified theorem.}
\GSfield{Assessment}{literal statement: not determinable; publication match: partial; verification: indirect; development: annotated; confidence: medium}
\end{GSitem}

\GStalk{207}{1971-02-08}{D. A. Kazhdan}{P000329-P000331}{Catalogued Kazhdan entry without a securely attributable surviving block}
The catalog names Kazhdan on 8 February, but these pages do not preserve an independent, securely readable Kazhdan lecture block. P000330 begins with a cancelled name above the quartic oscillator; P000331 has a different cancelled Ka\ldots dan-like name inside the continuation of Fedosov's Bott discussion. Neither cancellation licenses assigning a new mathematical theorem to this catalog record. The entry is retained with its attribution uncertainty and no invented mathematical items.


\GStalk{208}{1971-02-08}{B. V. Fedosov}{P000330-P000331}{Bott’s residue formula for holomorphic vector fields}
The Fedosov block starts at the clearly marked name below the oscillator on P000330. It gives the zero linearization and the residue expression for Chern numbers, then continues at the top of P000331 with its infinitesimal-Lefschetz interpretation and a real-analogue remark. It ends at the separator before the next-week announcement and the 15 February heading. The manuscript's tentative 1965/66 date is corrected by the explicit 1967 bibliography in record 210.

\begin{GSitem}{1970-208.1}{E}{Linearization at an isolated holomorphic zero}
\GSfield{Kind; source}{construction; \GSpage{P000330}}
\GSfield{The notes}{\(\Theta(X)Y=[X,Y]\); at a zero \(p\), this depends only on \(Y_p\), and nondegeneracy means \(\Theta_p\) is invertible.}
\GSfield{Statement}{If \(X\) is a holomorphic vector field and \(X(p)=0\), the map \(\Theta_p:T_pM\to T_pM\), \(Y_p\mapsto[X,Y]_p\), is well defined. With the usual Lie bracket it equals \(-dX_p\). A nondegenerate zero has \(\det dX_p\ne0\), and is therefore isolated. Reversing the sign convention does not change the degree-\(m\) residue quotient.}
\GSfield{Fate}{The Lie-derivative endomorphism is the construction used in \citet[section~3.4]{bruzzorubtsov2012localization}, which recovers \citet{bott1967vector}. The notebook has the right fibrewise dependence at a zero.}
\GSfield{Present status}{As of 9 November 2026: Established with the hypotheses in the statement.}
\GSfield{Assessment}{literal statement: correct with missing hypotheses; publication match: exact; verification: read the result; development: none; confidence: high}
\end{GSitem}
\begin{GSitem}{1970-208.2}{E}{The isolated-zero Bott formula computes Chern numbers}
\GSfield{Kind; source}{theorem; \GSpage{P000330}}
\GSfield{The notes}{For compact complex \(M\), a holomorphic field with finitely many nondegenerate zeros computes all Chern classes; the displayed formula is \(\overline\psi=\sum_p\psi(\Theta_p)/\det\Theta_p\).}
\GSfield{Statement}{Let \(M\) be a compact complex manifold of complex dimension \(m\), and \(X\) a holomorphic vector field with only isolated nondegenerate zeros. For a symmetric homogeneous polynomial \(\psi\) of degree \(m\) in the Chern roots, let \(\Psi(c_1,\ldots,c_m)\) be its expression in elementary symmetric functions. Then \[\int_M\Psi(c_1(TM),\ldots,c_m(TM))=\sum_{X(p)=0}\frac{\psi(\Theta_p)}{\det\Theta_p}.\] This computes characteristic numbers, not all cohomology classes themselves. Symmetry of \(\psi\) is required to make the expression independent of ordering of eigenvalues.}
\GSfield{Reconstruction}{The root polynomial is required to be symmetric, and the source phrase Chern classes is narrowed to the integrated Chern numbers that its formula computes.}
\GSfield{Fate}{This is the isolated-zero result of \citet{bott1967vector}; the same-year extension is \citet{bott1967residue}. Publication metadata and the original abstracts were opened. The exact formula was read in \citet[section~3.4]{bruzzorubtsov2012localization}, rather than inferred solely from a title.}
\GSfield{Present status}{As of 9 November 2026: Established before the seminar. Compactness, holomorphicity, symmetric polynomial degree, and isolated nondegenerate zeros delimit this displayed formula.}
\GSfield{Assessment}{literal statement: correct with missing hypotheses; publication match: exact; verification: read the result; development: annotated; confidence: high}
\end{GSitem}
\begin{GSitem}{1970-208.3}{E}{Lower-degree residue sums vanish}
\GSfield{Kind; source}{theorem; \GSpage{P000330}}
\GSfield{The notes}{The parenthesis says that if the polynomial degree is less than \(m\), the sum is zero.}
\GSfield{Statement}{With the same compact holomorphic setup and a symmetric homogeneous polynomial \(\psi\) of degree \(d<m\), \(\sum_p\psi(\Theta_p)/\det\Theta_p=0\). This is a lower-degree localization identity. No corresponding vanishing for all degrees \(d>m\) is asserted.}
\GSfield{Fate}{The vanishing is the lower-degree part of Bott localization; \citet[Theorem~3.2 and section~3.4]{bruzzorubtsov2012localization} supplies the general localization mechanism, and its equation~(15) displays the constant-polynomial instance. The source correctly states only the lower-degree case.}
\GSfield{Present status}{As of 9 November 2026: Established with the hypotheses in the statement.}
\GSfield{Assessment}{literal statement: correct with missing hypotheses; publication match: exact; verification: read the result; development: none; confidence: high}
\end{GSitem}
\begin{GSitem}{1970-208.4}{E}{The infinitesimal fixed-point interpretation}
\GSfield{Kind; source}{remark; \GSpage{P000331}}
\GSfield{The notes}{The formula is called an infinitesimal analogue of the Atiyah--Bott theorem because exponentiating the field gives a transformation of the manifold.}
\GSfield{Statement}{A holomorphic vector field generates a local holomorphic flow. Linearizing the fixed-point formulas for the flow at its stationary points relates the infinitesimal data \(dX_p\) to the characteristic-number residues. The notebook gives this conceptual relationship, not a complete limiting argument or a new fixed-point theorem.}
\GSfield{Fate}{The named fixed-point theory is \citet{atiyahbott1967lefschetz}; the original journal metadata was opened. Bott's residue formula is independently verified through \citet[section~3.4]{bruzzorubtsov2012localization}. The relationship is reported at the level of the notebook's remark and is not presented as a proof that was read in the original articles. The same named fixed-point theory appeared earlier in item 1965-56.9.}
\GSfield{Present status}{As of 9 November 2026: Established with the hypotheses in the statement.}
\GSfield{Assessment}{literal statement: correct; publication match: partial; verification: indirect; development: none; confidence: high}
\end{GSitem}
\begin{GSitem}{1970-208.5}{E}{The real analogue and the comparison with Morse inequalities}
\GSfield{Kind; source}{remark; \GSpage{P000331}}
\GSfield{The notes}{A real analogue is mentioned, with exact numerical values rather than estimates as in Morse theory; the preceding name is cancelled.}
\GSfield{Statement}{Bott-type characteristic-number formulas also have a real metric-preserving vector-field setting. The required real hypotheses are additional structure, not a statement about arbitrary smooth vector fields determining all characteristic numbers. The comparison with Morse theory concerns the exact sums in localization formulas; it does not replace all Morse inequalities by equalities or compute every Betti number.}
\GSfield{Fate}{\citet[introduction]{bruzzorubtsov2012localization} explicitly distinguishes the holomorphic and Riemannian-metric-preserving cases of \citet{bott1967vector}. The real case is treated more concretely in the subsequent Fedosov record 212; this early remark is retained with a partial publication match and no attribution to a cancelled speaker.}
\GSfield{Present status}{As of 9 November 2026: Established with the hypotheses in the statement.}
\GSfield{Assessment}{literal statement: correct with missing hypotheses; publication match: partial; verification: read the result; development: none; confidence: medium}
\end{GSitem}

\GStalk{209}{1971-02-15}{I. M. Gelfand}{P000331}{Differential geometry and Lie-algebra cohomology: announcement}
The dated heading announces work with Kazhdan on differential-geometric notions and Lie-algebra cohomology. No theorem or construction is written under it. The holomorphic-vector-field discussion above the date belongs to the preceding session.

\begin{GSitem}{1970-209.1}{F}{Announcement of work with Kazhdan}
\GSfield{Kind; source}{remark; \GSpage{P000331}}
\GSfield{The notes}{The heading announces \ru{Некоторые понятия дифф. геометрии и их применение к вопросам} followed by an abbreviated reference to Lie-algebra cohomology.}
\GSfield{Statement}{Only the subject and joint authorship are recoverable. No particular formal-coordinate bundle, cochain map, or cohomology calculation is stated.}
\GSfield{Reconstruction}{Identification with the later Gelfand--Kazhdan paper is bibliographical, not a reconstruction of an unwritten theorem.}
\GSfield{Fate}{The closely matching joint title appears in \citet[pp.~269--272]{gelfandkazhdan1971geometry}; the notes do not permit an exact result-level match.}
\GSfield{Present status}{As of November 2026, the topic and subsequent publication are identified, but the announced mathematical content remains unrecoverable from this entry.}
\GSfield{Assessment}{literal statement: not determinable; publication match: thematic only; verification: publisher metadata checked; development: none; confidence: medium}
\end{GSitem}

\GStalk{210}{1971-02-15}{B. V. Fedosov}{P000331}{Bibliography for vector-field residue formulas}
Fedosov supplies four references, including Bott\textquoteright s two 1967 papers and Baum--Bott\textquoteright s 1970 meromorphic-vector-field paper, and recalls the compact complex-manifold setting. Bernstein then develops the meromorphic version under his own heading.

\begin{GSitem}{1970-210.1}{E}{Bibliographical route from holomorphic to meromorphic residues}
\GSfield{Kind; source}{method; \GSpage{P000331}}
\GSfield{The notes}{Fedosov lists Michigan Math. J. 14 (1967), 231--244; J. Differential Geom. 1 (1967), 311--330; a 1970 Matematika translation; and Baum--Bott, pp.~29--47 in the de Rham volume.}
\GSfield{Statement}{The first two entries are Bott\textquoteright s vector-field characteristic-number and holomorphic-residue papers. The fourth is Baum--Bott\textquoteright s meromorphic-vector-field extension. These are sources for, rather than new assertions beyond, the adjacent formulas.}
\GSfield{Reconstruction}{Titles, authors, and full publication details are supplied by bibliographical matching. The unspecified Matematika entry is not assigned a separate result.}
\GSfield{Fate}{The titles and pagination are verified in the authors\textquoteright curated bibliographies \citep{bott1967vector,bott1967residue,baumbott1970zeroes}.}
\GSfield{Present status}{As of November 2026, the reconstructed statement is established; no open problem is asserted.}
\GSfield{Assessment}{literal statement: correct; publication match: exact; verification: publisher metadata checked; development: none; confidence: high}
\end{GSitem}

\GStalk{211}{1971-02-15}{I. N. Bernstein}{P000331-P000332}{Polar divisors and regularized Chern numbers}
The polar divisor of a meromorphic vector field produces a degree-two class \(\rho\). The notes replace the Chern classes by the coefficients of \(c(TX)/(1-\rho)\) before applying a zero-residue formula. Isolated zero and nondegeneracy assumptions are largely shorthand.

\begin{GSitem}{1970-211.1}{E}{Polar divisor of a meromorphic vector field}
\GSfield{Kind; source}{construction; \GSpage{P000331}}
\GSfield{The notes}{The pole set is written as irreducible hypersurfaces \(Y_i\), with an effective cycle \(\sum\alpha_iY_i\). A local equation removes the pole of \(V\).}
\GSfield{Statement}{For a meromorphic section \(V\) of \(TX\) on a complex manifold, choose an effective Cartier divisor \(D=\sum m_iY_i\) clearing its poles. Then \(V\) is a holomorphic section of \(TX\otimes\mathcal O(D)\), or equivalently a morphism \(\mathcal O(-D)\to TX\). Set \(\rho=c_1(\mathcal O(D))\in H^2(X;\mathbb Z)\).}
\GSfield{Reconstruction}{The divisor is a real codimension-two cycle; the notes\textquoteright wording \emph{two-dimensional cycles} confuses degree and dimension. Local multiplicities and the line-bundle formulation are made explicit.}
\GSfield{Fate}{This is the setup of the 1970 meromorphic residue theorem \citep[pp.~29--47]{baumbott1970zeroes}; its abstract explicitly uses a tangent field twisted by a line bundle. The exact virtual-bundle formulation is independently verified in \citet[\S4.5, Proposition 4.6]{bruzzorubtsov2012localization}; its notation \(L=\mathcal O(D)\) gives \(TX-L^*\).}
\GSfield{Present status}{As of November 2026, the reconstructed statement is established; no open problem is asserted.}
\GSfield{Assessment}{literal statement: correct with missing hypotheses; publication match: exact; verification: read the result; development: none; confidence: high}
\end{GSitem}
\begin{GSitem}{1970-211.2}{E}{Regularized Chern classes of the meromorphic field}
\GSfield{Kind; source}{definition; \GSpage{P000332}}
\GSfield{The notes}{\(a_1=c_1+\rho\), \(a_2=c_2+c_1\rho+\rho^2\), and \(a_i=c_i+c_{i-1}\rho+\cdots+\rho^i\).}
\GSfield{Statement}{These are the degree-\(2i\) components of \(c(TX-\mathcal O(-D))=c(TX)/(1-\rho)\). They are classes of a virtual bundle of rank \(n-1\), not the ordinary Chern classes of \(TX\otimes\mathcal O(D)\).}
\GSfield{Reconstruction}{The virtual-bundle identification is editorial; it prevents replacing the notebook\textquoteright s actual formula with a different tensor-product formula.}
\GSfield{Fate}{Baum--Bott\textquoteright s meromorphic formulation is the prior source named in the notes \citep[pp.~29--47]{baumbott1970zeroes}. The algebraic identity is checked directly from the displayed notebook coefficients. The exact virtual-bundle formulation is independently verified in \citet[\S4.5, Proposition 4.6]{bruzzorubtsov2012localization}; its notation \(L=\mathcal O(D)\) gives \(TX-L^*\).}
\GSfield{Present status}{As of November 2026, the reconstructed statement is established; no open problem is asserted.}
\GSfield{Assessment}{literal statement: correct; publication match: exact; verification: read the result; development: none; confidence: high}
\end{GSitem}
\begin{GSitem}{1970-211.3}{E}{Residues for meromorphic one-dimensional foliations}
\GSfield{Kind; source}{theorem; \GSpage{P000332}}
\GSfield{The notes}{The notes ask that the preceding holomorphic residue formula be applied using the regularized \(a_i\).}
\GSfield{Statement}{Let \(X\) be compact complex of dimension \(n\), and \(L=\mathcal O(-D)\to TX\) a holomorphic line-bundle morphism with isolated nondegenerate zeros. For an invariant polynomial \(\phi\) of degree \(n\), its characteristic number on \(TX-L\) is the sum of local residues. In a local frame, with derivative matrix \(A_p\), the nondegenerate residue is \(\phi(A_p)/\det A_p\). Rescaling the local frame cancels because both numerator and denominator have degree \(n\).}
\GSfield{Reconstruction}{Compactness, isolated zeros of the extended morphism, and nondegeneracy are supplied. Zeros created on the polar divisor must also be included; counting only zeros away from poles need not give the characteristic number.}
\GSfield{Fate}{The prior source named in the notes is \citet[pp.~29--47]{baumbott1970zeroes}. The exact virtual-bundle formula and isolated-zero hypothesis were read in \citet[\S4.5, Proposition 4.6]{bruzzorubtsov2012localization}, which explicitly identifies it as the Baum--Bott theorem. The original 1970 proof was not retrieved.}
\GSfield{Present status}{As of November 2026, the reconstructed statement is established; no open problem is asserted.}
\GSfield{Assessment}{literal statement: correct with missing hypotheses; publication match: partial; verification: read the result; development: none; confidence: high}
\end{GSitem}

\GStalk{212}{1971-02-15}{B. V. Fedosov}{P000332}{Pontryagin numbers from zeros of a Killing field}
A real analogue of Bott\textquoteright s formula uses a Killing vector field and the Pfaffian of its skew linearization. Compactness, nondegenerate isolated zeros, orientation, and characteristic-form normalization must be supplied.

\begin{GSitem}{1970-212.1}{E}{Killing field and skew linearization}
\GSfield{Kind; source}{definition; \GSpage{P000332}}
\GSfield{The notes}{The metric identity is \(([X,Y],Z)+(Y,[X,Z])=X(Y,Z)\); at a zero, the derivative \(\Theta_p\) is skew.}
\GSfield{Statement}{A smooth vector field \(X\) is Killing when \(\mathcal L_Xg=0\). At a zero \(p\), \(A_p=(\nabla X)_p\) is skew-adjoint on \(T_pM\), and a nondegenerate zero means \(\det A_p\ne0\).}
\GSfield{Reconstruction}{The displayed identity is read as a formula for the Riemannian inner product, not a Lie-algebra bilinear form.}
\GSfield{Fate}{This is the infinitesimal-isometry framework used by \citet[pp.~173--193]{baumcheeger1969isometries}.}
\GSfield{Present status}{As of November 2026, the reconstructed statement is established; no open problem is asserted.}
\GSfield{Assessment}{literal statement: correct with missing hypotheses; publication match: exact; verification: publisher metadata checked; development: none; confidence: high}
\end{GSitem}
\begin{GSitem}{1970-212.2}{E}{Real localization formula with a signed Pfaffian}
\GSfield{Kind; source}{theorem; \GSpage{P000332}}
\GSfield{The notes}{For an oriented manifold of dimension \(4k\), the Pontryagin characteristic number is a sum of \(\psi(\Theta_p)/\operatorname{Pf}(\Theta_p)\).}
\GSfield{Statement}{For a closed oriented Riemannian \(4k\)-manifold and a Killing field with isolated nondegenerate zeros, an invariant polynomial representing a Pontryagin number localizes as \(\int_M\psi(R/(2\pi))=\sum_p\psi(A_p)/\operatorname{Pf}(A_p)\), with matching Chern--Weil conventions. The Pfaffian uses the given orientation; it is not an unsigned square root of the determinant.}
\GSfield{Reconstruction}{Closedness and isolated nondegenerate zeros are essential. A consistent normalization is specified; numerator and denominator have the same total degree, so common \(2\pi\) factors cancel.}
\GSfield{Fate}{The prior infinitesimal-isometry result is \citet[pp.~173--193]{baumcheeger1969isometries}; Bott\textquoteright s 1967 characteristic-number work is also explicitly on the reading list \citep{bott1967vector}. Original result-level access was incomplete.}
\GSfield{Present status}{As of November 2026, the reconstructed statement is established; no open problem is asserted.}
\GSfield{Assessment}{literal statement: correct with missing hypotheses; publication match: exact; verification: publisher metadata checked; development: none; confidence: medium}
\end{GSitem}

\GStalk{213}{1971-02-22}{V. I. Arnold}{P000333-P000334}{Versal deformations of matrices}
Matrices depending on parameters are studied up to parameter-dependent conjugation. Jordan blocks explain singularities in the quotient, while transverse parameter families give versal deformations. Gelfand proposes a broader local picture for group actions and tensors, but the concluding conjectural language does not specify a mathematical category or equivalence.

\begin{GSitem}{1970-213.1}{A}{Equivalence and induced matrix deformations}
\GSfield{Kind; source}{definition; \GSpage{P000333}}
\GSfield{The notes}{A matrix \(A(\lambda)\) depending on parameters is considered after changes of parameters and parameter-dependent similarity transformations.}
\GSfield{Statement}{A deformation of \(A_0\in\operatorname{Mat}_n(\mathbb C)\) is a smooth or holomorphic germ \(A:(\mathbb C^r,0)\to(\operatorname{Mat}_n,A_0)\). Another deformation is induced if \(B(\mu)=S(\mu)^{-1}A(\varphi(\mu))S(\mu)\), where \(\varphi(0)=0\) and \(S(0)\) centralizes \(A_0\) (one may require \(S(0)=I\)). Versality means every deformation is so induced in the chosen category.}
\GSfield{Reconstruction}{The smooth and holomorphic categories are separated; the notes do not authorize a topological classification under arbitrary homeomorphisms.}
\GSfield{Fate}{Arnold\textquoteright s 1971 paper gives the parameter-dependent normal-form theory \citep[pp.~29--43]{arnold1971matrices}. The lecture precedes the journal issue; no earlier preprint date was established.}
\GSfield{Present status}{As of November 2026, the reconstructed statement is established; no open problem is asserted.}
\GSfield{Assessment}{literal statement: correct with missing hypotheses; publication match: exact; verification: publisher metadata checked; development: none; confidence: high}
\end{GSitem}
\begin{GSitem}{1970-213.2}{A}{Transverse slices as versal parameter families}
\GSfield{Kind; source}{method; \GSpage{P000333-P000334}}
\GSfield{The notes}{Jordan normal forms organize the possible matrices, but smooth parameter-dependent normal forms require extra entries near repeated eigenvalues.}
\GSfield{Statement}{At \(A_0\), the conjugacy orbit has tangent space \([A_0,\operatorname{Mat}_n]\). If \(V\) is a linear complement, then \(A_0+V\) is miniversal locally: the map \((S,v)\mapsto S^{-1}(A_0+v)S\) is a submersion at \((I,0)\). Its parameter dimension is \(\dim\ker\operatorname{ad}_{A_0}\).}
\GSfield{Reconstruction}{The complement and submersion argument are a modern precise reconstruction. Jordan canonical form chosen independently at each parameter does not itself prove smooth versality.}
\GSfield{Fate}{The normal-form framework matches \citet[pp.~29--43]{arnold1971matrices}; the detailed continuation and parameter count occur in record 218.}
\GSfield{Present status}{As of November 2026, the reconstructed statement is established; no open problem is asserted.}
\GSfield{Assessment}{literal statement: correct with missing hypotheses; publication match: exact; verification: publisher metadata checked; development: none; confidence: high}
\end{GSitem}
\begin{GSitem}{1970-213.3}{F}{Gelfand’s proposed local model for general group actions}
\GSfield{Kind; source}{conjecture; \GSpage{P000334}; I. M. Gelfand}
\GSfield{The notes}{Gelfand discusses high-order tensors and proposes that locally one has a nonsingular fibration plus a remaining adjoint/Jordan-form problem.}
\GSfield{Statement}{The group, action, locality, topology, allowable equivalence, and meaning of the remaining adjoint problem are not fixed. This is a research heuristic, not a quantified slice theorem.}
\GSfield{Reconstruction}{No universal theorem is supplied: a modern slice result for a specified proper or reductive action would add substantial assumptions not present here.}
\GSfield{Fate}{Searches under Arnold\textquoteright s matrix paper and the Gelfand--Kazhdan bibliography give related normal-form/cohomology topics, not an exact publication of this general conjecture \citep{arnold1971matrices,gelfandkazhdan1971geometry}.}
\GSfield{Present status}{As of November 2026, no definite proposition is recoverable; no open problem or absence-of-development claim is entered.}
\GSfield{Assessment}{literal statement: not determinable; publication match: thematic only; verification: unresolved lead; development: none; confidence: low}
\end{GSitem}
\begin{GSitem}{1970-213.4}{F}{Invariant distributions as motivation}
\GSfield{Kind; source}{remark; \GSpage{P000334}}
\GSfield{The notes}{The closing remarks mention invariant generalized functions together with the geometry of orbits.}
\GSfield{Statement}{No distribution space, invariance equation, support condition, or claimed classification is given.}
\GSfield{Reconstruction}{This occurrence is retained separately because it suggests an analytic application, but no theorem about invariant distributions is inferred.}
\GSfield{Fate}{The matrix-deformation publication is only a thematic context \citep{arnold1971matrices}; the remark cannot be matched to a result.}
\GSfield{Present status}{As of November 2026, this remains an unidentified remark, not a documented unsolved problem.}
\GSfield{Assessment}{literal statement: not determinable; publication match: thematic only; verification: unresolved lead; development: none; confidence: low}
\end{GSitem}

\GStalk{214}{1971-03-01}{A. M. Gabrielov}{P000335}{Griffiths on periods: bibliographical heading}
The entire entry is a heading and a reference to Griffiths\textquoteright s Annals papers, volume 90, number 3. The substantive exposition resumes in records 220 and 222; it is not retrospectively placed under this heading.

\begin{GSitem}{1970-214.1}{E}{Reference to Griffiths’s rational-period papers}
\GSfield{Kind; source}{remark; \GSpage{P000335}}
\GSfield{The notes}{The heading gives Griffiths, Annals of Mathematics 90, no.~3.}
\GSfield{Statement}{The bibliographical reference is to the two 1969 papers, pp.~460--495 and 496--541. No independent assertion accompanies this entry.}
\GSfield{Reconstruction}{The title and page ranges are supplied from the author\textquoteright s publication list; the later Gabrielov exposition has its own occurrences.}
\GSfield{Fate}{The author\textquoteright s bibliography verifies both papers \citep{griffiths1969periods1,griffiths1969periods2}. Its attached PDF is a different 1968 paper and was not treated as the 1969 full text.}
\GSfield{Present status}{As of November 2026, the reconstructed statement is established; no open problem is asserted.}
\GSfield{Assessment}{literal statement: correct; publication match: exact; verification: publisher metadata checked; development: none; confidence: high}
\end{GSitem}

\GStalk{215}{1971-03-01}{I. N. Bernstein}{P000335-P000336}{Quasihomogeneous singularities and a torsion preprint}
Bernstein reports Milnor--Orlik\textquoteright s formulas for the Milnor number and the monodromy divisor of an isolated quasihomogeneous hypersurface singularity. A brief comment connects polynomial-power distributions with cycles, and the page closes with news of the Ray--Singer preprint on analytic and combinatorial torsion.

\begin{GSitem}{1970-215.1}{E}{Weighted homogeneity and the isolated singularity}
\GSfield{Kind; source}{definition; \GSpage{P000335}}
\GSfield{The notes}{For the weights \(w_i\), monomials satisfy \(\sum i_j/w_j=1\); \(f=0\) has an isolated critical point at the origin.}
\GSfield{Statement}{Let \(f\in\mathbb C[z_1,\ldots,z_m]\) be weighted homogeneous of weighted degree one with positive rational coordinate weights \(q_i=1/w_i\), and suppose \(\nabla f\) vanishes only at \(0\). The Milnor fibre has reduced homology concentrated in degree \(m-1\); its rank is the Milnor number \(\mu\).}
\GSfield{Reconstruction}{The notebook\textquoteright s \(w_i\) are reciprocals of a widespread modern weight convention. Reduced homology is required also for \(m=1\).}
\GSfield{Fate}{This is the setting of \citet[pp.~385--393]{milnororlik1970weighted}; the modern formulation is checked in \citet[definition preceding Theorem~1.5]{hertlingmase2022weights}.}
\GSfield{Present status}{As of November 2026, the reconstructed statement is established; no open problem is asserted.}
\GSfield{Assessment}{literal statement: correct with missing hypotheses; publication match: exact; verification: read the result; development: none; confidence: high}
\end{GSitem}
\begin{GSitem}{1970-215.2}{E}{Milnor number of a weighted homogeneous hypersurface}
\GSfield{Kind; source}{theorem; \GSpage{P000335-P000336}}
\GSfield{The notes}{\(\mu=\prod_i(w_i-1)\); the diagonal example is \(z_1^{a_1}+\cdots+z_m^{a_m}\).}
\GSfield{Statement}{Under the isolated-critical-point hypothesis, \(\mu=\dim_{\mathbb C}\mathbb C[z]/(\partial f)=\prod_i(q_i^{-1}-1)\). For the Brieskorn polynomial with integers \(a_i\ge2\), this is \(\prod_i(a_i-1)\).}
\GSfield{Reconstruction}{The finite-dimensional Jacobian algebra is supplied as the precise meaning of multiplicity; arbitrary positive rational weights without an isolated singularity do not suffice.}
\GSfield{Fate}{The formula was published before the talk in \citet{milnororlik1970weighted}; its Jacobian-algebra interpretation is checked against \citet[Theorem~6.4(b)]{hertlingmase2022weights}.}
\GSfield{Present status}{As of November 2026, the reconstructed statement is established; no open problem is asserted.}
\GSfield{Assessment}{literal statement: correct with missing hypotheses; publication match: exact; verification: read the result; development: none; confidence: high}
\end{GSitem}
\begin{GSitem}{1970-215.3}{E}{The group-ring divisor calculation for monodromy}
\GSfield{Kind; source}{construction; \GSpage{P000336}}
\GSfield{The notes}{In \(\mathbb Q[\mathbb C^*]\), \(E_k=k^{-1}\sum_{\zeta^k=1}[\zeta]\) is an idempotent. For reduced \(w_i=u_i/v_i\), the monodromy divisor is \(\prod_i(w_iE_{u_i}-1)\).}
\GSfield{Statement}{Write \(\Lambda_d=\sum_{\zeta^d=1}[\zeta]\), with multiplication induced by multiplication of complex numbers. Then \(\Lambda_a\Lambda_b=\gcd(a,b)\Lambda_{\operatorname{lcm}(a,b)}\), and \(\operatorname{div}\det(tI-h_*)=\prod_i(\Lambda_{u_i}/v_i-[1])\). For an isolated weighted homogeneous singularity the final coefficients are nonnegative integers.}
\GSfield{Reconstruction}{The symbol \(1\) denotes the identity group-ring element \([1]\). Rational intermediate coefficients do not mean the characteristic polynomial has fractional multiplicities.}
\GSfield{Fate}{The monodromy formula is exactly the Milnor--Orlik theorem, explicitly restated in \citet[Theorem~6.4(a)]{hertlingmase2022weights}; the original paper is \citet{milnororlik1970weighted}.}
\GSfield{Present status}{As of November 2026, the reconstructed statement is established; no open problem is asserted.}
\GSfield{Assessment}{literal statement: correct with missing hypotheses; publication match: exact; verification: read the result; development: none; confidence: high}
\end{GSitem}
\begin{GSitem}{1970-215.4}{F}{Polynomial-power distributions and cycles}
\GSfield{Kind; source}{remark; \GSpage{P000336}}
\GSfield{The notes}{Bernstein connects poles or residues of \(P^\lambda\) with cycles and their transformations.}
\GSfield{Statement}{The note does not specify the polynomial, real or complex distribution, integration cycle, pole, or claimed relation. Meromorphic continuation and monodromy are distinguishable questions; neither a precise residue formula nor a monodromy eigenvalue theorem is recorded.}
\GSfield{Reconstruction}{No theorem is reconstructed from the general association.}
\GSfield{Fate}{Searches for Bernstein, polynomial powers, residues, and monodromy locate the later Bernstein--Sato theory, but do not identify this unquantified remark with one exact result. The preceding Milnor--Orlik formulas supply context only \citep{milnororlik1970weighted}.}
\GSfield{Present status}{As of November 2026, the general research direction is recognizable, but this occurrence does not define a new question.}
\GSfield{Assessment}{literal statement: not determinable; publication match: thematic only; verification: unresolved lead; development: none; confidence: low}
\end{GSitem}
\begin{GSitem}{1970-215.5}{E}{The Ray–Singer preprint and analytic torsion}
\GSfield{Kind; source}{construction; \GSpage{P000336}}
\GSfield{The notes}{The notebook mentions a preprint by Singer \emph{and someone else} on an analytic expression for torsion.}
\GSfield{Statement}{For a closed Riemannian manifold \(M\) and a unitary flat bundle \(E\), the positive Hodge-Laplacian spectrum defines \(\zeta_q(s)\). Analytic torsion is \(T(M,E)=\exp\{\frac12\sum_q(-1)^q q\zeta_q^{\prime}(0)\}\), with a stated torsion convention. For nonacyclic \(E\), comparison uses the corresponding determinant-line metric or compatible cohomology bases.}
\GSfield{Reconstruction}{Daniel B. Ray is the missing coauthor. Zeta regularization is not an ordinary convergent product of all eigenvalues; zero modes are omitted.}
\GSfield{Fate}{The preprint was circulating by this 1 March 1971 report; the paper appeared later in 1971 as \citet{raysinger1971torsion}. This identifies the construction, not a proof then available of the general equality. The original Introduction, p.~145, explicitly says the equality was not proved there; Definition 1.6 on p.~149 gives the zeta-regularized torsion.}
\GSfield{Present status}{As of November 2026, the reconstructed statement is established; no open problem is asserted.}
\GSfield{Assessment}{literal statement: correct with missing hypotheses; publication match: partial; verification: read the result; development: none; confidence: high}
\end{GSitem}
\begin{GSitem}{1970-215.6}{B}{Equality of analytic and Reidemeister torsion}
\GSfield{Kind; source}{question; \GSpage{P000336}}
\GSfield{The notes}{The preprint is described as expressing combinatorial torsion analytically; the status of a general equality is not separated from the proposal.}
\GSfield{Statement}{For a closed smooth manifold and an acyclic unitary flat bundle, does Ray--Singer analytic torsion equal the corresponding positive Reidemeister torsion? With cohomology present the comparison is one of metrics on the determinant line.}
\GSfield{Reconstruction}{The equality is identified as the Ray--Singer conjecture; it must not be reported as a theorem already proved by the 1971 preprint.}
\GSfield{Fate}{Cheeger announced a proof in 1977 \citep[p.~2651, abstract]{cheeger1977torsion}; M\"uller\textquoteright s independent full paper appeared in 1978 \citep{muller1978torsion}. The original Ray--Singer Introduction, p.~145, independently confirms its conjectural status in 1971 \citep{raysinger1971torsion}.}
\GSfield{Present status}{As of November 2026, the unitary comparison is the established Cheeger--M"uller theorem. The notes do not formulate later nonunitary anomaly questions.}
\GSfield{Assessment}{literal statement: not determinable; publication match: partial; verification: indirect; development: none; confidence: medium; problem \hyperref[prob:GP-1970-50]{GP-1970-50}}
\end{GSitem}

\GStalk{216}{1971-03-01}{I. M. Gelfand}{P000337}{Block determinants and positivity}
Gelfand recalls the Schur-complement factorization and a positivity criterion. The same subject occurred in item 1964-17.2. The notebook conflates nonnegative and positive-definite inequalities, so both precise versions are separated in the commentary.

\begin{GSitem}{1970-216.1}{E}{Schur determinant factorization}
\GSfield{Kind; source}{theorem; \GSpage{P000337}}
\GSfield{The notes}{For a block matrix, \(\det\begin{pmatrix}A&B\\C&D\end{pmatrix}=\det D\det(A-BD^{-1}C)\).}
\GSfield{Statement}{For square \(A,D\) over \(\mathbb C\), conformable off-diagonal blocks, and invertible \(D\), the stated determinant identity follows from block Gaussian elimination.}
\GSfield{Reconstruction}{Finite dimensions and invertibility are made explicit; no infinite-dimensional regularized determinant is meant. The same elementary identity already appears as 1964-17.2.}
\GSfield{Fate}{This is the classical Schur-complement identity; the prior occurrence 1964-17.2 is reused as provenance. A modern reference is \citet[Chapter~1]{zhang2005schur}.}
\GSfield{Present status}{As of November 2026, the reconstructed statement is established; no open problem is asserted.}
\GSfield{Assessment}{literal statement: correct with missing hypotheses; publication match: exact; verification: publisher metadata checked; development: none; confidence: high}
\end{GSitem}
\begin{GSitem}{1970-216.2}{E}{Positivity through the Schur complement}
\GSfield{Kind; source}{theorem; \GSpage{P000337}}
\GSfield{The notes}{The notebook uses \(D\ge0\) and \(A-BD^{-1}C\ge0\) while concluding positive definiteness.}
\GSfield{Statement}{If \(M=\begin{pmatrix}A&B\\B^*&D\end{pmatrix}\) is Hermitian and \(D>0\), then \(M>0\) iff \(A-BD^{-1}B^*>0\), and \(M\ge0\) iff the latter Schur complement is \(\ge0\). If \(D\) is singular, the inverse formula must be replaced by a pseudoinverse together with \(\operatorname{ran}B^*\subseteq\operatorname{ran}D\).}
\GSfield{Reconstruction}{Hermitian symmetry, strict versus weak positivity, and invertibility are supplied. These are corrections of the mathematical shorthand, not new assumptions silently attributed to the speaker.}
\GSfield{Fate}{The standard criterion is summarized in \citet[Chapter~1]{zhang2005schur}; the strict/weak distinction follows directly from the block congruence.}
\GSfield{Present status}{As of November 2026, the reconstructed statement is established; no open problem is asserted.}
\GSfield{Assessment}{literal statement: false as written; publication match: exact; verification: publisher metadata checked; development: none; confidence: high}
\end{GSitem}

\GStalk{217}{1971-03-01}{S. P. Novikov}{P000337-P000339}{Determinants of complexes and lens-space torsion}
An acyclic based finite complex has an alternating determinant, independent of the auxiliary boundary splittings. Twisted cellular complexes of lens spaces illustrate Reidemeister torsion. A nontrivial character is essential: the untwisted universal-cover complex is not acyclic.

\begin{GSitem}{1970-217.1}{E}{Determinant of an exact based complex}
\GSfield{Kind; source}{definition; \GSpage{P000337-P000338}}
\GSfield{The notes}{Novikov assigns a determinant to \(0\to E_0\to\cdots\to E_n\to0\) by choosing bases in images and lifting them to neighboring terms.}
\GSfield{Statement}{For a finite acyclic complex of finite-dimensional vector spaces over a field, fixed bases \(c_i\), boundary bases \(b_i\), and lifts of \(b_{i-1}\), torsion is the alternating product of the transition determinants \([b_i\widetilde b_{i-1}/c_i]^{(-1)^{i+1}}\), up to the explicitly chosen reciprocal/sign convention.}
\GSfield{Reconstruction}{The bases in the complex itself are part of the data. Independence concerns auxiliary boundary bases and lifts, not arbitrary changes of all \(c_i\).}
\GSfield{Fate}{The algebraic construction and its change-of-basis rules precede the lecture; see \citet[\S3]{milnor1966torsion}.}
\GSfield{Present status}{As of November 2026, the reconstructed statement is established; no open problem is asserted.}
\GSfield{Assessment}{literal statement: correct with missing hypotheses; publication match: exact; verification: read the result; development: none; confidence: high}
\end{GSitem}
\begin{GSitem}{1970-217.2}{E}{Cancellation of auxiliary determinants}
\GSfield{Kind; source}{claim; \GSpage{P000338}}
\GSfield{The notes}{The alternating product is said not to depend on the arbitrary intermediate choices.}
\GSfield{Statement}{Changing a boundary basis contributes equal determinants in adjacent degrees with opposite exponents; changing its lift gives a block-unitriangular change. Therefore torsion is independent of these choices. Changing \(c_i\) multiplies it by the corresponding alternating determinant.}
\GSfield{Reconstruction}{The elementary cancellation argument is expanded to specify precisely which basis-independence is true.}
\GSfield{Fate}{This is the standard based-complex torsion invariance \citep[\S3]{milnor1966torsion}.}
\GSfield{Present status}{As of November 2026, the reconstructed statement is established; no open problem is asserted.}
\GSfield{Assessment}{literal statement: correct with missing hypotheses; publication match: exact; verification: read the result; development: none; confidence: high}
\end{GSitem}
\begin{GSitem}{1970-217.3}{E}{Twisted cellular complex of a lens space}
\GSfield{Kind; source}{construction; \GSpage{P000338-P000339}}
\GSfield{The notes}{The lens-space complex has differentials \(1-\xi\), \(1+\xi+\cdots+\xi^{p-1}\), and \(1-\xi^\beta\).}
\GSfield{Statement}{For \(L(p,q)\), with \(p>1\), \(\gcd(p,q)=1\), and a nontrivial complex character taking a generator to \(\xi\), the twisted cellular chain complex has one one-dimensional term in degrees zero through three. Up to units/signs and the chosen generator, its maps are \(\xi-1,0,\xi^\beta-1\), with \(q\beta\equiv1\pmod p\). Both nonzero maps are invertible, so the twisted complex is acyclic.}
\GSfield{Reconstruction}{The character is nontrivial and \(\xi\) need not be primitive provided \(\xi^p=1\); coprimality ensures \(\xi^\beta\ne1\). The universal-cover complex before twisting is not acyclic.}
\GSfield{Fate}{The classical lens-space example and the role of a character are treated in \citet[\S12]{milnor1966torsion}.}
\GSfield{Present status}{As of November 2026, the reconstructed statement is established; no open problem is asserted.}
\GSfield{Assessment}{literal statement: correct with missing hypotheses; publication match: exact; verification: read the result; development: none; confidence: high}
\end{GSitem}
\begin{GSitem}{1970-217.4}{E}{Lens-space determinant and the torsion convention}
\GSfield{Kind; source}{claim; \GSpage{P000339}}
\GSfield{The notes}{The written answer is \((1-\xi)(1-\xi^\beta)\).}
\GSfield{Statement}{With Novikov\textquoteright s alternating determinant convention one may obtain this product; with the usual chain-torsion convention \(\tau(L(p,q),\xi)\) is its reciprocal, modulo the units caused by cell lifts and orientations. The invariant must be quoted with its convention.}
\GSfield{Reconstruction}{The notebook is not declared wrong merely because another convention inverts torsion. Its product is distinguished from an unconditional equality of canonically normalized complex numbers.}
\GSfield{Fate}{The reciprocal convention and unit ambiguity are explicit in the lens-space torsion calculation \citep[\S12]{milnor1966torsion}.}
\GSfield{Present status}{As of November 2026, the reconstructed statement is established; no open problem is asserted.}
\GSfield{Assessment}{literal statement: correct with missing hypotheses; publication match: exact; verification: read the result; development: none; confidence: high}
\end{GSitem}

\GStalk{218}{1971-03-01}{V. I. Arnold}{P000339-P000340}{Miniversal matrix families and centralizers}
The exposition of record 213 continues with a categorical description of induced deformations, an explicit nine-parameter family, and the centralizer-dimension formula. Transversality supplies versality; the notebook\textquoteright s description as a final object is stronger than the actual property.

\begin{GSitem}{1970-218.1}{A}{Versality is not a categorical final object}
\GSfield{Kind; source}{definition; \GSpage{P000339}}
\GSfield{The notes}{A deformation inducing all nearby deformations is described as a final object.}
\GSfield{Statement}{Versality requires existence of inducing parameter maps and conjugacies, not their uniqueness. Miniversality additionally minimizes the tangent parameter dimension. Nontrivial stabilizers generally obstruct the claimed final-object universal property.}
\GSfield{Reconstruction}{The distinction between a versal object and a universal object is an editorial correction. It continues item 1970-213.1.}
\GSfield{Fate}{Arnold\textquoteright s matrix-deformation framework is the precise match \citep[pp.~29--43]{arnold1971matrices}; metadata and abstract were checked, not the complete proof.}
\GSfield{Present status}{As of November 2026, the reconstructed statement is established; no open problem is asserted.}
\GSfield{Assessment}{literal statement: false as written; publication match: exact; verification: publisher metadata checked; development: none; confidence: high}
\end{GSitem}
\begin{GSitem}{1970-218.2}{A}{Nine parameters for Jordan blocks of sizes three and two}
\GSfield{Kind; source}{construction; \GSpage{P000339-P000340}}
\GSfield{The notes}{A five-by-five Jordan matrix with a common eigenvalue and blocks of lengths three and two is supplied with nine starred perturbation entries.}
\GSfield{Statement}{For \(A_0=J_3(\lambda)\oplus J_2(\lambda)\), the centralizer dimension is \(3+2+2+2=9\). A complement to \([A_0,\operatorname{Mat}_5]\) therefore yields a nine-parameter miniversal deformation. A particular star pattern is valid only if its coordinate subspace is such a complement.}
\GSfield{Reconstruction}{The exact count is verified algebraically; the drawn placement of stars is not promoted into a uniquely determined normal form.}
\GSfield{Fate}{This is a concrete instance of \citet{arnold1971matrices}, illustrated by the notebook itself.}
\GSfield{Present status}{As of November 2026, the reconstructed statement is established; no open problem is asserted.}
\GSfield{Assessment}{literal statement: correct with missing hypotheses; publication match: exact; verification: publisher metadata checked; development: none; confidence: high}
\end{GSitem}
\begin{GSitem}{1970-218.3}{A}{Centralizer codimension and the dimension of a miniversal base}
\GSfield{Kind; source}{theorem; \GSpage{P000340}}
\GSfield{The notes}{For Jordan blocks, the parameter count is \(n_1+3n_2+5n_3+\cdots\), summed over eigenvalues.}
\GSfield{Statement}{If the Jordan-block sizes for eigenvalue \(\lambda\) are \(n_{\lambda,1}\ge n_{\lambda,2}\ge\cdots\), then \(\dim Z(A)=\sum_\lambda\sum_{i,j}\min(n_{\lambda,i},n_{\lambda,j})=\sum_\lambda\sum_i(2i-1)n_{\lambda,i}\). This is the codimension of the conjugacy orbit and the dimension of a miniversal base.}
\GSfield{Reconstruction}{The ordering of block sizes and separate summation over distinct eigenvalues are supplied.}
\GSfield{Fate}{The formula and transverse-family interpretation are the theory announced in \citet{arnold1971matrices}; the centralizer count itself is elementary Jordan-module algebra.}
\GSfield{Present status}{As of November 2026, the reconstructed statement is established; no open problem is asserted.}
\GSfield{Assessment}{literal statement: correct with missing hypotheses; publication match: exact; verification: publisher metadata checked; development: none; confidence: high}
\end{GSitem}

\GStalk{219}{1971-03-15}{L. D. Faddeev}{P000341-P000344}{Inverse scattering and Hamiltonian KdV}
The one-dimensional inverse-scattering solution of KdV is followed by its Hamiltonian structure and scattering action--angle variables, subsequently published by Zakharov and Faddeev. The final part concerns multidimensional scattering and a Volterra factorization of the scattering operator already announced by Faddeev in 1966. The notebook changes KdV normalizations without consistently changing all coefficients.

\begin{GSitem}{1970-219.1}{E}{KdV through an evolving Schrödinger operator}
\GSfield{Kind; source}{method; \GSpage{P000341}}
\GSfield{The notes}{The notes write \(u_t+2u_{xxx}+6uu_x=0\), then associate \(-\partial_x^2+u\) to the initial potential.}
\GSfield{Statement}{For real Schwartz initial data, one standard normalization is \(u_t-6uu_x+u_{xxx}=0\) with \(L=-\partial_x^2+u\). Inverse scattering evolves the reflection coefficient by a cubic phase, keeps the negative eigenvalues fixed, and evolves their norming constants exponentially; inverse reconstruction gives the KdV solution.}
\GSfield{Reconstruction}{The standard normalization is stated separately. The literal coefficients on P000341 do not agree with the Hamiltonian displayed on P000342; changing signs or rescaling must be done consistently.}
\GSfield{Fate}{The inverse-scattering method predates the talk; \citet[Introduction, equations (1)--(5)]{zakharovfaddeev1971kdv} explicitly reviews the earlier Gardner--Greene--Kruskal--Miura method.}
\GSfield{Present status}{As of November 2026, the reconstructed statement is established; no open problem is asserted.}
\GSfield{Assessment}{literal statement: correct with missing hypotheses; publication match: partial; verification: read the result; development: none; confidence: high}
\end{GSitem}
\begin{GSitem}{1970-219.2}{E}{Reflection, transmission, and complete scattering data}
\GSfield{Kind; source}{definition; \GSpage{P000341}}
\GSfield{The notes}{The scattering solution behaves as \(e^{ikx}+a(k)e^{-ikx}\) at the left and \(b(k)e^{ikx}\) at the right, with \(|a|^2+|b|^2=1\).}
\GSfield{Statement}{For real \(u\) with \(\int(1+|x|)|u(x)|\,dx<\infty\), real \(k\ne0\) defines reflection \(r(k)\) and transmission \(t(k)\) coefficients with \(|r|^2+|t|^2=1\). To reconstruct a potential with bound states, reflection data must be supplemented by the negative eigenvalues and positive norming constants.}
\GSfield{Reconstruction}{The notebook\textquoteright s \(a,b\) are reflected/transmitted amplitudes, whereas the later paper uses different letters. The discrete spectrum alone does not supply the norming constants.}
\GSfield{Fate}{The full data and inverse problem are stated in \citet[Introduction and \S1]{zakharovfaddeev1971kdv}.}
\GSfield{Present status}{As of November 2026, the reconstructed statement is established; no open problem is asserted.}
\GSfield{Assessment}{literal statement: correct with missing hypotheses; publication match: exact; verification: read the result; development: none; confidence: high}
\end{GSitem}
\begin{GSitem}{1970-219.3}{E}{Reflectionless potentials and solitons}
\GSfield{Kind; source}{claim; \GSpage{P000341}}
\GSfield{The notes}{When the reflected wave vanishes, one obtains solitons; the discrete spectrum remains relevant.}
\GSfield{Statement}{A real rapidly decaying reflectionless potential with finitely many simple negative eigenvalues and their norming constants produces a finite soliton solution under KdV evolution. With no bound states and zero reflection the potential is zero.}
\GSfield{Reconstruction}{Reflectionless does not mean absence of scattering data: the eigenvalues and norming constants encode the solitons.}
\GSfield{Fate}{The scattering evolution and discrete variables are reviewed in \citet[Introduction, equations (4)--(5)]{zakharovfaddeev1971kdv}.}
\GSfield{Present status}{As of November 2026, the reconstructed statement is established; no open problem is asserted.}
\GSfield{Assessment}{literal statement: correct with missing hypotheses; publication match: exact; verification: read the result; development: none; confidence: high}
\end{GSitem}
\begin{GSitem}{1970-219.4}{E}{Conservation laws and the absence of shock formation}
\GSfield{Kind; source}{claim; \GSpage{P000341-P000342}}
\GSfield{The notes}{The margin associates infinitely many first integrals with the observation that shock waves do not arise.}
\GSfield{Statement}{For a smooth decaying KdV solution, integrations by parts give conserved local integrals, beginning with \(\int u\), \(\int u^2\), and, in the stated normalization, \(\int(u^3+\tfrac12u_x^2)\). Smooth real Schwartz data yield global smooth solutions in this integrable setting. The mere existence of many formal integrals is not by itself a proof of regularity.}
\GSfield{Reconstruction}{The qualitative observation is separated from the invalid general implication that an unspecified collection of first integrals rules out singularities.}
\GSfield{Fate}{The established KdV integrals are listed in \citet[equation (2) and \S3]{zakharovfaddeev1971kdv}; the marginal explanatory claim is stronger than what that list alone proves.}
\GSfield{Present status}{As of November 2026, the reconstructed statement is established; no open problem is asserted.}
\GSfield{Assessment}{literal statement: correct with missing hypotheses; publication match: partial; verification: read the result; development: none; confidence: medium}
\end{GSitem}
\begin{GSitem}{1970-219.5}{E}{Gardner’s Poisson bracket and a symplectic leaf}
\GSfield{Kind; source}{construction; \GSpage{P000342}}
\GSfield{The notes}{A double-integral skew form on variations \(\delta u\) is written, and the Hamiltonian is displayed as \(Q=\int(u^3+u_x^2)\,dx\).}
\GSfield{Statement}{The Hamiltonian operator \(J=\partial_x\) defines \(\{F,G\}=\int(\delta F/\delta u)\partial_x(\delta G/\delta u)\,dx\). On real Schwartz variations of zero integral, its inverse is \(\partial_x^{-1}v(x)=\frac12\int\operatorname{sgn}(x-y)v(y)\,dy\), giving a weak symplectic form. With \(H=\int(u^3+\tfrac12u_x^2)\,dx\), one obtains \(u_t=6uu_x-u_{xxx}\).}
\GSfield{Reconstruction}{The fixed-mass leaf, skew inverse, and coefficient \(1/2\) specify a consistent convention. The literal notebook Hamiltonian yields \(6uu_x-2u_{xxx}\), not its earlier displayed equation.}
\GSfield{Fate}{The paper credits the bracket to Gardner and records the consistent Hamiltonian \citep[equations (6)--(7)]{zakharovfaddeev1971kdv}.}
\GSfield{Present status}{As of November 2026, the reconstructed statement is established; no open problem is asserted.}
\GSfield{Assessment}{literal statement: false as written; publication match: partial; verification: read the result; development: none; confidence: high}
\end{GSitem}
\begin{GSitem}{1970-219.6}{A}{Trace identities and scattering expressions for local integrals}
\GSfield{Kind; source}{method; \GSpage{P000342-P000343}}
\GSfield{The notes}{The local integrals are expressed by moments of \(\log(1-|a|^2)\); the displayed expression for \(Q\) contains a \(k^4\) integral.}
\GSfield{Statement}{For real Schwartz potentials, the high-energy expansion of the transmission coefficient gives a hierarchy of trace identities: each appropriate local KdV integral is a fixed multiple of a moment of \(\log(1-|r(k)|^2)\) plus a finite sum of powers of the bound-state parameters. The constants depend on the KdV and scattering conventions.}
\GSfield{Reconstruction}{The notebook\textquoteright s schematic integral is not a complete equality when eigenvalues occur. No unchecked numerical coefficient is inserted into that display.}
\GSfield{Fate}{The precise hierarchy and discrete terms appear in \citet[\S3, equations (25)--(30)]{zakharovfaddeev1971kdv}, received 14 June 1971.}
\GSfield{Present status}{As of November 2026, the reconstructed statement is established; no open problem is asserted.}
\GSfield{Assessment}{literal statement: correct with missing hypotheses; publication match: partial; verification: read the result; development: none; confidence: high}
\end{GSitem}
\begin{GSitem}{1970-219.7}{A}{Canonical scattering coordinates}
\GSfield{Kind; source}{construction; \GSpage{P000343}}
\GSfield{The notes}{The skew form is rewritten using \(k\,\delta\log(1-|a|^2)\) and \(\delta\arg a\), with additional terms when eigenvalues exist.}
\GSfield{Statement}{In the published scattering convention \(r=b/a\), canonical variables include \(P(k)=-(k/\pi)\log(1-|r(k)|^2)\) and \(Q(k)=\arg b(k)\), together with bound-state action--angle pairs. On charts where the phases are defined, the symplectic form is the continuous canonical integral plus the discrete sum. Phase variables are singular at zeros of the corresponding coefficient.}
\GSfield{Reconstruction}{The notebook reflection coefficient is not automatically the paper\textquoteright s \(b\); replacing \(\arg b\) by \(\arg r\) without the induced corrections is unjustified. The paper explicitly calls its calculations formal derivations rather than a complete analytic manifold construction.}
\GSfield{Fate}{This is the central new result of \citet[\S2, equations (22)--(24)]{zakharovfaddeev1971kdv}; the paper\textquoteright s methodological qualification is on p.~282.}
\GSfield{Present status}{As of November 2026, the reconstructed statement is established; no open problem is asserted.}
\GSfield{Assessment}{literal statement: correct with missing hypotheses; publication match: partial; verification: read the result; development: none; confidence: high}
\end{GSitem}
\begin{GSitem}{1970-219.8}{E}{Three-dimensional scattering amplitude and the unitary operator}
\GSfield{Kind; source}{definition; \GSpage{P000343}}
\GSfield{The notes}{The asymptotic wave is \(e^{i\mathbf k\cdot x}+r^{-1}e^{ikr}f(k;\alpha,\beta)+o(r^{-1})\); \(f\) has five variables, and the notes write \(S(k)=I+f\), calling it unitary.}
\GSfield{Statement}{For a real sufficiently regular short-range potential in \(\mathbb R^3\), positive energy \(k^2\) and incident/outgoing directions \(\alpha,\beta\in S^2\) define the amplitude. The on-shell scattering operator on \(L^2(S^2)\), with its normalization factor and identity term, is unitary. The amplitude kernel alone is not a unitary operator.}
\GSfield{Reconstruction}{The factor relating the amplitude integral operator to \(S(k)-I\) is convention dependent. The dimensional count is one energy variable and two variables for each direction.}
\GSfield{Fate}{The amplitude/operator distinction is explicit in \citet[p.~70, equations (5)--(8)]{faddeev1966factorization}.}
\GSfield{Present status}{As of November 2026, the reconstructed statement is established; no open problem is asserted.}
\GSfield{Assessment}{literal statement: false as written; publication match: partial; verification: read the result; development: none; confidence: high}
\end{GSitem}
\begin{GSitem}{1970-219.9}{E}{Volterra factorization of the multidimensional scattering operator}
\GSfield{Kind; source}{method; \GSpage{P000343-P000344}}
\GSfield{The notes}{The notes factor \(S(k)\) by \(Q_\gamma^{\pm}\), with kernels supported on one side of \((\alpha-\beta)\cdot\gamma=0\), and mention analytic continuation of diagonal data.}
\GSfield{Statement}{Faddeev\textquoteright s factorization uses modified Green functions and a fixed direction \(\gamma\). Provided the corresponding Fredholm equation is uniquely solvable, the resulting factors are triangular with respect to the angular order and recover \(S\); their diagonal data have the specified complex-direction continuation. The 1966 note assumes unique solvability and explicitly proves it only in a sufficiently small-potential regime.}
\GSfield{Reconstruction}{The notebook signs and factor order depend on incoming/outgoing conventions. The uniqueness hypothesis is material; the construction is not an unconditional global inverse theorem for all decaying potentials.}
\GSfield{Fate}{The precise prior source is \citet[pp.~71--72, equations (10)--(13)]{faddeev1966factorization}, received 3 June 1965 and published in 1966.}
\GSfield{Present status}{As of November 2026, the reconstructed statement is established; no open problem is asserted.}
\GSfield{Assessment}{literal statement: correct with missing hypotheses; publication match: partial; verification: read the result; development: none; confidence: high}
\end{GSitem}
\begin{GSitem}{1970-219.10}{F}{The discrete-spectrum obstruction in the inverse problem}
\GSfield{Kind; source}{remark; \GSpage{P000344}}
\GSfield{The notes}{The Green function and factorization had been written down earlier, but \ru{Обратной задачи не было из-за дискретного спектра}.}
\GSfield{Statement}{This is a historical report about the method\textquoteright s unresolved analytic continuation and exceptional parameters. It does not specify a modern inverse-data problem, a uniqueness class, or a claimed obstruction for every potential.}
\GSfield{Reconstruction}{It is not converted into a present-day open inverse-scattering conjecture.}
\GSfield{Fate}{The earlier paper explicitly left singularities of the continued modified solutions uncontrolled \citep[p.~72]{faddeev1966factorization}; the notebook does not identify a stronger exact question.}
\GSfield{Present status}{As of November 2026, the historical limitation is identifiable, but no definite surviving open problem follows from this sentence.}
\GSfield{Assessment}{literal statement: not determinable; publication match: partial; verification: read the result; development: none; confidence: medium}
\end{GSitem}

\GStalk{220}{1971-03-15}{A. M. Gabrielov}{P000344-P000345}{Rational periods and tubes about a hypersurface}
Following Griffiths, rational differential forms are integrated over cycles avoiding their polar hypersurface. Projective compactification and a tube construction explain which cycles are relevant. Record 222 continues the homological and pole-reduction argument.

\begin{GSitem}{1970-220.1}{E}{Periods of rational differential forms}
\GSfield{Kind; source}{definition; \GSpage{P000344}}
\GSfield{The notes}{\(\omega=P(z)Q(z)^{-1}dz_1\wedge\cdots\wedge dz_n\), \(V=\{Q=0\}\), and \(\int_\Gamma\omega\) for \(n\)-cycles in the complement.}
\GSfield{Statement}{A closed rational top form on \(U=\mathbb C^n\setminus V\) pairs with \(H_n(U;\mathbb Z)\) by integration. After projective compactification all polar components, including a possible hyperplane at infinity, must be included in the divisor.}
\GSfield{Reconstruction}{Closedness of a holomorphic top form is automatic on \(U\); a projective denominator must account for infinity rather than just the affine zero set.}
\GSfield{Fate}{The bibliography identifies \citet[Part~I]{griffiths1969periods1} as the prior source; the exact subject is stated in the notebook itself.}
\GSfield{Present status}{As of November 2026, the reconstructed statement is established; no open problem is asserted.}
\GSfield{Assessment}{literal statement: correct with missing hypotheses; publication match: exact; verification: publisher metadata checked; development: none; confidence: high}
\end{GSitem}
\begin{GSitem}{1970-220.2}{E}{One-variable pole reduction and residues}
\GSfield{Kind; source}{method; \GSpage{P000344}}
\GSfield{The notes}{Homology is generated by small loops around poles; higher-order poles are removed by subtracting exact differentials.}
\GSfield{Statement}{For a rational one-form on the punctured complex line, partial fractions split it into an exact rational differential plus a sum of simple-pole forms. Its period around a cycle is \(2\pi i\) times the integral linear combination of the enclosed residues.}
\GSfield{Reconstruction}{The conventional \(2\pi i\) is restored; polynomial terms and all poles of order at least two contribute exact differentials.}
\GSfield{Fate}{This elementary prototype is the starting point of the rational-period exposition \citep{griffiths1969periods1}.}
\GSfield{Present status}{As of November 2026, the reconstructed statement is established; no open problem is asserted.}
\GSfield{Assessment}{literal statement: correct with missing hypotheses; publication match: exact; verification: publisher metadata checked; development: none; confidence: high}
\end{GSitem}
\begin{GSitem}{1970-220.3}{E}{Middle-dimensional complement cycles vanish in projective space}
\GSfield{Kind; source}{theorem; \GSpage{P000345}}
\GSfield{The notes}{\(H_n(\mathbb P^n\setminus V;\mathbb Z)\to H_n(\mathbb P^n;\mathbb Z)\) is zero; the even-dimensional case uses Bézout.}
\GSfield{Statement}{For a nonempty smooth hypersurface \(V\subset\mathbb P^n(\mathbb C)\) of degree \(q>0\), inclusion sends every middle-dimensional complement cycle to zero in projective homology. For odd \(n\), the target is zero. For \(n=2m\), intersecting with \(V\cap\mathbb P^{m+1}\), whose class is \(q[\mathbb P^m]\), forces its coefficient to vanish.}
\GSfield{Reconstruction}{Positive degree and the torsion-free ambient group are explicit. The notebook\textquoteright s intersection proof is retained, rather than inferred from vanishing of a different cohomology group.}
\GSfield{Fate}{The projective-complement framework belongs to \citet{griffiths1969periods1}; the complete elementary argument is present in the notes and repeated in record 222.}
\GSfield{Present status}{As of November 2026, the reconstructed statement is established; no open problem is asserted.}
\GSfield{Assessment}{literal statement: correct with missing hypotheses; publication match: exact; verification: publisher metadata checked; development: none; confidence: high}
\end{GSitem}
\begin{GSitem}{1970-220.4}{E}{The tube map and its surjectivity in middle dimension}
\GSfield{Kind; source}{construction; \GSpage{P000345}}
\GSfield{The notes}{If \(\Gamma=\partial R\), take \(\gamma=R\cap V\); the circle tube \(\tau(\gamma)\) is homologous to \(\Gamma\).}
\GSfield{Statement}{For smooth \(V\), the oriented real rank-two normal bundle gives a tube homomorphism \(\tau:H_k(V;\mathbb Z)\to H_{k+1}(\mathbb P^n\setminus V;\mathbb Z)\). In degree \(k=n-1\), it is surjective by the preceding vanishing and the long exact sequence of the pair plus the Thom isomorphism.}
\GSfield{Reconstruction}{The tube is the normal circle bundle over the cycle, not an arbitrary full neighborhood of the cycle in projective space. Chain orientations fix a possible conventional boundary sign.}
\GSfield{Fate}{This is the residue/tube construction underlying \citet{griffiths1969periods1}; record 222 supplies the exact sequence.}
\GSfield{Present status}{As of November 2026, the reconstructed statement is established; no open problem is asserted.}
\GSfield{Assessment}{literal statement: correct with missing hypotheses; publication match: exact; verification: publisher metadata checked; development: none; confidence: high}
\end{GSitem}

\GStalk{221}{1971-03-15}{V. I. Arnold}{P000346}{Real plane curves and the Gudkov congruence}
Harnack and Petrovskii bounds lead to Gudkov\textquoteright s modulo-eight conjecture for maximal curves. Arnold reports the modulo-four result and links it to characteristic elements of integral intersection forms on a branched double cover. The stronger congruence was proved by Rokhlin in a note received in December 1971.

\begin{GSitem}{1970-221.1}{E}{Harnack’s bound and real sextic configurations}
\GSfield{Kind; source}{theorem; \GSpage{P000346}}
\GSfield{The notes}{For degree \(n=2k\), the number of ovals is at most \(g+1\), \(g=(n-1)(n-2)/2\). A separate note says degree six has “50 cases (Gudkov)” and sketches an eleven-oval arrangement.}
\GSfield{Statement}{A nonsingular real projective plane curve of even degree \(2k\) has at most \((2k-1)(k-1)+1\) connected components, all ovals. For degree six the maximum is eleven. The unqualified number fifty does not specify whether empty curves, isotopy schemes, or a partial classification are counted.}
\GSfield{Reconstruction}{The numerical shorthand is not substituted for the number of all real schemes. The maximal-oval illustration is separated from that unresolved bibliographical count.}
\GSfield{Fate}{Arnold recalls Harnack and Gudkov\textquoteright s maximal sextics in \citet[\S1]{arnold1971ovals}. The exact interpretation of the notebook\textquoteright s “50” was not established.}
\GSfield{Present status}{As of November 2026, the reconstructed statement is established; no open problem is asserted.}
\GSfield{Assessment}{literal statement: correct with missing hypotheses; publication match: partial; verification: read the result; development: none; confidence: medium}
\end{GSitem}
\begin{GSitem}{1970-221.2}{E}{Petrovskii’s inequalities for even-degree plane curves}
\GSfield{Kind; source}{theorem; \GSpage{P000346}}
\GSfield{The notes}{Writing \(p,m\) for ovals inside an even or odd number of other ovals gives \(|2(p-m)-1|\le3k^2-3k+1\).}
\GSfield{Statement}{For a nonsingular real projective plane curve of degree \(2k\), let \(p\) and \(m\) denote the numbers of even and odd ovals. Petrovskii\textquoteright s inequalities are equivalent to the displayed double inequality. Maximality is not needed for this bound.}
\GSfield{Reconstruction}{“Positive oval” in the notes means even nesting parity, not orientation of the oval.}
\GSfield{Fate}{The prior inequalities are recalled in \citet[\S1]{arnold1971ovals}.}
\GSfield{Present status}{As of November 2026, the reconstructed statement is established; no open problem is asserted.}
\GSfield{Assessment}{literal statement: correct with missing hypotheses; publication match: exact; verification: read the result; development: none; confidence: high}
\end{GSitem}
\begin{GSitem}{1970-221.3}{B}{Gudkov’s congruence modulo eight}
\GSfield{Kind; source}{conjecture; \GSpage{P000346}}
\GSfield{The notes}{\(p-m\equiv k^2\pmod8\), following the discussion of curves with the maximal number of ovals.}
\GSfield{Statement}{For a nonsingular real projective \(M\)-curve of degree \(2k\), meaning it has \((2k-1)(k-1)+1\) components, one has \(p-m\equiv k^2\pmod8\). Without the \(M\)-curve hypothesis this is not the asserted theorem.}
\GSfield{Reconstruction}{The maximality hypothesis is carried forward from the preceding paragraph and stated explicitly.}
\GSfield{Fate}{Rokhlin proved the exact congruence \citep[\S1]{rokhlin1972gudkov}; his note was received 29 December 1971 and published in 1972.}
\GSfield{Present status}{As of November 2026, solved: this is the Gudkov--Rokhlin congruence, not an unresolved part of Hilbert' s sixteenth problem.}
\GSfield{Assessment}{literal statement: correct with missing hypotheses; publication match: exact; verification: read the result; development: none; confidence: high; problem \hyperref[prob:GP-1970-51]{GP-1970-51}}
\end{GSitem}
\begin{GSitem}{1970-221.4}{A}{Arnold–Varchenko’s modulo-four result}
\GSfield{Kind; source}{theorem; \GSpage{P000346}}
\GSfield{The notes}{Arnold and Varchenko have proved \(p-m\equiv k^2\pmod4\) using quadratic forms.}
\GSfield{Statement}{For every nonsingular real projective \(M\)-curve of even degree \(2k\), the weaker congruence modulo four holds.}
\GSfield{Reconstruction}{The published paper is authored by Arnold; its final discussion acknowledges Varchenko\textquoteright s initial argument and the development of the proof. Joint spoken attribution is not silently changed into joint publication authorship.}
\GSfield{Fate}{\citet[\S1 and final acknowledgments]{arnold1971ovals}, received 22 March 1971, publishes the result after the 15 March lecture.}
\GSfield{Present status}{As of November 2026, the reconstructed statement is established; no open problem is asserted.}
\GSfield{Assessment}{literal statement: correct with missing hypotheses; publication match: exact; verification: read the result; development: none; confidence: high; problem \hyperref[prob:GP-1970-51]{GP-1970-51}}
\end{GSitem}
\begin{GSitem}{1970-221.5}{E}{Characteristic squares are well defined modulo eight}
\GSfield{Kind; source}{claim; \GSpage{P000346}}
\GSfield{The notes}{For a characteristic \(u\), \((u+2x)^2=u^2+4(u,x)+4x^2\), and the two latter terms have matching parity.}
\GSfield{Statement}{For an integral symmetric bilinear form, if \((u,x)\equiv(x,x)\pmod2\) for every \(x\), then \((u+2x,u+2x)\equiv(u,u)\pmod8\). This calculation proves independence of the characteristic lift, not yet the equality with the signature.}
\GSfield{Reconstruction}{The distinction between lift independence and the full van der Blij lemma prevents treating the displayed two-line calculation as a proof of both assertions.}
\GSfield{Fate}{The arithmetic tool is used in \citet{arnold1971ovals} and developed explicitly in record 225.}
\GSfield{Present status}{As of November 2026, the reconstructed statement is established; no open problem is asserted.}
\GSfield{Assessment}{literal statement: correct with missing hypotheses; publication match: exact; verification: read the result; development: none; confidence: high}
\end{GSitem}
\begin{GSitem}{1970-221.6}{A}{Branched double covers connect real curves with intersection forms}
\GSfield{Kind; source}{method; \GSpage{P000346}}
\GSfield{The notes}{The notes introduce \(z^2=f(x,y)\) and identify \(p-m\) with an Euler characteristic of a region of the real plane.}
\GSfield{Statement}{For a smooth real plane curve of even degree, the corresponding double cover of \(\mathbb P^2(\mathbb C)\), branched along its complexification, is a compact smooth complex surface. Real involutions and their fixed loci connect nesting data to its integral intersection form; characteristic-vector congruences then constrain the ovals.}
\GSfield{Reconstruction}{The affine equation is compactified via the square root line bundle \(\mathcal O(k)\). Smoothness of the branch curve is required. The full topology is sourced, not claimed proved by the single notebook line.}
\GSfield{Fate}{This method is developed in \citet[\S\S2--8]{arnold1971ovals}; Rokhlin\textquoteright s stronger proof explicitly builds on it \citep[\S2]{rokhlin1972gudkov}.}
\GSfield{Present status}{As of November 2026, the reconstructed statement is established; no open problem is asserted.}
\GSfield{Assessment}{literal statement: correct with missing hypotheses; publication match: exact; verification: read the result; development: none; confidence: high}
\end{GSitem}

\GStalk{222}{1971-03-22}{A. M. Gabrielov}{P000347-P000349}{The tube exact sequence and Griffiths reduction}
The tube-map kernel is identified with the image of ambient projective homology. Homogeneous coordinates then give rational projective forms and a Jacobian-ideal formula reducing pole order. Two displayed formulas need explicit corrections: kernel means image of the intersection map, and the numerator degree depends on the pole order.

\begin{GSitem}{1970-222.1}{E}{The integral tube exact sequence and the kernel question}
\GSfield{Kind; source}{theorem; \GSpage{P000347-P000348}}
\GSfield{The notes}{The question \(\ker\tau=?\) is answered by \(H_{n+1}(\mathbb P^n)\to H_{n-1}(V)\to H_n(\mathbb P^n\setminus V)\to H_n(\mathbb P^n)\), followed by \(\ker\tau=H_{n+1}(\mathbb P^n)\).}
\GSfield{Statement}{For a smooth positive-degree hypersurface, the pair sequence and the oriented Thom isomorphism give an exact sequence over \(\mathbb Z\). Precisely, \(\ker\tau=\operatorname{im}(H_{n+1}(\mathbb P^n;\mathbb Z)\xrightarrow{\cap V}H_{n-1}(V;\mathbb Z))\). Thus the kernel is zero for even \(n\); for \(n=2m+1\), it is generated by the hyperplane-section class \([V\cap\mathbb P^{m+1}]\).}
\GSfield{Reconstruction}{The image, not the domain as a literal subgroup equality, is the kernel. Exactness needs no passage modulo torsion. The ambient source is injective here because intersecting its image with a suitable power of the hyperplane class has nonzero degree.}
\GSfield{Fate}{This continues the tube construction in \citet{griffiths1969periods1}; the proof is reconstructed directly from the exact sequence already written in the notes.}
\GSfield{Present status}{As of November 2026, the reconstructed statement is established; no open problem is asserted.}
\GSfield{Assessment}{literal statement: false as written; publication match: partial; verification: publisher metadata checked; development: none; confidence: high}
\end{GSitem}
\begin{GSitem}{1970-222.2}{E}{A closed chain need not bound}
\GSfield{Kind; source}{method; \GSpage{P000348}}
\GSfield{The notes}{The proof says \(\partial(T\gamma-R)=0\), hence \(T\gamma-R\) is a boundary.}
\GSfield{Statement}{The closed chain \(T\gamma-R\) defines an ambient homology class. Its intersection with \(V\) gives \(\gamma\), establishing membership in the image of the intersection map. There is no need, and in general no justification, to declare the ambient cycle a boundary.}
\GSfield{Reconstruction}{This repairs the notebook\textquoteright s proof of exactness without changing the intended sequence.}
\GSfield{Fate}{The corrected chain argument is the ordinary proof of the pair/Thom exact sequence used in the rational-period construction \citep{griffiths1969periods1}.}
\GSfield{Present status}{As of November 2026, the reconstructed statement is established; no open problem is asserted.}
\GSfield{Assessment}{literal statement: false as written; publication match: partial; verification: publisher metadata checked; development: none; confidence: high}
\end{GSitem}
\begin{GSitem}{1970-222.3}{E}{Homogeneous rational forms on projective space}
\GSfield{Kind; source}{construction; \GSpage{P000348-P000349}}
\GSfield{The notes}{\(\Omega=\sum(-1)^j\xi_j\,d\xi_0\wedge\cdots\widehat{d\xi_j}\cdots\wedge d\xi_n\), and \(P\Omega/Q^\ell\) descends to projective space.}
\GSfield{Statement}{If \(Q\) is homogeneous of degree \(q\) and \(P\) of degree \(\ell q-(n+1)\), then \(P\Omega/Q^\ell\) is horizontal and invariant under \(\mathbb C^*\), hence descends to a rational \(n\)-form on \(\mathbb P^n\). The formula \(\deg P=q-(n+1)\) applies only to a simple pole.}
\GSfield{Reconstruction}{The pole-order factor \(\ell\) is missing from the final repeated degree formula on P000349; the earlier simple-pole formula on P000348 is correct.}
\GSfield{Fate}{This is the homogeneous form setup of \citet{griffiths1969periods1}; its degree is independently checked by scaling every coordinate.}
\GSfield{Present status}{As of November 2026, the reconstructed statement is established; no open problem is asserted.}
\GSfield{Assessment}{literal statement: false as written; publication match: exact; verification: publisher metadata checked; development: none; confidence: high}
\end{GSitem}
\begin{GSitem}{1970-222.4}{E}{Lower-degree basic forms by contraction with the Euler field}
\GSfield{Kind; source}{construction; \GSpage{P000348}}
\GSfield{The notes}{A more general alternating expression omits several \(d\xi_j\), and gives \(\deg B=\deg A+n+2-\ell\).}
\GSfield{Statement}{Let \(E=\sum\xi_i\partial_{\xi_i}\) and let \(\eta\) be a polynomial \((r+1)\)-form with homogeneous coefficients of degree \(a\). If \(B\) has degree \(a+r+1\), then \((\iota_E\eta)/B\) is a horizontal, scaling-invariant rational \(r\)-form. With \(r=n+1-\ell\), its degree condition is exactly the one in the notes.}
\GSfield{Reconstruction}{The lengthy coordinate formula is interpreted as Euler contraction; this states a clean sufficient descent construction rather than an unchecked claim about every possible presentation.}
\GSfield{Fate}{The basic-form construction is part of the projective rational-form calculus cited as \citet{griffiths1969periods1}; scaling and \(\iota_E^2=0\) check the asserted property.}
\GSfield{Present status}{As of November 2026, the reconstructed statement is established; no open problem is asserted.}
\GSfield{Assessment}{literal statement: correct with missing hypotheses; publication match: exact; verification: publisher metadata checked; development: none; confidence: high}
\end{GSitem}
\begin{GSitem}{1970-222.5}{E}{Griffiths reduction of a Jacobian numerator}
\GSfield{Kind; source}{method; \GSpage{P000349}}
\GSfield{The notes}{Subtracting a coboundary yields \(P=\ell\sum A_jQ_j-Q\sum\partial_jA_j\).}
\GSfield{Statement}{For homogeneous \(Q\) of degree \(q\), \(\ell\ge1\), and homogeneous \(A_j\) of degree \(\ell q-n\), one has \(\frac{(\sum A_j\partial_jQ)\Omega}{Q^{\ell+1}}\equiv\frac{1}{\ell}\frac{(\sum\partial_jA_j)\Omega}{Q^\ell}\pmod{d\Omega^{n-1}(\mathbb P^n\setminus V)}\). Thus a numerator in the Jacobian ideal can be reduced by one pole order. A general numerator need not be in that ideal.}
\GSfield{Reconstruction}{The two numerator degrees are distinguished, and the reduction is a cohomological congruence. Smoothness of \(V\) supports the subsequent finite Jacobian-ring description, although the displayed differentiation identity itself is algebraic.}
\GSfield{Fate}{This is the Griffiths pole-reduction mechanism, explicitly attributed in the notes to \citet{griffiths1969periods1}. Exact differentiation supplies the stated reconstruction; the original 1969 full text was not successfully retrieved.}
\GSfield{Present status}{As of November 2026, the reconstructed statement is established; no open problem is asserted.}
\GSfield{Assessment}{literal statement: correct with missing hypotheses; publication match: exact; verification: publisher metadata checked; development: none; confidence: high}
\end{GSitem}

\GStalk{223}{1971-03-22}{V. P. Palamodov}{P000349}{Characteristics of differential systems}
Palamodov argues that differential consequences must be included before assigning a characteristic variety to a system. The first example is corrupted as a mathematical formula even though the transcription follows the scan. The underdetermined divergence equation illustrates the difference between analytic existence and unique determination by Cauchy data.

\begin{GSitem}{1970-223.1}{E}{Characteristics must use the differential module of the system}
\GSfield{Kind; source}{method; \GSpage{P000349}}
\GSfield{The notes}{\ru{Надо выписать все следствия и с ними иметь дело}: all differential consequences must be included.}
\GSfield{Statement}{For an analytic linear differential system represented by a coherent module \(\mathcal M\) over the sheaf of differential operators, the characteristic support is the support of \(\operatorname{gr}\mathcal M\) for a good filtration. Taking only the symbols of an arbitrary initial generating list may omit principal symbols of differential consequences.}
\GSfield{Reconstruction}{The coherent differential-module formulation is a modern reconstruction of the method, not terminology asserted to occur in the notes.}
\GSfield{Fate}{Palamodov\textquoteright s earlier work on differential operators in coherent analytic sheaves is a thematic source \citep{palamodov1968coherent}; the notebook does not identify a particular theorem from it.}
\GSfield{Present status}{As of November 2026, the reconstructed statement is established; no open problem is asserted.}
\GSfield{Assessment}{literal statement: correct with missing hypotheses; publication match: thematic only; verification: publisher metadata checked; development: none; confidence: medium}
\end{GSitem}
\begin{GSitem}{1970-223.2}{F}{The displayed two-equation example is inconsistent}
\GSfield{Kind; source}{claim; \GSpage{P000349}}
\GSfield{The notes}{The system \(D_xu=f\), \((D_y+\lambda)u=g\) is followed by \(u=D_yg-D_xf\). The scan confirms the written \(\lambda\).}
\GSfield{Statement}{For constant \(\lambda\) and commuting partial derivatives, the actual compatibility is \(D_xg-D_yf=\lambda f\); it does not recover \(u\). An instructive different example would use \((D_y+x)u=g\), when \([D_x,D_y+x]u=u\) gives \(u=D_xg-(D_y+x)f\).}
\GSfield{Reconstruction}{The latter system is explicitly an editorial illustration, not a proposed transcription of \(\lambda\). Both the notebook coefficient and its consequent fail to yield the claimed relation.}
\GSfield{Fate}{The general differential-consequence principle matches Palamodov\textquoteright s work \citep{palamodov1968coherent}; the intended specific example cannot be identified from the surviving formula.}
\GSfield{Present status}{As of November 2026, the literal countercalculation is decisive, but the author' s intended example is not recoverable.}
\GSfield{Assessment}{literal statement: false as written; publication match: thematic only; verification: unresolved lead; development: none; confidence: high}
\end{GSitem}
\begin{GSitem}{1970-223.3}{E}{Underdetermined divergence separates existence from uniqueness}
\GSfield{Kind; source}{claim; \GSpage{P000349}}
\GSfield{The notes}{For \(D_xu+D_yv=f\), Holmgren\textquoteright s theorem works in no direction, while Cauchy--Kovalevskaya works in every direction.}
\GSfield{Statement}{For analytic \(f\) in two variables, in any chosen nonzero normal direction one can solve for a normal derivative of one component after prescribing the other analytic component, giving local analytic solutions with compatible data. The equation alone does not uniquely determine both \(u\) and \(v\) from data on a hypersurface; hence a scalar Holmgren uniqueness conclusion cannot simply be transferred to this underdetermined system.}
\GSfield{Reconstruction}{The statement is an existence/uniqueness contrast, not a contradiction between the two theorems. The arbitrary free component must be included.}
\GSfield{Fate}{The example illustrates the theory of systems underlying \citet{palamodov1968coherent}; no exact publication match for this pedagogical example was found.}
\GSfield{Present status}{As of November 2026, the reconstructed statement is established; no open problem is asserted.}
\GSfield{Assessment}{literal statement: correct with missing hypotheses; publication match: thematic only; verification: publisher metadata checked; development: none; confidence: medium}
\end{GSitem}

\GStalk{224}{1971-03-29}{S. Rolewicz}{P000350}{Linear orbits and the hypercyclic weighted shift}
Finite-dimensional orbit behavior is contrasted with the dense orbit of a multiple of the unilateral backward shift. The concatenation construction is sketched, with additional remarks on topological equivalence and algebraic operators unbounded in every norm.

\begin{GSitem}{1970-224.1}{E}{Finite-dimensional orbit alternatives}
\GSfield{Kind; source}{claim; \GSpage{P000350}}
\GSfield{The notes}{The limit behavior of \(x,Ax,A^2x,\ldots\) is described as zero, infinity, or a compact set excluding zero.}
\GSfield{Statement}{For a linear map on a finite-dimensional real or complex normed space and a fixed vector, either \(A^nx\to0\), or \(\|A^nx\|\to\infty\), or the orbit is bounded and its omega-limit set is nonempty compact and avoids zero. This follows by separating generalized eigenspaces with spectral radius below, equal to, or above one, and the polynomial growth of nontrivial Jordan chains.}
\GSfield{Reconstruction}{“Limit points” means asymptotic behavior; infinity is not an ordinary point of the vector space, and the full orbit closure also contains initial transient points.}
\GSfield{Fate}{This contrast supplies the finite-dimensional background to Rolewicz\textquoteright s previously published example \citep{rolewicz1969orbits}.}
\GSfield{Present status}{As of November 2026, the reconstructed statement is established; no open problem is asserted.}
\GSfield{Assessment}{literal statement: correct with missing hypotheses; publication match: exact; verification: publisher metadata checked; development: none; confidence: high}
\end{GSitem}
\begin{GSitem}{1970-224.2}{E}{Rolewicz’s hypercyclic multiple of the backward shift}
\GSfield{Kind; source}{theorem; \GSpage{P000350}}
\GSfield{The notes}{\(S(x_1,x_2,\ldots)=(x_2,x_3,\ldots)\), \(A=aS\), \(a>1\), can have an everywhere dense orbit on \(\ell^p\) or \(c^0\).}
\GSfield{Statement}{On \(\ell^p(\mathbb N)\), \(1\le p<\infty\), or \(c_0(\mathbb N)\), over \(\mathbb R\) or \(\mathbb C\), the bounded operator \(\lambda S\) is hypercyclic when \(|\lambda|>1\). The separability restriction excludes \(\ell^\infty\).}
\GSfield{Reconstruction}{The index range and complex analogue \(|\lambda|>1\) are supplied. The Hilbert-space example is \(p=2\); no theorem about every Hilbert space dimension is intended.}
\GSfield{Fate}{\citet[pp.~17--22]{rolewicz1969orbits} appeared in 1969; the author\textquoteright s own institutional account identifies it as his hypercyclic example.}
\GSfield{Present status}{As of November 2026, the reconstructed statement is established; no open problem is asserted.}
\GSfield{Assessment}{literal statement: correct with missing hypotheses; publication match: exact; verification: publisher metadata checked; development: none; confidence: high; problem \hyperref[prob:GP-1970-52]{GP-1970-52}}
\end{GSitem}
\begin{GSitem}{1970-224.3}{E}{Concatenating small blocks to obtain a dense orbit}
\GSfield{Kind; source}{construction; \GSpage{P000350}}
\GSfield{The notes}{A list of finite rational blocks is concatenated after long zero gaps, with block \(j\) divided by a suitable power of \(a\).}
\GSfield{Statement}{Choose a dense sequence of finite-support vectors \(v_j\) with each target repeated, and increasing positions \(N_j\). Put the block \(\lambda^{-N_j}v_j\) at position \(N_j\), choosing later positions so that the tail after applying \((\lambda S)^{N_j}\) has norm below \(2^{-j}\) and all scaled block norms are summable. Then \((\lambda S)^{N_j}x-v_j\to0\), proving density.}
\GSfield{Reconstruction}{The long gaps must control amplified later tails as well as place blocks without overlap; this quantitative choice fills the notebook\textquoteright s proof sketch. For complex spaces use rational real and imaginary parts.}
\GSfield{Fate}{The construction explains the prior result \citep{rolewicz1969orbits}; the detailed tail condition is an editorial unpacking of the notebook, not a new theorem.}
\GSfield{Present status}{As of November 2026, the reconstructed statement is established; no open problem is asserted.}
\GSfield{Assessment}{literal statement: correct with missing hypotheses; publication match: exact; verification: publisher metadata checked; development: none; confidence: high; problem \hyperref[prob:GP-1970-52]{GP-1970-52}}
\end{GSitem}
\begin{GSitem}{1970-224.4}{E}{Topological conjugacy preserves a dense orbit}
\GSfield{Kind; source}{definition; \GSpage{P000350}}
\GSfield{The notes}{\(A=UBU^{-1}\) for a homeomorphism \(U\); orbit density is a topological property.}
\GSfield{Statement}{If \(U:X\to Y\) is a homeomorphism with \(UA=BU\), then \(U(\{A^nx\})=\{B^nU(x)\}\). Therefore hypercyclicity is invariant under topological conjugacy, even when \(U\) is nonlinear.}
\GSfield{Reconstruction}{No claim of linear similarity follows from topological equivalence.}
\GSfield{Fate}{The invariant is an immediate consequence of the definition and contextualizes \citet{rolewicz1969orbits}.}
\GSfield{Present status}{As of November 2026, the reconstructed statement is established; no open problem is asserted.}
\GSfield{Assessment}{literal statement: correct; publication match: exact; verification: publisher metadata checked; development: none; confidence: high}
\end{GSitem}
\begin{GSitem}{1970-224.5}{F}{An algebraic operator unbounded in every norm}
\GSfield{Kind; source}{claim; \GSpage{P000350}}
\GSfield{The notes}{One can construct an operator unbounded in every norm.}
\GSfield{Statement}{On the algebraic vector space of finite-support sequences, define \(Te_n=ne_n\). For every norm on this space, a bound \(\|Tx\|\le C\|x\|\) would imply \(n\le C\) for all \(n\), so no such bound exists.}
\GSfield{Reconstruction}{The example is editorial and proves a natural literal existential reading. The notes give no ambient space, completeness requirement, or specific operator, so they do not identify this particular construction.}
\GSfield{Fate}{The sentence was searched together with Rolewicz\textquoteright s algebraic-operator and orbit work; no exact publication or intended ambient space was established. The institutional bibliography records \emph{Equations in linear spaces} (1968), but that broad title is not an exact match.}
\GSfield{Present status}{As of November 2026, the existential algebraic reading is elementary and true; the intended generality is unidentified, not an open problem.}
\GSfield{Assessment}{literal statement: correct with missing hypotheses; publication match: none found; verification: unresolved lead; development: none; confidence: low}
\end{GSitem}

\GStalk{225}{1971-03-29}{V. I. Arnold}{P000350-P000351}{Unimodular integral forms and characteristic vectors}
Arnold continues the arithmetic used for real curves: characteristic elements, the signature congruence modulo eight, and the classification of indefinite unimodular forms. Definite forms obey a finiteness theorem but are not classified merely by rank and signature. Serre\textquoteright s arithmetic text is recommended.

\begin{GSitem}{1970-225.1}{E}{Unimodular lattices, parity, and characteristic elements}
\GSfield{Kind; source}{definition; \GSpage{P000350-P000351}}
\GSfield{The notes}{\(\varphi:\mathbb Z^n\times\mathbb Z^n\to\mathbb Z\) is unimodular when it identifies the lattice with its dual. Rank, signature, and types I/II are listed.}
\GSfield{Statement}{An integral symmetric nondegenerate form is unimodular when its determinant is \(\pm1\). It is even if \((x,x)\) is even for every \(x\), and odd otherwise. Because reduction modulo two is nondegenerate, the linear functional \(x\mapsto(x,x)\bmod2\) has a unique representing class \(u\in L/2L\), called characteristic.}
\GSfield{Reconstruction}{Symmetry and nondegeneracy are supplied. “Type I” denotes odd and “type II” even; the characteristic vector is unique only modulo \(2L\).}
\GSfield{Fate}{The notebook\textquoteright s cited arithmetic text treats these forms in Chapter~V \citep{serre1970arithmetic}. Earlier lectures underlying the book date to 1962 and 1964.}
\GSfield{Present status}{As of November 2026, the reconstructed statement is established; no open problem is asserted.}
\GSfield{Assessment}{literal statement: correct with missing hypotheses; publication match: exact; verification: publisher metadata checked; development: none; confidence: high}
\end{GSitem}
\begin{GSitem}{1970-225.2}{E}{Van der Blij’s characteristic-signature congruence}
\GSfield{Kind; source}{theorem; \GSpage{P000351}}
\GSfield{The notes}{\(\varphi(u,u)\bmod8\) is independent of the lift and equals \(\operatorname{sgn}\varphi\bmod8\); a Fresnel integral is mentioned.}
\GSfield{Statement}{For an integral unimodular symmetric bilinear lattice and any characteristic vector \(u\), \((u,u)\equiv\sigma(L)\pmod8\). In particular, an even unimodular lattice has signature divisible by eight.}
\GSfield{Reconstruction}{Unimodularity is needed for the signature assertion; the simpler lift-independence calculation is separately recorded in 1970-221.5. The mention of a Fresnel integral is a proof method, not a supplied analytic proof.}
\GSfield{Fate}{The classical arithmetic result is used explicitly by \citet{arnold1971ovals} and belongs to the unimodular-form treatment cited as \citet[Chapter~V]{serre1970arithmetic}.}
\GSfield{Present status}{As of November 2026, the reconstructed statement is established; no open problem is asserted.}
\GSfield{Assessment}{literal statement: correct with missing hypotheses; publication match: exact; verification: read the result; development: none; confidence: high}
\end{GSitem}
\begin{GSitem}{1970-225.3}{E}{Classification of indefinite integral unimodular forms}
\GSfield{Kind; source}{theorem; \GSpage{P000350-P000351}}
\GSfield{The notes}{Rank, signature, and parity are complete invariants in the indefinite case; odd forms diagonalize, and even forms use \(U\) and \(\Gamma_8\).}
\GSfield{Statement}{An odd indefinite unimodular lattice of signature \((p,q)\), \(p,q>0\), is \(I_{p,q}\). An even indefinite unimodular lattice exists exactly when \(p-q\equiv0\pmod8\) and is unique for that signature; it is an orthogonal sum of hyperbolic planes and copies of signed \(E_8\).}
\GSfield{Reconstruction}{The sign of \(E_8\) is chosen to realize either signature. The indefinite hypothesis is indispensable; rank and signature alone do not distinguish parity.}
\GSfield{Fate}{The exposition points to \citet[Chapter~V]{serre1970arithmetic}; the publisher\textquoteright s description confirms this chapter\textquoteright s subject, but its full theorem text was not retrieved.}
\GSfield{Present status}{As of November 2026, the reconstructed statement is established; no open problem is asserted.}
\GSfield{Assessment}{literal statement: correct with missing hypotheses; publication match: exact; verification: publisher metadata checked; development: none; confidence: high}
\end{GSitem}
\begin{GSitem}{1970-225.4}{E}{Finiteness in the positive-definite case}
\GSfield{Kind; source}{theorem; \GSpage{P000351}}
\GSfield{The notes}{Positive-definite unimodular forms are finite in number in every dimension.}
\GSfield{Statement}{For each fixed rank \(n\), there are finitely many isometry classes of positive-definite integral unimodular lattices. This is finiteness of classes, not uniqueness or an explicit classification by rank and signature.}
\GSfield{Reconstruction}{Equivalence means integral lattice isometry, and the determinant is fixed at one. Without fixing determinant the assertion would be false.}
\GSfield{Fate}{This is the definite reduction-theory counterpart of the Chapter~V discussion in \citet{serre1970arithmetic}.}
\GSfield{Present status}{As of November 2026, the reconstructed statement is established; no open problem is asserted.}
\GSfield{Assessment}{literal statement: correct with missing hypotheses; publication match: exact; verification: publisher metadata checked; development: none; confidence: high}
\end{GSitem}
\begin{GSitem}{1970-225.5}{F}{What kind of stabilization erases parity?}
\GSfield{Kind; source}{remark; \GSpage{P000351}}
\GSfield{The notes}{The final line says that stably only rank and signature remain, followed by “Dear Misha you are wrong.”}
\GSfield{Statement}{If stabilization means adding only hyperbolic even planes \(U\), parity remains invariant, so rank and signature do not suffice. If both forms may be augmented by the same odd indefinite plane \(\langle1\rangle\oplus\langle-1\rangle\), then forms of equal rank and signature become odd indefinite and hence isometric by the classification theorem.}
\GSfield{Reconstruction}{The notes do not identify the permitted stabilizations or the target of the handwritten disagreement. The two precise readings are kept separate.}
\GSfield{Fate}{The arithmetic classification in \citet[Chapter~V]{serre1970arithmetic} explains both readings, but does not identify which was intended here.}
\GSfield{Present status}{As of November 2026, the ambiguity is terminological and does not produce a new unsolved classification problem.}
\GSfield{Assessment}{literal statement: not determinable; publication match: thematic only; verification: publisher metadata checked; development: none; confidence: medium}
\end{GSitem}

\GStalk{226}{1971-04-05}{A. A. Kirillov}{P000352-P000353}{Sites, sheaves, and representations}
Kirillov presents material attributed to Geronimus: coverings in a category, the sheaf condition, and cohomology. For a finite group the site of finite group-sets recovers ordinary representations. Geronimus\textquoteright s short 1971 announcement and detailed article identify the Lie-group site and its comparisons with classifying-space and differentiable cohomology. The notebook confuses general morphisms with the submersions used in covering families.

\begin{GSitem}{1970-226.1}{E}{Fiber products and their universal property}
\GSfield{Kind; source}{definition; \GSpage{P000352}}
\GSfield{The notes}{The object \(B\times_A C\) completes a commutative square and every compatible pair of maps to \(B,C\) factors through it uniquely.}
\GSfield{Statement}{For maps \(B\to A\leftarrow C\), a fiber product represents pairs \((X\to B,X\to C)\) whose composites to \(A\) agree. Its structure maps and unique factorization are part of the definition.}
\GSfield{Reconstruction}{Compatibility with the maps to \(A\) is explicit; an arbitrary pair of maps does not suffice. This reprises item 1967-143.1.}
\GSfield{Fate}{This standard prerequisite is used in \citet[\S1]{geronimus1971topology} and the site axioms \citep[Tag 00VH]{stacks2026sites}.}
\GSfield{Present status}{As of November 2026, the reconstructed statement is established; no open problem is asserted.}
\GSfield{Assessment}{literal statement: correct; publication match: exact; verification: read the result; development: none; confidence: high}
\end{GSitem}
\begin{GSitem}{1970-226.2}{E}{Covering families define a Grothendieck pretopology}
\GSfield{Kind; source}{definition; \GSpage{P000352}}
\GSfield{The notes}{The conditions are composition of coverings, stability under base change, and \(U\to U\) being a cover; the last axiom number is overwritten.}
\GSfield{Statement}{A pretopology specifies covering families, contains identity covers, is stable under pullback, and is closed under composing a cover with covers of its members. Required pullbacks must exist; these families generate a Grothendieck topology by covering sieves.}
\GSfield{Reconstruction}{The family axioms describe a pretopology. The overwritten numeral is immaterial. Related definitions occur in 1967-143.2 and 1968-151.1.}
\GSfield{Fate}{Geronimus uses this framework \citep[\S1]{geronimus1971topology}; precise modern family axioms are in \citet[Tag 00VH, Definition 7.6.2]{stacks2026sites}.}
\GSfield{Present status}{As of November 2026, the reconstructed statement is established; no open problem is asserted.}
\GSfield{Assessment}{literal statement: correct with missing hypotheses; publication match: exact; verification: read the result; development: none; confidence: high}
\end{GSitem}
\begin{GSitem}{1970-226.3}{E}{The sheaf axiom is an equalizer condition}
\GSfield{Kind; source}{definition; \GSpage{P000352}}
\GSfield{The notes}{A contravariant functor \(F:C\to K\) is followed by \(0\to F(A)\to\prod F(U_\alpha)\rightrightarrows\prod F(U_\alpha\times_AU_\beta)\).}
\GSfield{Statement}{For a set-valued presheaf, \(F(A)\) must be the equalizer of the restriction maps to overlaps for every cover. In an abelian target this is \(0\to F(A)\to\prod F(U_\alpha)\xrightarrow{d}\prod F(U_\alpha\times_AU_\beta)\), with \(d\) the difference of restrictions. General target categories need the stated limits.}
\GSfield{Reconstruction}{Zero and exact-sequence notation cannot literally be imposed on an arbitrary category \(K\).}
\GSfield{Fate}{The matching condition is recorded by \citet[Tag 00VM, Definition 7.7.1]{stacks2026sites} and used by \citet[\S1]{geronimus1971topology}.}
\GSfield{Present status}{As of November 2026, the reconstructed statement is established; no open problem is asserted.}
\GSfield{Assessment}{literal statement: correct with missing hypotheses; publication match: exact; verification: read the result; development: none; confidence: high}
\end{GSitem}
\begin{GSitem}{1970-226.4}{E}{Čech cohomology and associated sheaves}
\GSfield{Kind; source}{construction; \GSpage{P000353}}
\GSfield{The notes}{For a cover \(\mathcal U\), one forms \(H(\mathcal U,F)\); a marked nontrivial theorem says that the preceding construction produces a sheaf.}
\GSfield{Statement}{For an abelian presheaf, overlap products give a Čech complex for each cover. The matching-family plus construction, applied twice, gives the associated sheaf. Čech cohomology of one cover need not equal derived sheaf cohomology without acyclicity hypotheses; a spectral sequence relates them.}
\GSfield{Reconstruction}{The unspecified construction is matched cautiously with Geronimus\textquoteright s sheafification discussion. No claim that one cover computes all cohomology is inserted.}
\GSfield{Fate}{\citet[Propositions 1.3 and 1.6]{geronimus1971topology} supplies the associated-sheaf construction and cohomological spectral sequence.}
\GSfield{Present status}{As of November 2026, the reconstructed statement is established; no open problem is asserted.}
\GSfield{Assessment}{literal statement: correct with missing hypotheses; publication match: partial; verification: read the result; development: none; confidence: medium}
\end{GSitem}
\begin{GSitem}{1970-226.5}{E}{Geronimus’s site of smooth group actions}
\GSfield{Kind; source}{construction; \GSpage{P000353}}
\GSfield{The notes}{Objects are \(G\)-manifolds and morphisms are called regular \(G\)-maps with surjective tangent maps.}
\GSfield{Statement}{For a Lie group \(G\), take smooth \(G\)-manifolds with all smooth equivariant maps as morphisms. A family is a cover when its members are equivariant submersions whose images jointly cover the target. Covers are stable under base change and composition.}
\GSfield{Reconstruction}{Submersion is a condition on covering arrows, not every morphism.}
\GSfield{Fate}{The short announcement \citet[\S0]{geronimus1971lie}, received 28 August 1970 and published in the first 1971 issue, precedes this April exposition. The detailed construction is \citet[\S3.1]{geronimus1971topology}, received 22 September 1970 and published later in 1971.}
\GSfield{Present status}{As of November 2026, the reconstructed statement is established; no open problem is asserted.}
\GSfield{Assessment}{literal statement: false as written; publication match: exact; verification: read the result; development: none; confidence: high}
\end{GSitem}
\begin{GSitem}{1970-226.6}{E}{Finite-group representations as sheaves on the group-set site}
\GSfield{Kind; source}{theorem; \GSpage{P000353}}
\GSfield{The notes}{For finite \(G\), take finite \(G\)-sets and vector spaces; a representation \(T\) gives tuples satisfying \(\xi_{g i}=T_g\xi_i\). All sheaves arise this way.}
\GSfield{Statement}{On finite left \(G\)-sets with jointly surjective equivariant covers, sheaves of vector spaces over a fixed field are equivalent to linear representations of \(G\). A module \(E\) gives \(F_E(X)=\operatorname{Map}_G(X,E)\); evaluation at the regular \(G\)-set recovers \(E\) and its action.}
\GSfield{Reconstruction}{The topology and coefficient field are supplied. The five-point illustration does not restrict the equivalence to five-element sets.}
\GSfield{Fate}{Geronimus identifies this older finite/discrete-group construction and credits Artin in the introduction and \S2 \citep{geronimus1971topology}; \citet[Theorem 1]{geronimus1971lie} records its cohomological comparison.}
\GSfield{Present status}{As of November 2026, the reconstructed statement is established; no open problem is asserted.}
\GSfield{Assessment}{literal statement: correct with missing hypotheses; publication match: exact; verification: read the result; development: none; confidence: high}
\end{GSitem}
\begin{GSitem}{1970-226.7}{E}{A Lie-group representation gives an equivariant-function sheaf}
\GSfield{Kind; source}{construction; \GSpage{P000353}}
\GSfield{The notes}{From \(X,T_g\) construct vector-valued functions \(\xi\), with \(T_g\xi(xg^{-1})=\xi(x)\). There are sheaves not obtained from representations.}
\GSfield{Statement}{For a continuous Banach representation \(E\) of a Lie group, smooth equivariant maps \(X\to E\) form a sheaf on the submersion site. Its value on \(G\) is the space \(E^\infty\) of smooth vectors. The finite-group equivalence does not justify an equivalence with all continuous Banach representations in the Lie-group case.}
\GSfield{Reconstruction}{Left and right conventions are reconciled by equivariance. The wider non-equivalence claim depends on the category of representations, which the notes omit; it is not asserted for every generalized notion of representation.}
\GSfield{Fate}{\citet[\S7.1 and Proposition 7.3]{geronimus1971topology} gives the function sheaf and recovers the Gårding space. The comparison is announced in \citet[Theorem 3]{geronimus1971lie}.}
\GSfield{Present status}{As of November 2026, the reconstructed statement is established; no open problem is asserted.}
\GSfield{Assessment}{literal statement: correct with missing hypotheses; publication match: partial; verification: read the result; development: none; confidence: medium}
\end{GSitem}
\begin{GSitem}{1970-226.8}{E}{Classifying-space and differentiable cohomology comparisons}
\GSfield{Kind; source}{theorem; \GSpage{P000353}}
\GSfield{The notes}{Sometimes the construction gives cohomology of a topological space, sometimes cohomology as in Gelfand--Fuchs.}
\GSfield{Statement}{For the group-action site of a connected Lie group \(G\), the sheaf of locally constant functions with values in a fixed discrete abelian group \(A\), with trivial action on the coefficient group, gives \(H^*(BG;A)\). For the smooth equivariant-function sheaf of a continuous Banach representation, site cohomology agrees with the differentiable group cohomology denoted \(H_{\rm alg}^*(G;E^\infty)\) by Geronimus.}
\GSfield{Reconstruction}{The coefficient group is abelian; the two comparisons use distinct coefficient sheaves. Connectedness is retained for the announced Lie-group comparison. No general identification of Lie-algebra and classifying-space cohomology is inserted.}
\GSfield{Fate}{The comparisons appear in \citet[Theorems 2 and 3]{geronimus1971lie}; the detailed classifying-space statement is \citet[Proposition 4.3]{geronimus1971topology}.}
\GSfield{Present status}{As of November 2026, the reconstructed statement is established; no open problem is asserted.}
\GSfield{Assessment}{literal statement: correct with missing hypotheses; publication match: partial; verification: read the result; development: none; confidence: high}
\end{GSitem}
\begin{GSitem}{1970-226.9}{F}{Representation theory and the classification of particles}
\GSfield{Kind; source}{remark; \GSpage{P000353}}
\GSfield{The notes}{Elementary-particle theory is incomplete because not all representations are known; spinors and electrons are mentioned.}
\GSfield{Statement}{This motivates connecting particle descriptions with representation theory. No symmetry group, representation class, or mathematical completeness assertion is specified.}
\GSfield{Reconstruction}{The remark is preserved without replacing it by a claim that physical theories are determined solely by a list of representations.}
\GSfield{Fate}{The surrounding Geronimus construction supplies context \citep{geronimus1971topology}, but no precise mathematical claim or publication match follows.}
\GSfield{Present status}{As of November 2026, this is a research motivation, not a formulated mathematical problem.}
\GSfield{Assessment}{literal statement: not determinable; publication match: thematic only; verification: unresolved lead; development: none; confidence: low}
\end{GSitem}

\GStalk{227}{1971-04-05}{I. M. Gelfand}{P000354-P000355}{The elliptic-curve modularity conjecture}
The \emph{Weil conjecture} here is modularity of elliptic curves over \(\mathbb Q\), expressed by a modular parametrization, local Euler factors, and integrals of a weight-two modular form. It is not the Weil conjectures for varieties over finite fields. The handwritten cubic exponent is retained rather than silently normalized, and the treatment of conductors explicitly starts with \(p\ne2,3\).

\begin{GSitem}{1970-227.1}{E}{The elliptic curve and the handwritten cubic exponent}
\GSfield{Kind; source}{definition; \GSpage{P000354}}
\GSfield{The notes}{The scan reads \(y^3=x^3-g_1x-g_2\), with rational coefficients, and calls the curve elliptic.}
\GSfield{Statement}{A smooth projective genus-one curve over \(\mathbb Q\) with a rational point is an elliptic curve. The written homogeneous cubic \(Y^3=X^3-g_1XZ^2-g_2Z^3\), if smooth, has the rational point \([1:1:0]\), so it can indeed define an elliptic curve. The familiar short Weierstrass equation instead has \(y^2\).}
\GSfield{Reconstruction}{The exponent three is confirmed by the scan and is not automatically a mathematical error. Smoothness is missing; subsequent reduction is interpreted after passage to a minimal Weierstrass equation. No silent replacement of three by two is made.}
\GSfield{Fate}{The Weierstrass-model conventions used for the subsequent arithmetic are in \citet[\S3.1]{cremona1997algorithms}. The intended equation is not determined beyond the surviving cubic.}
\GSfield{Present status}{As of November 2026, the reconstructed statement is established; no open problem is asserted.}
\GSfield{Assessment}{literal statement: correct with missing hypotheses; publication match: partial; verification: read the result; development: none; confidence: medium}
\end{GSitem}
\begin{GSitem}{1970-227.2}{E}{The congruence subgroup and compact modular curve}
\GSfield{Kind; source}{construction; \GSpage{P000354}}
\GSfield{The notes}{\(\Gamma_0(A)\) consists of matrices \(\left(\begin{smallmatrix}\alpha&\beta\\A\gamma&\delta\end{smallmatrix}\right)\), but its displayed action is \((\alpha z+A\gamma)/(\beta z+\delta)\).}
\GSfield{Statement}{For \(N\ge1\), \(\Gamma_0(N)=\{\left(\begin{smallmatrix}a&b\\c&d\end{smallmatrix}\right)\in\mathrm{SL}_2(\mathbb Z):N\mid c\}\) acts on \(\mathfrak H\) by \(z\mapsto(az+b)/(cz+d)\). Adding cusps gives the compact modular curve \(X_0(N)\).}
\GSfield{Reconstruction}{The off-diagonal entries are interchanged in the notebook action. The transposed subgroup is conjugate, but the written matrix and action cannot both be kept under the standard convention. A fundamental region is not yet the compact curve.}
\GSfield{Fate}{The conventions are in \citet[\S\S2.1.1 and 2.15.1]{cremona1997algorithms}.}
\GSfield{Present status}{As of November 2026, the reconstructed statement is established; no open problem is asserted.}
\GSfield{Assessment}{literal statement: false as written; publication match: exact; verification: read the result; development: none; confidence: high}
\end{GSitem}
\begin{GSitem}{1970-227.3}{B}{Modularity of every elliptic curve over the rationals}
\GSfield{Kind; source}{conjecture; \GSpage{P000354-P000355}}
\GSfield{The notes}{The Weil conjecture asserts a nontrivial analytic map \(X_A\to T\), with level \(A\) chosen as the conductor.}
\GSfield{Statement}{For every elliptic curve \(E/\mathbb Q\) of conductor \(N\), there is a normalized weight-two cuspidal newform on \(\Gamma_0(N)\) with \(L(E,s)=L(f,s)\). Equivalently, \(E\) admits a nonconstant modular parametrization \(X_0(N)\to E\) over \(\mathbb Q\).}
\GSfield{Reconstruction}{Nonconstant, compactification at cusps, and the arithmetic field are explicit. This is elliptic-curve modularity, not the Weil conjectures for varieties over finite fields.}
\GSfield{Fate}{\citet[Theorem A and Introduction]{breuilconraddiamondtaylor2001modularity} proves modularity for all \(E/\mathbb Q\), completing the earlier semistable case. The integral parametrization is explained in \citet[\S\S2.6 and 2.15.1]{cremona1997algorithms}.}
\GSfield{Present status}{As of November 2026, solved by the modularity theorem; the full result was published in 2001.}
\GSfield{Assessment}{literal statement: correct with missing hypotheses; publication match: exact; verification: read the result; development: none; confidence: high; problem \hyperref[prob:GP-1970-53]{GP-1970-53}}
\end{GSitem}
\begin{GSitem}{1970-227.4}{F}{There are not finitely many unrestricted maps to a torus}
\GSfield{Kind; source}{claim; \GSpage{P000354}}
\GSfield{The notes}{A parenthesis says there are finitely many maps from every curve to a torus.}
\GSfield{Statement}{Without normalization this is false: translating a map by arbitrary target points gives infinitely many maps, even constant ones. If a nonconstant \(\varphi:C\to E\) exists, the maps \([m]\circ\varphi\), \(m\ge1\), remain distinct after fixing a base point sent to zero, since their degrees grow as \(m^2\deg\varphi\).}
\GSfield{Reconstruction}{The intended restriction, perhaps fixed degree or a quotient by operations on the target, is unspecified. These counterexamples do not require identifying that intention.}
\GSfield{Fate}{The modular-parametrization construction \citep[\S2.15.1]{cremona1997algorithms} supplies examples of maps but does not support the unrestricted finiteness sentence.}
\GSfield{Present status}{As of November 2026, the literal assertion is false; an intended historical finiteness theorem cannot be uniquely recovered.}
\GSfield{Assessment}{literal statement: false as written; publication match: none found; verification: indirect; development: none; confidence: high}
\end{GSitem}
\begin{GSitem}{1970-227.5}{E}{Reduction types and conductor exponents}
\GSfield{Kind; source}{definition; \GSpage{P000354-P000355}}
\GSfield{The notes}{For \(p\ne2,3\), smooth reduction contributes \(A_p=1\), distinct tangents give \(A_p=p\), and coincident tangents give \(A_p=p^2\); \(A=\prod A_p\).}
\GSfield{Statement}{For a minimal integral Weierstrass model over \(\mathbb Q_p\), good, multiplicative, and additive reduction have conductor exponents \(0,1,2\), respectively, for \(p\ge5\). Multiplicative reduction is split or nonsplit according to whether its tangent directions are defined over \(\mathbb F_p\). At \(p=2,3\), wild ramification can increase the additive exponent.}
\GSfield{Reconstruction}{Minimality is essential. The global conductor includes exceptional primes, whose exponents the simple table does not determine. Split/nonsplit replaces real/imaginary over finite fields.}
\GSfield{Fate}{Reduction and conductor computations are in \citet[\S\S3.1--3.2]{cremona1997algorithms}, including Tate\textquoteright s algorithm at exceptional primes. This modern source documents prior background, not a work claimed to precede the talk.}
\GSfield{Present status}{As of November 2026, the reconstructed statement is established; no open problem is asserted.}
\GSfield{Assessment}{literal statement: correct with missing hypotheses; publication match: exact; verification: read the result; development: none; confidence: high}
\end{GSitem}
\begin{GSitem}{1970-227.6}{E}{Local Euler factors of the elliptic-curve L-function}
\GSfield{Kind; source}{construction; \GSpage{P000355}}
\GSfield{The notes}{At good primes \(L_p(s)=(1-c_pp^{-s}+p^{1-2s})^{-1}\), with \(\#E(\mathbb F_p)=1-c_p+p\); multiplicative reduction gives signs \(\pm1\), and additive reduction the factor one.}
\GSfield{Statement}{At a good prime set \(a_p=p+1-\#E(\mathbb F_p)\), so \(L_p(E,s)=(1-a_pp^{-s}+p^{1-2s})^{-1}\). Split and nonsplit multiplicative primes give \((1-p^{-s})^{-1}\) and \((1+p^{-s})^{-1}\); additive primes give one. Their product defines \(L(E,s)\) in its initial half-plane of absolute convergence.}
\GSfield{Reconstruction}{The factors use a minimal model and the correct reduction type. This elliptic \(L\)-function is not the full zeta function of a curve.}
\GSfield{Fate}{The local data and matching coefficients are in \citet[\S\S2.6 and 3.9]{cremona1997algorithms}; equality with the modular \(L\)-series follows from \citet{breuilconraddiamondtaylor2001modularity}.}
\GSfield{Present status}{As of November 2026, the reconstructed statement is established; no open problem is asserted.}
\GSfield{Assessment}{literal statement: correct with missing hypotheses; publication match: exact; verification: read the result; development: none; confidence: high}
\end{GSitem}
\begin{GSitem}{1970-227.7}{B}{From Dirichlet coefficients to a weight-two modular form}
\GSfield{Kind; source}{construction; \GSpage{P000355}}
\GSfield{The notes}{Expand \(\zeta(s)=\sum c_n/n_s\), then write \(\omega(z)=\sum c_ne^{inz}\).}
\GSfield{Statement}{Write \(L(E,s)=\sum_{n\ge1}a_nn^{-s}\). Modularity states that \(f(z)=\sum_{n\ge1}a_ne^{2\pi inz}\) is the normalized weight-two newform of level \(N(E)\). Its invariant differential is \(2\pi i f(z)\,dz\).}
\GSfield{Reconstruction}{The visible \(n_s\) is mathematically emended to \(n^s\). The factor \(2\pi\) uses the standard unit-period coordinate; \(e^{inz}\) requires rescaling coordinate and action. Forming the Fourier series does not prove modularity.}
\GSfield{Fate}{The newform assertion is part of \citet[Introduction, Theorem A]{breuilconraddiamondtaylor2001modularity}; normalization is in \citet[\S2.6]{cremona1997algorithms}.}
\GSfield{Present status}{As of November 2026, the reconstructed statement is established; no open problem is asserted.}
\GSfield{Assessment}{literal statement: false as written; publication match: exact; verification: read the result; development: none; confidence: high; problem \hyperref[prob:GP-1970-53]{GP-1970-53}}
\end{GSitem}
\begin{GSitem}{1970-227.8}{B}{Integrating the modular differential gives the parametrization}
\GSfield{Kind; source}{method; \GSpage{P000355}}
\GSfield{The notes}{\(\int\omega(z)\,dz\) is the sought map, and \(\int_z^{gz}\omega(z)\,dz\) should lie in a lattice for \(g\in\Gamma_0(A)\).}
\GSfield{Statement}{The rational-coefficient newform attached to \(E\) has a rank-two period lattice \(\Lambda_f\). Integration from a base cusp defines \(X_0(N)\to\mathbb C/\Lambda_f\). This target is isogenous to \(E\); composition with an appropriate isogeny gives a parametrization of \(E\). Under \(\Gamma_0(N)\), the integral changes by a period.}
\GSfield{Reconstruction}{The base point, invariant differential, quotient, and possible isogeny are supplied. The assertion does not hold for an arbitrary integer-coefficient Fourier series.}
\GSfield{Fate}{\citet[\S2.15.1]{cremona1997algorithms} gives the period construction. Its applicability to every rational elliptic curve follows from \citet{breuilconraddiamondtaylor2001modularity}.}
\GSfield{Present status}{As of November 2026, the reconstructed statement is established; no open problem is asserted.}
\GSfield{Assessment}{literal statement: correct with missing hypotheses; publication match: exact; verification: read the result; development: none; confidence: high; problem \hyperref[prob:GP-1970-53]{GP-1970-53}}
\end{GSitem}

\GStalk{228}{1971-04-12}{A. M. Perelomov}{P000356-P000359}{Theta functions, coherent lattices, and uncertainty}
Integral symplectic lattices and quadratic refinements lead to theta-function identities and completeness of lattice coherent states. The critical-density deletion result was already published in Perelomov\textquoteright s February 1971 paper. A question about bounded centered position--momentum uncertainty for an orthogonal basis has a positive general answer; restrictions preserving a Gabor lattice change the answer. The final leaf P000360 is an adjacent out-of-scope record: Kirillov on 6 September 1971, not a continuation of this April session.

\begin{GSitem}{1970-228.1}{E}{A positive compatible complex structure and a principal lattice}
\GSfield{Kind; source}{definition; \GSpage{P000356}}
\GSfield{The notes}{\(B\) is nondegenerate skew and integral on \(L\); \(J^2=-1\), \(J^TBJ=B\), \(BJ>0\). The Hermitian form is \(H(z,w)=B(z,Jw)+iB(z,w)\).}
\GSfield{Statement}{Let \((V,B)\) be real symplectic of dimension \(2N\), \(L\) a full integral lattice, and \(J\) a compatible complex structure with \(B(v,Jv)>0\) for nonzero \(v\). Then \(H\) is positive Hermitian, conjugate-linear in its first slot. Principality means \(B:L\to L^*\) is unimodular.}
\GSfield{Reconstruction}{Volume one uses the symplectic normalization. Compatibility alone does not imply positivity; this is a separate condition, and nondegeneracy makes nonnegative positivity strict.}
\GSfield{Fate}{These conventions are in \citet[\S\S1.1 and 1.2(c)]{cartier1968theta}; the notebook also refers to Cartier\textquoteright s earlier exposition \citep{cartier1966theta}.}
\GSfield{Present status}{As of November 2026, the reconstructed statement is established; no open problem is asserted.}
\GSfield{Assessment}{literal statement: correct with missing hypotheses; publication match: exact; verification: read the result; development: none; confidence: high}
\end{GSitem}
\begin{GSitem}{1970-228.2}{E}{Quadratic refinements of the integral symplectic pairing}
\GSfield{Kind; source}{construction; \GSpage{P000356}}
\GSfield{The notes}{\(F(\alpha+\beta)=F(\alpha)+F(\beta)+B(\alpha,\beta)\pmod2\); adding linear forms gives \(4^N\) solutions.}
\GSfield{Statement}{Refinements of the alternating pairing on \(L/2L\) form an affine space over \(\operatorname{Hom}(L,\mathbb F_2)\). In a symplectic basis, \(F(m,n)=m\cdot n+\varepsilon\cdot m+\widetilde\varepsilon\cdot n\), with \(\varepsilon,\widetilde\varepsilon\in\mathbb F_2^N\); there are \(2^{2N}=4^N\) choices.}
\GSfield{Reconstruction}{Only values modulo two matter; integer lifts give the same signs. The extra prime on one source coefficient does not affect the count.}
\GSfield{Fate}{\citet[\S1.1]{cartier1968theta} develops this description.}
\GSfield{Present status}{As of November 2026, the reconstructed statement is established; no open problem is asserted.}
\GSfield{Assessment}{literal statement: correct; publication match: exact; verification: read the result; development: none; confidence: high}
\end{GSitem}
\begin{GSitem}{1970-228.3}{E}{Theta automorphy and its one-dimensional solution space}
\GSfield{Kind; source}{theorem; \GSpage{P000356}}
\GSfield{The notes}{Entire \(\Theta\) satisfies \(\Theta(z+\alpha)=(-1)^{F(\alpha)}e^{\pi H(\alpha,\alpha)/2+\pi H(\alpha,z)}\Theta(z)\), uniquely up to a constant for a principal lattice.}
\GSfield{Statement}{For the positive compatible principally polarized data in 1970-228.1--2, the entire functions satisfying this law form a one-dimensional vector space. Nonzero solutions are unique up to nonzero scalar.}
\GSfield{Reconstruction}{Positivity and principality are essential. Linearity of \(H\) in its second slot makes the factor holomorphic in \(z\); no conjugate of \(z\) should be inserted.}
\GSfield{Fate}{The theta realization appears in \citet[\S1.3]{cartier1968theta} and is used in \citet[\S5]{perelomov1971completeness}.}
\GSfield{Present status}{As of November 2026, the reconstructed statement is established; no open problem is asserted.}
\GSfield{Assessment}{literal statement: correct; publication match: exact; verification: read the result; development: none; confidence: high}
\end{GSitem}
\begin{GSitem}{1970-228.4}{F}{The idempotent anecdote attributed to Pauli}
\GSfield{Kind; source}{remark; \GSpage{P000356}; I. M. Gelfand}
\GSfield{The notes}{Gelfand writes \(P^2=P\), \(\lambda^2=\lambda\), then \(\lambda=0,+1,-1\), with \(-1\) discarded for physical reasons.}
\GSfield{Statement}{Over \(\mathbb R\) or \(\mathbb C\), an eigenvalue of an idempotent satisfies \(\lambda(\lambda-1)=0\), hence is zero or one. Minus one is excluded algebraically.}
\GSfield{Reconstruction}{This is an anecdote or joke, not silently changed into a different operator equation. No Pauli volume or passage is identified.}
\GSfield{Fate}{No precise Pauli publication match was recovered from the book-and-idempotent wording. The correction follows directly from the written polynomial.}
\GSfield{Present status}{As of November 2026, the algebraic statement is settled; the anecdote' s target is unidentified.}
\GSfield{Assessment}{literal statement: false as written; publication match: none found; verification: unresolved lead; development: none; confidence: high}
\end{GSitem}
\begin{GSitem}{1970-228.5}{E}{Coherent states, the Bargmann space, and its evaluation kernel}
\GSfield{Kind; source}{construction; \GSpage{P000357-P000358}}
\GSfield{The notes}{\(a|\alpha\rangle=\alpha|\alpha\rangle\), labels fill \(\mathbb C\), and \(|\langle\alpha|\beta\rangle|^2=e^{-|\alpha-\beta|^2}\). Orthogonality means an entire function vanishes at the chosen labels.}
\GSfield{Statement}{In \(\mathcal F^2=\{f\text{ entire}:\pi^{-1}\int_{\mathbb C}|f(z)|^2e^{-|z|^2}\,dA<\infty\}\), normalized kernels are \(k_\alpha(z)=e^{-|\alpha|^2/2+\bar\alpha z}\). With the inner product linear in the second slot, \(\langle k_\alpha,f\rangle=e^{-|\alpha|^2/2}f(\alpha)\). Their span is dense exactly when the labels form a uniqueness set. A finite accumulation point of distinct labels suffices by the identity theorem.}
\GSfield{Reconstruction}{The conjugation convention and positive area measure are fixed. The Gaussian overlap kernel is positive definite, also by multiplying a Gram kernel entrywise with its conjugate; the marginal Neumann--Schoenberg reference supplies no distinct new theorem.}
\GSfield{Fate}{The reduction and coherent-state formulas appear in \citet[\S\S1--2 and 4]{perelomov1971completeness}, published in February 1971.}
\GSfield{Present status}{As of November 2026, the reconstructed statement is established; no open problem is asserted.}
\GSfield{Assessment}{literal statement: correct with missing hypotheses; publication match: exact; verification: read the result; development: none; confidence: high; problem \hyperref[prob:GP-1970-54]{GP-1970-54}}
\end{GSitem}
\begin{GSitem}{1970-228.6}{E}{The critical area and the exact deletion threshold}
\GSfield{Kind; source}{theorem; \GSpage{P000357}}
\GSfield{The notes}{Cell area greater than \(\pi\) gives incompleteness, smaller area overcompleteness. At area \(\pi\), deleting one state leaves completeness, deleting two destroys it.}
\GSfield{Statement}{For a full lattice \(\Lambda\subset\mathbb C\) with fundamental area \(S\), coherent states are incomplete for \(S>\pi\), and complete after every finite deletion for \(S<\pi\). At \(S=\pi\), the full system and every one-point deletion are complete, whereas every two-point deletion is incomplete.}
\GSfield{Reconstruction}{The abbreviation means deleting points in context, not selecting only one or two. The critical excess-one statement is more precise than a density threshold.}
\GSfield{Fate}{\citet[Assertion 1, \S4]{perelomov1971completeness} proves these exact alternatives. Received 20 August 1970, it appeared in February 1971; its 2002 arXiv date is a later archival posting.}
\GSfield{Present status}{As of November 2026, the reconstructed statement is established; no open problem is asserted.}
\GSfield{Assessment}{literal statement: correct; publication match: exact; verification: read the result; development: none; confidence: high; problem \hyperref[prob:GP-1970-54]{GP-1970-54}}
\end{GSitem}
\begin{GSitem}{1970-228.7}{E}{Perelomov’s critical-lattice theta identity}
\GSfield{Kind; source}{theorem; \GSpage{P000357-P000358}}
\GSfield{The notes}{For \(\lambda_{mn}=m\omega_1+n\omega_2\) at area \(\pi\), \(\sum_{m,n}(-1)^{mn+m+n}e^{-|\lambda_{mn}|^2/2+\lambda_{mn}z}=0\); the notes call this one relation.}
\GSfield{Statement}{If \(|\operatorname{Im}(\overline{\omega_1}\omega_2)|=\pi\), this series converges absolutely and uniformly on compact subsets of \(\mathbb C\) and vanishes identically. It is an identity of entire functions. Its coefficient signs are not square summable, so it is not automatically a norm-convergent dependence in the coherent-state Hilbert space.}
\GSfield{Reconstruction}{Reflecting the lattice reconciles kernel conjugations. The informal one-relation statement concerns the completeness defect, not the kernel of an unspecified bounded synthesis operator.}
\GSfield{Fate}{This is \citet[Eq.~(39) and the following theta argument]{perelomov1971completeness}. Whether it follows from a particular Jacobi identity concerns the proof route, not an unresolved identity.}
\GSfield{Present status}{As of November 2026, the reconstructed statement is established; no open problem is asserted.}
\GSfield{Assessment}{literal statement: correct with missing hypotheses; publication match: exact; verification: read the result; development: none; confidence: high; problem \hyperref[prob:GP-1970-54]{GP-1970-54}}
\end{GSitem}
\begin{GSitem}{1970-228.8}{B}{Can an orthonormal basis have a uniformly bounded uncertainty product?}
\GSfield{Kind; source}{question; \GSpage{P000357-P000358}}
\GSfield{The notes}{Von Neumann\textquoteright s question asks whether \(\delta_x\delta_p\) stays bounded after orthogonalizing the critical coherent lattice. The notes then ask the general question for an orthogonal system; Gelfand thinks a positive answer unlikely.}
\GSfield{Statement}{For unit vectors \(e_j\in L^2(\mathbb R)\), use centered standard deviations of position and Fourier momentum. There exists an orthonormal basis with \(\sup_j\Delta_x(e_j)\Delta_p(e_j)<\infty\). There even exists one with both centered second moments uniformly bounded. For a regular single-window Gabor basis, simultaneous finite second moments are impossible.}
\GSfield{Reconstruction}{The scan expressly asks about the product. General basis existence, a particular Gram--Schmidt ordering of coherent states, and preservation of a fixed Gabor lattice are different formulations.}
\GSfield{Fate}{Meyer\textquoteright s Schwartz wavelets give bounded product \citep[\S2, Theorem 1]{meyer1987uncertainty}, lectured February 1986 and published 1987. Bourgain\textquoteright s November 1986 preprint, published 1988, gives separate bounds \citep[Theorem, preprint p.~3]{bourgain1988uncertainty}. The distinction from Balian--Low and higher-moment obstructions appears in \citet[Eq.~(1.1), Theorem 1.1, pp.~115--116]{grochenigmalinnikova2013localization}.}
\GSfield{Present status}{As of November 2026, the unrestricted basis question has a positive established answer. An unspecified orthogonalization algorithm is not settled by this existence theorem.}
\GSfield{Assessment}{literal statement: correct with missing hypotheses; publication match: exact; verification: read the result; development: annotated; confidence: high; problem \hyperref[prob:GP-1970-59]{GP-1970-59}}
\end{GSitem}
\begin{GSitem}{1970-228.9}{E}{Odd theta characteristics give the higher-dimensional identities}
\GSfield{Kind; source}{construction; \GSpage{P000358-P000359}}
\GSfield{The notes}{The count is \(2^{N-1}(2^N-1)\), but the closing paragraph says half the \(4^N\) choices vanish. A Gaussian theta series equals \(c\Theta\).}
\GSfield{Statement}{For a principal lattice, odd refinements number \(2^{N-1}(2^N-1)\), and even ones \(2^{N-1}(2^N+1)\). The Gaussian Poincaré series \(\sum_{\alpha\in L}(-1)^{F(\alpha)}e^{-\pi H(\alpha,\alpha)/2+\pi H(\alpha,z)}\) has the theta automorphy and is a scalar multiple of its theta function. It is even in \(z\); for an odd characteristic the theta function is odd, forcing the scalar to vanish.}
\GSfield{Reconstruction}{The word half is false even for \(N=1\). Parity proves the odd-characteristic identities but does not exclude extra vanishing of even theta constants on special higher-dimensional tori.}
\GSfield{Fate}{The refinement and parity framework appears in \citet[\S\S1.1 and 1.3]{cartier1968theta}; the theta-series argument is used in \citet[\S5]{perelomov1971completeness}.}
\GSfield{Present status}{As of November 2026, the reconstructed statement is established; no open problem is asserted.}
\GSfield{Assessment}{literal statement: false as written; publication match: exact; verification: read the result; development: none; confidence: high}
\end{GSitem}
\begin{GSitem}{1970-228.10}{E}{Weyl operators and the central lattice cocycle}
\GSfield{Kind; source}{construction; \GSpage{P000358-P000359}}
\GSfield{The notes}{Heisenberg matrices are followed by \(D=e^{i\varphi}e^{i(\alpha^*a-\alpha a^*)}\). The subgroup uses lattice \(\alpha+i\beta\) and central coordinate \(k\pi\).}
\GSfield{Statement}{With \([a,a^\dagger]=1\), unitary displacement operators are \(D(\alpha)=\exp(\alpha a^\dagger-\bar\alpha a)\), up to replacing \(\alpha\) by \(-\alpha\). Multiplication carries the symplectic phase cocycle. A lattice lift must incorporate it, equivalently a quadratic refinement; restricting a raw triangular central entry to \(\pi\mathbb Z\) need not give a subgroup for an arbitrary plane lattice.}
\GSfield{Reconstruction}{The extra inner \(i\) turns a skew-adjoint exponent into a self-adjoint one and does not give the claimed unitary operator. Triangular and symmetric group coordinates have different central cocycles.}
\GSfield{Fate}{The displacement and cocycle formulas are in \citet[\S5, Eqs.~(60)--(71)]{perelomov1971completeness}; the lattice lift is described by \citet[\S1.2]{cartier1968theta}.}
\GSfield{Present status}{As of November 2026, the reconstructed statement is established; no open problem is asserted.}
\GSfield{Assessment}{literal statement: false as written; publication match: exact; verification: read the result; development: none; confidence: high}
\end{GSitem}
\begin{GSitem}{1970-228.11}{E}{The lattice model is irreducible after fixing the central character}
\GSfield{Kind; source}{theorem; \GSpage{P000359}}
\GSfield{The notes}{Induction from the signs \((-1)^{mn+m+n}\) on a discrete subgroup is said to give an irreducible representation.}
\GSfield{Statement}{For the Heisenberg group with circle center \(T\), let \(\Gamma\) be the compatible lift of a principal lattice, with \(T\cap\Gamma=\{1\}\). Inducing from \(T\Gamma\) by the character \(\chi(t\gamma)=t\) gives the irreducible lattice realization of the Schrödinger representation with \(U_t=tI\). Induction from the discrete subgroup alone, or replacement by an arbitrary nontrivial central character, does not have this stated multiplicity-one conclusion.}
\GSfield{Reconstruction}{The full center and its tautological character are essential. Principality is relative to this central weight; it does not give multiplicity one for every character \(t\mapsto t^m\).}
\GSfield{Fate}{The exact formulation is \citet[\S1.2(b), pp.~420--421]{cartier1968theta}; the notebook abbreviates the essential central-character condition.}
\GSfield{Present status}{As of November 2026, the reconstructed statement is established; no open problem is asserted.}
\GSfield{Assessment}{literal statement: correct with missing hypotheses; publication match: exact; verification: read the result; development: none; confidence: high}
\end{GSitem}
\begin{GSitem}{1970-228.12}{C}{Completeness for an arbitrary sequence of coherent states}
\GSfield{Kind; source}{question; \GSpage{P000357-P000358}}
\GSfield{The notes}{\ru{Для каких послед.} \(\lambda_n\) \ru{полная система?} The proof asks whether a nonzero Fock function can vanish at all labels.}
\GSfield{Statement}{For a locally finite set \(\Lambda\subset\mathbb C\), characterize when no nonzero \(f\in\mathcal F^2\) vanishes on \(\Lambda\). Equivalently, characterize completeness of its coherent kernels. This is uniqueness, not the stronger stable sampling property.}
\GSfield{Reconstruction}{Distinct labels and local finiteness isolate the nontrivial question. Finite accumulation is covered by 1970-228.5; area alone suffices only in the lattice case.}
\GSfield{Fate}{A verified partial criterion is \citet[Definition 1.1, Theorem 1.2 and Remark 1.3]{mitkovskiwick2013uniqueness}, circulated 3 June 2013. In normalization \(e^{-\pi|z|^2}\), a \(d\)-regular perturbation \(\gamma_{mn}\) of \((m+in)/d\), with \(|m+in-d\gamma_{mn}|\le\phi(|m+in|)\), \(|\gamma_{mn}-\gamma_{m^{\prime}n^{\prime}}|\ge c|(m-m^{\prime})+i(n-n^{\prime})|\) for some \(c>0\), \(\phi(t)/t\) nonincreasing and \(\int_1^\infty\phi(t)t^{-2}dt<\infty\), is unique for \(d>1\), nonunique for \(d<1\); \(d=1\) admits both. A September 2026 zero-set preprint is an unverified lead \citep{houma2026zeros}.}
\GSfield{Present status}{As of November 2026, the lattice and stated perturbative cases are verified. Those results leave arbitrary irregular sets untreated; this audit has not established whether a complete modern characterization is available. The 2026 lead could not be opened, so no global claim of continued openness or resolution is made.}
\GSfield{Assessment}{literal statement: correct; publication match: partial; verification: read the result; development: none; confidence: medium; problem \hyperref[prob:GP-1970-54]{GP-1970-54}}
\GSfield{Worth developing}{A precise critical-density extension could merit work, but further development should first check the inaccessible 2026 lead and other recent classification results; no Stage 2 novelty claim is warranted here.}
\end{GSitem}
\begin{GSitem}{1970-228.13}{F}{A proposed group-orbit generalization}
\GSfield{Kind; source}{remark; \GSpage{P000359}; D. A. Kazhdan (provisional attribution)}
\GSfield{The notes}{An intervention provisionally attributed to Kazhdan starts with \(T:G\to U(H)\), \(h_0\in H\), and discrete \(\Gamma\subset G\). For a reasonable subgroup, \(\{T(\gamma)h_0\}\) should be complete with finitely many relations.}
\GSfield{Statement}{No group or representation class, cyclicity assumption, meaning of reasonable, or permitted topology and coefficients for relations is given. Completeness already fails if \(h_0\) lies in a proper invariant subspace.}
\GSfield{Reconstruction}{The abbreviation after \(G\) is not expanded into a theorem about all semisimple groups. The coherent-state prototype does not determine the missing quantifiers.}
\GSfield{Fate}{The Heisenberg model \citep[\S1.2]{cartier1968theta} and exact coherent-lattice result \citep{perelomov1971completeness} identify the prototype. No exact general-group theorem follows from this intervention.}
\GSfield{Present status}{As of November 2026, the generalization is too underspecified for a definite problem-index entry; no unsolved universal theorem is asserted.}
\GSfield{Assessment}{literal statement: not determinable; publication match: thematic only; verification: unresolved lead; development: none; confidence: low}
\end{GSitem}
\begin{GSitem}{1970-228.14}{E}{Finite-deletion overcompleteness}
\GSfield{Kind; source}{definition; \GSpage{P000357}}
\GSfield{The notes}{Overcomplete means that throwing away any finite number of vectors leaves a complete system.}
\GSfield{Statement}{A total family in a Hilbert space has the finite-deletion property if the complement of every finite subset still has dense linear span. This is stronger than some uses of overcomplete meaning merely redundant or nonminimal.}
\GSfield{Reconstruction}{The notebook explicitly fixes its convention; no uniform frame bounds are implied.}
\GSfield{Fate}{The lattice alternatives in \citet[Assertion 1]{perelomov1971completeness} distinguish this property from critical excess one.}
\GSfield{Present status}{As of November 2026, the reconstructed statement is established; no open problem is asserted.}
\GSfield{Assessment}{literal statement: correct; publication match: exact; verification: read the result; development: none; confidence: high; problem \hyperref[prob:GP-1970-54]{GP-1970-54}}
\end{GSitem}
\begin{GSitem}{1970-228.15}{E}{The unnormalized complex-shift Gaussian series needs correction}
\GSfield{Kind; source}{claim; \GSpage{P000357}}
\GSfield{The notes}{After the entire-function identity the notes write an alternative \(\sum_{m,n}(-1)^{mn+m+n}e^{-(x-\lambda_{mn})^2/2}\), with complex lattice labels.}
\GSfield{Statement}{The literal complex-shift series need not even have terms tending to zero. For normalized oscillator coherent states write \(\alpha=(q+ip)/\sqrt2\) and \(\psi_\alpha(x)=\pi^{-1/4}\exp(-(x-q)^2/2+ip(x-q/2))\). The normalization and phase are essential when transporting a Bargmann identity to the position representation.}
\GSfield{Reconstruction}{The second display is not justified by substituting complex centers into real normalized Gaussians. Its interpretation as a formal vector identity also needs the convergence qualifications in 1970-228.7.}
\GSfield{Fate}{The normalized displacement representation in \citet[\S5, Eqs.~(60)--(71)]{perelomov1971completeness} supplies the correct coherent states. For a nonreal lattice direction with imaginary part larger in absolute value than its real part, the notebook terms grow along its multiples, directly disproving ordinary convergence.}
\GSfield{Present status}{As of November 2026, the reconstructed statement is established; no open problem is asserted.}
\GSfield{Assessment}{literal statement: false as written; publication match: partial; verification: read the result; development: none; confidence: high}
\end{GSitem}

\section{Developments}\label{dev:1969}
\subsection{Interval support for a current group: a precise positive construction}
\label{dev:current-support}

\paragraph{The question and its category.}
On P000291, item 1969-180.2, Gelfand asks for an irreducible unitary
representation of an $SU(2)$ current group whose support is an interval
(problem GP-1969-02). The displayed matrices are functions of
$t\in[0,1]$. The notes do not specify their regularity, the topology on
the group, or the continuity required of a representation. Support is
defined through the implication that a current equal to the identity
near a closed set acts trivially. These omissions prevent a single
category-independent answer. They are not silently repaired here.

We give a complete construction for
\[
 G^1=C^1([0,1],SU(2)),\qquad
 G^\infty=C^\infty([0,1],SU(2)),
\]
with their usual $C^1$ and smooth topologies, respectively, and
\emph{strongly continuous} unitary representations. For either group,
say that a closed set $F$ carries a representation $\pi$ if
\[
 g=1\text{ on a neighbourhood of }F\quad\Longrightarrow\quad\pi(g)=1.
\]
Our representation has no proper closed carrying set, so its support in
this precise sense is $[0,1]$. Its restriction to $G^\infty$ remains
irreducible. It is not strongly continuous for the uniform topology.
This is a bounded contribution to the explicitly reconstructed question,
not a classification of all representations, a claim of historical
priority, or an identification of an unstated topology in the notebook.

\paragraph{What the checked literature establishes.}
For operator-norm-continuous irreducible representations of the smooth
current group, the finite-evaluation classification of
\citet[arXiv version, Thms.~4.11 and 4.15]{janssensneeb2015norm}
gives finite support. The preprint circulated on 11 February 2013;
online publication was in December 2014 and the journal issue in 2015.
The group theorem is numbered 4.16 in the published version.
Ismagilov's energy-representation irreducibility theorem requires base
dimension at least five \citep[\S\S1--2]{ismagilov1976unitary}; his
later segment note proves reducibility in dimension one
\citep[pp.~73--74 in the Russian original]{ismagilov1981segment}.
Neither statement rules out the strongly continuous representation
constructed below. Positive-energy localization theorems impose
additional projective and spacetime-action hypotheses
\citep[Introduction]{janssensneeb2024positive}.

These sources document why the continuity and representation class
matter. They do not document why this particular seminar question was
abandoned, nor justify saying that all strongly continuous interval
representations have now been classified. The search included the
original Math-Net notes, the norm-continuous classification and its
references, and the 2024 localization monograph. No historical claim
that the following construction is new is made.

\begin{lemma}[Uniform smooth interpolation at finitely many points]\ours{}
\label{lem:interpolation}
Let $s_1,\ldots,s_q$ be distinct points of $(0,1)$ and set
\[
 r=\frac13\min\bigl(\{s_i,1-s_i:1\leq i\leq q\}
       \cup\{|s_i-s_j|:i\ne j\}\bigr),\qquad L=12/r.
\]
For any $h_1,\ldots,h_q\in SU(2)$ there is a smooth current $g$,
equal to $1$ near the endpoints, with $g(s_i)=h_i$ and
$\|g'\|_\infty\leq L$. The norm is the operator norm on $2\times2$
matrices. If $q=1$, omit the empty set of pairwise distances.
\end{lemma}

\begin{proof}
Every $h_i$ is $\exp X_i$ for a skew-Hermitian traceless matrix $X_i$
of norm at most $\pi$, by unitary diagonalization of an $SU(2)$ matrix.
Use the smooth bump
\[
 \eta(u)=\begin{cases}
 \exp\bigl(1-(1-u^2)^{-1}\bigr),&|u|<1,\\
 0,&|u|\geq1.
 \end{cases}
\]
It satisfies $\eta(0)=1$ and $\|\eta'\|_\infty<3$: putting
$v=(1-u^2)^{-1}\geq1$, its derivative has absolute value at most
$2v^2e^{1-v}\leq8/e<3$. The intervals $(s_i-r,s_i+r)$ are disjoint
and contained in $(0,1)$. Thus
\[
 g(t)=\prod_{i=1}^q\exp\bigl(\eta((t-s_i)/r)X_i\bigr)
\]
has at most one nonidentity factor at any point. It has the required
values, is smooth also at the bump boundaries, and
$\|g'\|_\infty\leq3\pi/r<12/r$.
\end{proof}

\begin{lemma}[Dense pairs with an interpolation tail bound]\ours{}
\label{lem:pairs}
There are pairwise distinct points $a_n,b_n\in(0,1)$, with
$\{a_n:n\geq1\}$ dense, such that, writing
$\epsilon_n=|a_n-b_n|$ and letting $L_N$ be the preceding lemma's
bound for $a_1,b_1,\ldots,a_N,b_N$,
\[
 \sum_n\epsilon_n<\infty,
 \qquad L_N\sum_{n>N}\epsilon_n\leq 2^{-2N}.
 \tag{1}\label{eq:pairtail}
\]
\end{lemma}

\begin{proof}
Enumerate a countable basis $J_1,J_2,\ldots$ of nonempty open
intervals in $(0,1)$. Choose $a_1,b_1\in J_1$ distinct. Inductively
choose $a_m\in J_m$ outside the preceding finite set, and then
$b_m\in J_m$, also new, so close to $a_m$ that
\[
 0<\epsilon_m\leq 2^{-m}
       \min_{1\leq j<m}\frac{2^{-j}}{1+L_j}\quad(m\geq2).
 \tag{2}\label{eq:pairchoice}
\]
All the bounds on the right have already been fixed. Such a choice
is possible in an open interval, even if those bounds are very small.
For fixed $N$ and $m>N$, the $j=N$ term gives
\[
 \sum_{m>N}\epsilon_m
 \leq\frac{2^{-N}}{1+L_N}\sum_{m>N}2^{-m}
 =\frac{2^{-2N}}{1+L_N}.
\]
This proves (1) and summability. Meeting every basic interval proves
density. There is no assumption of a positive separation for the
infinite point set; separation is used only at each finite stage.
\end{proof}

\begin{theorem}[An irreducible representation with interval support]\ours{}
\label{thm:interval-support}
For either $G^1$ or $G^\infty$ as specified above, there is a
strongly continuous unitary representation on a separable,
infinite-dimensional Hilbert space that is irreducible and has support
$[0,1]$. One representation of $G^1$ has all these properties also
upon restriction to $G^\infty$.
\end{theorem}

\begin{proof}
\emph{The Hilbert space and convergence.}
Put $H_n=\mathbb C^2\otimes\overline{\mathbb C^2}$ and
\[
 \Omega_n=2^{-1/2}(e_1\otimes\bar e_1+e_2\otimes\bar e_2).
\]
Here the overline denotes the conjugate Hilbert space and conjugate
representation. The vector $\Omega_n$ is fixed by
$h\otimes\bar h$ for every $h\in SU(2)$. Let $H$ be the completion
of the increasing union $H_1\otimes\cdots\otimes H_N$, embedded at
the next stage by $v\mapsto v\otimes\Omega_{N+1}$. Write
$\Omega=\bigotimes_n\Omega_n$. This defines the incomplete tensor
product directly and makes separability and infinite dimension clear.

For $g\in G^1$ set
\[
 U_n(g)=g(a_n)\otimes\overline{g(b_n)},\qquad
 R_N(g)=\bigotimes_{n\leq N}U_n(g)\otimes1.
\]
Diagonal invariance and the fundamental theorem of calculus give
\[
 \|U_n(g)\Omega_n-\Omega_n\|
 \leq\|g(a_n)-g(b_n)\|\leq\|g'\|_\infty\epsilon_n.
 \tag{3}\label{eq:vacuum-bound}
\]
For a vector $v$ with excitations only in its first $M$ factors,
telescoping a tensor product of unit vectors gives, for $P>Q\geq M$,
\[
 \|R_P(g)v-R_Q(g)v\|
 \leq\|v\|\,\|g'\|_\infty\sum_{Q<n\leq P}\epsilon_n.
\]
The same estimate holds for an arbitrary vector of the first $M$
tensor factors, since that head factors from the vacuum tail.
Summability therefore gives a strong limit $\pi(g)$ first on this
dense union and then on all of $H$, because all $R_N(g)$ have norm
one. Finite products satisfy $R_N(gh)=R_N(g)R_N(h)$.
Strong limits of uniformly bounded products preserve this identity.
Applying it also to $g^{-1}$ proves that $\pi(g)$ is unitary and
$\pi$ is a representation.

\emph{Strong continuity.}
For a vector supported in the first $M$ factors,
\[
 \|\pi(g)v-R_M(g)v\|
 \leq\|v\|\,\|g'\|_\infty\sum_{n>M}\epsilon_n.
 \tag{4}\label{eq:headtail}
\]
As $g\to1$ in $C^1$, the finite head tends to the identity and
$\|g'\|_\infty\to0$. Thus $\pi(g)v\to v$ on the dense union;
unitarity extends this to every vector. Continuity at other group
elements follows from the homomorphism property. The inclusion
$G^\infty\to G^1$ is continuous, so the restriction is strongly
continuous as well.

\emph{Irreducibility.}
Fix $k$ and arbitrary, independently chosen $h_{a_n},h_{b_n}\in SU(2)$
for $n\leq k$. Consider the local unitary
\[
 V=\bigotimes_{n\leq k}
       (h_{a_n}\otimes\overline{h_{b_n}})\otimes1.
\]
For every $N\geq k$, Lemma~\ref{lem:interpolation} supplies a smooth
$g_N$ taking these prescribed values at the first $2k$ points, and
the value $1$ at the remaining points through $2N$, with
$\|g_N'\|_\infty\leq L_N$. For $v$ supported in the first $M$
factors and $N\geq\max(k,M)$, (4) and (1) give
\[
 \|\pi(g_N)v-Vv\|
 \leq\|v\|L_N\sum_{n>N}\epsilon_n
 \leq\|v\|2^{-2N}\longrightarrow0.
\]
Hence $V$ is a strong limit of representation operators, already
from the smooth subgroup. Any operator commuting with the
representation commutes with every such $V$.

The complex linear span of $SU(2)$ in $M_2(\mathbb C)$ is all of
$M_2(\mathbb C)$: the identity and $i$ times the three Pauli matrices
belong to $SU(2)$ and span it. Independent choices at $a_n,b_n$
therefore span $B(H_n)$, and independent choices through $k$ span
the full matrix algebra on the first $k$ factors. If $T$ commutes
with all these local algebras, it commutes with the projections
\[
 Q_N=\bigotimes_{n\leq N}|\Omega_n\rangle\langle\Omega_n|\otimes1.
\]
On the dense union these projections converge strongly to the
rank-one projection onto $\mathbb C\Omega$; hence they do so on
$H$. Consequently $T\Omega=c\Omega$. The vectors $A\Omega$ for
local matrices $A$ span a dense subspace, and
$TA\Omega=AT\Omega=cA\Omega$. Thus $T=c1$. This proves
irreducibility for both groups.

\emph{Support.}
Let $F$ be a proper closed subset of $[0,1]$. Its complement
contains an open interval with closure disjoint from $F$ and lying
in $(0,1)$. Density supplies $a_n$ in that interval. Choose a smooth
bump current $g$, equal to $1$ outside the interval, with $g(a_n)$
noncentral. On the $n$th local matrix algebra, conjugation by
$\pi(g)$ is exactly conjugation by $U_n(g)$: this is true for
$R_N(g)$ when $N\geq n$ and passes to the strong limits of the
unitaries and their inverses. The matrix $U_n(g)$ is not scalar.
Indeed, if a tensor product of two invertible matrices is scalar,
conjugating $X\otimes1$ shows that its first matrix commutes with
every $X$, so that matrix is scalar. Here $g(a_n)$ is noncentral.
It follows that $\pi(g)\ne1$, although $g=1$ near $F$. No proper
closed set carries $\pi$; $[0,1]$ does. Notice that this argument
does not try to isolate one point from the infinite dense set.
\end{proof}

\begin{proposition}[The uniform-topology obstruction]\ours{}
\label{prop:not-czero}
The representation in Theorem~\ref{thm:interval-support}, even
restricted to smooth currents, is not strongly continuous for the
$C^0$ topology. In particular it has no strongly continuous
extension to $C([0,1],SU(2))$ with its uniform topology.
\end{proposition}

\begin{proof}
For each $N$ choose a smooth $f_N:[0,1]\to[0,1]$ with
$f_N(a_j)=1$ and $f_N(b_j)=0$ for $j\leq N$. Disjoint bumps around
the finitely many $a_j$, avoiding the finitely many $b_j$, provide
such a function. With $\sigma_3=\operatorname{diag}(1,-1)$, put
\[
 h_N(t)=\exp(iN^{-1/2}f_N(t)\sigma_3).
\]
Then $\|h_N-1\|_\infty\leq N^{-1/2}$. The first $N$ vacuum
coefficients are all $\cos(N^{-1/2})$, since
\[
 \langle(g(a)\otimes\overline{g(b)})\Omega_n,\Omega_n\rangle
 =\tfrac12\operatorname{Tr}(g(a)g(b)^*).
\]
Every remaining coefficient has absolute value at most one.
Taking limits of the finite tensor products yields
\[
 |\langle\pi(h_N)\Omega,\Omega\rangle|
 \leq\cos(N^{-1/2})^N\longrightarrow e^{-1/2}<1.
\]
Strong convergence $\pi(h_N)\Omega\to\Omega$ is impossible.
An extension continuous for the uniform topology would imply just
that convergence on this sequence.
\end{proof}

\begin{corollary}[Exact support in the norm-continuous class]\ours{}
For $G^\infty$, an irreducible operator-norm-continuous unitary
representation has finite support; conversely, every finite subset
occurs as the support of such a representation. An interval of
positive length cannot occur in this class.
\end{corollary}

\begin{proof}
The smooth mapping group is connected: $g(t)\mapsto g(st)$ joins
it to its constant subgroup, and $SU(2)$ is connected.
The cited classification therefore expresses an irreducible
norm-continuous representation as a finite tensor product
$\bigotimes_{x\in S}\rho_x(g(x))$, omitting trivial factors.
The finite set $S$ carries it. If a closed $F$ omits $x\in S$,
choose a smooth current supported in a small neighbourhood of $x$
disjoint from $F$ and from $S\setminus\{x\}$, taking at $x$ a value
on which the nontrivial $\rho_x$ is not the identity. It acts
nontrivially, so $F$ cannot carry the representation. Thus the
support is exactly $S$. Conversely, independent finite
interpolation and irreducibility of each factor show that the
tensor product of nontrivial finite-dimensional irreducible
$SU(2)$ representations at distinct points is irreducible and
has exactly that support. For $S=\varnothing$ use the trivial
one-dimensional representation.
\end{proof}

\paragraph{Independent check and remaining obstruction.}
The proof was checked by a separate reader and by re-deriving its
two potentially delicate steps. The tail estimate is selected
\emph{after} every finite interpolation constant is fixed, so very
small pair separations do not create a circular bound. Irreducibility
uses approximate realization of finite local matrix algebras, not
impossible exact separation of the dense evaluation set. Support
was checked separately through conjugation of a local algebra.
The $C^0$ obstruction supplies an independent diagnostic that the
topology cannot be weakened without changing the conclusion.
No finite computation is substituted for these infinite-dimensional
arguments.

What remains unresolved here is a classification of strongly
continuous irreducible representations, the identification of the
notebook's intended category, and the extent to which arbitrary
paired sequences yield equivalent representations. A precise next
question for this construction is: for distinct dense pairs
$(a_n,b_n)$ with $\sum|a_n-b_n|<\infty$, characterize when the
corresponding representation is irreducible, without imposing
the sufficient interpolation-tail condition (1). A further bounded
run could compare that condition with intrinsic continuity and
commutant criteria in the infinite-tensor-product literature,
and trace the historical priority of paired evaluation constructions.
It should not promise a classification from the present sufficient
criterion. The status of GP-1969-02 is therefore a positive answer
in the explicit $C^1$/strong category and a negative answer in the
smooth/operator-norm category, with the archival ambiguity retained.

\subsection{Four subspaces: an elementary proof of the requested deformation}
\label{dev:four-subspaces}

On P000317, item 1970-198.12 asks for a simple proof that
$S(5,2)\oplus S(1,2)$ deforms to $S(3,2)\oplus S(3,2)$.
The canonical systems are identified in
\citet[\S2]{gelfandponomarev1971quadruples}, received on
21 October 1970 and published in 1971. The original announcement
was read, as were the later normal forms reproduced in
\citet[\S9]{enomotowatatani2004subspaces}. Those checked texts
identify the systems but do not establish the exact requested
deformation in a located passage. The publication match is therefore
partial. The bounded mathematical question is answered below;
this is not a claim that it remained open after the seminar or
that the argument is new.

\paragraph{Precise formulation.}
Work over $\mathbb C$ with ordered quadruples of subspaces and
simultaneous change of basis in the ambient vector space. The
source's positive-defect notation here means
$\sum_i\dim E_i-2\dim P=2$ for each indecomposable summand;
the notebook's earlier displayed definition has the opposite sign.
Define $S(1,2)$ as four copies of $\mathbb C$ in $\mathbb C$.
For $R=S(3,2)$ take ambient basis $e_1,e_2,f_1$ and
\[
 \begin{aligned}
 E_1&=\langle e_1,e_2\rangle,&
 E_2&=\langle f_1,e_2\rangle,\\
 E_3&=\langle e_1,e_2+f_1\rangle,&
 E_4&=\langle f_1,e_1+e_2\rangle.
 \end{aligned}
\]
For $S(5,2)$ use basis $e_1,e_2,e_3,f_1,f_2$ and
\[
 \begin{aligned}
 E_1&=\langle e_1,e_2,e_3\rangle,&
 E_2&=\langle f_1,f_2,e_3\rangle,\\
 E_3&=\langle e_1,e_2+f_1,e_3+f_2\rangle,&
 E_4&=\langle f_1,e_1+f_2,e_2+e_3\rangle.
 \end{aligned}
\]
Both direct sums in the question consequently define points of
$X=\operatorname{Gr}(4,6)^4$. We prove the stronger statement
that $R\oplus R$ is the unique dense open orbit in $X$; thus every
point of $X$, including the displayed special sum, admits an
analytic deformation to that orbit. The field, topology and
ordered-subspace convention are explicit reconstruction choices.

\begin{lemma}[Endomorphisms of the three-dimensional system]\ours{}
Every endomorphism of $R$ preserving its four subspaces is scalar.
\end{lemma}

\begin{proof}
In the dual basis the four annihilator lines are represented by
\[
 \ell_1=(0,0,1),\quad\ell_2=(1,0,0),\quad
 \ell_3=(0,1,-1),\quad\ell_4=(1,-1,0).
\]
The first three form a basis, and
$\ell_4=-\ell_1+\ell_2-\ell_3$ has all three coefficients nonzero.
An endomorphism preserving each primal plane has transpose
preserving these four dual lines. It is diagonal in the first
three lines, and preserving the fourth forces all three diagonal
entries to agree. This proves the assertion, including
noninvertible endomorphisms.
\end{proof}

\begin{theorem}[The notebook's deformation]\ours{}
The $\mathrm{GL}_6(\mathbb C)$ orbit of $R\oplus R$ is open and
dense in $X$. There is a holomorphic one-parameter family of ordered
quadruples $x(t)\in X$, defined near $0$, with
$x(0)=S(5,2)\oplus S(1,2)$ and $x(t)\cong R\oplus R$ for all
sufficiently small nonzero $t$.
\end{theorem}

\begin{proof}
An endomorphism of a direct sum is a block matrix of morphisms
between its summands. The lemma gives
$\operatorname{End}(R\oplus R)=M_2(\mathbb C)$, of dimension four.
Its invertible elements form the stabilizer, also of dimension
four. The orbit therefore has dimension $36-4=32$. Meanwhile
$\dim X=4\cdot4(6-4)=32$, and $X$ is irreducible. An algebraic
group orbit is locally closed. A locally closed subset of full
dimension in this irreducible variety has closure $X$ and is open
in that closure; the orbit is consequently open dense. Two
disjoint open dense orbits in an irreducible variety cannot exist.

For completeness, the closure statement gives an actual
one-parameter deformation, not just a dimension comparison.
Choose a product of standard affine Grassmann charts through
$x(0)$; in graph coordinates it is $\mathbb C^{32}$ with $x(0)=0$.
The complement of the open orbit is a proper closed algebraic
subset of this chart. Choose a nonzero polynomial $f$ vanishing
on that complement and a point $v$ with $f(v)\ne0$.
Then $f(tv)$ is a nonzero one-variable polynomial, so for every
sufficiently small $t\ne0$ it is nonzero. The graph coordinates
$tv$ give the asserted holomorphic family. Finally, the common
intersection of the four subspaces is one-dimensional for
$S(5,2)\oplus S(1,2)$ and zero-dimensional for $R\oplus R$;
the special fibre is genuinely a different orbit.
\end{proof}

\begin{example}[Independent endomorphism calculation]\oursComputed{}
The script \nolinkurl{code/verify_four_subspaces.py} imposes
$\ell^T A v=0$ for every subspace basis vector $v$ and every
annihilating covector $\ell$, then row-reduces the resulting
linear system over $\mathbb Q$. It independently finds
\[
\begin{array}{c|rrrrr}
\text{system}&S(1,2)&S(3,2)&S(5,2)&S(5,2)\oplus S(1,2)&R\oplus R\\
\hline
\dim\operatorname{End}&1&1&1&7&4
\end{array}
\]
The last two orbit dimensions are $29$ and $32$. The output is
\nolinkurl{output/verify_four_subspaces.json}. This calculation
checks the normal forms and stabilizer dimensions independently
of the projective-frame proof. The first two nontrivial single
summands have full-dimensional orbits in their own products of
Grassmannians, confirming the source's stability claim for these
particular summands, without endorsing its separate false slogan
about every parameter-free indecomposable.
\end{example}

The proof was independently reviewed through both the projective
frame and the orbit-dimension argument. No obstruction remains
for this particular deformation over $\mathbb C$. What remains
unverified is the exact historical publication of the requested
elementary proof, not a mathematical openness claim. One further
bounded run could locate that attribution or produce a convenient
closed formula for a chosen one-parameter change of bases;
neither is needed for the existence theorem proved here.

\subsection{Audit annotation: commuting quantum marginals need not have product eigenbases}
\label{dev:bell}

This section corrects item 1970-197.2 on P000312--P000313. It is an
annotation of a false literal claim, not a newly discovered open
problem selected from class C or D. The notes first discuss the
simple-spectrum situation, where a matrix commuting with all
one-site marginals is diagonal in their uniquely determined tensor
eigenbasis, and then assert that a suitable tensor eigenbasis can
always be chosen in the degenerate case. The following counterexample
checks the stronger condition that \emph{all} finite marginals
commute after embedding in their common tensor algebra.

\begin{proposition}[A faithful consistent Bell-block counterexample]\ours{}
For each $0<p<1$ there is a state on the qubit chain indexed by
$\mathbb Z$ whose finite density matrices are strictly positive,
consistent under partial trace, and pairwise commuting when
embedded into a common finite tensor algebra, but whose two-site
density matrix on a designated adjacent pair is not diagonal in
any product of one-site orthonormal bases.
\end{proposition}

\begin{proof}
Partition the sites into pairs $\{2k,2k+1\}$ and put on each pair
\[
 D_p=p|\Phi\rangle\langle\Phi|+\frac{1-p}{4}1_4,
 \qquad \Phi=\frac{|00\rangle+|11\rangle}{\sqrt2}.
\]
Its eigenvalues are $(1+3p)/4$ on $\Phi$ and $(1-p)/4$ on its
three-dimensional orthogonal complement. Thus it is positive
definite and has trace one. Each one-site partial trace is $1_2/2$.
For a finite set $I\subset\mathbb Z$, define $\rho_I$ as the
tensor product of $D_p$ for complete pairs in $I$ and $1_2/2$ for
single sites of partially included pairs. Factors are reordered
into site order by the canonical tensor permutation.

Tracing out one member of a complete pair replaces $D_p$ by
$1_2/2$; tracing out the remaining member gives one. It follows
that $\operatorname{Tr}_{I\setminus J}\rho_I=\rho_J$ for every
$J\subset I$. These normalized positive consistent functionals
define a state on the increasing union of local matrix algebras;
their norm is one, so the state extends to its norm completion.
Every finite density matrix is strictly positive.

For finite $I,J$, embed $\rho_I$ and $\rho_J$ in the tensor
algebra on $I\cup J$ by tensoring with identity operators. On each
fixed two-site block, each embedded matrix is either $D_p$ or a
scalar multiple of the identity (including the case with only one
site in the union). These block factors commute, hence so do the
two full matrices. This checks every finite $I,J$, not just nested
intervals.

If $D_p$ were diagonal in a product of one-site orthonormal bases,
its simple eigenvalue $(1+3p)/4$ would have a product eigenvector.
Its eigenspace is instead $\mathbb C\Phi$. A two-qubit product
vector has a rank-one coefficient matrix, whereas $\Phi$ has
coefficient matrix $2^{-1/2}1_2$ of determinant $1/2$. Thus
$\Phi$ is not a product vector, a contradiction.
\end{proof}

The example has period two. No one-site translation-invariance
hypothesis is stated at this point in the notebook, and none is
needed for the claim being refuted. Strict positivity shows that
zero eigenvalues are not responsible. The obstruction is
degeneracy of the one-site marginals, which leaves room for a
nonproduct eigenvector of a joint density matrix.

\begin{proposition}[The valid simple-spectrum repair]\ours{}
Let $\rho$ be a density matrix on
$H_1\otimes\cdots\otimes H_m$, with finite-dimensional $H_i$.
Suppose that each one-site marginal $\rho_i$ has simple spectrum
and that $\rho$ commutes with each embedded $\rho_i$. Then $\rho$
is diagonal in the tensor product of the one-site eigenbases.
\end{proposition}

\begin{proof}
For a simple-spectrum self-adjoint matrix, each spectral projection
is a polynomial in that matrix by Lagrange interpolation on its
distinct eigenvalues. Thus $\rho$ commutes with every one-site
rank-one spectral projection and with every product of these
projections. Those products are the rank-one projections of the
tensor eigenbasis, so each corresponding line is invariant under
$\rho$. This is the desired diagonalization.
\end{proof}

\begin{example}[Independent exact finite checks]\oursComputed{}
The file \nolinkurl{code/verify_annotations.py} uses rational
arithmetic from the Python standard library. On four sites it
computes every one of the $16$ subset marginals by an explicit
basis definition of partial trace, independently compares them
with the block formula in the proof, and checks all $256$ ordered
commutators after embedding. It repeats the entire calculation
for $p=1/2$ and $p=2/3$. Both runs pass exactly. The first parameter
gives spectrum $5/8,1/8,1/8,1/8$, and the second gives
$3/4,1/12,1/12,1/12$. As an additional diagnostic, the partial
transpose has a negative eigenvalue $-1/8$ and $-1/4$, respectively.
The full output is \nolinkurl{output/verify_annotations.json}.
The infinite-chain conclusion follows from the block proof above;
four-site computations are an independent check of its algebra,
not evidence substituted for an all-volume theorem.
\end{example}

\subsection{Audit annotation: the product uncertainty question and wavelet scaling}
\label{dev:wavelet}

On P000358 the general orthogonal-system question asks for a bound
on the \emph{product} $\delta_x\delta_p$; it is distinct from the
proposed Gram--Schmidt construction from a coherent-state lattice.
The item and GP-1970-59 retain that distinction. This is a class-B
historical resolution, not an open problem manufactured for
development. Meyer's Schwartz orthonormal wavelet construction
\citep[\S2, Thm.~1]{meyer1987uncertainty} already gives the
following elementary implication. The lecture was delivered in
February 1986 and published in 1987.

\begin{proposition}[Uniform uncertainty products from a Schwartz wavelet]\ours{}
Let $\psi\in\mathcal S(\mathbb R)$ have norm one and suppose that
\[
 \psi_{j,k}(x)=2^{j/2}\psi(2^jx-k),\qquad j,k\in\mathbb Z,
\]
is an orthonormal basis of $L^2(\mathbb R)$. With
$Qf=xf$, $Pf=-if'$ and $\hbar=1$, let
$\Delta_f A=\|(A-\langle Af,f\rangle)f\|$ for a normalized vector
in the domain of $A$. Then
\[
 \Delta_{\psi_{j,k}}Q=2^{-j}\Delta_\psi Q,
 \qquad\Delta_{\psi_{j,k}}P=2^j\Delta_\psi P,
 \qquad
 \Delta_{\psi_{j,k}}Q\,\Delta_{\psi_{j,k}}P
 =\Delta_\psi Q\,\Delta_\psi P<\infty.
\]
\end{proposition}

\begin{proof}
All vectors are Schwartz, so the required first and second moments
exist. In the position integrals substitute $y=2^jx-k$ to obtain
\[
 \langle Q\psi_{j,k},\psi_{j,k}\rangle
 =2^{-j}(\langle Q\psi,\psi\rangle+k),
\]
and the centered second moment is multiplied by $2^{-2j}$.
Differentiation gives
$P\psi_{j,k}(x)=2^{3j/2}(P\psi)(2^jx-k)$, so the momentum mean is
multiplied by $2^j$ and its centered second moment by $2^{2j}$.
Taking square roots proves the identities.
\end{proof}

As a second check, the dilation--translation unitary
$U_{j,k}f(x)=2^{j/2}f(2^jx-k)$ satisfies on Schwartz functions
\[
 U_{j,k}^*QU_{j,k}=2^{-j}(Q+k),\qquad
 U_{j,k}^*PU_{j,k}=2^jP.
\]
These conjugation identities reproduce the variance calculation
without repeating the moment integrals. They also show its exact
limitation: neither individual deviation stays bounded as
$j\to\pm\infty$.

Bourgain proved the stronger existence of unrestricted orthonormal
bases with both centered second moments uniformly bounded
\citep[main theorem; IHES preprint, November 1986]{bourgain1988uncertainty}.
The journal paper appeared in 1988. The distinction between such
bases, regular Gabor systems, and higher-moment obstructions is
explicit in \citet[pp.~115--116 and Thm.~1.1]{grochenigmalinnikova2013localization}.
Neither the wavelet argument nor that general existence statement
proves the notebook's particular Gram--Schmidt procedure works in
an unspecified ordering. The chapter accordingly separates the
settled existence question from that incompletely specified method.

\section{Problem index}\label{problems:1969}
Earlier-year occurrence identifiers are retained as text; current-year identifiers link to their entries.

\subsection{GP-1966-13: Entire minimal-graph Bernstein problem}\label{prob:GP-1966-13}
For which dimensions \(m\) must every smooth entire solution \(u:\mathbb R^m\to\mathbb R\) of the minimal-surface equation be affine?
\par\noindent\textit{Occurrences:} 1966-114.6, \hyperref[item:1969-174.4]{1969-174.4}.
\par\noindent\textit{Status (2026-11):} As of November 2026: resolved. Every such graph is affine for \(m\leq7\); non-affine entire graphs exist for \(m\geq8\). The counterexamples were published in 1969 by \citet{bombieridegiorgigiusti1969minimal}. The later Russian translation is not the result\textquotesingle s discovery date.

\subsection{GP-1969-01: Irreducible representations of an interval current group}\label{prob:GP-1969-01}
After fixing the regularity and topology of a group of maps \([0,1]\to SU(2)\), classify its continuous irreducible unitary representations up to equivalence. Distinguish operator-norm continuity from strong continuity; the notebook specifies neither the map category nor this convention.
\par\noindent\textit{Occurrences:} \hyperref[item:1969-180.1]{1969-180.1}.
\par\noindent\textit{Status (2026-11):} As of November 2026: partly developed, with an archival ambiguity. In the smooth compact-semisimple setting, norm-continuous irreducibles are finite evaluation tensors \citep[Theorem 4.11; group theorem 4.15 in the arXiv version]{janssensneeb2015norm}; preprint 2013, online 2014, print 2015. Positive-energy and measurable-current results impose other hypotheses \citep{janssensneeb2024positive,araki1970factorizable}. The checked literature does not supply a full strong-continuity classification for every interpretation of the notebook.

\subsection{GP-1969-02: Interval support of an irreducible current-group representation}\label{prob:GP-1969-02}
In a specified continuous unitary representation category for \(SU(2)\)-valued interval currents, does there exist an irreducible representation nontrivial on currents supported in every nonempty open subinterval? Equivalently in the full-support case, no proper closed subset carries the representation in the sense recorded on P000291.
\par\noindent\textit{Occurrences:} \hyperref[item:1969-180.2]{1969-180.2}.
\par\noindent\textit{Status (2026-11):} As of November 2026: the chapter proves existence for strong continuity on \(C^1([0,1],SU(2))\), with irreducibility also on the smooth subgroup and full interval support. This constructed representation has no strongly continuous uniform-topology extension. Smooth norm-continuous irreducibles instead have finite support \citep{janssensneeb2015norm}. The high-dimensional energy theorem does not apply to an interval, whose energy representation is reducible \citep{ismagilov1976unitary,ismagilov1981segment}. These are precise positive and negative variants; the unrecorded archival category and broader classification remain unresolved here. No historical novelty is claimed for the construction.

\subsection{GP-1969-03: Analytic approximation of formal solutions}\label{prob:GP-1969-03}
Given analytic equations \(F(x,y)=0\) over \(\mathbb C\) and a formal solution \(\widehat y(x)\) with the prescribed constant term, does every finite jet of \(\widehat y\) extend to a convergent solution of the same equations?
\par\noindent\textit{Occurrences:} \hyperref[item:1969-175.3]{1969-175.3}, \hyperref[item:1969-177.1]{1969-177.1}.
\par\noindent\textit{Status (2026-11):} As of November 2026: resolved affirmatively by \citet{artin1968analytic}, published December 1968. For every \(c\), a convergent solution can be chosen congruent to \(\widehat y\) modulo \((x)^c\). Raynaud presented the result at Bourbaki in June 1969; that exposition appeared in 1971 \citep[Theorem 1]{raynaud1969artin}. The given formal series itself need not converge.

\subsection{GP-1970-20: Rota's finite-field vector-space Ramsey conjecture}\label{prob:GP-1970-20}
For fixed \(q,r,\ell,k\), \(1\leq\ell<k\), does every \(r\)-coloring of the \(\ell\)-subspaces of a sufficiently large \(F_q^N\) admit a \(k\)-subspace monochromatic in dimension \(\ell\)?
\par\noindent\textit{Occurrences:} \hyperref[item:1970-191.1]{1970-191.1}, \hyperref[item:1970-191.2]{1970-191.2}.
\par\noindent\textit{Status (2026-11):} Solved by the Graham--Leeb--Rothschild theorem; January 1972 announcement. No identical earlier problem was found in the supplied previous index.

\subsection{GP-1970-21: Quadratic pairs in characteristic three}\label{prob:GP-1970-21}
Classify nontrivial finite groups \(G\) and faithful irreducible \(F_3G\)-modules generated by elements acting with \((g-I)^2=0\).
\par\noindent\textit{Occurrences:} \hyperref[item:1970-192.3]{1970-192.3}.
\par\noindent\textit{Status (2026-11):} Later classification work includes Chermak's 2002 preprint/2004 article, Theorems A--B, for the component case, together with its cited component-free treatment. The notebook's separate characteristic-two slogan has insufficient restrictions for the same identification.

\subsection{GP-1970-22: Factor criteria for infinite quantum spin states}\label{prob:GP-1970-22}
When does a state of the spin UHF algebra have a factor GNS representation, and which types can occur?
\par\noindent\textit{Occurrences:} \hyperref[item:1970-195.8]{1970-195.8}, \hyperref[item:1970-197.4]{1970-197.4}, \hyperref[item:1970-200.1]{1970-200.1}, \hyperref[item:1970-200.2]{1970-200.2}, \hyperref[item:1970-200.3]{1970-200.3}, \hyperref[item:1970-200.4]{1970-200.4}.
\par\noindent\textit{Status (2026-11):} Powers\' 1967 relative-commutant correlation criterion answers the factor question; product states and Araki's one-dimensional finite-range Gibbs states supply established families. The type depends on the state. No identical spin-chain problem was found in the supplied previous index.

\subsection{GP-1970-23: Elementary proof of the four-subspace deformation example}\label{prob:GP-1970-23}
Give a direct proof that \(S(5,2)\oplus S(1,2)\) has arbitrarily small deformations of type \(S(3,2)\oplus S(3,2)\), although the two individual summands are rigid.
\par\noindent\textit{Occurrences:} \hyperref[item:1970-198.12]{1970-198.12}.
\par\noindent\textit{Status (2026-11):} The normal forms are identified in the 1971 classification. The requested deformation over \(\mathbf C\) is proved in Section~\ref{dev:four-subspaces} by the dense open orbit of \(S(3,2)^{\oplus2}\); an analytic one-parameter family follows in a Grassmann chart. The exact historical publication of the elementary proof was not verified. No mathematical openness is claimed.

\subsection{GP-1970-24: The unit-sheaf cohomology of the punctured elliptic curve}\label{prob:GP-1970-24}
For \(A=\mathbf C[x,y]/(y^2-x^3+1)\), show \(H^1(\operatorname{Spec}A,\mathcal O^\times)\ne0\).
\par\noindent\textit{Occurrences:} \hyperref[item:1970-199.3]{1970-199.3}.
\par\noindent\textit{Status (2026-11):} Solved by the Picard group of a smooth elliptic curve with its point at infinity removed.

\subsection{GP-1970-25: Punctured affine cohomology and residue duality}\label{prob:GP-1970-25}
Compute \(H^*(\mathbf A^n\setminus\{0\},\mathcal O)\) and identify the module and residue pairing in its top nonzero degree.
\par\noindent\textit{Occurrences:} \hyperref[item:1970-199.4]{1970-199.4}, \hyperref[item:1970-201.2]{1970-201.2}, \hyperref[item:1970-201.4]{1970-201.4}, \hyperref[item:1970-201.5]{1970-201.5}, \hyperref[item:1970-201.6]{1970-201.6}.
\par\noindent\textit{Status (2026-11):} Solved classical localization/Koszul calculation. For \(n\geq2\) the top group is the module of Laurent monomials negative in every variable; its residue dual is restricted linear, not ring-linear. Earlier affine-cover and non-affineness entries 1965-70.6 and 1965-58.7 are related background, not an identical unresolved question.

\subsection{GP-1970-26: Homotoping a distribution to a foliation}\label{prob:GP-1970-26}
Is every smooth constant-rank subbundle of a tangent bundle homotopic through subbundles to the tangent distribution of a nonsingular foliation?
\par\noindent\textit{Occurrences:} \hyperref[item:1970-206.1]{1970-206.1}, \hyperref[item:1970-206.2]{1970-206.2}.
\par\noindent\textit{Status (2026-11):} Negative in general, as reported from Bott's pre-seminar integrability obstruction. Normal Pontryagin-ring classes in degrees above twice the codimension obstruct such a homotopy. The compressed CP² example is not securely identified and is not needed to state the obstruction. Earlier item 1968-149.3 records a related foliation-existence slogan, but its class-F title alone does not establish identical quantifiers and no earlier problem ID is supplied; the link is thematic.

\subsection{GP-1970-50: Analytic and combinatorial torsion}\label{prob:GP-1970-50}
Equality of Ray--Singer and Reidemeister torsion for closed manifolds with an acyclic unitary flat bundle, with the determinant-line interpretation when cohomology is present.
\par\noindent\textit{Occurrences:} \hyperref[item:1970-215.6]{1970-215.6}.
\par\noindent\textit{Status (2026-11):} As of November 2026, solved: Cheeger\textquoteright s 1977 announcement and M\"uller\textquoteright s 1978 proof establish the unitary comparison. The 1971 report is a preprint occurrence, not its solution.

\subsection{GP-1970-51: Gudkov congruence for maximal real plane curves}\label{prob:GP-1970-51}
For every nonsingular real projective \(M\)-curve of degree \(2k\), prove \(p-m\equiv k^2\pmod8\), with \(p,m\) the even/odd nesting counts.
\par\noindent\textit{Occurrences:} \hyperref[item:1970-221.3]{1970-221.3}, \hyperref[item:1970-221.4]{1970-221.4}.
\par\noindent\textit{Status (2026-11):} As of November 2026, solved by Rokhlin: note received 29 December 1971, published 1972. Arnold\textquoteright s modulo-four result was the March 1971 intermediate step.

\subsection{GP-1970-52: A bounded linear operator with a dense orbit}\label{prob:GP-1970-52}
Does there exist a separable infinite-dimensional Banach space and a bounded linear operator on it with an orbit dense in the entire space? The notebook realizes this on sequence spaces by \(\lambda S\), \(|\lambda|>1\).
\par\noindent\textit{Occurrences:} \hyperref[item:1970-224.2]{1970-224.2}, \hyperref[item:1970-224.3]{1970-224.3}.
\par\noindent\textit{Status (2026-11):} As of November 2026, the sequence-space existence question was already answered by Rolewicz in 1969. No universal assertion about all nonseparable or finite-dimensional spaces is included.

\subsection{GP-1970-53: Modular parametrization of a rational elliptic curve}\label{prob:GP-1970-53}
For every elliptic curve \(E/\mathbb Q\) of conductor \(N\), construct a nonconstant morphism \(X_0(N)\to E\) over \(\mathbb Q\), equivalently identify its \(L\)-series with a weight-two newform.
\par\noindent\textit{Occurrences:} \hyperref[item:1970-227.3]{1970-227.3}, \hyperref[item:1970-227.7]{1970-227.7}, \hyperref[item:1970-227.8]{1970-227.8}.
\par\noindent\textit{Status (2026-11):} As of November 2026, solved: the full modularity theorem was proved by Breuil, Conrad, Diamond, and Taylor and published in 2001.

\subsection{GP-1970-54: Which sequences of coherent states are complete?}\label{prob:GP-1970-54}
Characterize sets \(\Lambda\subset\mathbb C\) for which the normalized coherent states span the Bargmann--Fock space densely, equivalently its uniqueness sets.
\par\noindent\textit{Occurrences:} \hyperref[item:1970-228.5]{1970-228.5}, \hyperref[item:1970-228.6]{1970-228.6}, \hyperref[item:1970-228.7]{1970-228.7}, \hyperref[item:1970-228.12]{1970-228.12}, \hyperref[item:1970-228.14]{1970-228.14}.
\par\noindent\textit{Status (2026-11):} As of November 2026, Perelomov\textquoteright s lattice criterion and the stated Mitkovski--Wick perturbation criteria are verified. An unrestricted characterization is not established by this audit; a September 2026 zero-set preprint remains an unverified lead. Sampling density theorems alone do not answer the exact question.

\subsection{GP-1970-59: Uniform centered uncertainty for an orthonormal basis}\label{prob:GP-1970-59}
Does \(L^2(\mathbb R)\) admit an orthonormal basis \((e_j)\) with \(\sup_j\Delta_x(e_j)\Delta_p(e_j)<\infty\)? The notebook also mentions an unspecified orthogonalization of a critical coherent lattice.
\par\noindent\textit{Occurrences:} \hyperref[item:1970-228.8]{1970-228.8}.
\par\noindent\textit{Status (2026-11):} As of November 2026, yes for general bases: Meyer wavelets give the product bound and Bourgain gives separate uniform centered second-moment bounds. Balian--Low prohibits the corresponding regular single-window Gabor basis. No assertion covers every ordering or an unspecified Gram--Schmidt procedure.

\section{Corrections to the transcription and translation}\label{corr:1969}
The evidence distinguishes a transcription correction from a mathematical emendation of wording faithfully preserved in the notebook.
\subsection*{P000282 (both; provenance): opening carryover}
\noindent\textit{Transcription reads.} Catalog record 174 includes P000282\par
\noindent\textit{Critical reading.} Attribute P000282 to the continuation of record 173, already indexed in the previous chapter\par
\noindent\textit{Evidence.} Source continuity and \nolinkurl{previous_index.md} identify the Buchstaber torus and film/Morse discussion as 1968-173.1--3.\par
\subsection*{P000283 (both; provenance): heading before spectral measures}
\noindent\textit{Transcription reads.} Kazhdan (catalog attribution)\par
\noindent\textit{Critical reading.} Kazhdan and Bernstein are jointly named in the notebook heading\par
\noindent\textit{Evidence.} The scan clearly reads Каждан и Бернштейн.; the catalogue speaker remains a record identifier, while item attribution is joint.\par
\subsection*{P000285 (both; faithful error): classical admissible displacements}
\noindent\textit{Transcription reads.} \(C^1(\Sigma)\)\par
\noindent\textit{Critical reading.} \(C^1(\overline A)\), with the trace inequality on \(\Sigma\)\par
\noindent\textit{Evidence.} The energy contains interior strain derivatives. The scan faithfully writes the boundary symbol.\par
\subsection*{P000285 (both; faithful error): surface-work term}
\noindent\textit{Transcription reads.} Integral over Sigma\par
\noindent\textit{Critical reading.} Integral over Sigma* for the prescribed load\par
\noindent\textit{Evidence.} The preceding source assigns the prescribed traction to Sigma*, and the scalar energy on the same page uses Sigma*. The scan faithfully has Sigma in the first formula.\par
\subsection*{P000287 (both; provenance): Soviet-priority author}
\noindent\textit{Transcription reads.} Rybkin / Рыбкин\par
\noindent\textit{Critical reading.} Rybin / Рыбин (bibliographic identification)\par
\noindent\textit{Evidence.} Math-Net rm6746 identifies P. P. Rybin at exactly UMN 15(4), 1960, pp.169--172. This corrects the identification, not a certainty about the handwritten letters.\par
\subsection*{P000288 (both; faithful error): sphere two-form and primitive}
\noindent\textit{Transcription reads.} \(\Omega=a\,dz\wedge d\bar z/(1+|z|^2)^2\), \(\omega=(a/i)(z\,d\bar z-\bar z\,dz)/(1+|z|^2)\)\par
\noindent\textit{Critical reading.} \(d\omega=(2a/i)\,dz\wedge d\bar z/(1+|z|^2)^2\); choose a consistent real symplectic normalization\par
\noindent\textit{Evidence.} Direct differentiation of the two formulas visible in the scan; the discrepancy is faithfully transcribed.\par
\subsection*{P000289 (both; faithful error): integral parameter statement}
\noindent\textit{Transcription reads.} A measure exists only for integral a\par
\noindent\textit{Critical reading.} A global equivariant prequantum line bundle requires an integral curvature class; existence of a positive weighted measure alone does not\par
\noindent\textit{Evidence.} The scan confirms the measure wording. The radial integral is finite for every real a>1.\par
\subsection*{P000293 (both; faithful error): Fourier p-norm display}
\noindent\textit{Transcription reads.} \(\|\widehat f\|_p=(\sum_\sigma d_\sigma\|\widehat f(\sigma\|_p^p)\)\par
\noindent\textit{Critical reading.} \(\|\widehat f\|_p=(\sum_\sigma d_\sigma\|\widehat f(\sigma)\|_p^p)^{1/p}\)\par
\noindent\textit{Evidence.} The scan has no secure final exponent; homogeneity of a norm requires the p-th root. This emends the formula rather than silently changing the transcription.\par
\subsection*{P000294 (both; faithful error): monodromy parenthesis}
\noindent\textit{Transcription reads.} \(T\) is unipotent\par
\noindent\textit{Critical reading.} \(T\) is quasi-unipotent; \(T^N\) is unipotent\par
\noindent\textit{Evidence.} Both the Russian scan and the transcription have T without the exponent in the parenthesis. The preceding displayed formula gives the correct qualified assertion.\par
\subsection*{P000295 (both; ambiguity): flat-section expansion lower index}
\noindent\textit{Transcription reads.} \(\sum_{n\leq0}\)\par
\noindent\textit{Critical reading.} For the regular-singular expansion use \(\sum_{n\geq n_0}\) after finite ramification\par
\noindent\textit{Evidence.} The scan has crowded index strokes and the transcription explicitly marks the reading as uncertain. Mathematics requires a finite principal part; the proposed index is a reconstruction, not a definitive paleographic correction.\par
\subsection*{P000297 (both; faithful error): localization applied to all principal opens}
\noindent\textit{Transcription reads.} \(0\notin S\) is required, followed by \(S=\{f^n:n\geq0\}\) for every \(f\)\par
\noindent\textit{Critical reading.} Allow the zero-ring localization when f is nilpotent\par
\noindent\textit{Evidence.} The empty principal open has zero-ring sections; the preceding restricted localization definition does not cover this case.\par
\subsection*{P000297 (both; ambiguity): sheaf equalizer symbols}
\noindent\textit{Transcription reads.} \(\bigoplus\) with the caption product of sets\par
\noindent\textit{Critical reading.} Interpret the two symbols as \(\prod\), not direct sums\par
\noindent\textit{Evidence.} The scan visibly has circled-plus signs but explicitly labels the first a product of sets. An arbitrary sheaf cover requires products; the source meaning is clearer than its notation.\par
\subsection*{P000297 (both; faithful error): uniqueness of a sheaf with prescribed stalks}
\noindent\textit{Transcription reads.} a sheaf with the same stalks, and unique\par
\noindent\textit{Critical reading.} the canonical sheafification with its stalk-isomorphism map, unique by its universal property\par
\noindent\textit{Evidence.} Abstract stalks alone do not determine a sheaf; the classical exercise concerns canonical sheafification.\par
\subsection*{P000298 (both; provenance): prescheme and scheme terminology}
\noindent\textit{Transcription reads.} prescheme; scheme when the diagonal is closed\par
\noindent\textit{Critical reading.} modern scheme; separated scheme when the diagonal is closed\par
\noindent\textit{Evidence.} This is the early EGA convention used consistently by the lecture; it is a terminology note rather than an error in transcription.\par
\subsection*{P000299 (both; faithful error): affine cohomology parenthesis}
\noindent\textit{Transcription reads.} for affine spaces cohomology is zero\par
\noindent\textit{Critical reading.} for affine schemes and quasi-coherent coefficients, cohomology in strictly positive degrees is zero\par
\noindent\textit{Evidence.} The scan has no degree or coefficient qualification. Degree zero consists of global sections and is generally nonzero.\par
\subsection*{P000299 (both; faithful error): empty inverse-image parenthesis}
\noindent\textit{Transcription reads.} the inverse image may be empty, then the set of sections is empty\par
\noindent\textit{Critical reading.} sections on the empty open form a singleton set, or the zero group\par
\noindent\textit{Evidence.} Both occurrences of empty are visible in the Russian scan. The sheaf axiom for the empty cover determines a terminal section object.\par
\subsection*{P000300 (both; provenance): inverse-image functor notation}
\noindent\textit{Transcription reads.} \(f^*\) for sheaves on topological spaces\par
\noindent\textit{Critical reading.} Modern \(f^{-1}\); reserve module pullback \(f^*\) for tensoring with the structure sheaf\par
\noindent\textit{Evidence.} The source discusses sets or abelian groups and uses the sheaf adjunction; this is a notation clarification, not a transcription error.\par
\subsection*{P000301 (both; faithful error): associated object after Hom display}
\noindent\textit{Transcription reads.} F is the associated presheaf\par
\noindent\textit{Critical reading.} F is the associated sheaf\par
\noindent\textit{Evidence.} The scan repeats predpuchok, but the heading, canonical map and universal property identify sheafification.\par
\subsection*{P000302 (both; faithful error): Rota conjecture, order of quantifiers}
\noindent\textit{Transcription reads.} \(\ell,k\) are fixed, then \(N\), with no bound on the number of partition classes.\par
\noindent\textit{Critical reading.} Fix the finite number \(r\) of colors before choosing \(N\).\par
\noindent\textit{Evidence.} The scan and Russian agree in omitting \(r\). A coloring with one color per \(\ell\)-space is a counterexample to unrestricted colors. The verified 1972 announcement concerns Ramsey partitions into fixed finitely many classes.\par
\subsection*{P000302 (both; ambiguity): elliptic-system qualifier}
\noindent\textit{Transcription reads.} on a manifold without lower-order terms\par
\noindent\textit{Critical reading.} the system has no lower-order terms; the source does not say without boundary\par
\noindent\textit{Evidence.} The Russian text places bez before mladshikh chlenov. This prevents a tempting but unsupported reconstruction as a closed-manifold problem.\par
\subsection*{P000302 (both; faithful error): regular centralizer condition}
\noindent\textit{Transcription reads.} centralizer rank n minus 1\par
\noindent\textit{Critical reading.} centralizer dimension equals rank G\par
\noindent\textit{Evidence.} The scan distinguishes rang from dim; the correct regularity condition is dimension, and the role of n is not specified. Steinberg 1965, p. 49.\par
\subsection*{P000302 (both; faithful error): subregular centralizer condition}
\noindent\textit{Transcription reads.} dimension n minus 2\par
\noindent\textit{Critical reading.} dimension rank G plus 2 in the usual subregular setting\par
\noindent\textit{Evidence.} For the stated two-equal-eigenvalues example in \(SL_n\), a block calculation gives n plus 1. The source n minus 2 cannot be the centralizer dimension.\par
\subsection*{P000303 (both; ambiguity): free presentation diagram}
\noindent\textit{Transcription reads.} \(A^m\to A^n\to M\) with the same finite-looking exponents for sheaves\par
\noindent\textit{Critical reading.} Use arbitrary direct-sum index sets, or assume M is finitely presented\par
\noindent\textit{Evidence.} The preceding lecture explicitly allowed arbitrary sets. Finite presentations do not characterize all quasi-coherent sheaves.\par
\subsection*{P000303-P000304 (both; provenance): catalogue source boundary}
\noindent\textit{Transcription reads.} Record 190 is sometimes assigned an overlapping catalogue page interval\par
\noindent\textit{Critical reading.} Record 190 runs from the Kazhdan heading on P000303 through the natural-map paragraph on P000304, before Thompson\par
\noindent\textit{Evidence.} The horizontal separator and new Thompson heading mark the end. Rota on P000302 and Thompson below the separator are independent catalogue records.\par
\subsection*{P000304 (both; faithful error): Conway group}
\noindent\textit{Transcription reads.} A finite simple group is the lattice automorphism group.\par
\noindent\textit{Critical reading.} Take the quotient of \(\operatorname{Aut}(\Lambda)\) by its central \(\{\pm I\}\).\par
\noindent\textit{Evidence.} The scan makes no quotient visible. Central inversion is an immediate obstruction to literal simplicity; the opened Chermak preprint identifies the automorphism group as \(2\cdot\mathrm{Co}_1\).\par
\subsection*{P000304 (both; faithful error): Thompson classification}
\noindent\textit{Transcription reads.} \(p>5\): only classical groups, except \(E_8\).\par
\noindent\textit{Critical reading.} \(p\geq5\); central-product reduction, then a specified list of Lie-type simple quotients, including exceptional types other than \(E_8\).\par
\noindent\textit{Evidence.} The scan faithfully has \ru{классич.}; Thompson's ICM Main theorem includes \(G_2,F_4,E_6,E_7\), and the preceding Central product theorem supplies an omitted hypothesis/reduction.\par
\subsection*{P000305 (both; faithful error): Leech lattice minimum}
\noindent\textit{Transcription reads.} \(\min(\lambda,\lambda)=4\).\par
\noindent\textit{Critical reading.} \(\min_{0\ne\lambda\in\Lambda}(\lambda,\lambda)=4\).\par
\noindent\textit{Evidence.} The scan omits the range. The zero vector otherwise makes the minimum zero.\par
\subsection*{P000305 (both; ambiguity): Proposed skyscraper example}
\noindent\textit{Transcription reads.} Fibre \(A\) at \(x\), zero elsewhere; automatically a non-quasi-coherent sheaf of modules.\par
\noindent\textit{Critical reading.} The \(A_{\mathfrak p}\)-action and closedness/support conditions must be specified; usual closed-point skyscrapers are quasi-coherent.\par
\noindent\textit{Evidence.} The scan contains no action or qualification. The correction concerns missing mathematical data, not an English mistranslation.\par
\subsection*{P000305 (both; faithful error): Signature remark}
\noindent\textit{Transcription reads.} Every positive-definite form has signature divisible by eight.\par
\noindent\textit{Critical reading.} The even integral unimodular lattice hypotheses must be retained.\par
\noindent\textit{Evidence.} The scan and Russian preserve only positivity in this sentence; a rank-one positive form is a counterexample to the unrestricted claim.\par
\subsection*{P000307 (both; faithful error): Cylinder probability}
\noindent\textit{Transcription reads.} \(\varphi(e_{\alpha_1}^{i_1}\otimes\cdots\otimes e_{\alpha_n}^{i_n})\).\par
\noindent\textit{Critical reading.} Apply \(\varphi\) to the tensor product of the associated rank-one projections.\par
\noindent\textit{Evidence.} \(\varphi\) is defined on matrices, not on Hilbert-space vectors. This is a typing correction to the mathematical abbreviation.\par
\subsection*{P000310 (both; faithful error): Čech differential target indices}
\noindent\textit{Transcription reads.} The left side of the differential formula omits the final target index.\par
\noindent\textit{Critical reading.} \((dc)_{i_0\ldots i_{q+1}}\).\par
\noindent\textit{Evidence.} The right side has \(q+2\) alternating faces; a degree-\(q+1\) cochain has \(q+2\) indices.\par
\subsection*{P000311 (both; faithful error): Middle equality of the primitive calculation}
\noindent\textit{Transcription reads.} An additional prefactor \(1/(f_{i_0}^Nf_{i_1}^N)\) multiplies already localized cocycle terms.\par
\noindent\textit{Critical reading.} Drop the additional prefactor, keeping the identity \(\sum_i L_i f_i^N(m_{i_0i}-m_{i_1i})=m_{i_0i_1}\).\par
\noindent\textit{Evidence.} Scan checked: the extra factor is present in the intermediate line, but disappears in the next line. Direct substitution into the displayed fractions detects the extra denominator.\par
\subsection*{P000311 (both; faithful error): Partition-of-unity coefficients}
\noindent\textit{Transcription reads.} \(L_i\in X_i\).\par
\noindent\textit{Critical reading.} \(L_i\in A\).\par
\noindent\textit{Evidence.} The scan visibly has \(X_i\), faithfully transcribed. \(X_i\) is an open set, whereas the identity \(\sum L_if_i^N=1\) requires ring coefficients.\par
\subsection*{P000312 (both; faithful error): Local Hilbert-space display}
\noindent\textit{Transcription reads.} \(H_I=\bigoplus_{i\in I}H_i\).\par
\noindent\textit{Critical reading.} \(H_I=\bigotimes_{i\in I}H_i\).\par
\noindent\textit{Evidence.} The scan has the direct-sum-like circled plus. The adjacent tensor basis, \(a\otimes1\) inclusion, and partial trace require tensor products.\par
\subsection*{P000313 (both; faithful error): Product eigenbasis assertion}
\noindent\textit{Transcription reads.} Commuting embedded densities always admit a product eigenbasis.\par
\noindent\textit{Critical reading.} This follows under simple one-site spectra; it is false in general with degeneracies.\par
\noindent\textit{Evidence.} The scan explicitly mentions simple spectrum as the obvious case, then asserts existence more generally. The checked faithful Bell-pair counterexample refutes that extension.\par
\subsection*{P000314 (both; faithful error): Quiver equivalence}
\noindent\textit{Transcription reads.} All vector-space diagrams of the displayed four-armed shape.\par
\noindent\textit{Critical reading.} The full subcategory with injective arrows into the central vertex.\par
\noindent\textit{Evidence.} Subspace inclusions are injective; an arbitrary diagram may have outer-vertex kernels and is not literally an ordered-subspace system.\par
\subsection*{P000315 (both; faithful error): Definition of defect}
\noindent\textit{Transcription reads.} \(\rho=2\dim P-\sum_i\dim E_i\), with subsequent labels \(S_3(2k,-1)\), \(S_3(2k,1)\).\par
\noindent\textit{Critical reading.} Either negate every displayed defect label, or use the published convention \(\delta=\sum_i\dim E_i-2\dim P\).\par
\noindent\textit{Evidence.} Scan confirms the boxed sign. Counting the dimensions of the very next two displayed bases gives the reverse labels; the original 1971 paper uses the latter convention.\par
\subsection*{P000316 (both; faithful error): Anharmonic orbit of the cross ratio}
\noindent\textit{Transcription reads.} Five listed transformations.\par
\noindent\textit{Critical reading.} Include \(\lambda/(\lambda-1)\) as the sixth, allowing coincidences at special values.\par
\noindent\textit{Evidence.} The scan has only five entries. The earlier P000314 generators include the missing transform. The 1971 announcement repeats the compressed five-entry list, so this is not solely an OCR error.\par
\subsection*{P000316 (both; faithful error): Odd-dimensional classification}
\noindent\textit{Transcription reads.} Defect zero does not occur in odd dimension.\par
\noindent\textit{Critical reading.} Odd-dimensional exceptional defect-zero indecomposables occur.\par
\noindent\textit{Evidence.} Scan checked. \((k;k,0,k,0)\) is an immediate counterexample; Claim 2 of the original announcement explicitly lists the family.\par
\subsection*{P000316-P000317 (both; faithful error): Genericity and stability slogans}
\noindent\textit{Transcription reads.} General position has no parameter; every parameter-free indecomposable is stable.\par
\noindent\textit{Critical reading.} Genericity and rigidity depend on the dimension vector and orbit; exceptional forms can deform to parameter families.\par
\noindent\textit{Evidence.} Four distinct lines already vary in cross ratio. The notebook's own displayed \(k=1\) defect-zero exceptional form is a coincident-line limit of distinct-line systems.\par
\subsection*{P000317 (both; faithful error): Trial primitive for a one-cocycle}
\noindent\textit{Transcription reads.} With \(m_{ij}=n_i-n_j\) and \(n_1=0\), take \(n_i=m_{1i}\).\par
\noindent\textit{Critical reading.} Take \(n_i=-m_{1i}=m_{i1}\).\par
\noindent\textit{Evidence.} Scan checked. The final P000318 formula \(n_i^{(k)}=m_{ik}^{(k)}\) already has the correct orientation.\par
\subsection*{P000319 (both; faithful error): Powers theorem quantifiers}
\noindent\textit{Transcription reads.} \(\forall x\,\forall\varepsilon>0\,\exists y\) with the estimate.\par
\noindent\textit{Critical reading.} \(\forall x\,\forall\varepsilon>0\,\exists n\,\forall y\in\mathcal A([-n,n])^{\prime}\cap\mathcal A\) with the estimate.\par
\noindent\textit{Evidence.} The scan visibly writes existential \(y\), making the condition vacuous. Powers Theorem 2.5 and its proof supply the uniform quantifiers.\par
\subsection*{P000320 (both; provenance): Catalog date of record 201}
\noindent\textit{Transcription reads.} \texttt{16\_Nov-70}; no ISO date.\par
\noindent\textit{Critical reading.} 1970-11-16.\par
\noindent\textit{Evidence.} The manuscript header clearly reads 16/XI; the catalog's underscore spelling prevented its supplied ISO rendering.\par
\subsection*{P000320 (both; ambiguity): General-base zero notation}
\noindent\textit{Transcription reads.} The notation \(\{0\}\) denotes the ideal \((x_1,\ldots,x_n)\) and, explicitly, all prime ideals containing it.\par
\noindent\textit{Critical reading.} In modern notation the removed closed subscheme is the zero section \(V(x_1,\ldots,x_n)\cong\operatorname{Spec}A\).\par
\noindent\textit{Evidence.} The Russian text and scan P320 explicitly include all containing primes and distinguish the field case. This is a clarification of notation already correctly explained in the source, not a mathematical error.\par
\subsection*{P000320 (both; faithful error): Intersection-localization display}
\noindent\textit{Transcription reads.} The list of inverted variables starts after the first indicated index.\par
\noindent\textit{Critical reading.} Invert every variable indexing the intersection, including \(x_{k_0}\).\par
\noindent\textit{Evidence.} By the preceding definition the intersection is \(D(x_{k_0}\cdots x_{k_q})\); omission would describe a different open set.\par
\subsection*{P000323 (both; faithful error): Top-degree module dual}
\noindent\textit{Transcription reads.} \(\operatorname{Hom}_{R}(R,k)\), or \(\operatorname{Hom}_{H^0}(H^0,A)\).\par
\noindent\textit{Critical reading.} The graded restricted linear dual over \(k\), respectively \(A\), with contragredient \(R\)-action and degree shift.\par
\noindent\textit{Evidence.} Scan checked: the ring subscript and evaluation-at-zero explanation are explicit. Ring-linear Hom is one-dimensional over \(k\), whereas the preceding negative-monomial module is infinite-dimensional.\par
\subsection*{P000324 (both; provenance): Date of the program}
\noindent\textit{Transcription reads.} Catalog record 202: 23-Nov-70; manuscript 22/XI.\par
\noindent\textit{Critical reading.} Retain the catalog date as catalog provenance and explicitly report the manuscript date 22 November.\par
\noindent\textit{Evidence.} Scan checked: both digits of 22 are visible. The surviving page does not determine whether the plan was written the previous day or the catalog date is mistaken.\par
\subsection*{P000325 (both; faithful error): Primes defining projective space}
\noindent\textit{Transcription reads.} Homogeneous prime ideals not contained in \((x_0,\ldots,x_n)\).\par
\noindent\textit{Critical reading.} Homogeneous prime ideals not containing \((x_0,\ldots,x_n)\).\par
\noindent\textit{Evidence.} Both transcriptions and the scan preserve the reversed inclusion; \(\operatorname{Proj}\) excludes primes containing the irrelevant ideal.\par
\subsection*{P000325 (both; ambiguity): Toeplitz index sign}
\noindent\textit{Transcription reads.} The index is the number of turns.\par
\noindent\textit{Critical reading.} For positive circle orientation and \(\operatorname{ind}=\dim\ker-\dim\operatorname{coker}\), use minus the winding number.\par
\noindent\textit{Evidence.} The convention is unstated in the source; \(T_z\) is the unilateral shift with index \(-1\). Gohberg 1967, Theorem 1, gives the same sign.\par
\subsection*{P000326 (both; ambiguity): Hopf sheaf convention}
\noindent\textit{Transcription reads.} \(\mathcal O(1)\) is the Hopf sheaf.\par
\noindent\textit{Critical reading.} With the displayed grading, \(\mathcal O(1)\) is the dual of the tautological \(\mathcal O(-1)\).\par
\noindent\textit{Evidence.} The name Hopf is convention-dependent; the section-degree convention is unambiguous.\par
\subsection*{P000326-P000328 (both; faithful error): Scope of coherence and Serre finiteness}
\noindent\textit{Transcription reads.} The coefficient ring A is not restricted in the displayed theorem.\par
\noindent\textit{Critical reading.} Assume A Noetherian for the stated coherence and finite-generation proof.\par
\noindent\textit{Evidence.} The proof uses coherent kernels and finite generation of submodules of finite modules, which are justified in the standard Noetherian formulation.\par
\subsection*{P000328 (both; faithful error): Euler finite-difference equation}
\noindent\textit{Transcription reads.} \(\chi(\mathcal F(k+1))-\chi(\mathcal F(k))=\overline{\mathcal F(k)}\).\par
\noindent\textit{Critical reading.} Put \(\chi(\overline{\mathcal F}(k))\) on the right.\par
\noindent\textit{Evidence.} Scan checked: no \(\chi\) is visible on the right. The equality is between integers and uses additivity of Euler characteristic.\par
\subsection*{P000328 (both; faithful error): Hyperplane injection}
\noindent\textit{Transcription reads.} An arbitrary coordinate multiplication \(\mathcal F\to\mathcal F(1)\) is injective.\par
\noindent\textit{Critical reading.} Choose a linear form regular on the sheaf, for example a general form over an infinite field.\par
\noindent\textit{Evidence.} A sheaf supported on \(x_i=0\) is annihilated by \(x_i\); Stacks Lemma 33.35.3 gives the missing avoidance condition.\par
\subsection*{P000328 (both; ambiguity): Twisted middle cohomology term}
\noindent\textit{Transcription reads.} \(H^{i_0}(\mathcal F(k))\) between terms twisted by \(k^{\prime}\).\par
\noindent\textit{Critical reading.} \(H^{i_0}(\mathcal F(k^{\prime}))\).\par
\noindent\textit{Evidence.} The preceding short exact sequence twists every term by the same new parameter; the scan lacks a clear prime in the middle term.\par
\subsection*{P000329 (both; ambiguity): Bundle and degree in the Bott bound}
\noindent\textit{Transcription reads.} \(\operatorname{Pontr}_i^k(\ldots)=0\) for \(k>2\dim\) of the bundle.\par
\noindent\textit{Critical reading.} The normal bundle \(TM/E\), real rank \(q\), has Pontryagin-ring classes zero in total cohomological degree greater than \(2q\).\par
\noindent\textit{Evidence.} The displayed submersion sequence identifies the intended quotient bundle. The source leaves the degree convention compressed; Bott's primary later exposition identifies the normal Pontryagin classes.\par
\subsection*{P000329-P000331 (both; provenance): Record 207 attribution}
\noindent\textit{Transcription reads.} Catalogued Kazhdan lecture dated 8 February over overlapping pages.\par
\noindent\textit{Critical reading.} Retain the catalog entry, but do not assign the oscillator or the cancelled intervention to it without additional evidence.\par
\noindent\textit{Evidence.} Scans P000330 and P000331 checked. The first name is cancelled and ambiguous; the second is within the continuing Bott/Fedosov block. No independent positive heading establishes the catalog boundary.\par
\subsection*{P000330 (both; ambiguity): Anharmonic-oscillator branch-order slogan}
\noindent\textit{Transcription reads.} Second-order branch points tend to a third-order branch point at zero.\par
\noindent\textit{Critical reading.} Distinguish infinitely accumulating singularities at a nonisolated origin from a global three-sheeted continuation; do not infer an isolated cubic Puiseux point.\par
\noindent\textit{Evidence.} Bender--Wu explicitly says the origin is nonisolated. Simon's rigorous follow-up is securely identified bibliographically but was not read in full.\par
\subsection*{P000330 (both; provenance): Bott bibliography date}
\noindent\textit{Transcription reads.} 1965, 66; Michigan.\par
\noindent\textit{Critical reading.} Bott, Vector fields and characteristic numbers, Michigan Math. J. 14(2) (1967), 231–244.\par
\noindent\textit{Evidence.} The complete citation on P000331 and the opened author bibliography agree on 1967. Earlier circulation was not established by this search.\par
\subsection*{P000330 (both; faithful error): Chern classes versus numbers}
\noindent\textit{Transcription reads.} Using the field one can compute all Chern classes.\par
\noindent\textit{Critical reading.} The displayed isolated-zero formula computes all Chern numbers.\par
\noindent\textit{Evidence.} The next lines define the corresponding Chern number by integration. The mathematical formula and its primary modern verification concern characteristic numbers.\par
\subsection*{P000330 (both; faithful error): Polynomial hypothesis}
\noindent\textit{Transcription reads.} \(\psi(X_1,\ldots,X_m)\) is homogeneous of degree \(m\).\par
\noindent\textit{Critical reading.} Require \(\psi\) symmetric, equivalently use an invariant polynomial of matrices; its degree in Chern roots is \(m\).\par
\noindent\textit{Evidence.} The source subsequently expresses the polynomial through elementary symmetric functions, which requires symmetry. The degree convention differs from weighted Chern-class degree only by that expression.\par
\subsection*{P000331 (both; provenance): 15 February heading and preceding text}
\noindent\textit{Transcription reads.} Catalog range begins at P000330 for record 209.\par
\noindent\textit{Critical reading.} Record 209 begins at the dated heading on P000331; the preceding holomorphic-vector-field argument belongs to record 208.\par
\noindent\textit{Evidence.} EN/RU headings and P000331 scan; checked jointly with the adjoining segment.\par
\subsection*{P000331 (both; faithful error): polar divisor cohomology}
\noindent\textit{Transcription reads.} The divisor produces two-dimensional cycles.\par
\noindent\textit{Critical reading.} The divisor has real codimension two and defines a degree-two cohomology class; its homology dimension is \(2n-2\).\par
\noindent\textit{Evidence.} P000331 scan plus the local hypersurface description.\par
\subsection*{P000332 (both; ambiguity): real fixed-point denominator}
\noindent\textit{Transcription reads.} Square root of determinant of the skew linearization.\par
\noindent\textit{Critical reading.} Use the signed Pfaffian with the chosen orientation and consistent characteristic-class normalization.\par
\noindent\textit{Evidence.} The square root alone loses the sign required by the real localization formula.\par
\subsection*{P000332 (both; ambiguity): regularized Chern classes}
\noindent\textit{Transcription reads.} Apply the regularized coefficients without explaining the bundle.\par
\noindent\textit{Critical reading.} Interpret \(a_i\) as coefficients of \(c(TX)/c(\mathcal O(-D))\), not \(c(TX\otimes\mathcal O(D))\).\par
\noindent\textit{Evidence.} Notebook expansion and Bruzzo--Rubtsov Proposition 4.6.\par
\subsection*{P000335 (both; provenance): record 214 heading}
\noindent\textit{Transcription reads.} Catalog suggests substantive text over neighboring pages.\par
\noindent\textit{Critical reading.} Record 214 is the Griffiths bibliographical heading only; Bernstein begins the subsequent singularity exposition.\par
\noindent\textit{Evidence.} Speaker headings in EN/RU P000335.\par
\subsection*{P000336 (both; ambiguity): torsion preprint authors}
\noindent\textit{Transcription reads.} Singer and someone else.\par
\noindent\textit{Critical reading.} Ray and Singer.\par
\noindent\textit{Evidence.} Ray--Singer 1971 article metadata; historical notebook attribution is incomplete.\par
\subsection*{P000337 (both; faithful error): positive definiteness}
\noindent\textit{Transcription reads.} Nonnegative \(D\) and nonnegative Schur complement imply positive definiteness.\par
\noindent\textit{Critical reading.} For strict positivity require Hermitian symmetry, \(D>0\), and strictly positive Schur complement. Semidefinite and singular cases have separate conditions.\par
\noindent\textit{Evidence.} P000337 scan and block congruence; item 1970-216.2.\par
\subsection*{P000339 (both; faithful error): lens-space complex}
\noindent\textit{Transcription reads.} The complex is acyclic without stating the twist.\par
\noindent\textit{Critical reading.} The complex becomes acyclic after a nontrivial character twist; the untwisted universal cover retains ordinary homology.\par
\noindent\textit{Evidence.} P000339 scan and Milnor Section 12.\par
\subsection*{P000339 (both; faithful error): versal deformation terminology}
\noindent\textit{Transcription reads.} Final object.\par
\noindent\textit{Critical reading.} Versal object: inducing maps exist but need not be unique.\par
\noindent\textit{Evidence.} The displayed deformation equivalence permits stabilizers; Arnold matrix-deformation framework.\par
\subsection*{P000341-P000342 (both; faithful error): KdV equation and Hamiltonian}
\noindent\textit{Transcription reads.} \(u_t+2u_{xxx}+6uu_x=0\) together with \(Q=\int(u^3+u_x^2)\,dx\).\par
\noindent\textit{Critical reading.} With \(J=\partial_x\), the displayed \(Q\) yields \(u_t=6uu_x-2u_{xxx}\). A standard consistent normalization uses \(H=\int(u^3+\tfrac12u_x^2)\,dx\), giving \(u_t=6uu_x-u_{xxx}\).\par
\noindent\textit{Evidence.} Direct variation of \(Q\); Zakharov--Faddeev equations (1), (6)--(7).\par
\subsection*{P000342-P000343 (both; ambiguity): scattering angle}
\noindent\textit{Transcription reads.} Argument of the reflection coefficient as the canonical angle, without convention changes.\par
\noindent\textit{Critical reading.} The cited published canonical pair uses \(\arg b(k)\), while the physical reflection coefficient is \(b(k)/a(k)\); do not identify these phases without the corresponding transformation.\par
\noindent\textit{Evidence.} Zakharov--Faddeev equation (23), p.285.\par
\subsection*{P000344 (both; provenance): Faddeev/Gabrielov boundary}
\noindent\textit{Transcription reads.} Broad catalog ranges overlap both subjects.\par
\noindent\textit{Critical reading.} Faddeev ends in the upper part of P000344; Gabrielov starts at his explicit lower heading and continues through P000345.\par
\noindent\textit{Evidence.} P000344 scan and both transcriptions.\par
\subsection*{P000344 (both; faithful error): unitarity of scattering}
\noindent\textit{Transcription reads.} Scattering amplitude is unitary.\par
\noindent\textit{Critical reading.} The full scattering operator on the energy sphere, including its identity term and normalization, is unitary. The amplitude alone is not.\par
\noindent\textit{Evidence.} Faddeev 1966 equations (5)--(8).\par
\subsection*{P000345 (both; ambiguity): tube notation}
\noindent\textit{Transcription reads.} Full neighborhood \(T_\varepsilon(\Gamma)\) and ambiguous \(n/m\) index.\par
\noindent\textit{Critical reading.} Use the normal circle bundle over \(\gamma\) in \(V\); for \(n=2m\) intersect \(V\) with \(\mathbb P^{m+1}\).\par
\noindent\textit{Evidence.} P000345 argument and its repeated exact-sequence formulation on P000347--348.\par
\subsection*{P000346 (both; ambiguity): Gudkov congruence}
\noindent\textit{Transcription reads.} \(p-m\equiv k^2\pmod8\) without a repeated maximality condition.\par
\noindent\textit{Critical reading.} Carry forward the \(M\)-curve hypothesis from the maximal-oval paragraph.\par
\noindent\textit{Evidence.} Arnold Section 1 and Rokhlin Theorem in Section 1.\par
\subsection*{P000346 (both; ambiguity): Gudkov count}
\noindent\textit{Transcription reads.} 50 cases for sextics.\par
\noindent\textit{Critical reading.} Retain as an unverified, underspecified count; do not replace it by the number of all real isotopy schemes.\par
\noindent\textit{Evidence.} The scan shows 50 but does not specify what is counted; the exact intended enumeration was not found.\par
\subsection*{P000348 (both; faithful error): kernel of tube map}
\noindent\textit{Transcription reads.} \(\ker\tau=H_{n+1}(\mathbb P^n)\).\par
\noindent\textit{Critical reading.} \(\ker\tau\) is the image of the intersection map \(H_{n+1}(\mathbb P^n)\to H_{n-1}(V)\).\par
\noindent\textit{Evidence.} The preceding exact sequence on the same leaf.\par
\subsection*{P000348 (both; faithful error): proof of exactness}
\noindent\textit{Transcription reads.} A closed ambient chain is therefore a boundary.\par
\noindent\textit{Critical reading.} A closed chain defines an ambient homology class; intersect that class with \(V\) to obtain the kernel element.\par
\noindent\textit{Evidence.} Closed does not imply bounding in projective space; pair-sequence argument.\par
\subsection*{P000349 (both; provenance): Gabrielov/Palamodov boundary}
\noindent\textit{Transcription reads.} Broad catalog ranges attach the whole leaf to both records.\par
\noindent\textit{Critical reading.} Gabrielov occupies the upper pole-reduction part; Palamodov starts at the lower dated speaker heading.\par
\noindent\textit{Evidence.} P000349 scan and both transcriptions.\par
\subsection*{P000349 (both; faithful error): Palamodov two-equation example}
\noindent\textit{Transcription reads.} \(D_xu=f\), \((D_y+\lambda)u=g\) imply \(u=D_yg-D_xf\).\par
\noindent\textit{Critical reading.} For constant \(\lambda\) the compatibility is \(D_xg-D_yf=\lambda f\), and it does not determine \(u\). The intended example is unidentified.\par
\noindent\textit{Evidence.} P000349 scan confirms \(\lambda\); direct commuting-derivative calculation. The alternative coefficient \(x\) in the commentary is explicitly editorial.\par
\subsection*{P000349 (both; faithful error): numerator degree for pole order ell}
\noindent\textit{Transcription reads.} \(\deg P=q-(n+1)\).\par
\noindent\textit{Critical reading.} \(\deg P=\ell q-(n+1)\); the written formula is only the simple-pole case.\par
\noindent\textit{Evidence.} Scaling of \(P\Omega/Q^\ell\) and the earlier simple-pole formula.\par
\subsection*{P000350 (both; provenance): Rolewicz/Arnold boundary}
\noindent\textit{Transcription reads.} The leaf appears in both record ranges.\par
\noindent\textit{Critical reading.} Rolewicz is the upper orbit discussion; Arnold begins the lower unimodular-form discussion and continues on P000351.\par
\noindent\textit{Evidence.} P000350 EN/RU headings.\par
\subsection*{P000350 (both; ambiguity): weighted shift space}
\noindent\textit{Transcription reads.} \(\ell^p\) without an index range.\par
\noindent\textit{Critical reading.} \(1\le p<\infty\), or \(c_0\); \(\ell^\infty\) is nonseparable and cannot have a countable dense orbit.\par
\noindent\textit{Evidence.} Rolewicz example and separability.\par
\subsection*{P000351 (both; ambiguity): stable invariants}
\noindent\textit{Transcription reads.} Stably only rank and signature remain.\par
\noindent\textit{Critical reading.} Specify stabilization: adding even hyperbolic planes preserves parity; adding the same odd indefinite plane makes the equal-rank/signature forms isometric.\par
\noindent\textit{Evidence.} Unimodular classification; the handwritten disagreement does not specify the intended stabilization.\par
\subsection*{P000352 (both; ambiguity): sheaf valued in arbitrary category K}
\noindent\textit{Transcription reads.} Exact-sequence notation starting with \(0\).\par
\noindent\textit{Critical reading.} Use an equalizer in a category with the needed limits; use \(0\) and a difference map only for the abelian formulation.\par
\noindent\textit{Evidence.} Sheaf equalizer axiom; Stacks Tag 00VM.\par
\subsection*{P000353 (both; faithful error): Lie-group site morphisms}
\noindent\textit{Transcription reads.} Every morphism is a regular equivariant map with surjective derivative.\par
\noindent\textit{Critical reading.} All smooth equivariant maps are morphisms; jointly surjective families of submersions are covers.\par
\noindent\textit{Evidence.} Geronimus Section 3.1, full Russian primary article; scan confirms the notebook wording.\par
\subsection*{P000354 (both; ambiguity): elliptic curve exponent}
\noindent\textit{Transcription reads.} \(y^3=x^3-g_1x-g_2\).\par
\noindent\textit{Critical reading.} Retain \(y^3\) as the scan reading; if smooth, its projective cubic with rational point can be elliptic. A Weierstrass \(y^2\) may have been intended but is not established.\par
\noindent\textit{Evidence.} P000354 scan, homogeneous cubic and rational point \([1:1:0]\).\par
\subsection*{P000354 (both; faithful error): finitely many maps}
\noindent\textit{Transcription reads.} There are finitely many maps from every curve to a torus.\par
\noindent\textit{Critical reading.} False without restrictions: translations and compositions with multiplication maps yield infinitely many.\par
\noindent\textit{Evidence.} Translations provide a direct counterexample; nonconstant maps have unbounded degrees after multiplication.\par
\subsection*{P000354 (both; faithful error): fractional-linear action}
\noindent\textit{Transcription reads.} Matrix \(\left(\begin{smallmatrix}\alpha&\beta\\A\gamma&\delta\end{smallmatrix}\right)\) acts by \((\alpha z+A\gamma)/(\beta z+\delta)\).\par
\noindent\textit{Critical reading.} For that matrix the standard action is \((\alpha z+\beta)/(A\gamma z+\delta)\).\par
\noindent\textit{Evidence.} P000354 scan and the defining matrix-action convention.\par
\subsection*{P000354-P000355 (both; ambiguity): conductor table}
\noindent\textit{Transcription reads.} Product of factors \(1,p,p^2\) over all primes.\par
\noindent\textit{Critical reading.} Use a minimal Weierstrass model and include exceptional-prime conductor exponents at 2 and 3; the displayed table explicitly concerns \(p\ge5\).\par
\noindent\textit{Evidence.} The source brackets \(p\ne2,3\); Cremona Section 3.2.\par
\subsection*{P000355 (both; faithful error): Dirichlet denominator}
\noindent\textit{Transcription reads.} \(c_n/n_s\).\par
\noindent\textit{Critical reading.} \(c_n/n^s\) in the intended Dirichlet series.\par
\noindent\textit{Evidence.} The scan puts \(s\) below \(n\); this is a mathematical emendation of the faithful transcription.\par
\subsection*{P000355 (both; ambiguity): Fourier expansion}
\noindent\textit{Transcription reads.} \(\sum_z c_ne^{inz}\).\par
\noindent\textit{Critical reading.} Use \(\sum_{n\ge1}a_ne^{2\pi inz}\) in the standard coordinate; an exponent \(inz\) needs a compatible rescaling.\par
\noindent\textit{Evidence.} P000355 scan and Cremona Section 2.6.\par
\subsection*{P000356 (both; faithful error): idempotent eigenvalues}
\noindent\textit{Transcription reads.} \(\lambda=0,+1,-1\) from \(\lambda^2=\lambda\).\par
\noindent\textit{Critical reading.} Only 0 and 1; minus one is excluded algebraically.\par
\noindent\textit{Evidence.} P000356 scan and factorization \(\lambda(\lambda-1)=0\).\par
\subsection*{P000357 (both; faithful error): complex-shift Gaussian identity}
\noindent\textit{Transcription reads.} \(\sum_{m,n}(-1)^{mn+m+n}\exp(-(x-\lambda_{mn})^2/2)\).\par
\noindent\textit{Critical reading.} Supply normalized coherent wavefunctions with the required momentum phase; the raw complex-shift terms can grow and do not define the claimed convergent sum.\par
\noindent\textit{Evidence.} P000357 scan and the explicit normalized displacement formula; item 1970-228.15.\par
\subsection*{P000358 (both; ambiguity): uncertainty scope}
\noindent\textit{Transcription reads.} \(\delta_x\delta_p\) is bounded after orthogonalization.\par
\noindent\textit{Critical reading.} Preserve the product. Separate general orthonormal-basis existence from a prescribed Gram--Schmidt order and from a Gabor-covariant construction.\par
\noindent\textit{Evidence.} P000358 scan, Meyer Theorem 1, Bourgain main theorem, Grochenig--Malinnikova pp.115--116.\par
\subsection*{P000358-P000359 (both; ambiguity): induced Heisenberg lattice model}
\noindent\textit{Transcription reads.} Induce only from the discrete subgroup and its signs.\par
\noindent\textit{Critical reading.} Induce from \(T\Gamma\) with the tautological central character \(\chi(t\gamma)=t\), hence \(U_t=tI\). Principality alone does not give multiplicity one for every nontrivial central character.\par
\noindent\textit{Evidence.} Cartier Section 1.2(b), pp.420--421.\par
\subsection*{P000359 (both; faithful error): number of vanishing theta characteristics}
\noindent\textit{Transcription reads.} Half of the \(4^N\) choices vanish.\par
\noindent\textit{Critical reading.} Odd characteristics number \(2^{N-1}(2^N-1)\), not half; additional even theta-null vanishing may occur for special higher-dimensional data.\par
\noindent\textit{Evidence.} Correct count already written on P000358; Cartier quadratic-refinement parity.\par
\subsection*{P000359 (both; faithful error): unitary displacement exponent}
\noindent\textit{Transcription reads.} \(e^{i\varphi}e^{i(\bar\alpha a-\alpha a^*)}\).\par
\noindent\textit{Critical reading.} Remove the extra inner \(i\) for the standard skew-adjoint displacement generator, with an overall sign fixed by convention.\par
\noindent\textit{Evidence.} P000359 scan and Perelomov equations (60)--(71).\par
\subsection*{P000360 (both; provenance): adjacent next-year leaf}
\noindent\textit{Transcription reads.} Supplied record 228 range extends through P000360.\par
\noindent\textit{Critical reading.} P000360 begins Kirillov on 6 September 1971, record 229 of the next academic year; exclude its quantization discussion from the April record.\par
\noindent\textit{Evidence.} P000360 has a Kirillov heading but no written date. P000359 announces the next meeting on the first Monday of September. The parent independently checked catalog record229, Kirillov on6September1971. The date and next-year assignment are a supported inference from that chain, not a date visibly written on P000360.\par
